\documentclass[11pt]{article}
% Preamble for the eiPartialID technical report (pdflatex).
\usepackage[T1]{fontenc}
\usepackage[utf8]{inputenc}
\usepackage{lmodern}
\usepackage[a4paper,margin=2.4cm]{geometry}
\usepackage{amsmath,amssymb,amsthm,mathtools}
\usepackage{longtable,booktabs,array,calc,tabularx}
\usepackage{enumitem}
\usepackage{xcolor}
\usepackage{newunicodechar}
\usepackage{etoolbox}
\usepackage{footnote}
\makesavenoteenv{longtable}
\usepackage[numbers,sort&compress]{natbib}
\usepackage[hidelinks]{hyperref}
\usepackage{bookmark}
\usepackage{microtype}
\setlength{\emergencystretch}{3em}
\providecommand{\tightlist}{\setlength{\itemsep}{0pt}\setlength{\parskip}{0pt}}
\providecommand{\real}[1]{#1}
\def\LTcaptype{none}
% pandoc's longtables: allow page breaks and avoid overfull boxes
\setlength{\LTleft}{0pt}\setlength{\LTright}{0pt}
\makeatletter
\patchcmd\longtable{\par}{\if@noskipsec\mbox{}\fi\par}{}{}
\makeatother

% ---------- unicode characters used in the repository's Markdown ----------
\newunicodechar{Δ}{\ensuremath{\Delta}}
\newunicodechar{Θ}{\ensuremath{\Theta}}
\newunicodechar{α}{\ensuremath{\alpha}}
\newunicodechar{β}{\ensuremath{\beta}}
\newunicodechar{γ}{\ensuremath{\gamma}}
\newunicodechar{δ}{\ensuremath{\delta}}
\newunicodechar{ε}{\ensuremath{\varepsilon}}
\newunicodechar{η}{\ensuremath{\eta}}
\newunicodechar{κ}{\ensuremath{\kappa}}
\newunicodechar{π}{\ensuremath{\pi}}
\newunicodechar{σ}{\ensuremath{\sigma}}
\newunicodechar{τ}{\ensuremath{\tau}}
\newunicodechar{χ}{\ensuremath{\chi}}
\newunicodechar{ᵃ}{\textsuperscript{a}}
\newunicodechar{ᵇ}{\textsuperscript{b}}
\newunicodechar{ᶜ}{\textsuperscript{c}}
\newunicodechar{ᵈ}{\textsuperscript{d}}
\newunicodechar{ᵉ}{\textsuperscript{e}}
\newunicodechar{ᶠ}{\textsuperscript{f}}
\newunicodechar{′}{\ensuremath{'}}
\newunicodechar{₁}{\ensuremath{_1}}
\newunicodechar{↑}{\ensuremath{\uparrow}}
\newunicodechar{→}{\ensuremath{\rightarrow}}
\newunicodechar{⇐}{\ensuremath{\Leftarrow}}
\newunicodechar{⇒}{\ensuremath{\Rightarrow}}
\newunicodechar{⇔}{\ensuremath{\Leftrightarrow}}
\newunicodechar{∈}{\ensuremath{\in}}
\newunicodechar{∎}{\ensuremath{\blacksquare}}
\newunicodechar{−}{\ensuremath{-}}
\newunicodechar{∩}{\ensuremath{\cap}}
\newunicodechar{∪}{\ensuremath{\cup}}
\newunicodechar{≈}{\ensuremath{\approx}}
\newunicodechar{≠}{\ensuremath{\neq}}
\newunicodechar{≤}{\ensuremath{\le}}
\newunicodechar{≥}{\ensuremath{\ge}}
\newunicodechar{≲}{\ensuremath{\lesssim}}
\newunicodechar{⊂}{\ensuremath{\subset}}
\newunicodechar{⊃}{\ensuremath{\supset}}
\newunicodechar{⊆}{\ensuremath{\subseteq}}
\newunicodechar{⊇}{\ensuremath{\supseteq}}
\newunicodechar{⊕}{\ensuremath{\oplus}}
\newunicodechar{⚠}{\textbf{(!)}}
\newunicodechar{✓}{\ensuremath{\checkmark}}
\newunicodechar{²}{\ensuremath{^2}}
\newunicodechar{½}{\ensuremath{\tfrac12}}
\newunicodechar{·}{\ensuremath{\cdot}}
\newunicodechar{×}{\ensuremath{\times}}
\newunicodechar{±}{\ensuremath{\pm}}
\newunicodechar{†}{\textsuperscript{\dag}}
\newunicodechar{…}{\ldots}
\newunicodechar{≡}{\ensuremath{\equiv}}

% ---------- notation (docs/extension/notation.tex) ----------
\newcommand{\Simplex}{\Delta}
\newcommand{\Bset}{\mathfrak B}
\newcommand{\Kspace}{\mathcal K}
\newcommand{\ThI}{\Theta_I}
\newcommand{\ThX}{\Theta_X}
\newcommand{\ThO}{\Theta_O}
\newcommand{\E}{\mathbb E}
\newcommand{\vertA}{\operatorname{vert}(A)}
\newcommand{\conv}{\operatorname{conv}}
\newcommand{\supp}{\operatorname{supp}}
\newcommand{\Var}{\operatorname{Var}}
\newcommand{\Cov}{\operatorname{Cov}}
\newcommand{\sym}{\operatorname{sym}}
\newcommand{\skw}{\operatorname{skew}}

% ---------- theorem environments ----------
\theoremstyle{plain}
\newtheorem{result}{Result}[section]
\theoremstyle{definition}
\newtheorem{assumption}{Assumption}
\newtheorem{definitionx}[result]{Definition}
\newenvironment{plain}{\par\smallskip\noindent\textbf{In plain terms.}\ }{\par\smallskip}
\newenvironment{sketch}{\par\smallskip\noindent\textit{Sketch.}\ }{\par\smallskip}
\newenvironment{revnote}{\par\smallskip\noindent\begin{minipage}{\linewidth}\color{black!75}\small\textbf{Revision note.}\ }{\end{minipage}\par\smallskip}
\newcommand{\src}[1]{\textup{[\emph{#1}]}}
\newcommand{\status}[1]{\textsf{\small #1}}
\newcommand{\proofref}[1]{\par\noindent\textit{Proof:} #1.\par\smallskip}


\title{Partial Identification for R$\times$C Ecological Inference:\\
Theory, Inference, Referee Revisions and the \texttt{eiPartialID} Package\\[4pt]
\large Technical report}
\author{Compiled from the \texttt{mdweaver/eiid} repository\\
\small branch \texttt{claude/wizardly-volta-lbh4bm}, commit \texttt{b270285}}
\date{November 2026}

\begin{document}
\maketitle

\begin{abstract}
\noindent This report collects, in one document, the partial-identification
programme for ecological inference developed in the \texttt{eiid} repository and
implemented in the R package \texttt{eiPartialID} (version 0.3.1). It extends
Jiang, King, Schmaltz and Tanner's $2\times2$ bounds under linear contextual
effects to general R$\times$C tables (the sets $\Theta_I\subseteq\Theta_X\subseteq\Theta_O$),
to units nested in districts (the rung ladder DD, M2, M1($\delta$), M1, L4, L5,
L6 beside the pooled M0), to sign, reinforcement and monotonicity beliefs from
Monotone Ecological Inference (N1, N2 and the M-SIGN module), and to confidence
intervals that are valid when many candidate constraints are nearly tied. Each
result is stated with its source label in the repository, explained in plain
terms, sketched, and proved in full in the appendices, which reproduce the
repository's theory documents. Every change made in response to three external
referees over two rounds is attached to the result it changed and collected in a
revision table. A status table separates proved results, results proved under
stated conditions, cited theorems, conjectures, algorithms and open problems.
\end{abstract}

\tableofcontents
\clearpage

% Section 1: purpose and scope
\section{Purpose and scope}\label{sec:purpose}

\subsection{The problem}

Ecological inference (EI) asks how groups behaved when only group shares and
outcome shares are observed for each unit. A unit is a polling place, a
precinct, a county or a riding. For unit $i$ we see the row shares
$X_i\in\Delta^{R-1}$ (for example, how its electorate voted in election~1, or
its racial composition) and the column shares $T_i\in\Delta^{C-1}$ (for
example, how it voted in election~2). Here $\Delta^{k-1}$ is the probability
simplex in $\mathbb R^k$. We want the latent $R\times C$ table of rates
$\beta_i=[\beta_{irc}]$, where $\beta_{irc}$ is the share of row-$r$ members
of unit $i$ who chose column $c$, and above all its electorate-wide average,
the transition rate $B_{rc}$ (Section~\ref{sec:setting}).

The data link rows and columns only through the accounting identity
$T_i=\beta_i^\top X_i$. In each unit there are $RC$ unknown rates and only
$C-1$ free equations. Without further structure the rates are not
point-identified.

\subsection{Why partial identification}

The classical answers either assume enough to obtain a point (Goodman's
constancy~\citep{goodman1953}; King's truncated normal~\citep{king1997}; hierarchical R$\times$C models) or give the
widest possible range (the Duncan--Davis ``method of bounds''). The first
kind of answer can be confidently wrong when its assumption fails. The second
is always valid but often too wide to be useful.

Partial identification takes the middle road. It states an assumption, derives
the set of values of the target that are consistent with the data and that
assumption (the \emph{identified set}), and reports that set. When the
assumption is weak, the set is wide but credible. When it is strong, the set is
narrow but rests on more. The researcher can see exactly what each assumption
buys. Jiang, King, Schmaltz and Tanner~\citep{jkst2019} (JKST below) did this
for $2\times2$ tables, in the tradition of partial identification~\citep{molinari2020}, under \emph{linear contextual effects} (LCE): each
group's mean rate may depend linearly on the composition of the unit. The
present work extends that programme to general R$\times$C tables, to units
nested in districts, to sign and monotonicity beliefs from Monotone EI
\citep{egghm2025}, and to valid inference for all of these.

\subsection{The plan of the work: a ladder of assumptions}

The work organises assumptions as a ladder. Each rung adds an assumption, and
under the conditions stated in each result, its identified set lies inside the
one before it. The rungs are not all on one line: some are non-nested and are
reported side by side. Table~\ref{tab:ladder} lists them with what each adds.

\begin{longtable}{@{}L{0.17\linewidth}L{0.42\linewidth}L{0.33\linewidth}@{}}
\caption{The rungs and what each adds.}\label{tab:ladder}\\
\toprule
Rung & Assumption added (plain) & What it buys \\
\midrule\endfirsthead
\toprule Rung & Assumption added (plain) & What it buys \\ \midrule\endhead
DD (Duncan--Davis) & None beyond the accounting identity (A1). & The model-free range; the baseline everyone can accept. Districts add nothing here (Prop.~D0). \\
M2 & Each district has its own linear contextual effects. & Usually little: district bounds add up (Prop.~D5). \\
M1($\delta$) & District contextual effects differ from a shared set by at most $\delta$. & A sensitivity dial from M2 ($\delta$ large) to M1 ($\delta=0$) (Prop.~D6). \\
M1 & Shared contextual effects, district baselines (``riding dummies''). & One district with lopsided units and an extreme baseline tightens the bound for all (Prop.~D3). \\
L4($r$) & Row $r$ has no contextual effect. & Pins $(C-1)(R-1)$ unknown directions; point identification when $R=2$ (Prop.~D7(b)). \\
L5 & Constancy within each district (Goodman by district). & A point (Prop.~D7(a)); best used as a diagnostic. \\
L6 & Constancy everywhere (Goodman). & A point; historical reference. \\
M0 (pooled, beside the ladder) & One set of linear contextual effects for all units, no districts. & JKST's model in R$\times$C; non-nested with M1 (Prop.~D4). \\
M-SIGN (any linear rung) & The sign of the between or within association, or that the two reinforce each other (CR). & Large gains in simulation; a false belief gives narrow, wrong intervals (PHASE5a~\S3.4). \\
N1 & The same sign beliefs on realised unit tables, no linear model. & Bounds sharp for continuous unit tables (Prop.~M2). Not in the package. \\
N2 & Each group's mean rate is monotone in composition, any shape. & An outer LP that contains LCE with signed slopes (Prop.~N2.2). Not in the package. \\
\bottomrule
\end{longtable}

What each rung adds, in one sentence each:
\begin{itemize}[nosep]
\item \textbf{LCE (M0) over DD.} The regression of $T$ on $X$ is quadratic under
LCE and identifies the symmetric part $Q$ of the coefficient matrices. Only the
skew part $K$, of dimension $(C-1)R(R-1)/2$, is left unknown, and requiring the
implied rates to be probabilities at the extreme units bounds it.
\item \textbf{Districts (M1 and relatives) over M0.} District baselines protect
against between-district aggregation bias, and shared contextual effects carry
the information of the most informative district to all districts.
\item \textbf{Signs and reinforcement (M-SIGN) over a linear rung.} One or two
extra linear rows in $K$ can cut the bound substantially; the price is a
belief that the data cannot always check.
\item \textbf{Inference.} Every bound is computed from an estimated regression.
The plug-in bound is too narrow, so confidence intervals that treat each end
as an intersection of many estimated candidates are needed (Phase~3).
\end{itemize}

\subsection{Method map: where each method is identified and where its inference is proved}

Table~\ref{tab:methodmap} lists, for each method, the results that identify its
bound, the results that cover its confidence intervals, its status, and where the
proofs are. Appendix key: \ref{app:phase1} = PHASE1, \ref{app:phase2a} = PHASE2a,
\ref{app:phase2b} = PHASE2b, \ref{app:phase2c} = PHASE2c, \ref{app:phase3} = PHASE3,
\ref{app:phase5a} = PHASE5a, \ref{app:pkgdesign} = PACKAGE\_DESIGN.

{\footnotesize
\setlength{\tabcolsep}{3pt}
\begin{longtable}{@{}L{0.13\linewidth}L{0.24\linewidth}L{0.29\linewidth}L{0.22\linewidth}L{0.07\linewidth}@{}}
\caption{Method map.}\label{tab:methodmap}\\
\toprule Method & Identification & Inference & Status & App. \\ \midrule\endfirsthead
\toprule Method & Identification & Inference & Status & App. \\ \midrule\endhead
DD & 2b Prop.~D0; 2a Prop.~8 (cell bounds) & none needed: no estimated quantity; CIs for $B$ are intersected with DD & Proved & \ref{app:phase2b}, \ref{app:phase2a} \\
JKST ($2\times2$ M0) & 2a Prop.~4 (equivalence), Props~1--2 & P3.5 (JKST's $\mathrm{CI}_x$ is (b) with one candidate); P3.4 under tie-breaking & Proved & \ref{app:phase2a}, \ref{app:phase3} \\
M0 (pooled LCE) & 2a Lemmas~1--3, Thm~1, Props~1--2, 7, 7$'$; Cor.~1 & P3.1, P3.2, P3.4/P3.4$'$; P3.5(a), (b2), (c1); P3.6(a); P3.8. CLR = Conjecture P3.5(b3) & Identification and (a), (b2), (c1) proved in every regime; CLR algorithmic & \ref{app:phase2a}, \ref{app:phase3} \\
M2 & 2b Thm~1$'$, Thm~2, Prop.~D5 & P3.5(a), (b2) in regime R1; P3.8 in any regime; P3.6(a) & Proved (R1); R3 conjecture; R2: P3.8 only & \ref{app:phase2b}, \ref{app:phase3} \\
M1($\delta$) & 2b Thm~2(c), Prop.~D6 & as M2 & as M2 & \ref{app:phase2b}, \ref{app:phase3} \\
M1 & 2b Lemmas~D1--D2, Thm~2, Props~D3--D4 & as M2; plus P3.6(b) (M1 vs M2) and P3.6(c) (pooling), wild bootstrap; P3.7. CLR in R2: simulation finding & Proved (R1); bootstrap simulation-supported & \ref{app:phase2b}, \ref{app:phase3} \\
L4($r$), L5, L6 & 2b Prop.~D7, Thm~2(c) & P3.5 (near point identification: two-sided split; P3 \S4) & Proved & \ref{app:phase2b}, \ref{app:phase3} \\
Change across periods (M1) & 2b Prop.~D8 & via the M1 interval for the change & Proved (identification) & \ref{app:phase2b} \\
N1 (signs on realised tables) & 2c Lemma~M1, Props~M2--M4, Rem.~M2.1 & none needed (no estimated quantity) & Proved for continuous tables; integer R15 open & \ref{app:phase2c} \\
N2 (monotone means) & 2c Lemma~N2.1, Prop.~N2.2, Lemma~N2.4; Conj.~N2.2$'$ & P3.5 conditional on a prespecified degree (2c \S6); R7 open & Outer set proved; sharp set conjecture & \ref{app:phase2c} \\
M-SIGN rows & 2c Prop.~M5; 5a Thm~R3, Prop.~R2 & 5a Thm~R10(a), by interval version & Proved for pooled M0; grouped rungs R16 open & \ref{app:phase2c}, \ref{app:phase5a} \\
CR & 2c Prop.~M3; 5a Thm~R3(c)--(d) & 5a Thm~R10(b) (slab, conservative), (c) (branch union); CLR union algorithmic; R10(e) (plug-in bracket invalid) & Proved for pooled M0 & \ref{app:phase2c}, \ref{app:phase5a} \\
Ordering / M-LOC & PD Props~P-S, P-R, P-N; 2c M5(b) & 5a Thm~R10(a) as added rows (pooled M0) & P-S proved; C1 conjecture; M-LOC not built & \ref{app:pkgdesign}, \ref{app:phase2c} \\
Sensitivity $\eta$, $\delta$ & P3 Prop.~S1; 2b Prop.~D6 & CI version of the bound (P3 \S7) & Outer bound proved & \ref{app:phase3}, \ref{app:phase2b} \\
\bottomrule
\end{longtable}
}

\subsection{What this report contains and how to read it}

This report collects the theory, the package mapping and two rounds of
external referee review into one document. It is pinned to commit
\texttt{ec3644f} of branch \texttt{claude/wizardly-volta-lbh4bm} of the
repository \texttt{mdweaver/eiid}.

\begin{itemize}[nosep]
\item Sections~\ref{sec:setting}--\ref{sec:did} state each result, give its
source label in the repository (for example ``Proposition P3.4$'$ (PHASE3
\S4)''), explain it in plain terms, sketch the argument and point to the full
proof in the appendices. Where a referee changed a result, a \emph{revision
note} follows it.
\item Section~\ref{sec:evidence} collects the numerical evidence. It is
evidence, not proof.
\item Section~\ref{sec:package} maps results to the R package
\texttt{eiPartialID} 0.3.1.
\item Section~\ref{sec:status} is the status table of every result.
Section~\ref{sec:open} lists what is open.
\item Appendix~\ref{app:revisions} is the consolidated table of every referee
point, the result it changed and the change made.
\item Appendices~\ref{app:phase1}--\ref{app:pkgdesign} reproduce the
repository's theory documents verbatim (converted mechanically from Markdown).
They contain every detailed proof, with the justification tag on each step.
\end{itemize}

\paragraph{Status labels.} The referees insisted that established results,
conjectures and algorithmic procedures be kept apart. Every result carries one
of these labels:
\begin{description}[nosep,font=\normalfont\sffamily]
\item[Proved] A proof is written in the repository, in numbered steps with
justification tags, and checked numerically where possible.
\item[Proved under conditions] Proved, but only under a stated scope (for
example, the pooled model M0 only, or a fixed number of districts).
\item[Cited] The result applies a published theorem whose conditions are
checked here but whose proof is not repeated.
\item[Sketch] Only an outline of a proof exists in the repository.
\item[Conjecture] Stated precisely, supported by checks or simulation, not
proved.
\item[Algorithm] A procedure whose validity is not proved (for example, the
CLR interval as implemented). Simulation evidence cannot promote it.
\item[Open] Not settled.
\item[Proposal] Written in a Phase-1 or Phase-2 proposal document that has not
been through the review rounds (PHASE1b, PHASE2d).
\end{description}
Verification checks are tagged as in the repository: \textbf{[F]} fails on an
earlier statement and passes on the corrected one; \textbf{[V]} verifies a
corrected or new statement; \textbf{[I]} illustrates or strengthens a
statement that was already true.

% Section 2: setting, notation, targets, assumptions
\section{Setting, notation and assumptions}\label{sec:setting}

Source: PHASE1 \S\S1--3 (Appendix~\ref{app:phase1}) and PHASE2a \S0
(Appendix~\ref{app:phase2a}). Notation follows
\texttt{docs/extension/notation.tex}; Table~\ref{tab:notation} reproduces it.

\subsection{Data, latent rates and the accounting identity}

There are units $i=1,\dots,p$ (in Phase~3, $n$), nested in districts
$d=1,\dots,D$ (written $D^\star$ for the number of districts in PHASE2b and
PHASE3). Unit $i$ has size $N_i$, district $D_i$, row shares
$X_i\in\Delta^{R-1}$ and column shares $T_i\in\Delta^{C-1}$. The latent rate
$\beta_{irc}=N_{irc}/N_{ir}$ is the share of the $N_{ir}$ row-$r$ members of
unit $i$ who fall in column $c$, and $\beta_i\in\Bset:=(\Delta^{C-1})^R$. The
accounting identity is
\begin{equation}\label{eq:acc}
T_{ic}=\sum_{r=1}^RX_{ir}\beta_{irc},\qquad\text{i.e.}\quad T_i=\beta_i^\top X_i .
\end{equation}
The unit Duncan--Davis polytope $S(x,t)=\{b\in\Bset:b^\top x=t\}$ is the set of
rate matrices consistent with one unit's margins. Its support function is
$h_S(u;x,t)=\max_{b\in S(x,t)}\langle u,b\rangle$.

\subsection{Three targets}

PHASE1 \S1.7 and PHASE3 \S0.2 distinguish three targets. Each result says which
one it is about.
\begin{itemize}[nosep]
\item $B_{rc}=\sum_iN_iX_{ir}\beta_{irc}/\sum_iN_iX_{ir}$: the \emph{actual}
(realised, finite-population) rate in this electorate. This is JKST's target.
\item $\tilde B_{rc}=\sum_iw_{ir}m_{rc}(V_i)$, with $w_{ir}=N_iX_{ir}/\sum_jN_jX_{jr}$,
$V_i=(X_i,N_i,D_i)$ and $m=\E[\beta\mid V]$: the \emph{model-implied} rate at
the observed design. It is the conditional mean of $B$ given the design.
\item $B^*_{rc}=\E[NX_r\beta_{rc}]/\E[NX_r]$: the \emph{superpopulation} rate
used in the identification results.
\end{itemize}
$B=\tilde B+E$, where $E_{rc}=\sum_iw_{ir}e_{irc}$ and $e_i=\beta_i-m(V_i)$ is
the unit residual. Only $\tilde B$ and $B^*$ are affine in the unknown
parameter. All inference is conditional on the design and covers $B$ or
$\tilde B$; inference for $B^*$ is not pursued. Any linear target is written
$\varphi=\sum_{rc}a_{rc}B_{rc}$ (loyalty, defection, demobilisation, a gap
between groups).

\paragraph{In plain terms.} $B$ is ``what these voters actually did''.
$\tilde B$ is ``what the systematic pattern implies these units did''; it
ignores the local luck of each unit. The confidence interval for $B$ is wider
because it must allow for that luck.

\subsection{Linear contextual effects in homogeneous form}

Because $\mathbf 1^\top X=1$, an affine function of $X$ can be written as a
homogeneous linear one. The pooled LCE model (M0) says that for each column $c$
there is $H_c\in\mathbb R^{R\times R}$ with
\begin{equation}\label{eq:lce}
\E[\beta_{irc}\mid X_i,N_i]=e_r^\top H_cX_i\qquad\text{on }\{X_i\in A\},
\end{equation}
where $A\subseteq\Delta^{R-1}$ is the region on which the model is assumed
(by default the population hull $\conv\supp P_X$). Write
$H_c=Q_c+K_c$ with $Q_c=\tfrac12(H_c+H_c^\top)$ symmetric and
$K_c=\tfrac12(H_c-H_c^\top)$ skew. Then
\[
\tau_c(x):=\E[T_c\mid X=x]=x^\top H_cx=x^\top Q_cx,
\]
because $x^\top K_cx=0$. The regression identifies $Q$; it never sees $K$.
The unknown $K=(K_1,\dots,K_C)$ lives in
$\Kspace=\{K_c^\top=-K_c,\ \sum_cK_c=0\}$, of dimension $(C-1)R(R-1)/2$: one
number in $2\times2$, two in $2\times3$, six in $3\times3$.
$M_H(x)=[e_r^\top H_cx]_{rc}$ is the implied mean transition matrix.

\paragraph{What $H$ and $K$ mean (PHASE1 \S7.2).} $H_{c,rs}$ is the rate at which
row-$r$ voters go to column $c$ in a hypothetical unit made up entirely of
row-$s$ voters. $Q_{c,rs}$ is the average of ``$r$-voters in pure-$s$ units''
and ``$s$-voters in pure-$r$ units''; the data identify it. $K_{c,rs}$ is half
the gap between the two; the data never identify it. It is the R$\times$C version
of JKST's question ``is turnout low because Black residents vote less, or
because white residents in Black areas vote less?''

\subsection{Three identified sets and three hulls}

PHASE1 \S1.7 names three nested sets, $\ThI\subseteq\ThX\subseteq\ThO$
(PHASE2a Prop.~1):
\begin{itemize}[nosep]
\item $\ThI$: the sharp set, using the full law of $(N,X,T)$ and conditioning
on $(X,N)$ (PHASE2a Thm~1). Not a finite LP.
\item $\ThX$: the same, conditioning on $X$ only (JKST Prop.~1).
\item $\ThO$: the outer set, using only the regression $\E[T\mid X]$ and
probability validity (PHASE2a Prop.~7). A finite LP; sharp \emph{for the
regression function} only (PHASE2a Prop.~2).
\end{itemize}
The rule of the documents is that $\ThO$ is never called ``sharp'' without
``for the regression function''. There are also three hulls: the population hull
$A=\conv\supp P_X$ (the default in identification), the hull of the empirical
law (exact for finite-sample identification), and the sample hull used as an
estimate of $A$. Since the sample hull lies inside $A$, estimating it only
widens the outer bound (PHASE3 Prop.~P3.2(c)). Every ``sharp'' statement is
relative to continuous (real-valued) latent rates, not integer cell counts
(PHASE2a Remark~4.2).

\subsection{The assumptions, and what each implies}\label{sec:assumptions}

Each assumption is stated as in PHASE1 \S2, followed by a plain-language
reading: what it rules in or out, when it is credible, and what goes wrong if it
fails. Which rungs use which assumption is shown in Table~\ref{tab:assumption-use}.

\begin{assumption}[A1: closed population / accounting identity]
$T_i=\beta_i^\top X_i$ exactly, with $N_i$ the same people in both margins.
\end{assumption}
\begin{plain}
Both margins describe the same people: the same electorate in both elections,
or the same eligible population for group and outcome. It rules out turnover,
moving, deaths, new registrants and boundary changes. It is credible when both
margins come from one register and the row and column categories include
abstention and, where relevant, ``entered'' or ``left''. If it fails, the
Duncan--Davis sets and the regression identity both shift, and every bound can
exclude the truth. Bounded turnover is an open extension: it cannot be handled
by widening the final interval (PHASE1 A1, as corrected after review). Every
rung, including DD, needs A1.
\end{plain}

\begin{assumption}[A2: linear contextual effects on $A$]
Equation~\eqref{eq:lce}: each group's mean rate is linear in the unit's
composition on $A$.
\end{assumption}
\begin{plain}
A group's behaviour may depend on its surroundings, but only along a straight
line: for example, in-party loyalty rises steadily as the in-party share of the
poll rises. It rules out thresholds and tipping (a machine that mobilises only
above 50\%) and U-shapes (competitive polls mobilising both sides). It is the
\emph{only} identifying restriction of the LCE rungs. If it fails, the
quadratic regression is misspecified and the bound can miss the truth by an
amount that does not shrink with more data. Some violations are invisible:
nonlinear departures that cancel in the aggregate leave the data unchanged
while moving the target (PHASE2a Prop.~9). A cubic term or an empty bound can
reveal some failures; goodness of fit can never validate A2.
\end{plain}

\begin{assumption}[A3: mean independence from size]
$\E[\beta\mid X,N]=\E[\beta\mid X]$ (with $D$ added to both sides under the
district rungs).
\end{assumption}
\begin{plain}
Large and small units with the same mix behave alike. It matters because the
target weights units by size. If large urban polls differ in loyalty from small
rural ones with the same partisan mix, the size-weighted target is biased by
$\Cov(NX_r,\beta_{rc}-M(X))/\E[NX_r]$. District dummies help when size differs
by district (PHASE2b Prop.~D4). It is partly testable through $X\cdot\log N$
interactions.
\end{plain}

\begin{assumption}[A4: support and extrapolation]
$\supp(X)\subseteq A$, compact, and A2 holds on all of $A$.
\end{assumption}
\begin{plain}
The straight lines are trusted over the whole region $A$, whose corners (the
most extreme units) set the constraints. A larger $A$ gives tighter bounds but
assumes more: using the whole simplex assumes the lines extend to pure units
that may not exist (PHASE2a Prop.~5). Outlier units at the hull's corners can
make constraints bind spuriously. Check by trimming the hull. PHASE2a adds
the condition (S3), probability validity on all of $A$; it is automatic for
the population hull (PHASE2a Lemma~3(c)).
\end{plain}

\begin{assumption}[A5: rank / richness]
For symmetric $D$, $x^\top Dx=0$ for all $x\in\supp_AP_X$ implies $D=0$, where
$\supp_AP_X$ is the support of $P_X$ restricted to $A$.
\end{assumption}
\begin{plain}
The compositions vary enough to fit the quadratic regression. For two groups it
means at least three distinct compositions. For three groups it needs six
compositions not on a common conic: five are never enough (PHASE2a
Remark~1.1). It fails, or nearly fails, when shares barely vary (JKST's
``gender'' problem). Then $Q$ is weakly identified and the bound is wide and
unstable. Check the condition number of the design matrix.
\end{plain}

\begin{assumption}[A6: source of randomness and independence]
Given the design $\{(N_i,X_i,D_i)\}$ and the district effects, the unit rates
$\beta_i$ are random with mean $m(V_i)$, and the residuals $e_i$ are
independent across units (P3-S1a).
\end{assumption}
\begin{plain}
Every unit is observed, so the uncertainty is not from sampling units. It is
model-based: given its characteristics, each unit's behaviour has a random part.
District shocks are fixed features of this electorate (part of the mean under
M1/M2; a misspecification under M0). Independence is a substantive assumption.
Neighbourhood spillovers within a district violate it, and then the variances
and residual allowances are too small. Observing every unit does \emph{not}
remove dependence (PHASE3 \S0.4, as corrected in round~2). The pooling test
does not test it.
\end{plain}

\begin{assumption}[A7: regularity for inference]
A unique active constraint at the optimum, $\mu_{rs}>0$, and $\ThO$ not a
singleton.
\end{assumption}
\begin{plain}
This is JKST's tie-breaking condition. With it the bound's endpoint is a smooth
function of the regression coefficients and the delta method applies. In
R$\times$C, and especially with many districts, near-ties are the norm, so the
inference of Phase~3 is built not to need it.
\end{plain}

\begin{assumption}[A8: homogeneous contextual effects across districts (M1)]
$K$ (and the symmetric slope parts) are common to all districts; only
intercepts vary.
\end{assumption}
\begin{plain}
The way composition shifts behaviour is the same in every district, while each
district may be more or less loyal across the board (parallel lines). It is
what produces the intersection gain of M1. It fails if, say, a strong
incumbent changes the contextual gradient in one district, or if ceiling and
floor effects make an additive shift impossible near 0 or 1, which is most
likely in the lopsided districts that tighten the bound most. The bound can then
exclude the truth and still look narrow. The symmetric part is testable
(PHASE2b Lemma~D1(b)); the shared $K$ is not.
\end{plain}

\begin{assumption}[A9: within-district richness]
Under M2, A5 within each district. Under M1, each district's support contains
$R$ affinely independent compositions and the stacked design has full column
rank (PHASE2b Lemma~D1(d)).
\end{assumption}
\begin{plain}
Districts with few, similar units identify almost nothing. Under M1 the bar is
lower than under M2: a district's baseline needs only $R$ affinely
independent compositions. Weak identification per district, combined with a
max/min over districts, biases plug-in bounds inward.
\end{plain}

\begin{assumption}[A10: asymptotic regime (a choice)]
(R1) $D$ fixed, $p_d\to\infty$; (R2) $D\to\infty$, $p_d$ bounded; (R3) both
grow.
\end{assumption}
\begin{plain}
Which limit theory applies depends on the design. With ten polls in each of
ninety ridings (R2), district intercepts are incidental parameters and the
plug-in intersection is badly too narrow. PHASE3 \S0.5 says which result holds
in which regime (Table~\ref{tab:regimes} in Section~\ref{sec:inference}).
\end{plain}

\paragraph{Assumptions explicitly not made.} No distributional form for $\beta$
(unlike King's truncated normal), no independence between $\beta$ and $X$
(contextual effects are the point), and no point-identifying nonlinearity.

\begin{longtable}{@{}p{0.30\linewidth}p{0.62\linewidth}@{}}
\caption{Which rungs use which assumption (PLAIN\_LANGUAGE\_GUIDE \S2; PHASE2c \S0).}\label{tab:assumption-use}\\
\toprule Rung & Assumptions used \\ \midrule\endfirsthead
\toprule Rung & Assumptions used \\ \midrule\endhead
DD & A1 only. \\
N1 (sign rows on realised tables) & A1; continuous unit tables; no regression and no sampling inference. \\
M0 & A1--A5; inference adds A6, and A7 only for the normal limit of P3.4(i). \\
M1, M1($\delta$), M2, L4, L5, L6 & A1--A5 within districts (A5 replaced by A9), A8 for M1 and its relatives; inference adds A6 and the regime A10. \\
N2 & A1, A3 and the finite-support convention of PHASE2c \S4. \\
M-SIGN rows & Those of the rung they are added to. \\
\bottomrule
\end{longtable}

\begin{table}[ht]
\centering\small
\begin{tabular}{@{}ll@{}}
\toprule
Symbol & Meaning \\
\midrule
$i=1,\dots,p$ & unit (precinct) index \\
$d=1,\dots,D$;\ $D_i$;\ $\mathcal I_d$;\ $p_d$ & district; district of unit $i$; units in $d$; their number \\
$r,s=1,\dots,R$ & row categories (election-1 choice / group) \\
$c=1,\dots,C$ & column categories (election-2 choice) \\
$N_i,\ N_{ir},\ N_{irc}$ & unit size; row count; latent cell count \\
$X_i\in\Delta^{R-1}$;\ $T_i\in\Delta^{C-1}$ & observed row shares; observed column shares \\
$\beta_i=[\beta_{irc}]\in\mathfrak B=(\Delta^{C-1})^R$ & latent transition matrix of unit $i$ \\
$W_i=(N_i,X_i,T_i,Z_i,D_i)$;\ $V=(X,N,D)$ & observables of unit $i$; conditioning variables \\
$P$;\ $\mathbb Q$ & law of $W$ (observable); law of $(W,\beta)$ \\
$\mathcal G=\sigma(V)$ & conditioning $\sigma$-algebra of the contextual model \\
$B_{rc}$,\ $\tilde B_{rc}$,\ $B^*_{rc}$ & actual, model-implied, superpopulation rate \\
$\varphi=\sum_{rc}a_{rc}B_{rc}$ & linear target \\
$A\subseteq\Delta^{R-1}$,\ $\operatorname{vert}(A)$ & region where LCE is assumed; its vertices \\
$H_c$;\ $Q_c=\tfrac12(H_c+H_c^\top)$;\ $K_c=\tfrac12(H_c-H_c^\top)$ & LCE coefficients; identified part; unidentified part \\
$\mathcal K$;\ $k$ & space of $K$, $\dim=(C-1)R(R-1)/2$; its free coordinates \\
$\kappa_c=K_{c,12}$ & unidentified scalar per column when $R=2$ \\
$M_H(x)=[e_r^\top H_cx]_{rc}$;\ $\tau_c(x)=E[T_c\mid X=x]$ & implied mean transition matrix; regression function \\
$\mu_r=E[NX_r]$,\ $\mu_{rs}=E[NX_rX_s]$ & weighting moments \\
$S(x,t)$;\ $h_S(u;x,t)$ & unit Duncan--Davis polytope; its support function \\
$\Theta_I\subseteq\Theta_X\subseteq\Theta_O$ & sharp set; JKST Prop.~1 set; outer set \\
$\zeta_{c,d}\in\mathbb R^R$ & district intercept shifts (M1), identified \\
$Q^{(d)}_c,\ K^{(d)}_c,\ A_d$ & district-level identified part, unidentified part, support hull \\
$\delta_B,\ \delta_W,\ \Delta$ & between, within (W1) and pooled-within (W2) associations (PHASE2c) \\
$\mathrm{CI}_x,\ S_1$ & JKST confidence interval; residual (Popoviciu) standard-error term \\
$\Phi$ & standard normal CDF \\
\bottomrule
\end{tabular}
\caption{Notation (after \texttt{docs/extension/notation.tex}, extended with symbols used later).}\label{tab:notation}
\end{table}

% Section 3: the 2x2 original
\section{The $2\times2$ original: Duncan--Davis and JKST}\label{sec:2x2}

Source: PHASE1 \S0 (Appendix~\ref{app:phase1}), PHASE2a \S6.1, and the
original paper and package~\citep{jkst2019,duncan1953}.

\subsection{Duncan--Davis (the method of bounds)}

With two rows (Black, white; or group 1 and the rest) and a binary outcome, the
accounting identity in unit $i$ is $T_i=X_i\beta^b_i+(1-X_i)\beta^w_i$. Each
unit's rate for one group then lies in an interval:
\[
\beta^w_i\in\Big[\max\Big\{0,\tfrac{T_i-X_i}{1-X_i}\Big\},\ \min\Big\{1,\tfrac{T_i}{1-X_i}\Big\}\Big],
\]
and the same for $\beta^b_i$. Aggregating with the target weights gives the
Duncan--Davis (DD) bound on $B$ \citep{duncan1953}. In R$\times$C the per-cell
version is the Fréchet--Hoeffding bound
$n_{rc}\in[\max\{0,x_r+t_c-1\},\min\{x_r,t_c\}]$ (PHASE2a Prop.~8;
PHASE2b Prop.~D0).

\begin{plain}
DD uses nothing but arithmetic. It is always valid when A1 holds, never empty,
and often wide: in the BC $3\times3$ illustration of PHASE2a \S8 the DD bound
on loyalty is $[0.131,0.973]$.
\end{plain}

\subsection{JKST: linear contextual effects in $2\times2$}

JKST's Assumption~1 is that $(\beta^b_i,\beta^w_i,X_i,N_i)$ are iid with
$\E(\beta^w_i\mid X_i,N_i)=w_0+w_1X_i$ and
$\E(\beta^b_i\mid X_i,N_i)=b_0+b_1X_i$. The accounting identity then gives
\[
\E(T_i\mid X_i)=w_0+c_1X_i+d_1X_i^2 ,
\]
which identifies $(w_0,c_1,d_1)$ and leaves one free parameter, $w_1$.
\begin{itemize}[nosep]
\item \textbf{JKST Prop.~1} bounds $w_1$ through
$\E(L_i\mid X_i)\le w_0+w_1X_i\le \E(U_i\mid X_i)$ on a set $A$, where $L_i,U_i$
are the unit DD ends. It uses the conditional law of $T$ given $X$.
\item \textbf{JKST Prop.~2} relaxes Prop.~1 by Jensen's inequality and uses the
quadratic regression only. \emph{The package implements only Prop.~2}, with
$A=[\min X_i+10^{-5},\max X_i-10^{-5}]$, so the package bound is an outer set,
not the sharp one (PHASE1 \S0, fact~2).
\item \textbf{JKST Prop.~3} is the version with $A=[0,1]$ (extrapolation to pure
units).
\item \textbf{JKST Prop.~4 / App.~A (inference)}:
$\mathrm{CI}_x=[\hat B_L-xS_L,\ \hat B_U+xS_U]\cap DD$, with an HC1 sandwich
variance and the Popoviciu bound $\Var(\beta)\le1/4$ for the finite-population
term $S_1$. It needs tie-breaking (one active constraint) and covers the
parameter $B$ with probability at least $\Phi(x)$, in the spirit of
Imbens--Manski~\citep{imbensmanski2004}.
\item \textbf{App.~B}: the package reverts to DD when $\mathrm{CI}_0\cap DD$ is
empty; under the model $\mathrm{CI}_0\ne\emptyset$ in the large-$p$ limit.
\item \textbf{App.~C (covariates)}: covariates enter the intercept additively,
with a common slope; only $\tilde w_1$ is unidentified. Not in the package.
\end{itemize}

\begin{plain}
JKST's insight is that LCE leaves exactly one number unknown in $2\times2$: how
the curvature of the aggregate regression splits between the two groups'
contextual slopes. The requirement that both groups' implied rates be
probabilities at the two extreme compositions bounds that number, and hence $B$.
\end{plain}

\subsection{How the R$\times$C work contains JKST}

PHASE1 \S0 (fact~6) rewrote Prop.~2 in the homogeneous form of
Section~\ref{sec:setting} and checked numerically that it reproduces the
package's \texttt{generateBounds} to $10^{-16}$
(\path{docs/extension/check_2x2_equivalence.py}). PHASE2a Prop.~4
(Result~\ref{res:p2a-p4} below) proves the equivalence: the outer set $\ThO$ in
$2\times2$ is JKST Prop.~2 (with $A=[l,u]$) or Prop.~3 (with $A=\Delta^1$), the
dictionary is $K=c_1/2-w_1$, and $B^*(K)$ is JKST eq.~(16). PHASE2a Thm~1
(Result~\ref{res:p2a-t1}) confirms JKST's ``tightest bound'' claim for Prop.~1
when conditioning on $X$ only, and shows that conditioning on $(X,N)$ can be
strictly tighter.

\begin{revnote}
The Phases~1--3 referee asked that ``sharp'' never be applied to $\ThO$ without
``for the regression function'', and that the three sets, three hulls and three
targets be named in every statement. PHASE1 \S1.7 was added for this
(Appendix~\ref{app:revisions}, rows R-10, R-11 and R-33).
\end{revnote}

% Section 4: R x C under LCE (PHASE2a)
\section{R$\times$C under linear contextual effects}\label{sec:rxc}

Source: PHASE2a (Appendix~\ref{app:phase2a}), which builds on PHASE1
(Appendix~\ref{app:phase1}). Every proof in PHASE2a is written in numbered
steps with a justification tag; the tag legend is PHASE2a \S0.3 and the citation
check is \S0.4. Numerical and symbolic checks: \texttt{verify\_phase2a.py},
32/32 PASS. The model $\mathcal M(H;A)$ is the set of laws of $(W,\beta)$ with
(S1) $\beta\in\Bset$ and $T=\beta^\top X$; (S2) $\E[\beta\mid V]=M_H(X)$ on
$\{X\in A\}$; (S3) $M_H(x)\in\Bset$ for every $x\in A$.

\subsection{The unit polytope and what the data identify}

\begin{result}[Lemma 0 --- PHASE2a \S1]\label{res:p2a-l0}
\status{Proved.} For all $x\in\Delta^{R-1}$, $t\in\Delta^{C-1}$: (a)
$S(x,t)\ne\emptyset$; (b) it is compact and convex; (c) $(x,t)\mapsto S(x,t)$ has
closed graph, is upper hemicontinuous and Borel; (d) $h_S$ is sublinear and
$\sqrt R$-Lipschitz in $u$, concave in $t$, Borel; (e) $S$ is the intersection
of its supporting half-spaces.
\end{result}
\begin{plain}
Each unit's set of possible rate tables is never empty (``everyone in the unit
behaves like the unit average'' always fits), and it is a well-behaved convex
polytope. These are the regularity facts every later proof uses.
\end{plain}
\proofref{Appendix~\ref{app:phase2a}, \S1, steps 1--13}

\begin{result}[Lemma 1, identification decomposition --- PHASE2a \S2]\label{res:p2a-l1}
\status{Proved.} Under the model and A5: (a) $\E[T_c\mid X=x]=x^\top Q_cx$ and
$Q_c$ is unique; (b) $\sum_cH_c=J$, hence $\sum_cQ_c=J$ and $\sum_cK_c=0$;
(c) every $K'\in\Kspace$ gives the same regression, and
$\dim\Kspace=(C-1)R(R-1)/2$.
\end{result}
\begin{plain}
The data reveal the symmetric part $Q$ of the contextual-effect matrices, and
nothing about the skew part $K$. For each pair of groups and each outcome, the
data identify the sum ``effect of $s$-share on group $r$ plus effect of $r$-share
on group $s$'' but not how it splits. A $2\times2$ table has one unknown, a
$3\times3$ table six.
\end{plain}
\begin{sketch}
Take conditional expectations of the accounting identity: $\E[T_c\mid X]=X^\top
H_cX$. A skew matrix vanishes in a quadratic form, so only $Q_c$ appears. A5
makes the quadratic form determine $Q_c$. Row sums of $\beta$ equal one, which
under A5 forces $\sum_cH_c=J$.
\end{sketch}
\proofref{Appendix~\ref{app:phase2a}, \S2, (a) steps 1--8, (b) steps 1--7, (c) steps 1--3}
\begin{revnote}
(R-01, R-02, P-B39.) The derivation of $\sum_cH_c=J$ now uses A5 rather than an
unassumed open set. A5 is a matrix rank condition stated on
$\supp_AP_X$, the support of $P_X$ restricted to $A$; the earlier
$\supp P_X\cap A$ is wrong when $A$ carries less than full mass
($X\sim U[0,1]$, $A=\{1/2\}$).
\end{revnote}

\paragraph{Remarks 1.0--1.2 (PHASE2a \S2).} \emph{Remark~1.0}: A2 cannot be
validated by goodness of fit; Prop.~9 exhibits nonlinear departures that leave
the data unchanged (moved forward after review, R-16). \emph{Remark~1.1}
(proved): for $R=2$, A5 holds iff there are at least three distinct
compositions; for $R=3$ six points off a common conic are needed, and five are
never enough. \emph{Remark~1.2} (proved): column-by-column OLS on the quadratic
monomials gives $\sum_c\hat Q_c=J$ exactly, so no constrained regression is
needed.

\begin{result}[Lemma 2, the estimand --- PHASE2a \S3]\label{res:p2a-l2}
\status{Proved.} Under (S1)--(S2) with $P(X\in A)=1$,
\[
B^*_{rc}=\frac1{\mu_r}\Big[\sum_sQ_{c,rs}\mu_{rs}+\sum_{s\ne r}K_{c,rs}\mu_{rs}\Big],
\]
affine in $K$, with coefficient $\mu_{rs}/\mu_r\ge0$ on the ambient entry $K_{c,rs}$.
\end{result}
\begin{plain}
The target is a known number plus a weighted sum of the unknown $K$ entries. The
weights $\mu_{rs}/\mu_r$ measure how much group $r$ lives alongside group $s$:
the more mixed the units, the more the unknown split matters.
\end{plain}
\proofref{Appendix~\ref{app:phase2a}, \S3, steps 1--7; $2\times2$ check against JKST (16)}
\begin{revnote}
(R-03, R-05.) The coefficient is stated in ambient coordinates, with the free
coordinates in Prop.~7. $P(X\in A)=1$ is now explicit; Remark~2.1 and
Prop.~7$'$ treat units outside $A$.
\end{revnote}

\subsection{The sharp identified set}

\begin{result}[Theorem 1, characterisation of $\ThI$ --- PHASE2a \S4]\label{res:p2a-t1}
\status{Proved} (measure-theoretic steps rest on the cited theorems of PHASE2a
\S0.4). Assume A1, A5, standing integrability, $H=Q(P)+K$ with $K\in\Kspace$.
Then $K\in\ThI(P;A)$ iff (S3) holds and, outside one $P_V$-null set, for every
$v=(x,n)$ with $x\in A$,
\[
\langle u,M_H(x)\rangle\le\psi(u;v):=\int h_S(u;x,t)\,P_{T\mid V=v}(dt)\quad\text{for all }u\in\mathbb R^{R\times C}. \tag{T1}
\]
$\ThI$ is convex and closed; it is compact when $\operatorname{affdim}(A)=R-1$,
which A5 implies.
\end{result}
\begin{plain}
A value of the unknown $K$ is consistent with the data exactly when, at every
composition and size, the implied average rate table could be the average of
rate tables that each fit some unit's margins. It is the strongest possible
use of the data under the model. It needs the whole conditional distribution of
outcomes given composition and size, which in real data must be smoothed, so it
is not what the package computes.
\end{plain}
\begin{sketch}
Necessity: each unit's $\beta$ lies in its polytope, so its conditional mean lies
in the conditional selection expectation; one common null set for all $u$ is
obtained from a countable dense set and Lipschitz continuity. Sufficiency: the
set of measurable selections is weakly compact in $L^1$ (Dunford--Pettis,
Mazur); its image under $\E[\cdot\mid V]$ is weakly compact and convex; if the
target mean were outside it, Hahn--Banach would separate it, contradicting
(T1) evaluated at the separating direction (a measurable argmax attains the
support function). No auxiliary randomisation is needed (Remark~4.1). The random-set formulation follows Beresteanu, Molchanov and Molinari~\citep{bmm2011}.
\end{sketch}
\proofref{Appendix~\ref{app:phase2a}, \S4: (i)$\Rightarrow$(ii) steps 1--11; facts F1--F7; (ii)$\Rightarrow$(i) steps 1--22; convexity and compactness steps 1--4}
\begin{revnote}
(R-07, R-08, R-09.) Compactness needs $\operatorname{affdim}(A)=R-1$, not
``nonempty relative interior''. The sufficiency proof was expanded into
facts F1--F7 with checked citations, and a paragraph explains why (T1) must hold
for all $u$ outside one null set.
\end{revnote}

\paragraph{Remark 4.2 (continuous versus integer rates; PHASE2a \S4).}
Sharpness is relative to real-valued latent rates. Part~(1), that every vertex of
the transportation polytope is an integer table, is a standard fact checked on
200 cases. Part~(2), that randomising over integer vertices would leave the
sharp set unchanged, is \emph{only a sketch} (not verified) and is not used
anywhere (R-12).

\begin{result}[Corollary 1, sharp bounds for a linear target --- PHASE2a \S4]
\status{Proved.} If $P(X\in A)=1$ and $\ThI\ne\emptyset$, the sharp set for
$\varphi$ is $[\min_{\ThI}\varphi(K),\max_{\ThI}\varphi(K)]$.
\end{result}
\begin{plain}
Because the unknown enters the target linearly and the consistent set is convex
and compact, the answer is always an interval.
\end{plain}
\proofref{Appendix~\ref{app:phase2a}, \S4, Corollary~1 steps 1--3}

\subsection{The outer set: reduction, nesting, sharpness for the regression}

\begin{result}[Lemma 3 --- PHASE2a \S5]\label{res:p2a-l3}
\status{Proved.} (a) For a polytope $A$, linear constraints on $A$ reduce to its
vertices. (b) If $\operatorname{affdim}(A)=R-1$, $\ThO$ is bounded (a polytope);
otherwise it is a closed convex polyhedron. (c) Under (S1)--(S2), if
$A\subseteq\conv(\supp_AP_X)$, (S3) holds automatically, in particular for the
population hull.
\end{result}
\begin{plain}
Only the most extreme units matter for the constraints. If the region $A$ is
the range of the data, probability validity of the straight lines is automatic;
if $A$ extrapolates beyond the data, validity there is an extra assumption.
\end{plain}
\proofref{Appendix~\ref{app:phase2a}, \S5, steps 1--9; Remarks~5.1--5.2 give the counterexamples}
\begin{revnote}
(R-04, R-07, P-B39.) (c) was restated: (S3) is automatic only when
$A\subseteq\conv(\supp_AP_X)$. Remark~5.1 (a singleton inside the hull with no
mass) and Remark~5.2 (a singleton $A$ makes $\ThO$ unbounded; A5 rules it out)
were added. Checks in \texttt{verify\_phase2a.py} fail on the earlier statements.
\end{revnote}

\begin{result}[Proposition 1, nesting --- PHASE2a \S5]\label{res:p2a-p1}
\status{Proved.} $\ThI(P;A)\subseteq\ThX(P;A)\subseteq\ThO(P;A)$.
\end{result}
\begin{plain}
Using more of the data can only narrow the set. The outer set uses only the
regression, so it is the widest of the three.
\end{plain}
\begin{sketch}
Conditioning on $(X,N)$ implies conditioning on $X$ (tower property). Conditional
expectations of valid rate tables are valid rate tables, which is JKST's Jensen
step.
\end{sketch}
\proofref{Appendix~\ref{app:phase2a}, \S5, Proposition~1 steps 1--5}

\begin{result}[Proposition 2, $\ThO$ is sharp for the regression function --- PHASE2a \S5]\label{res:p2a-p2}
\status{Proved.} Let $\mathcal P(P)$ be the laws with the same law of $(N,X)$
and the same regression $\E[T\mid X]$. If $P(X\in A)=1$,
$\ThO(P;A)=\bigcup_{P'\in\mathcal P(P)}\ThI(P';A)$.
\end{result}
\begin{plain}
The outer bound is exactly what can be learned from the regression alone. It is
the sharp set for any data set in which units with the same composition have
identical outcomes. Calling it ``sharp'' without that qualifier is wrong.
\end{plain}
\begin{sketch}
Given $K\in\ThO$, set every unit's rates equal to the implied mean $M_H(X)$;
this is a valid law with the same regression, and $K$ is in its sharp set.
\end{sketch}
\proofref{Appendix~\ref{app:phase2a}, \S5, Proposition~2 steps 1--10}

\begin{result}[Corollary 2 and Proposition 3 --- PHASE2a \S5]
\status{Proved.} If $T=\E[T\mid X]$ a.s., then $\ThI=\ThO$ (Cor.~2). More
generally (Prop.~3, a \emph{sufficient} condition), if $\E[T\mid V]=\E[T\mid X]$
and, for a.e.\ $v$ and every $u$, $t\mapsto h_S(u;x,t)$ is affine on the
convex hull of the conditional support of $T$, then $\ThI=\ThO$.
\end{result}
\begin{plain}
The regression-only bound loses nothing unless units with similar composition
are so different that some sit at a Duncan--Davis corner and others do not
(Remark~3.1, the $2\times2$ reading). In 400 random valid $2\times2$ designs the
inclusion was strict in 46.
\end{plain}
\proofref{Appendix~\ref{app:phase2a}, \S5, Corollary~2 steps 1--3; Proposition~3 steps 1--7}

\subsection{Closed forms and the linear program}

\begin{result}[Proposition 4, equivalence with JKST --- PHASE2a \S6.1]\label{res:p2a-p4}
\status{Proved.} In $2\times2$ with $A=[l,u]$: (a) $Q_{1,22}=w_0$,
$Q_{1,12}=w_0+c_1/2$, $Q_{1,11}=w_0+c_1+d_1$, $K=c_1/2-w_1$; (b) $\ThO$ is JKST
Prop.~2; (c) with $A=\Delta^1$ it is JKST Prop.~3; (d) with $P(l\le X\le u)=1$,
$B^*(K)$ is JKST eq.~(16).
\end{result}
\begin{plain}
The R$\times$C machinery contains the published $2\times2$ method exactly; the
package's existing bounds are a special case.
\end{plain}
\proofref{Appendix~\ref{app:phase2a}, \S6.1, derivations (a)--(d); numerically to $10^{-16}$}

\begin{result}[Proposition 5, $A$ the whole simplex --- PHASE2a \S6.2]
\status{Proved.} With $A=\Delta^{R-1}$, $\ThO=\{|K_{c,rs}|\le Q_{c,rs}\}$; it is
non-empty iff every $Q_{c,rs}\ge0$; each $K_{c,rs}$ ranges over
$[-m_{c,rs},m_{c,rs}]$ with $m_{c,rs}=\min\{Q_{c,rs},1-Q_{c,rs}\}$, and $B^*$
has a closed form.
\end{result}
\begin{plain}
If the straight lines are trusted all the way to pure units, two hypothetical
probabilities with known average $Q$ can split it by at most
$\pm\min(Q,1-Q)$. This is narrower than the hull-based bound but assumes
extrapolation.
\end{plain}
\proofref{Appendix~\ref{app:phase2a}, \S6.2, steps 1--8; checked against an LP in 200 tables}

\begin{result}[Lemma 4, box $\cap$ hyperplane --- PHASE2a \S6.3]
\status{Proved.} For $\mathcal P=\{\ell\le\kappa\le h,\sum_c\kappa_c=0\}$:
(a) non-empty iff $\sum\ell\le0\le\sum h$; (b) coordinate projections in closed
form; (c) linear objectives are maximised by a greedy rule.
\end{result}
\proofref{Appendix~\ref{app:phase2a}, \S6.3, steps 1--6}

\begin{result}[Proposition 6, two rows, any number of columns --- PHASE2a \S6.4]
\status{Proved.} For $R=2$ and $A=[l,u]$ with $0<l<u<1$, $\ThO=\{\ell\le\kappa\le
h,\ \sum_c\kappa_c=0\}$ with closed-form $\ell_c,h_c$; target bounds follow by
Lemma~4.
\end{result}
\proofref{Appendix~\ref{app:phase2a}, \S6.4, steps 1--6; checked against an LP in 300 cases}

\begin{result}[Corollary 6.1, joint versus collapsed --- PHASE2a \S6.4]
\status{Proved.} For $R=2$: (a) the joint interval for $\kappa_c$ lies inside the
collapsed one; (b) they are equal when, for each row, the $Q_{c',rr}$ ($c'\ne c$)
share a sign; (c) strict tightening therefore requires some negative
extrapolated pure-poll rate. For $R\ge3$ the joint bound can be strictly tighter
with all rates valid; a characterisation is open.
\end{result}
\begin{plain}
With two groups, fitting all outcome columns jointly helps only if the straight
lines would give an impossible rate at some pure unit. With three or more groups,
joint fitting can help even when everything is valid.
\end{plain}
\proofref{Appendix~\ref{app:phase2a}, \S6.4, Corollary~6.1 steps 1--7}
\begin{revnote}
(R-14; PHASE2a \S10 item~1.) PHASE1's general claim of strict joint gain was
withdrawn; the exact sign condition replaced ``strong contextual effects''.
\end{revnote}

\begin{result}[Proposition 7, the general LP --- PHASE2a \S6.5]\label{res:p2a-p7}
\status{Proved.} With free coordinates $k_{c,rs}$ ($c<C$, $r<s$),
$\ThO=\{k:Gk\le g\}$, one row per (vertex, row, column), and
$\varphi(k)=\varphi_0+\gamma^\top k$ with explicit $\varphi_0$ and
$\gamma_{c,rs}=\mu_{rs}[(a_{rc}-a_{rC})/\mu_r-(a_{sc}-a_{sC})/\mu_s]$. Imposing
the rows at every observed $X_i$ equals imposing them at the hull vertices.
\end{result}
\begin{plain}
The computable bound is a small linear program: one variable per unknown
direction, one constraint per extreme unit, group and outcome. This is the
\emph{outer} set, not an LP form of the sharp set.
\end{plain}
\proofref{Appendix~\ref{app:phase2a}, \S6.5, steps 1--5}
\begin{revnote}
(R-06.) The closing paragraph on regularity was wrong: the value function is
concave and piecewise linear in $Q$ for fixed $\mu$, convex in $\gamma$ for
fixed $Q$, and in general neither jointly ($V=\mu Q$). Lemma~P3.1 is the
fixed-design case. Check: \texttt{verify\_phase2a.py} (midpoint test).
\end{revnote}

\begin{result}[Proposition 7$'$, units outside $A$ --- PHASE2a \S6.6]
\status{Proved.} If $P(X\in A)>0$ only, the bound is the LP over units in $A$
plus the support-function (Duncan--Davis) terms of the units outside $A$; each
endpoint is attained. This is what \texttt{range = "trim"} computes for single
cells.
\end{result}
\begin{plain}
Units outside the region where straight lines are trusted contribute their
model-free range; nothing links them to the model.
\end{plain}
\proofref{Appendix~\ref{app:phase2a}, \S6.6, steps 1--6}
\begin{revnote}
(R-05; decision D2.) New result written after review to document the existing
\texttt{trim} option.
\end{revnote}

\subsection{Refutability and undetectable violations}

\begin{result}[Proposition 8, non-emptiness --- PHASE2a \S7]
\status{Proved.} If $P$ is generated by the model with $P(X\in A)=1$, then the
true $K_0\in\ThI\subseteq\ThO$ and $B^*_{rc}\in[\varphi_L,\varphi_U]\cap DD^*_{rc}$.
An empty $\ThO$, or an empty intersection with $DD^*$, refutes A1--A5 jointly.
\end{result}
\begin{plain}
If the model is right, the population bound is never empty. An empty bound is
evidence against the model (in samples it can also be noise; Phase~3 turns it
into a test).
\end{plain}
\proofref{Appendix~\ref{app:phase2a}, \S7, steps 1--5}

\begin{result}[Proposition 9, what the data can never rule out --- PHASE2a \S7]\label{res:p2a-p9}
\status{Proved.} If $\Delta(x)$ satisfies $x^\top\Delta(x)=0$, $\Delta(x)\mathbf
1=0$ and keeps $M_H+\Delta$ valid on the support, then $\beta=M_H(X)+\Delta(X)$
and $\beta'=M_H(X)$ give the same $P$ but different $B^*$. Under A5 the linear
such $\Delta$ are exactly $M_{K'}$, $K'\in\Kspace$.
\end{result}
\begin{plain}
There are nonlinear contextual effects that cancel in the aggregate and leave
the data identical, yet change the transition rates. The linearity assumption is
therefore a genuine leap that no fit statistic can confirm. Phase~3 bounds its
consequences with the sensitivity parameter $\eta$ (Prop.~S1).
\end{plain}
\proofref{Appendix~\ref{app:phase2a}, \S7, Proposition~9 steps 1--4}

\paragraph{Illustration (PHASE2a \S8).} In a noise-free BC $3\times3$ design with
2000 polls, the $\ThO$ bounds on loyalty, defection, demobilisation and
mobilisation have width ratios 0.52--0.64 relative to DD and all contain the
truth. They remain wide because six numbers are unidentified.

\paragraph{Relation to Cross--Manski (PHASE2a \S12).} Cross and
Manski~\citep{crossmanski2002} treat Goodman-type constancy (an instrument) and
obtain a point; LCE is weaker, and the R$\times$C bounds sit between their
data-alone region and their point.

% Section 5: districts and the rung ladder (PHASE2b)
\section{Districts and the rung ladder}\label{sec:districts}

Source: PHASE2b (Appendix~\ref{app:phase2b}). Checks:
\texttt{verify\_phase2b.py}, 22/22 PASS. ``District'' means any partition of the
units into groups (ridings, counties, schools, or election periods). The district
model replaces (S2) by $\E[\beta\mid V]=M_{H^{(D)}}(X)$ and (S3) by validity on
each district hull $A_d$. The rungs are restrictions on
$\mathbf H=(H^{(1)},\dots,H^{(D^\star)})$:
\begin{itemize}[nosep]
\item \textbf{M2}: none (each district its own linear effects).
\item \textbf{M1($\delta$)}: $\exists\bar k$ with $\|k^{(d)}-\bar k-c_d\|_\infty\le\delta$ for all $d$.
\item \textbf{M1}: $H^{(d)}_c=H_c+\zeta_{c,d}\mathbf 1^\top$, $\sum_c\zeta_{c,d}=0$, $\zeta_{c,1}=0$.
\item \textbf{L4($r$)}: M1 and row $r$ has no contextual effect.
\item \textbf{L5}: $H^{(d)}_c=\alpha_{c,d}\mathbf 1^\top$ (constancy within districts).
\item \textbf{L6}: L5 with common $\alpha$ (Goodman).
\item \textbf{M0}: one $H$ for all units on $\conv\bigcup_dA_d$ (not nested).
\end{itemize}
$\delta$ is in probability units: $\delta=0.05$ lets the gap between ``$r$-voters
in pure-$s$ polls'' and ``$s$-voters in pure-$r$ polls'' differ across districts
by up to 10 points.

\begin{result}[Proposition D0, grouping adds nothing without a model --- PHASE2b \S1]\label{res:d0}
\status{Proved.} Under (S1) alone the sharp set for $B^*_{rc}$ is the
Duncan--Davis interval $DD^*_{rc}$, whatever the law of $D$ or of any
covariate. The joint sharp set for the whole matrix $B^*$ is convex and compact
but not the product of the cell intervals.
\end{result}
\begin{plain}
District dummies cannot narrow the model-free bound: it is already computed unit
by unit. Every gain from districts comes from an assumption that links districts.
\end{plain}
\proofref{Appendix~\ref{app:phase2b}, \S1, steps 1--8 (Gale's theorem)}
\begin{revnote}(R-18.) The per-cell sharp interval is now separated from the joint set.\end{revnote}

\begin{result}[Lemma D1, what each rung identifies --- PHASE2b \S2]\label{res:d1}
\status{Proved.} Under A5 within districts: (a) every rung identifies
$Q^{(d)}$; each $K^{(d)}\in\Kspace$ is unidentified. (b) Under M1,
$\zeta_{c,d}=\operatorname{diag}(Q^{(d)}_c-Q^{(1)}_c)$ is identified; the off-diagonal entries
satisfy the overidentifying restrictions
$(Q^{(d)}_c-Q^{(1)}_c)_{rs}=\tfrac12(\zeta_{c,d,r}+\zeta_{c,d,s})$; the unknown
dimension is the same as pooled. (c) Under L5, $K^{(d)}$ is identified and the
structure of $Q^{(d)}$ is testable. (d) Under M1 a weaker stacked rank condition
(A5-M1) suffices; e.g.\ A5 in one district and $R$ affinely independent points in
every other.
\end{result}
\begin{plain}
District baselines are identified from each district's own regression. Under M1
the curvature of the regression must be the same in every district: that is the
testable part. The shared contextual split $K$ is not testable. A district
needs only a few compositions to fix its baseline, but they must not lie on a
line (for three groups).
\end{plain}
\proofref{Appendix~\ref{app:phase2b}, \S2, proofs of (a), (b), (c), (d)}
\begin{revnote}
(R-19, R-30, P-B39.) Part (d) and Remark~D1.1 were added: ``$R$ distinct
compositions'' is not enough, they must be affinely independent. The number of
overidentifying restrictions is $(D^\star-1)(C-1)R(R-1)/2$. A5-d is stated on
the support of the restricted measure.
\end{revnote}

\begin{result}[Lemma D2, the estimand under districts --- PHASE2b \S2]
\status{Proved.} $B^*_{rc}=\mu_r^{-1}\sum_d[\sum_sQ^{(d)}_{c,rs}\mu^{(d)}_{rs}+\sum_{s\ne
r}K^{(d)}_{c,rs}\mu^{(d)}_{rs}]$; under M1 its sensitivity to the common $K_{c,rs}$
is $\mu_{rs}/\mu_r$, as pooled.
\end{result}
\proofref{Appendix~\ref{app:phase2b}, \S2, Lemma~D2 steps 1--4}

\begin{result}[Theorem 1$'$, sharpness for any affine mean model --- PHASE2b \S3]\label{res:t1prime}
\status{Proved} (by reduction to PHASE2a Thm~1; measure-theoretic content not
machine-checked). If $m_\theta(v)$ is Borel (H1), bounded (H2), affine and
uniformly continuous in $\theta$ (H3), and the selection set is unaffected
(H4), then $\theta\in\ThI$ iff $m_\theta$ is valid on $\mathcal V_A$ and
$\langle u,m_\theta(v)\rangle\le\psi(u;v)$ for a.e.\ $v$ and all $u$.
\end{result}
\begin{plain}
The sharpness argument of the pooled model carries over to any model in which
the mean is affine in the unknowns, including all the district rungs.
\end{plain}
\proofref{Appendix~\ref{app:phase2b}, \S3, Theorem~1$'$ (step-by-step correspondence with PHASE2a Thm~1)}
\begin{revnote}(R-20.) Hypotheses H1--H4 listed explicitly.\end{revnote}

\begin{result}[Theorem 2, the rungs --- PHASE2b \S3]\label{res:t2}
\status{Proved.} Under A1, A5-d and $P(X\in A_D)=1$: (a) the sharp set of each rung
is the rung restriction intersected with each district's own sharp condition;
(b) the outer set is a polyhedron, every target bound is an LP, and for M1
$\Theta^{M1}_O=\bigcap_d\Theta^{(d)}_O$ in the common $k$; (c) the ladder
\[
\mathcal B^{L6}\subseteq\mathcal B^{L5}\subseteq\mathcal B^{L4(r)}\subseteq\mathcal B^{M1}\subseteq\mathcal B^{M1(0)}\subseteq\mathcal B^{M1(\delta)}\subseteq\mathcal B^{M1(\delta')}\subseteq\mathcal B^{M2}\subseteq DD^*
\]
holds for $0\le\delta\le\delta'$, sharp and outer, whenever the smaller rung's
model is non-empty.
\end{result}
\begin{plain}
Each stronger assumption gives a bound inside the weaker one's, so the ladder can
be reported as a table and readers can see what each assumption buys. M0 is not
on this chain.
\end{plain}
\begin{sketch}
Each rung is an affine mean model, so Thm~1$'$ applies district by district;
the outer version replaces the sharp condition by vertex non-negativity;
nesting follows because each step restricts the same polyhedral system further.
\end{sketch}
\proofref{Appendix~\ref{app:phase2b}, \S3, proofs of (a) steps 1--2, (b) steps 1--4, (c) steps 1--7}
\begin{revnote}
(R-21.) M1(0) is a \emph{sensitivity benchmark}: it keeps the common skew part
but drops the testable symmetric restrictions of D1(b).
\end{revnote}

\begin{result}[Proposition D3, M1 with two rows --- PHASE2b \S4]\label{res:d3}
\status{Proved.} For $R=2$ and vertices $t\in\{l_d,u_d\}$ with $0<l_d\le u_d<1$,
$\Theta^{M1}_O=\{\ell\le\kappa\le h,\sum_c\kappa_c=0\}$ with
$\ell_c=\max_d\max_t[\lambda_c(t)-\zeta_{c,d,1}/(1-t)]$ and
$h_c=\min_d\min_t[\rho_c(t)+\zeta_{c,d,2}/t]$. Pure-composition vertices give
testable restrictions and are dropped from the max/min as stated.
\end{result}
\begin{plain}
A district raises the lower bound when its baseline for group 1 is low and it has
polls that are almost all group 1: there group 1's rate is nearly observed, and
the shared contextual effect carries that information to every district. One
riding with lopsided polls and an extreme baseline can tighten the province-wide
bound.
\end{plain}
\proofref{Appendix~\ref{app:phase2b}, \S4, derivation steps 1--5; pure-composition case; checked against the LP (126 and 195 instances)}
\begin{revnote}(R-22.) $0<l_d\le u_d<1$ made visible; pure-composition vertices derived.\end{revnote}

\begin{result}[Proposition D4, when dummies matter --- PHASE2b \S4]\label{res:d4}
\status{Proved.} (a) If $\zeta\equiv0$, M1's outer set equals the pooled outer set
on $\conv\bigcup_dA_d$. (b) Under M1, $\E[\beta\mid X,N]=M_H(X)+\E[Z_D\mid X,N]$,
and M0 holds iff $\E[Z_D\mid X,N]$ is a.s.\ a function of $X$ alone, affine on
the support; $D\perp(X,N)$ suffices; $D\perp X$ alone does not (Remark~D4.1).
\end{result}
\begin{plain}
The pooled model is safe when district baselines, given composition and unit
size, vary at most linearly with composition. When strongholds also have
different baselines (a popular local candidate in the strongest ridings), or when
district baselines correlate with unit size, the pooled fit mistakes baseline
differences for contextual effects and can miss the truth. This is an
ecological omitted-variable bias.
\end{plain}
\proofref{Appendix~\ref{app:phase2b}, \S4, (a) steps 1--2, (b) steps 1--6; Remark~D4.1}
\begin{revnote}
(R-23.) The earlier version conditioned on $X$ only and claimed $D\perp X$
suffices; the size-weighted estimand uses $\E[\beta\mid X,N]$. Counterexample
with district correlated with size; check in \texttt{verify\_phase2b.py} [F].
Conditional on the observed design, P3.6(c) tests the stronger condition
$\zeta=0$ (Section~\ref{sec:inference}).
\end{revnote}

\begin{result}[Propositions D5--D7 --- PHASE2b \S4]
\status{Proved.}
\textbf{D5}: under M2 the bound is the Minkowski sum of the district bounds, so
widths add.
\textbf{D6}: under M1($\delta$) the bounds are nested in $\delta$, equal M2 for
$\delta\ge\bar\delta$, have continuous piecewise-linear endpoints, and equal
the M1 outer bound at $\delta=0$ (computed with the same $Q^{(d)},\zeta$).
\textbf{D7}: L5 point-identifies $B^*_{rc}=\sum_d(\mu^{(d)}_r/\mu_r)\alpha_{c,d,r}$
(Goodman within districts); L4($r$) pins $K_{c,rs}=Q_{c,rr}-Q_{c,rs}$ for
$s\ne r$, i.e.\ $(C-1)(R-1)$ coordinates, and point-identifies when $R=2$.
\end{result}
\begin{plain}
District-specific lines (M2) are safe but rarely much tighter than DD. The dial
$\delta$ traces a smooth path between M2 and M1, so a reader can see how much
the shared-effects assumption is doing. The constancy rungs give single numbers;
reported beside the bounds they show what ``no contextual effects'' implies
(Remark~D7.1: in the $3\times3$ example L5 gives 0.870 against a truth of
0.836, inside M1 but wrong).
\end{plain}
\proofref{Appendix~\ref{app:phase2b}, \S4, Propositions~D5 (steps 1--3), D6 (steps 1--4), D7 (a) and (b)}

\paragraph{Remarks 6.1--6.3 (PHASE2b \S5).} Remark~6.1: everything above is
identification; with about ten polls per riding the plug-in intersection is
biased inward (72\% coverage in PHASE1 \S7.1), which Phase~3 corrects.
Remark~6.2: with election periods as districts, M1 is ``contextual effects
unchanged'' and leads to Prop.~D8 (Section~\ref{sec:did}). Remark~6.3: the
grouping need not be geographic.

\paragraph{Worked examples (PHASE2b \S6).} In the $2\times2$ two-riding example
(truth 0.819) DD is $[0.453,0.879]$, M0 is empty (and the package reverts to
DD), M2 is $[0.563,0.870]$ and M1 is $[0.618,0.866]$. In the $3\times3$
six-district example the ladder shrinks monotonically, contains the truth at every
rung, and M2$\to$M1 narrows loyalty from width 0.525 to 0.450.

% Section 6: monotone EI, N1, N2 (PHASE2c) and PACKAGE_DESIGN section 5
\section{Monotone ecological inference: signs, N1 and N2}\label{sec:monotone}

Source: PHASE2c (Appendix~\ref{app:phase2c}), following Elzayn, Goldin, Guage,
Ho and Morton~\citep{egghm2025} (EGGHM, arXiv v1 numbering), and PACKAGE\_DESIGN
\S\S2.2 and 5 (Appendix~\ref{app:pkgdesign}). Checks: \texttt{verify\_phase2c.py},
29/29 PASS, each Batch~B correction tagged [F], [V] or [I]. Two settings are used:
\begin{itemize}[nosep]
\item \textbf{Finite population (N1).} The design is fixed; unit weights
$p_i=N_i/\sum_jN_j$; the realised rates satisfy $\beta_i\in S(X_i,T_i)$; the
target is the actual rate $B$. All N1 results read the unit tables as
\emph{continuous} (real cell masses).
\item \textbf{Superpopulation (K-rows and N2).} $m(v)=\E[\beta\mid V=v]$; the
target is $\tilde B$; a superscript $^*$ marks an association computed with $m$.
\end{itemize}

\subsection{The associations}

Fix a column and write $w^r_i=p_iX_{ir}/\mu_r$.
\begin{itemize}[nosep]
\item \emph{Between, pairwise:} $\delta_B(r,s)=E_p[X_rX_s(\beta_r-\beta_s)]$;
\emph{one-vs-rest:} $\delta^r_B=\sum_{s\ne r}\delta_B(r,s)$.
\item \emph{Within, W2 (pooled):} $\Delta=\sum_s\mu_s\Cov_{w^s}(X,\beta_s)$, one
sum of row-specific covariances.
\item \emph{Within, W1 (one-vs-rest):}
$\delta^r_W=\mu_r\Cov_{w^r}(\beta_r,X_r)+(1-\mu_r)\Cov_{w^{\neg r}}(\beta^{\neg r},X_r)$,
with A's share as the predictor in both terms.
\end{itemize}
\begin{plain}
The \emph{between} association says that, on average within the same unit, group
A is more likely than the others to choose the outcome. It is a weighted average
with weights $p_iX_{ir}(1-X_{ir})$: pure units get no weight and a large mixed
unit can dominate. It is not an ordering in every unit or against each rival. The
\emph{within} association says that rates rise, on average, with A's share. W1 and
W2 can have opposite signs while every group's rate rises (EI referee's example:
$\delta^A_W=-0.02633$, $\Delta_A=+0.00110$, check C14). A sign on W2 signs one
pooled sum, not each group. Covariances also reflect sorting; a causal reading
needs a separate design.
\end{plain}
\begin{revnote}
(P-E1--P-E5, P-E8, P-B28.) W2 was earlier described as ``each group
separately''; the weighting of $\delta_B$, the predictor in W1, sorting, the
Goodman-bias relation and the status of W2, pair and CR as extensions of EGGHM
were added as remarks on Definition~1.
\end{revnote}

\begin{result}[Lemma M1, identities --- PHASE2c \S1]\label{res:m1}
\status{Proved.} Under A1 only (with the stated per-part conditions):
(a) $\mathbf 1^\top\Delta=0$ and Goodman's $g$ satisfies $g-B_{\cdot c}=M^{-1}\Delta$;
(b) $NM_r-B_r=-\delta^r_B/\mu_r$;
(c) the one-vs-rest collapse is a valid two-group population, so
$ER^r-D^r=\delta^r_W/\Var_p(X_r)$ and $(NM_r-NM_{\neg r})-D^r=-\delta^r_B/(\mu_r(1-\mu_r))$;
(d) $B$, every $\delta_B$, $\Delta$ and $\delta^r_W$ are affine in the unit rates.
\end{result}
\begin{plain}
These are accounting identities. The neighbourhood model (NM) misses the true
gap exactly by the between association; the aggregate regression slope (ER) misses
it exactly by the within association. So a sign belief about either association
is a one-sided bound relative to a number computed from the data.
\end{plain}
\proofref{Appendix~\ref{app:phase2c}, \S1, proofs of (a)--(d); check C1}
\begin{revnote}(P-B28.) The invertibility of $M$ is needed for (a) only.\end{revnote}

\subsection{N1: sign restrictions on realised tables}

\begin{result}[Proposition M2, N1 is sharp for continuous unit tables --- PHASE2c \S2]\label{res:m2}
\status{Proved} (continuous reading). For $F=\{\beta:\beta_i\in S(X_i,T_i)\cap
U_i,\ g_k(\beta)\ge-\varepsilon_k\}$: (a) the values of $\varphi$ attained by
continuous unit tables are $\varphi(F)=[\min_F\varphi,\max_F\varphi]$, sharp
given the margins; (b) $F=\emptyset$ iff the restrictions contradict the margins
in the continuous problem; (c) under CR the set is a union of at most two
intervals; (d) the sets are nested componentwise in $\varepsilon$ with continuous
piecewise-linear endpoints; the preserving set of slacks is in general not a
product of intervals, and its boundary is the breakdown frontier; a scalar
breakdown value needs a stated path.
\end{result}
\begin{plain}
With sign beliefs alone and no regression, the bound is an LP over all unit
tables that fit the margins, and it is exactly right for real-valued tables.
With whole-number ballots it can differ (Remark~M2.1).
\end{plain}
\proofref{Appendix~\ref{app:phase2c}, \S2, Proposition~M2 steps 1--6}
\begin{revnote}
(P-B1, P-B2, I.10.) ``Sharp in the finite-population sense'' became ``sharp for
continuous unit tables, conditional on the margins''; ``the breakdown value'' for
vector slacks became the breakdown frontier with a scalar value along a path; in
round~2, the set and its frontier were distinguished. Checks C7 [F], C7b [V].
\end{revnote}

\begin{result}[Remark M2.1, integer counts --- PHASE2c \S2]\label{res:m21}
\status{Proved} for $2\times2$ with margins only; \status{open} otherwise (R15).
In $2\times2$ with margins only and sign rows on one column, the integer gap set
is the DD interval intersected with a lattice of step $1/N_1+1/N_0=1/(N\mu(1-\mu))$;
endpoints differ from the continuous ones by less than one step (e.g.\
$4.8\times10^{-6}$ for $N=10^6$, $\mu=0.3$). With unit restrictions, joint
columns or $R\ge3$, the integer--continuous gap is not bounded by this argument.
\end{result}
\begin{plain}
Counting whole ballots matters only for tiny tables or near-ties, at least in
$2\times2$; the referees' four-voter examples show that CR can then be empty and
an endpoint can move.
\end{plain}
\proofref{Appendix~\ref{app:phase2c}, \S2, Remark~M2.1 (proposition and proof); checks C7, C7a}
\begin{revnote}(P-B1, P-E11, I.5.) Added after review; the lattice step quoted with both group totals in round~2.\end{revnote}

\begin{result}[Proposition M3, two rows --- PHASE2c \S2]\label{res:m3}
\status{Proved.} For $R=2$, $\Var_p(X)>0$, continuous tables, no unit
restrictions and no other sign rows: (a) $\delta_W\ge0\iff D\le D_{ER}$ and
$\delta_B\ge0\iff D\ge D_{NM}$; (b) $\{\delta_W\ge0\}$ gives
$[D^-,\min\{D^+,D_{ER}\}]$, empty iff $D_{ER}<D^-$; (c) $D_{NM}\in[D^-,D^+]$
always, so between signs are never refuted in this problem; (d)
$D_{NM}=\gamma D_{ER}$, $\gamma\in(0,1]$, and CR has exactly one non-empty branch
when $\gamma<1$ and $D_{ER}\ne0$; at $\gamma=1$ or $D_{ER}=0$ both branches are
non-empty with the same gap image. The CR set is $[D_{NM},\min\{D^+,D_{ER}\}]$
if $D_{ER}\ge0$ and $[\max\{D^-,D_{ER}\},D_{NM}]$ otherwise; under it
$\operatorname{sign}D=\operatorname{sign}D_{ER}$.
\end{result}
\begin{plain}
In this base problem (two groups, real-valued tables, nothing else imposed) only a
within sign can be contradicted by the margins; a between sign and contextual
reinforcement are always feasible, because ``every group votes like its unit''
fits them. That non-refutability belongs to this base system only: under the
linear rung CR can be rejected jointly (Remark~R3.2); with integer counts CR can be
rejected outright; known cells or unit orderings remove the witness. The
equal-rates table satisfies the CR \emph{union}, in R$\times$C too, but only the
branch matching its within sign.
\end{plain}
\proofref{Appendix~\ref{app:phase2c}, \S2, Proposition~M3 steps 1--8 and ``Reading''; checks C2, C8, C8b, C15}
\begin{revnote}
(P-B3--P-B5, P-E12, P-E13; round~2 I.1, I.9.) The hypotheses ``no unit-level
restrictions, no other sign rows'' were missing; (d) failed at $\gamma=1$;
``CR is never refuted'' now names its base system; branches and their gap images
are distinguished; in round~2 the equal-rates witness was corrected to give the
CR union, not both branches (three-row example, C15 [F]).
\end{revnote}

\begin{result}[Proposition M4, joint restrictions --- PHASE2c \S2]
\status{Proved.} Imposing signs on several columns or rows jointly gives a subset
of each single restriction's feasible set.
\end{result}
\begin{plain}
Joint beliefs are weakly tighter. In $2\times3$ simulations, signing two columns
jointly narrowed column 1 strictly in 75 of 139 populations. With two
complementary columns, the \emph{same} weak within sign on both forces $D=ER$.
\end{plain}
\proofref{Appendix~\ref{app:phase2c}, \S2, Proposition~M4; checks C3, C3b}
\begin{revnote}(I.8.) The round-1 text said ``opposite within signs''; it is the same sign on complementary columns.\end{revnote}

\subsection{Signs under linear contextual effects (K-rows)}

\begin{result}[Proposition M5 --- PHASE2c \S3]\label{res:m5}
\status{Proved}; (c) \status{for M0 only}. Write the rung's means as
$m_i=m_{\rm id}(V_i)+L_ik$. (a) Every association $\delta^*$ is affine in $k$ under
every rung, so each sign is one LP row and CR is a union of two LPs. (b) Mean-level
ordering in every district on its support implies $\delta^*_B(r,s)\ge0$, pooled.
(c) For $R=2$ under M0, $\delta^*_W=a(w_1+d_1)+bw_1$ with $a,b\ge0$, so both
slopes $\ge0$ imply $\delta^*_W\ge0$. Under district rungs the pooled
$\delta^*_W$ contains a between-district term with no sign, and the implication
fails (six-unit example: within districts $+0.00064$, pooled $-0.04927$).
\end{result}
\begin{plain}
Under a linear rung the sign beliefs become cheap extra constraints on the
unknown $K$. But a pooled sign under district models is one restriction across
all districts: positive slopes in every district do not make the pooled within
association positive, because district baselines can fall as composition rises.
\end{plain}
\proofref{Appendix~\ref{app:phase2c}, \S3, Proposition~M5 steps 1--5; checks C4 (rebuilt), C9}
\begin{revnote}
(P-B7; E2-4.) M5 was stated for every rung with the pooled formula. Now (a) uses
the stacked parameterisation of each rung, (b) needs district-specific ordering,
(c) holds for M0 only. Sharpness and inference for M-SIGN under district rungs
became proof item R16 (open). A district rung with sign rows has no coverage
theorem yet.
\end{revnote}

\paragraph{Remark M5.1 (which $\delta$).} N1 uses the realised associations;
K-rows use $\delta^*$. Their difference is a design-weighted sum of unit residuals
with mean zero; under P3-S1a its standard deviation is at most
$\tfrac12\{\sum_i(\sum_r|c_{ir}|)^2\}^{1/2}$. The earlier ``order $p^{-1/2}$''
claim needed conditions (P-B8).

\subsection{N2: monotone nonparametric contextual effects (two rows)}

N2 is defined on the finite empirical design law (support points
$x_1<\dots<x_{J_X}$). Model N2($\sigma$): (S1), $\E[\beta\mid V]=m(X)$, and for
designated $(r,c)$, $\sigma_{rc}m_{rc}$ is nondecreasing on the support.

\begin{result}[Lemma N2.1 and Proposition N2.2 --- PHASE2c \S4]\label{res:n22}
\status{Proved} (on the finite support). N2.1: $m(x)\in S(x,\tau(x))$ at every
support point. N2.2: (a) the LP over $\mathcal M_O$ (support-point means in
$S(x_j,\tau(x_j))$ plus adjacent order rows) is an \emph{outer} set for $\tilde
B$, not claimed sharp; (b) every point of $\mathcal M_O$ extends to a monotone
Borel $m$, so on the complete finite support nothing is lost (on a coarser grid
it is only a relaxation); (c) with quadratic $\tau$, M0 with signed slopes lies
inside N2; (d) both rows nondecreasing imply $\delta^*_W\ge0$; (e) for $R\ge3$ a
direction needs a simplex chart, and (d) has no general analogue.
\end{result}
\begin{plain}
N2 assumes only ``no reversals'': each group's average rate moves in one
direction as its share grows, in any shape. It is weaker than straight lines with
signed slopes, so its bound contains theirs. It is not in the package.
\end{plain}
\proofref{Appendix~\ref{app:phase2c}, \S4, Lemma~N2.1 steps 1--2; Proposition~N2.2 steps 1--5; checks C5, C10, C10b}
\begin{revnote}
(P-B9, P-B10, P-B29.) The support convention was added; (b) is false on a grid
(Remark~N2.2.2, C10 [F]); the proposed sharp version was not sharp when size is
informative (Remark~N2.2.1, C10b [F]); (e) gained a chart and weighted
covariances.
\end{revnote}

\begin{result}[Conjecture N2.2$'$ --- PHASE2c \S4, Remark N2.2.1]\label{res:n22prime}
\status{Conjecture.} The sharp N2 set is
$\{m:\ m(X)\in\E[S(X,T)\mid V]\ P_V\text{-a.s.}\}$ plus the shape restrictions.
A sufficiency argument under the selection-expectation reading (due to the
partial-identification referee, round~2) is recorded as a remark; necessity in
general is not written, and a procedure that uses only $\E[T\mid X]$ does not
attain the set when size is informative. The conjecture stays open.
\end{result}
\begin{plain}
When unit size carries information, conditioning on composition alone can admit
mean rates that the full model excludes (example: $(0.75,0.25)$ admitted,
$(0.5,0.5)$ forced). The same caution applies to the planned M-SHARP series
bound (R9): it is ``tighter than $\ThO$'', not sharp.
\end{plain}
\begin{revnote}(P-B10; I.11.) Added as a conjecture in round~1; the referee's sufficiency remark added in round~2, conjecture kept open by decision.\end{revnote}

\begin{result}[Proposition N2.3, structural local identity --- PHASE2c \S4]\label{res:n23}
\status{Proved} (as a calculus identity). For differentiable $m_{1c},m_{2c}$ and the
structural aggregate $\tau_m(z)=zm_1(z)+(1-z)m_2(z)$: (a)
$\tau'_m=D(x)+\delta_W(x)$; (b)--(c) local signs and local CR in terms of
$\tau'_m$; (d) if both rows are nondecreasing on a neighbourhood of $x$, the
local sign is $\ge0$; monotonicity on the finite support is not enough, and the
local sign everywhere neither implies nor is implied by the aggregate $\delta^*_W$
sign. On a finite support $\tau'$ is not identified, so N2.3 is \emph{not} an
observable restriction; the N2-local rows were dropped (decision U4(a)).
\end{result}
\begin{plain}
EGGHM's local results are a statement about smooth functions, but the data never
identify a derivative at isolated compositions. So ``local'' sign beliefs cannot be
imposed from aggregates without an identified regression family.
\end{plain}
\proofref{Appendix~\ref{app:phase2c}, \S4, Proposition~N2.3 steps 1--3; checks C6, C11, C13, C16}
\begin{revnote}
(P-B11, P-B12, P-B34; I.4.) N2.3 was stated as an observable LP row with an
estimated $\tau'$ (C11 [F]); the smooth witness replaced a kinked one (C13 [I]);
in round~2 (d)'s premise became monotonicity on a neighbourhood (C16 [F]).
\end{revnote}

\begin{result}[Lemma N2.4, person-weighted bins --- PHASE2c \S4]\label{res:n24}
\status{Proved.} With person-weighted bin rates
$\bar m_{1b}=E_b[Xm_1]/\bar x_b$, $\bar m_{2b}=E_b[(1-X)m_2]/(1-\bar x_b)$ and the
model-implied bin mean $\bar\tau_b=E_b\tau$: (i) $\bar m_b\in S(\bar x_b,\bar\tau_b)$;
(ii) the targets are linear in the bin rates; (iii) monotonicity is preserved
across ordered bins. The binned LP is an outer relaxation of global N2.
\end{result}
\begin{plain}
With many compositions, group them into bins, but average rates weighting by
people in each group, not by units; the ordinary bin average breaks the
accounting identity (0.3125 against 0.25 in the referee's example). In samples the
observed bin mean of $T$ estimates $\bar\tau_b$ with error.
\end{plain}
\proofref{Appendix~\ref{app:phase2c}, \S4, Lemma~N2.4 proof; check C12}
\begin{revnote}(P-B13; I.6.) Replaces an invalid [CE-cvx] argument; round~2 distinguished the model-implied bin mean from the observed one.\end{revnote}

\paragraph{What Phase~3 must supply for N2 (PHASE2c \S6).} A weighted polynomial of
\emph{prespecified} degree $q$ (not chosen by cross-validation: post-selection
coverage is about 0.68 in the referee's design, check C12b), conditional on the
specification; then P3.5 applies to the coefficients. (P-B14, I.6.)

\subsection{Results from the package design (PACKAGE\_DESIGN \S\S2.2, 5)}

These results are in Appendix~\ref{app:pkgdesign}. Write
$q_{rs}(x)=(Qx)_r-(Qx)_s$ (identified) and $g_{rs}(x;k)=(Kx)_r-(Kx)_s$.

\begin{result}[G2, JKST plus signs in $2\times2$ --- PACKAGE\_DESIGN \S2.2; register R13]\label{res:g2}
\status{Proved.} Under LCE in $2\times2$ with $\gamma<1$: $\delta_W\ge0$ gives
$[L_J,\min(U_J,D_{ER})]$; $\delta_B\ge0$ gives $[\max(L_J,D_{NM}),U_J]$; CR gives
$[L_J,U_J]\cap[\min,\max](D_{ER},D_{NM})$, where $[L_J,U_J]$ is the JKST interval
for the gap.
\end{result}
\begin{plain}
Adding a sign belief to JKST cuts its interval at one of two data-computed
numbers. In 2000 random LCE designs with true signs, the within sign tightened
JKST in 83\% of designs (median width ratio 0.70), the between sign in 89\%
(0.67). A sign is redundant exactly when the JKST-implied interval for that
association already excludes zero.
\end{plain}
\proofref{Appendix~\ref{app:pkgdesign}, \S2.2, steps 1--3; [V1], [V2]}

\begin{result}[Proposition P-S, sandwich lemma --- PACKAGE\_DESIGN \S5.1; register R12]\label{res:ps}
\status{Proved.} For any $R$ and one column, with the between-group gap
$O_{rs}(x)=(Hx)_r-(Hx)_s$ and the weighted contextual derivative
$L_{rs}(x)=(x^\top H)_r-(x^\top H)_s$: (a) $O_{rs}+L_{rs}=2q_{rs}=\partial_{e_r-e_s}\tau$;
(b) $O_{rs}=q_{rs}+g_{rs}$, $L_{rs}=q_{rs}-g_{rs}$; (c) $O_{rs}\ge0$ and
$L_{rs}\ge0$ at $x$ iff $|g_{rs}(x;k)|\le q_{rs}(x)$, possible only if
$q_{rs}(x)\ge0$ (a data test).
\end{result}
\begin{plain}
An ordering belief (``group $r$ at least as supportive as $s$'') and a local
contextual belief restrict the same unknown from opposite sides, so together
they confine it to an interval whose half-width is an identified slope.
\end{plain}
\proofref{Appendix~\ref{app:pkgdesign}, \S5.1, steps 1--4; [V4]}

\begin{result}[Proposition P-R, rank lemma --- PACKAGE\_DESIGN \S5.2; register R5]\label{res:pr}
(a) One pair restricts an $(R-1)$-dimensional subspace of $\Kspace_c$: \status{proved}.
(b) A star tree of pairs reaches all $\binom R2$ directions: \status{proved}.
(c) Any spanning tree does: \status{checked} for chains, $R=3,4,5$ [V7];
general proof open (R5). (d) M-ORD $\cap$ M-LOC on a star tree bounds $K_c$ without
DD rows: follows from (b) and P-S.
\end{result}
\proofref{Appendix~\ref{app:pkgdesign}, \S5.2}

\begin{result}[Proposition P-N and Conjecture C1, nesting in $2\times2$ under M0 --- PACKAGE\_DESIGN \S5.3; register R1]\label{res:pn}
\status{Proved} for (a)--(b); (c) \status{conjecture}. Under pooled M0: (a) signed
slopes imply weighted-local on the hull; (b) $\delta_W=(a+b)(w_1+d_1x^*)$, so
weighted-local implies $\delta_W\ge0$ whenever $x^*=a/(a+b)$ lies in the
composition hull; (c) Conjecture C1: $x^*$ always lies in the hull (no violation
in 2000 random designs and an adversarial search of 50{,}000).
\end{result}
\begin{revnote}
(P-B7.) The chain ``signed slopes $\subset$ weighted-local $\subset\delta_W\ge0$''
is scoped to M0; it fails for pooled M1 (check I2 of the plan).
\end{revnote}

% Section 7: M-SIGN (PHASE5a)
\section{The sign module M-SIGN: R3, R2 and R10}\label{sec:msign}

Source: PHASE5a (Appendix~\ref{app:phase5a}), proof-register items R3, R2 and
R10. Checks: \texttt{verify\_R3.py}, \texttt{verify\_R2.py},
\texttt{verify\_R10.py}. \textbf{Scope: every statement is for the pooled LCE
model M0.} The sign and CR rows are valid linear rows under the district rungs
too (Prop.~M5(a)), but sharpness and inference there are proof item R16, not
proved. The package labels the module \emph{experimental} in 0.3.0 and 0.3.1.

\paragraph{Setting.} $\ThI$ and $\ThO$ are viewed in the free coordinates $k$.
The target is a model-implied linear $\varphi(k)=\varphi_0+h^\top k$ (for example
$\tilde B_{rc}$ or the gap $\tilde D=\tilde B_1-\tilde B_2$). An \emph{admissible
restriction} is a set $C(P)$ that is (C1) closed and convex and (C2) depends on
$P$ only through the design and $Q(P)$. Examples: any sign on $\delta^*_B$,
$\Delta^*_r$ or $\delta^{r*}_W$; mean-level ordering on $A$; the M-LOC rows. A
\emph{union restriction} is $C_+\cup C_-$; contextual reinforcement (CR) for a pair
$(g_B,g_W)$ is $C_\pm=\{\pm g_B\ge0,\pm g_W\ge0\}$.

\subsection{R3: sharpness of LCE plus restrictions}

\begin{result}[Theorem R3 --- PHASE5a \S1]\label{res:r3}
\status{Proved under conditions} (pooled M0). Assume A1--A5, standing
integrability, $P(X\in A)=1$ and $\operatorname{affdim}(A)=R-1$.
\begin{enumerate}[label=(\alph*),nosep]
\item For admissible $C$, the sharp identified set for $\varphi$ under LCE $\cap\,C$
is $[\min_{\ThI\cap C}\varphi,\max_{\ThI\cap C}\varphi]$ or empty.
\item $\ThO\cap C=\bigcup_{P'\in\mathcal P(P)}(\ThI(P')\cap C(P'))$: the LP
interval is sharp for the regression and the design; a finite LP when $C$ is
polyhedral; equal to (a) when $\ThI=\ThO$; weakly outer otherwise.
\item For a union, the set is $\varphi(\Theta\cap C_+)\cup\varphi(\Theta\cap C_-)$,
at most two intervals; the hull is sharp iff they intersect; a common witness in
$\Theta\cap C_+\cap C_-$ certifies that (sufficient only).
\item One-vs-rest (W1) CR is one LP: $C_+\cup C_-=\{D^{r*}(k)\in\operatorname{hull}\{\gamma_rER^{r*},ER^{r*}\}\}$,
two linear rows; in $2\times2$ this is the only CR, and under it
$\operatorname{sign}\tilde D=\operatorname{sign}ER^*$.
\item Refutation is joint: an empty population LP rejects the conjunction LCE
$\cap\,C$, not either part alone.
\end{enumerate}
\end{result}
\begin{plain}
Adding a sign or ordering belief to the pooled linear rung keeps the bound
``as tight as the regression allows'': nothing narrower can be learned from the
regression and the design. The full data could rule out more (Remark~R3.1). For
the default one-vs-rest comparison, reinforcement becomes a single band in the
gap, so no union is needed. If the bound is empty, it is the combination that is
contradicted, not the belief on its own.
\end{plain}
\begin{sketch}
(a) The identified set is the image of the convex compact $\ThI\cap C$ under an
affine map. (b) Apply the construction of PHASE2a Prop.~2 to each $k\in\ThO\cap
C$; (C2) keeps $C$ unchanged across laws with the same regression. (d) Apply
Lemma~M1(c) to the model means: $\delta^{r*}_W=\Var_p(X_r)(ER^{r*}-D^{r*})$ and
$\delta^{r*}_B=\mu_r(1-\mu_r)(D^{r*}-\gamma_rER^{r*})$; the union of the two
same-sign cases is the hull.
\end{sketch}
\proofref{Appendix~\ref{app:phase5a}, \S1.3, steps 1--15; checks S1--S7 (\S1.5)}
\begin{revnote}
(M-1; P-B15, P-B16, P-B17; I.7.) The dimension hypothesis is
$\operatorname{affdim}(A)=R-1$, not ``nonempty relative interior''. ``The LP
interval'' is a finite LP only for polyhedral $C$, and in round~2 the
finite-LP sentence was narrowed to $\ThO\cap C$. (e) earlier read ``never $C$
alone''; it now says the conjunction is rejected and gives both directions of
counterexample. The scope was restricted to pooled M0 (P-B7).
\end{revnote}

\paragraph{Remark R3.1 (full law versus regression).} On $X\in\{1/4,1/2,3/4\}$ with
a fair binary $T$ at every composition, the full law forces $k=0$, but the
regression admits a range. For $C=\{k\in[.2,.4]\}$ the full-law set is empty
while the regression-only gap set is $[1/6,1/3]$ (check S6 [V]).

\paragraph{Remark R3.2 (LCE can exclude CR although LCE fits).} On
$X\in\{0,1/4,1\}$ with $m_0=1-x$, $m_1=(49/55)x$, the vertex rows force $k=0$, so
LCE fits with a single point, $ER=0$ and $\tilde D=-261/1925$; both CR branches
are empty. In the model-free continuous problem CR is still feasible. Any ``cannot
be refuted'' statement must name its base system (check S7 [F]).

\paragraph{\S1.4: what fails.} For CR pairings other than W1, (d) fails: both
branches can be non-empty with disjoint images (search S5 found gaps up to about
31\% of hull width). For the pairwise version only $\delta^*_B(r,s)$ is pairwise;
$\Delta^*_r$ pools all groups (E2-1). N2 shapes, N1 realised signs and targets
mixing $\tilde B$ with $B$ are not covered.

\subsection{R2: when does an added restriction tighten the bound?}

\begin{result}[Proposition R2 --- PHASE5a \S2]\label{res:r2}
\status{Proved.} Let $F$ be the current (non-empty, compact) feasible set,
$F^*_U,F^*_L$ its optimal faces for $\varphi$, and $N$ an admissible restriction
with $F\cap N\ne\emptyset$.
\begin{enumerate}[label=(\alph*),nosep]
\item $U_N<U$ iff $F^*_U\cap N=\emptyset$ (one feasibility LP, carrying the
target intercept: $h^\top k\ge U-\varphi_0$).
\item For one row $a^\top k\le b$: tightening iff $\min_{F^*_U}a^\top k>b$ iff
every optimal dual multiplier is positive. A positive multiplier from one solver
optimum does not certify tightening.
\item $N$ is implied iff it tightens no linear target (the witness need not be a
reported target).
\item In $2\times2$ under LCE with $\gamma<1$, each of the four sign rows is either
implied or tightens any target non-constant in $w_1$; closed forms are those of G2.
\item For stacks, report added-first and added-last gains; zero added-last gain
means zero incremental gain for this target, not that the row is implied.
\end{enumerate}
\end{result}
\begin{plain}
A belief narrows a quantity only if it removes every table that attains that end
of the bound. The package tests this with one feasibility LP per end and never
trusts solver dual values. A belief can narrow the set of possible tables without
narrowing the quantity you care about; with three or more groups this was common
in simulation (28\% of cases).
\end{plain}
\proofref{Appendix~\ref{app:phase5a}, \S2.3, steps 1--14; checks T1--T4, T1b, T1c, T4b}
\begin{revnote}
(P-B18--P-B22.) The face LP in step~4 omitted the target intercept (the code
was correct; check T1c [F]); the duplicate-row example shows a positive dual does
not certify tightening; (d) gained its qualifications; (e)'s ``redundant given
the others'' read as ``implied'' and was reworded (T4b [F]); width-gain
inference must treat both endpoints jointly.
\end{revnote}

\subsection{R10: inference with added rows, unions and near-empty sets}

\begin{result}[Theorem R10 --- PHASE5a \S3]\label{res:r10}
\status{Proved under conditions} (pooled M0; by interval version, below). Assume
P3.5's conditions for the augmented row system and that the restrictions hold at
the truth.
\begin{enumerate}[label=(\alph*),nosep]
\item \emph{Added rows.} For finitely many rows linear in $\beta$, P3.5 and
P3.6(a) hold for the augmented system, version by version. For actual contrasts
the residual allowance must cover the contrast (a contrast envelope, or the
per-cell sum the package uses).
\item \emph{Parallel disjunction (version (a) only).} Inflate each piece, then take
the maximum; this covers at the nominal level. W1 and $2\times2$ CR are of this
form (the ``slab''); no branch is selected from $\widehat{ER}$.
\item \emph{General union.} $\mathrm{CI}(C_+)\cup\mathrm{CI}(C_-)$, each at the
full level, covers the true target with probability $\ge1-\alpha$ asymptotically;
no Bonferroni split is needed. This is point coverage only.
\item \emph{Near-empty sets.} $P(\mathrm{CI}=\emptyset)\to\le\alpha_U+\alpha_L$
under the model, including point-identified cases; conditioning on non-emptiness
can only raise coverage. If the restrictions are false, no coverage statement
holds; a narrow non-empty interval is no evidence that they are true.
\item \emph{Plug-in brackets are not valid.} Imposing the estimated bracket
$\operatorname{hull}\{\gamma_r\widehat{ER},\widehat{ER}\}$ without inflation can
under-cover persistently: in a repeated three-point design
$\limsup_nP(\tilde D_0\in\mathrm{CI}_{\rm naive})\le\Phi(\tfrac7{12}z_{.995})=0.9335<0.95$.
\end{enumerate}
\end{result}

\begin{longtable}{@{}L{0.30\linewidth}L{0.42\linewidth}L{0.20\linewidth}@{}}
\caption{Status of R10 by interval version (PHASE5a \S3.2).}\label{tab:r10}\\
\toprule Interval version & What R10 claims & Status \\ \midrule\endfirsthead
\toprule Interval version & What R10 claims & Status \\ \midrule\endhead
P3.5(a), conservative & (a)--(d), including the slab (b) & Proved \\
P3.5(b1), CLR over all dual vertices & (a), (c), (d) under the full augmented CLR conditions & Cited \\
P3.5(b2)(i), fixed superset & branch union and crossing arguments & Proved \\
P3.5(b2)(ii), projection & both ends, any candidate set; $\Sigma$ p.d.\ or exact reduction & Proved \\
P3.5(b3), CLR on simulated candidates (package default for CR) & the inclusion and crossing steps are deterministic; their probabilities are conjectural & Algorithm \\
\bottomrule
\end{longtable}

\begin{plain}
Sign beliefs enter the same confidence-interval machinery as the rung's own
constraints. Reinforcement uses a band whose two pieces are each widened by their
own uncertainty before the larger is taken; this is valid without knowing which
branch is true. The tempting shortcut, plugging the estimated neighbourhood and
Goodman numbers into the bracket, is not a confidence interval, and more data do
not cure it. When the belief is false, the interval stays narrow and can be
reproducibly wrong (Section~\ref{sec:evidence}).
\end{plain}
\proofref{Appendix~\ref{app:phase5a}, \S3.3, steps 1--23 (step 18--23 is the referee's fixed-design proof of (e)); checks V9--V13}
\begin{revnote}
(M-2; P-B23--P-B27, P-E17, P-E21, P-E22, P-E33; I.2; E2-2, E2-3.) The status of
each claim is now given by interval version (table above) with CLR's
Condition~P spelled out; zero-variance candidates are exact only with zero
gradient; the contrast envelope was added; (b)--(d) gained the one-event,
point-coverage and triangular-array qualifications and the belief-selection
sentence; (e) became an analytic theorem (referee's proof, credited) instead of a
Monte Carlo sketch; ``a fortiori'' was replaced by a direct $\le\tfrac12$ argument
(V10 [F]); union lengths are occupied lengths (V11). In round~2 the R10 scope was
narrowed to pooled M0 (I.2), and the coverage figure was tied to the procedure:
the literal intersection covered 304/600, CR-naive 550/600 (V13 [V]).
\end{revnote}

\paragraph{Consequences for the package (PHASE5a \S\S1.6, 2.5, 3.5).} For pooled
M0 with sign, ordering or local rows the LP interval may be labelled ``sharp for
the regression''; under district rungs, ``outer for the regression''. W1 and
$2\times2$ CR use the slab for the conservative interval and the branch union for
CLR (the default, algorithmic). Other CR pairings use the union of branch CIs and
can print two or more pieces. Redundancy verdicts use the face test. Emptiness is
reported as a test result, naming the rejected combination and whether the rung
alone is rejected.

% Section 8: inference (PHASE3)
\section{Inference}\label{sec:inference}

Source: PHASE3 (Appendix~\ref{app:phase3}). Checks: \path{verify_phase3.py},
17/17 PASS. Monte Carlo: \path{phase3_montecarlo.py}, 1000 replications per
design (Section~\ref{sec:evidence}).

\paragraph{Sampling (PHASE3 \S0.1).} \textbf{(P3-S1a)} Given the design and the
district effects, the unit rates are independent with $\E[\beta_i\mid\text{design}]=m(V_i)$.
\textbf{(P3-S2)} Inference is conditional on the design, so the sample-hull
vertices are fixed. The uncertainty is model-based, conditional on the design
(\S0.4, following Abadie, Athey, Imbens and Wooldridge~\citep{aaiw2023}): every unit is observed, nothing is assigned, and district shocks are fixed
features of this electorate (in the mean under M1/M2; a misspecification under
M0). No clustered variance is used. Independence is a substantive assumption: it
does not follow from observing every unit, and within-district dependence with
no mean shift violates it without being detected by the pooling test.
\begin{revnote}
(R-28, R-36; A2-5.) Round~1: the earlier P3-S1 allowed arbitrary within-district
dependence for M0 while P3.5(a) assumed independence; it was replaced by P3-S1a,
the \texttt{cluster} option was removed and the pooling test added. Round~2: the
two arguments (``assume independence given district effects'' versus ``all units
observed, so no dependence'') were separated; only the first is made.
\end{revnote}

\paragraph{Estimators (PHASE3 \S1).} OLS of each of $C-1$ columns on the
homogeneous quadratic monomials (plus district $\times X_r$ shifts for M1; one
regression per district for M2), with HC3 covariance. Each constraint is
$s_j(\beta)+a_j^\top k\ge0$ with $s_j$ linear in the coefficients $\beta$; the
target is $\tilde B(k)=c(\beta)+\gamma^\top k$ with sample weights. The plug-in
bound is $c(\hat\beta)+[\min,\max]\{\gamma^\top k:s(\hat\beta)+Ak\ge0\}$.

\begin{longtable}{@{}L{0.12\linewidth}L{0.13\linewidth}L{0.42\linewidth}L{0.25\linewidth}@{}}
\caption{Which result holds in which regime (PHASE3 \S0.5).}\label{tab:regimes}\\
\toprule Model & Regime & Formal result & Status \\ \midrule\endfirsthead
\toprule Model & Regime & Formal result & Status \\ \midrule\endhead
any rung & any, any $n$ & Theorem P3.8 (finite-sample CI and emptiness test) & Theorem; uninformative in width \\
M0 & any & P3.2, P3.4, P3.5(a), (b1), (b2) & Theorem (asymptotic) \\
M1, M2 & R1 ($D^\star$ fixed) & P3.2, P3.4, P3.5(a), (b1), (b2) & Theorem (asymptotic) \\
M1, M2 & R3 & the same with a high-dimensional Gaussian approximation & Conjecture \\
M1 & R2 & coverage of the CLR interval & Simulation finding \\
any & any & the CLR interval as implemented & Conjecture P3.5(b3) \\
\bottomrule
\end{longtable}

\begin{result}[Lemma P3.1, the endpoint is a minimum of linear functions --- PHASE3 \S2]\label{res:p31}
\status{Proved} (fixed constraint system). If the LP is feasible and bounded at
$\beta$, $\varphi_U(\beta)=c(\beta)+\min_{\lambda\in\Lambda}\lambda^\top s(\beta)$
with $\Lambda$ the finite set of dual vertices; $\varphi_U$ is concave and
piecewise linear in $\beta$.
\end{result}
\begin{plain}
Each end of the bound is the tightest of finitely many candidate limits, each
linear in the regression coefficients. JKST's case has four candidates and assumes
one wins clearly; in R$\times$C with districts there are many and near-ties are
the norm.
\end{plain}
\proofref[pE-res-lemma-P3.1]{Appendix~\ref{app:phase3}, \S2, steps 1--5}
\begin{revnote}(R-25.) Stated for a fixed constraint matrix and objective; joint variation in $(Q,\mu)$ is not covered.\end{revnote}

\begin{result}[Proposition P3.2, consistency and rate --- PHASE3 \S3]\label{res:p32}
\status{Proved.} Under P3-S1--S2, A1--A5, $\sqrt n(\hat\beta-\beta_0)=O_p(1)$ and
Slater (strict feasibility): (a) $\hat\varphi_U-\varphi_U(\beta_0;\hat A)=O_p(n^{-1/2})$;
(b) $\varphi_U(\beta_0;\hat A)\to\varphi_U(\beta_0;A)$ for nested or boundedly
convergent hulls; (c) since $\hat A_d\subseteq A_d$, estimating the hull only widens
the bound.
\end{result}
\begin{plain}
The plug-in bound converges at the usual rate to the bound for the sample hull,
which is weakly wider than the population bound. Coverage of the true parameter
never depends on the hull, because the true parameter satisfies the constraints at
every observed composition. Slater fails when an implied rate is exactly 0 at a
vertex (a structural zero); the rate result is then open, but coverage results do
not use Slater.
\end{plain}
\proofref[pE-res-proposition-P3.2]{Appendix~\ref{app:phase3}, \S3, steps 1--9; Remarks}
\begin{revnote}(R-26, R-32.) Non-nested hulls; Slater as an interiority condition; the error split into estimation ($O_p(n^{-1/2})$) and hull ($\ge0$, rate not shown) parts.\end{revnote}

\begin{result}[Theorem P3.4, asymptotic distribution for a fixed constraint system --- PHASE3 \S4]\label{res:p34}
\status{Cited structure} (directional delta method~\citep{fangsantos2019}).
If the constraint system is fixed along the sequence and
$\sqrt n(\hat\beta-\beta_0)\Rightarrow\mathcal Z\sim N(0,\Sigma)$, then
$\sqrt n(\hat\varphi_U-\varphi_U)\Rightarrow\min_{\lambda\in\Lambda^*}g_\lambda^\top\mathcal Z$.
(i) A single active candidate gives a normal limit. (ii) If two active gradients
differ in a direction where the estimator varies, the limit has strictly negative
mean: the plug-in bound is too narrow. (iii) Then the standard bootstrap is
inconsistent.
\end{result}
\begin{plain}
When several candidate limits are nearly tied, the plug-in end picks the one that
looks tightest by chance, so the estimated bound is biased inward. This is why the
plug-in is not a confidence interval.
\end{plain}
\proofref[pE-res-theorem-P3.4]{Appendix~\ref{app:phase3}, \S4, steps 1--6; Remark~P3.4.1 (use $C-1$ columns)}
\begin{revnote}
(R-27; A2-2.) Round~1: (ii) needs gradients that differ where $\Sigma$ has
variance (a null-direction difference gives a normal limit). Round~2: the theorem
was stated ``with the sample hull fixed'', which is not coherent as the hull grows;
it now requires a fixed constraint system (referee's $f_n=\min\{\theta,-\theta+n^{-1/2}\}$
gives $\min\{Z,1-Z\}$).
\end{revnote}

\begin{result}[Proposition P3.4$'$, triangular arrays --- PHASE3 \S4; register R17]\label{res:p34prime}
\status{Proved.} If the constraint system changes with $n$, $|\Lambda_n|\le\bar L$,
gradients converge and the localised slacks
$c_{\lambda,n}=\sqrt n(\theta_{\lambda,n}(\beta_0)-\varphi_{U,n}(\beta_0))\to c_\lambda\in[0,\infty]$,
then $\sqrt n(\hat\varphi_{U,n}-\varphi_{U,n}(\beta_0))\Rightarrow\min_{\lambda:c_\lambda<\infty}(g_\lambda^\top\mathcal Z+c_\lambda)$.
P3.4's limit obtains under separation ($c_\lambda\in\{0,\infty\}$).
\end{result}
\begin{plain}
A candidate that is inactive only by a gap of order $1/\sqrt n$ still shapes the
limit. The result describes the sample-hull endpoint; it says nothing about the
population-hull error.
\end{plain}
\begin{sketch}
By Lemma~P3.1 the scaled error equals exactly $\min_\lambda(g_{\lambda,n}^\top\sqrt
n(\hat\beta-\beta_0)+c_{\lambda,n})$; candidates with $c\to\infty$ drop out;
continuous mapping.
\end{sketch}
\proofref[pE-res-proposition-P3.4p]{Appendix~\ref{app:phase3}, \S4, Proposition~P3.4$'$ steps 1--4; check ``P3.4$'$ [F]''}
\begin{revnote}(A2-2.) New in round~2.\end{revnote}

\begin{result}[Theorem P3.5, confidence intervals --- PHASE3 \S5]\label{res:p35}
Fix levels $\alpha_U,\alpha_L,\alpha_E$; candidate upper endpoints
$\theta_\lambda(\beta)=g_\lambda^\top\beta$ with standard errors $\hat\sigma_\lambda$.
\begin{enumerate}[label=(\alph*),nosep]
\item \textbf{Conservative} (\status{proved}). Inflate every constraint and the target
constant by a simultaneous Gaussian critical value $\hat k$, solve the LP over the
inflated set $\hat\Theta^+$, and add $\pm z_{\alpha_E}S_1$ with
$S_1=\tfrac12(\sum_iw_{ir}^2)^{1/2}$. Then
$\liminf P(B_{rc}\in\mathrm{CI})\ge1-\alpha_U-\alpha_L-\alpha_E$; without $S_1$ it
covers $\tilde B$.
\item[(b1)] \textbf{CLR with the full candidate set} (\status{cited}: Chernozhukov, Lee
and Rosen~\citep{clr2013}, Thm~4, under Condition~P). Same coverage.
\item[(b2)] \textbf{Any candidate set with a simultaneous critical value}
(\status{proved}). For any data-dependent set of dual-feasible points,
$\hat U=\min_{\lambda\in\hat\Lambda}(\hat\theta_\lambda+\hat k\hat\sigma_\lambda)$
covers $\varphi_U$ with probability $\ge1-\alpha_U$ asymptotically if (i) $\hat k$ is
computed over a fixed finite superset, or (ii) $\hat k=\sqrt{\chi^2_{q,1-\alpha}}$
(projection) with $\Sigma$ positive definite and $q=\dim\beta$, or an exact reduced
full-rank parameterisation. An exact reduction means documented identities
(for example, adding-up removes column $C$); a rank taken from a numerical cutoff
on $\hat\Sigma$ is not one, so then $q=\dim\beta$ must be used, and dropping
candidates that load on numerically null directions is only a safeguard (it can
only widen the interval).
\item[(b3)] \textbf{The CLR procedure as implemented} (\status{conjecture /
algorithm}). Candidates are collected from the optimal dual vertices at
$\hat\beta$ and at simulated draws. A reduction to (b1) is proved when the
collected set contains CLR's selected set and a minimiser; recovery of rare
vertices, the downward effect on $k_p$, and growing $|\Lambda|$ are open.
\item[(c1)] \textbf{Variance-adaptive allowance} (\status{proved}; register R18).
Replace $S_1$ by the Bhatia--Davis~\citep{bhatiadavis2000} bound $\hat S_{BD}$ computed from a P3.5(a) or
projection interval for $\bar m=\sum_i\pi_im_i$ on the rung's own constraints;
if $\max_iw_i^2/\bar s^2\to0$, coverage is $\ge1-\alpha_{\rm est}-\alpha_E-\alpha_v$.
\item[(c2)] With a CLR interval for $\bar m$: coverage conditional on that interval
covering; inherits (b3)'s status.
\end{enumerate}
\end{result}
\begin{plain}
The confidence interval widens each candidate limit by its own uncertainty before
taking the tightest, so noisy candidates (for example, districts with few polls)
rarely set the limit. Three versions carry proofs: the conservative interval (all
constraints widened at once), the projection interval (one ellipsoid for the
whole coefficient vector; valid for any set of candidates, but wide when there
are many coefficients) and the fixed-superset version. The package's default,
CLR with simulated candidates, is narrower in grouped designs and covered in every
simulation design, but its validity is not proved: it is an \emph{algorithm}.
For the actual rate $B$, an extra allowance covers each unit's local luck; a
tighter, variance-adaptive allowance is valid if its own input interval comes
from a proved method.
\end{plain}
\begin{sketch}
(a) On one simultaneous event the true $k_0$ is feasible for the inflated
constraints, so $\tilde B$ is inside the inflated LP range; the residual term is
handled by Chebyshev or Lindeberg along subsequences. (b2) By weak duality every
dual-feasible candidate is an upper bound for $\varphi_U(\beta_0)$; a simultaneous
event for all candidates (a fixed superset, or Cauchy--Schwarz in the metric of
$\hat\Sigma$ for projection) then covers. (c1) Concavity of $m(1-m)$ bounds the
residual variance by $\bar s^2$; on the event that the $\bar m$ interval covers,
$\hat S_{BD}\ge\bar s$; a union bound finishes.
\end{sketch}
\proofref[pE-res-theorem-P3.5]{Appendix~\ref{app:phase3}, \S5: (a) steps 1--7; (b2) steps 1--4; (b3) proved reduction; (c1) steps 1--5}
\begin{revnote}
(R-29; P-B40, P-B41; I.3; A2-3.) Round~1 (Phases~1--3 referee): the simulated-candidate
CLR was claimed as a theorem; P3.5(b) was split into (b1) cited, (b2) proved and
(b3) conjecture, and the package labels CLR algorithmic with projection alongside.
Batch~B: (b2)(ii) took $q=\operatorname{rank}\Sigma$ with a pseudoinverse, which
does not control errors in null directions (counterexample $(Z/\sqrt n,-1/n)$);
it now needs $\Sigma$ positive definite or an exact reduction. The residual step
of (a) applied a CLT under the weight condition alone; the referee's
Chebyshev/Lindeberg subsequence argument replaced it. Round~2: a numerical rank
cutoff is a safeguard, not case (ii-b), so the package uses $q=\dim\beta$; the
projection interval for $B$ was not covered when its Bhatia--Davis bound came from
CLR, which led to P3.5(c).
\end{revnote}

\begin{result}[Proposition P3.6, tests --- PHASE3 \S5]\label{res:p36}
(a) \textbf{Emptiness test} (\status{proved}): reject the rung if $\hat\Theta^+$ is
empty; level $\le\alpha_U+\alpha_L$. (b) \textbf{M1 against M2}: Wald test of the
$(D^\star-1)(C-1)R(R-1)/2$ overidentifying restrictions; $\chi^2$ calibration
\status{proved for fixed dimension} (R1). (c) \textbf{Pooling, M0 against M1}
(\status{proved}): conditional on the design, with $R$ affinely independent
compositions per district, M0's mean model holds iff all district shifts vanish;
test $\zeta=0$ with $(D^\star-1)(C-1)R$ restrictions. It tests mean shifts only.
\end{result}
\begin{plain}
An empty confidence set is a real rejection at the stated level. The curvature
test checks M1's testable part. The pooling test says whether the pooled model is
wrong for this design; rejection does not mean its intervals fail (design S1), and
the test has no power against within-district dependence.
\end{plain}
\begin{sketch}
(a) On the simultaneous event of P3.5(a) the true $k_0$ lies in the inflated set,
so emptiness requires that event to fail. (c) If M0 held for the observed units,
then in each district $(G_c-\zeta_{c,d}\mathbf 1^\top)X_i=0$ for all its units, with
$G=H'-H$. Because the district has $R$ affinely independent compositions, which
lie in $\{\mathbf 1^\top x=1\}$ and are therefore linearly independent, this forces
$G_c=\zeta_{c,d}\mathbf 1^\top$ in every district. With $\zeta_{c,1}=0$, $G_c=0$ and all
shifts vanish. So M0 holds iff $\zeta=0$, a linear hypothesis on the M1
regression.
\end{sketch}
\proofref[pE-res-proposition-P3.6]{Appendix~\ref{app:phase3}, \S5, Proposition~P3.6 proofs of (a) and (c)}

\paragraph{Calibration by restricted wild bootstrap (package 0.3.0; PHASE3 \S5).}
\status{Simulation-supported, not proved for growing dimension.} In moderate
dimension the HC3 Wald statistic against $\chi^2$ over-rejects (pooling test, 38
restrictions: 0.078 at nominal 5\%). The package fits the null model, draws
$Y^*=\hat Y_0+\hat e_0v$ with Rademacher $v_i$ shared across the $C-1$ equations,
refits the larger model and recomputes the HC Wald statistic;
$p=(1+\#\{W^*\ge W\})/(B+1)$, $B=399$ by default. Rejection rates at 5\% in design
S0: pooling 0.043 (SE 0.005; $\chi^2$ 0.078, F 0.072); M1 vs M2 0.054 (SE 0.007;
$\chi^2$ 0.057). Scripts: \path{benchmarks/pooling_test_size.R},
\path{benchmarks/m1_test_size.R}.
\begin{revnote}(R-30, R-37; WS6; A2-5.) Degrees of freedom corrected; bootstrap added after the over-rejection was found; the pooling test is stated as a test of mean shifts only.\end{revnote}

\begin{result}[Theorem P3.8, finite-sample CI and emptiness test --- PHASE3 \S5.1]\label{res:p38}
\status{Proved.} Under P3-S1a, a correctly specified mean model and a full-rank
design, with Hoeffding~\citep{hoeffding1963} radii for every constraint and the target constant (union
bound over $J+2$ events) and $t_E=(\tfrac12\log(2/\alpha_E)\sum_iw_{ir}^2)^{1/2}$:
for every $n$ and in every regime, (a) $P(B\in\mathrm{CI}^{\rm fin}\mid\text{design})\ge1-\alpha_s-\alpha_E$;
(b) without $t_E$ it covers $\tilde B$ at $1-\alpha_s$; (c) the emptiness test has
exact level $\le\alpha_s$; (d) intersecting with DD keeps (a).
\end{result}
\begin{plain}
This interval needs no large-sample theory and no fixed number of districts, so it
is the only formal guarantee with many small districts. In every design tried it is
as wide as the logical range $[0,1]$, mainly because of the
constraint-by-constraint union bound. It is valid but not yet useful, and is not
in the package.
\end{plain}
\begin{sketch}
Every constraint term and the target constant are linear functionals of the OLS
fit, and each is a sum over units of independent terms whose range is known from
the design (the outcome shares lie in the simplex). OLS is unbiased under the
mean model, so Hoeffding's inequality gives a finite-sample radius for each. A
union bound over the $J+2$ one-sided events puts the true $k_0$ in the widened
set; the residual term gets its own Hoeffding radius. No Gaussian approximation
is used, so nothing depends on the regime.
\end{sketch}
\proofref[pE-res-theorem-P3.8]{Appendix~\ref{app:phase3}, \S5.1, steps 1--7}
\begin{revnote}(R-31; A2-6.) New after round~1; round~2 agreed it is a fallback guarantee and left informative R2 inference open.\end{revnote}

\begin{result}[Proposition P3.7, many districts with few units --- PHASE3 \S6]
\status{Proved.} (a) Dropping districts' constraints is conservative. (b) If the
retained districts are chosen from the design only, any procedure keeps its
conditional coverage; asymptotic results transfer only if their regime conditions
hold for the retained districts.
\end{result}
\begin{plain}
With many small districts, ignoring some districts' constraints can only widen
the bound, so it never costs validity of the identified set. If the retained
districts are chosen from the design alone (for example, those with at least 30
polls or enough spread in composition), not from the outcomes, any interval that
is valid for a fixed set of districts stays valid. Selection does not, by itself,
move the problem into a regime where the asymptotic theory applies.
\end{plain}
\proofref[pE-res-proposition-P3.7]{Appendix~\ref{app:phase3}, \S6, steps 1--2}
\begin{revnote}(R-31.) (b) was weakened from ``every result holds conditionally'' to a conditional statement.\end{revnote}

\begin{result}[Proposition S1, bounded nonlinearity --- PHASE3 \S7]\label{res:s1}
\status{Proved} (an outer bound). If $\E[\beta\mid V]=M_H(X)+\Delta(X)$ with $\Delta$ in
the undetectable class of PHASE2a Prop.~9 and $\|\Delta\|_\infty\le\eta$, then
$\tilde B_{rc}\in[\min_{\ThO(\eta)}\varphi-\eta,\max_{\ThO(\eta)}\varphi+\eta]$ with
$\ThO(\eta)=\{k:s+Ak\ge-\eta\mathbf 1\}$.
\end{result}
\begin{plain}
If rates may depart from straight lines, invisibly, by up to $\eta$, the bound
widens in a controlled way. The breakdown value $\eta^*$ of a conclusion is the
largest departure under which it survives; in the $3\times3$ example,
$\eta^*(\text{loyalty}>0.30)=0.060$. The same logic applies to the dials $\delta$
(district heterogeneity, change between periods).
\end{plain}
\proofref[pE-res-proposition-S1]{Appendix~\ref{app:phase3}, \S7, steps 1--3}
\begin{revnote}(R-34.) Labelled an outer bound under bounded deviations, not the sharp set of the relaxed model.\end{revnote}

\paragraph{Boundary cases and JKST's $\mathrm{CI}_x$ (PHASE3 \S\S4--5).} When the
true $K$ is on the boundary, coverage is close to nominal; in the interior it tends
to one (Imbens--Manski). Near point identification, one-sided intervals per
endpoint under-cover unless the critical value adapts~\citep{stoye2009}; the
two-sided split used here is uniformly valid at some cost in width. JKST's
$\mathrm{CI}_x$ is version (b) with a single candidate per side.

\paragraph{Practical recipe (PHASE3 \S9).} Fit on $C-1$ columns with HC3; report
the plug-in as an estimate of the identified set; report CLR (algorithmic) for
$\tilde B$ and $B$ with the projection interval alongside (for $B$, with $\bar m$
from the projection interval on the rung alone, or Popoviciu); intersect with DD;
report the emptiness, M1 and pooling tests; report sensitivity in $\eta$ and
$\delta$ with breakdown values; with few units per district, prefer selection,
M1($\delta$) or coarser groups, and say that CLR's coverage there rests on
simulation.

% Section 9: DiD material
\section{Changes across periods and difference-in-differences}\label{sec:did}

Only one DiD result has been through external review: PHASE2b Proposition~D8,
which treats election periods as ``districts'' under M1. The rest of the DiD
material, in PHASE1b (\texttt{PHASE1b\_monotone\_did.md}) and PHASE2d
(\texttt{PHASE2d\_did.md}), is a \textbf{proposal that has not been reviewed}
(\texttt{MSIGN\_REVIEW\_LOG.md} item M-13; \texttt{EXTERNAL\_REVIEW\_PLAN.md}).
It is summarised here only so that the status table is complete; it is not
reproduced in the appendices, and nothing in Sections~\ref{sec:rxc}--\ref{sec:inference}
depends on it.

\subsection{The reviewed result: Proposition D8}

\begin{result}[Proposition D8, the change in rates under M1 with periods as districts --- PHASE2b \S5]\label{res:d8}
\status{Proved.} Let $D\in\{\text{pre},\text{post}\}$; assume M1, (S1)--(S2-d),
$P(X\in A_D)=1$ and A5-M1. With per-period moments $\mu^{[t]}$,
\[
B^*_{rc,\text{post}}-B^*_{rc,\text{pre}}=\text{(identified)}+L(K),\qquad
L(K)=\sum_{s\ne r}K_{c,rs}\,\Delta q_{rs},\quad
\Delta q_{rs}=\frac{\mu^{[\text{post}]}_{rs}}{\mu^{[\text{post}]}_r}-\frac{\mu^{[\text{pre}]}_{rs}}{\mu^{[\text{pre}]}_r}.
\]
(a) \emph{Identification by cancellation:} the change does not depend on $K$ for
every $K\in\Kspace$ iff $\Delta q_{rs}=0$ for all $s\ne r$ (for example when the
size-and-composition moments are equal across periods). (b) \emph{Identification
given the feasible set:} the change is point-identified iff
$\max_{K\in\Theta}L(K)-\min_{K\in\Theta}L(K)=0$, where $\Theta$ is the identified
set for $K$. This holds under (a), and also when $\Theta$ fixes $L$ although
$\Delta q\ne0$ (example: $2\times2$, $Q_{1,12}=0$, LCE on the whole simplex, so
$\Theta=\{0\}$).
\end{result}
\begin{plain}
If contextual effects are unchanged between two elections and only baselines move
(M1 across periods), the change in ecological transition rates is pinned down when
the electorate's size-weighted composition moments are the same in both periods,
or when the bounds happen to fix the unknown combination. Otherwise the change is
bounded by projecting the common identified set. This is a statement about the
difference of two ecological rates. A causal (difference-in-differences) reading
needs further structure, developed in PHASE2d, which is not yet reviewed.
\end{plain}
\proofref{Appendix~\ref{app:phase2b}, \S5, Proposition~D8 steps 1--5; check D8b [F]}
\begin{revnote}
(R-24; A2-4.) Round~1: promoted from Remark~6.2 to a proposition with the exact
moment condition. Round~2: the proposition said the change is point-identified,
whatever the identified set for $K$, \emph{iff} $\Delta q=0$; that is false as a
statement about identification (referee's example). It now separates (a)
cancellation, which is sufficient, from (b), necessary and sufficient given the
identified set, and drops the causal reading.
\end{revnote}

\subsection{Proposals not yet reviewed (PHASE1b, PHASE2d)}

PHASE2d develops DiD bounds for the Wisconsin and Iowa suffrage referenda (rows:
veterans and non-veterans; columns: yes, no, abstain). Its results are written
with proofs and checked in \texttt{verify\_phase2d.py} (16 PASS) and
\texttt{verify\_phase1\_mei\_did.py}, but they have \textbf{not been reviewed}.
In brief:
\begin{itemize}[nosep]
\item \textbf{Lemma G}: a general sharpness lemma for affine mean models with
convex selection sets (as PHASE2b Thm~1$'$).
\item \textbf{Proposition $\Delta0$}: with no link between periods, the DiD of the
gap lies in the difference of the two periods' DD ranges; widths add.
\item \textbf{Lemma FD1}: under first-difference LCE the regression of $\Delta T$
identifies $\Delta Q$ and leaves $\Delta K$ unidentified.
\item \textbf{Propositions FD2, FD3 and FD4}: single restrictions R1 (no differential
spillover onto non-veterans), R2 and R4 (Goodman on changes) each point-identify
the DiD, with R4 between R1 and R2; a bounded change $|\Delta w_1|\le\delta$ gives
an interval and a breakdown value; FD4 is the outer set with no restriction on
$\Delta K$.
\item \textbf{Proposition $\Delta4$}: Monotone EI applied to individual changes;
the DiD lies between Goodman-on-changes and the neighbourhood model on changes
under CR on changes.
\item \textbf{Proposition $\Delta6$}: with changing population counts, R1 gives a
short interval rather than a point.
\item \textbf{Propositions C1--C2, Remark C3}: composition change (stayers,
aging-in, exits).
\item \textbf{Proposition CL}: monotone treatment response, no decline absent
service and a selection assumption bound the ATT and are jointly refutable.
\item \textbf{Propositions P2--P3}: party information (a structural zero alone does
not help; selection on party only does).
\end{itemize}
PACKAGE\_DESIGN \S7 lists the change-versions of the modules (M-SIGN-$\Delta$,
M-ORD-$\Delta$, M-LOC-$\Delta$ and the sandwich DiD-S); DiD-S validity is
proved there by applying P-S to changes, and the rest of R11 is open. The DiD
rungs were renamed $\Delta0$--$\Delta8$ to avoid a clash with PHASE2b's
Propositions~D0--D8 (M-14). The PHASE1b and PHASE2d clustering language was
revised to ``one unit observed twice'' and to county effects in the mean (M-5).

% Section 10: numerical evidence
\section{Numerical evidence}\label{sec:evidence}

The results below are \emph{evidence, not proof}. They check that the theory
behaves as stated in designs where the truth is known, and they show where the
proved intervals are informative. A coverage of ``1'' means every replication
covered; with 1000 replications the exact one-sided 95\% lower bound is then
0.997. Figures in parentheses are Monte Carlo standard errors.

\subsection{Verification scripts}

Every phase has a script that checks its closed forms, inclusions and
counterexamples numerically or symbolically. The committed outputs report:

\begin{center}\small
\begin{tabular}{@{}lll@{}}
\toprule
Script & Result & Covers \\
\midrule
\texttt{check\_2x2\_equivalence.py} & agreement to $10^{-16}$ & JKST Prop.~2 in homogeneous form \\
\texttt{verify\_phase2a.py} & 32/32 PASS & PHASE2a (incl.\ [supp], Rems 5.1--5.2, $V=\mu Q$) \\
\texttt{verify\_phase2b.py} & 22/22 PASS & PHASE2b (incl.\ D4(b), D8b, M1 df) \\
\texttt{verify\_phase2c.py} & 29/29 PASS & PHASE2c (C1--C16, tagged [F]/[V]/[I]) \\
\texttt{verify\_phase3.py} & 17/17 PASS & PHASE3 (incl.\ P3.4$'$, P3.5(b2) null direction, P3.5(c)) \\
\texttt{verify\_R3.py} & S1--S7 & Theorem R3 \\
\texttt{verify\_R2.py} & T1--T4, T1b, T1c, T4b & Proposition R2 \\
\texttt{verify\_R10.py} & V9--V13 & Theorem R10 \\
\texttt{verify\_package\_design.py} & [V1]--[V7] & G2, P-S, P-R, P-N \\
\texttt{verify\_phase2d.py}, \texttt{verify\_phase1\_mei\_did.py} & 16, 11 PASS & DiD proposals (not reviewed) \\
\bottomrule
\end{tabular}
\end{center}
Not machine-checked: the measure-theoretic steps of PHASE2a Theorem~1 and PHASE2b
Theorem~1$'$, which rest on cited theorems, and PHASE2a Remark~4.2(2), a sketch.

\subsection{The Phase 3 Monte Carlo (PHASE3 \S8; 1000 replications)}

Unit rates are Beta/Dirichlet draws whose conditional means are exactly linear in
composition given the district, so M1 always holds and M0 holds only without
district shifts. All intervals are nominal 95\%. Cells are coverage / median
width (clipped to $[0,1]$).

\paragraph{Model holds (PHASE3 \S8.1, selected rows).}
\begin{center}\small
\begin{tabular}{@{}lccccc@{}}
\toprule
Design & DD & Plug-in cover / empty & CLR, $\tilde B$ & Projection, $\tilde B$ & Conservative, $B$ \\
\midrule
A: $2\times2$, M1, $10\times93$ & 0.574 & 0.982 / 0.000 & 1 / 0.432 & 1 / 0.554 & 1 / 0.696 \\
A: $2\times2$, M1, $93\times30$ & 0.567 & 0.656 / 0.002 & 1 / 0.414 & 1 / 0.932 & 1 / 0.694 \\
C: pooled, truth near boundary & 0.325 & 0.661 / 0.076 & 0.994 / 0.039 & 0.999 / 0.041 & 1 / 0.087 \\
E: pooled, Slater fails & 0.271 & 0.267 / 0.503 & \textbf{0.944} (0.007) / 0.143 & 0.995 / 0.194 & 1 / 0.249 \\
B: $3\times3$ pooled, loyalty & 0.835 & 1 / 0.000 & 1 / 0.557 & 1 / 0.581 & 1 / 0.675 \\
\bottomrule
\end{tabular}
\end{center}

\paragraph{Many small districts, R2 (PHASE3 \S8.2): M1, $2\times2$, 10 polls per district.}
\begin{center}\small
\begin{tabular}{@{}lcccccc@{}}
\toprule
Districts & $q$ & DD & Plug-in cover / empty & CLR, $\tilde B$ & Projection, $\tilde B$ & Thm P3.8, $\tilde B$ \\
\midrule
30 & 61 & 0.571 & 0.305 / 0.241 & 0.999 / 0.557 & 1 / 0.967 & 1 / 1.000 \\
93 & 187 & 0.566 & 0.042 / 0.463 & 0.995 / 0.467 & 1 / 1.000 & 1 / 1.000 \\
200 & 401 & 0.566 & 0.005 / 0.717 & 0.996 / 0.416 & 1 / 1.000 & 1 / 1.000 \\
\bottomrule
\end{tabular}
\end{center}

\paragraph{Riding shocks (PHASE3 \S8.3): $2\times2$, 20 districts $\times$ 50 polls.}
With no shifts (S0) every interval covers and the pooling test has size 0.090
(SE 0.009) against $\chi^2$ (0.043 with the wild bootstrap). With shifts affine in
the district's mean composition (S1) the pooling test rejects every time, but M0's
intervals still cover. With a common non-affine shock (S2), every M0 interval fails
(CLR for $\tilde B$ covers 0.009 and is empty 96\% of the time; projection for $B$
0.206), while M1 covers in every replication and the pooling test always rejects.

\paragraph{What the Monte Carlo shows (PHASE3 \S8.4).}
\begin{enumerate}[nosep]
\item No proved interval under-covers where its theorem applies: projection at least
0.995, conservative at least 0.990, Theorem~P3.8 always 1.
\item CLR (a simulation finding) covers at least 0.994 wherever Slater holds,
including R2 with 200 districts. At the Slater failure (design E) CLR for
$\tilde B$ covers 0.944 (SE 0.007), at the nominal level, as expected when the
identified set is a point.
\item The plug-in fails as Theorem~P3.4 (and P3.4$'$ for near-ties) predicts: in R2
its coverage falls to 0.305, 0.042 and 0.005.
\item The projection interval is informative only when $q$ is small (critical value
2.8 at $q=3$, 21.2 at $q=401$). In grouped designs with many districts CLR is the
only informative interval, on simulation support only.
\item Theorem~P3.8 is valid but uninformative: clipped width 1.000 in every M1 design.
\item The residual term is a large part of the width: targeting $\tilde B$ rather than
$B$ narrows CLR by 0.06--0.19.
\end{enumerate}

\paragraph{Round-2 rerun of the projection column for $B$.} With the Bhatia--Davis
bound taken from a projection interval on the rung alone (P3.5(c1)), coverage is
unchanged in every design and widths move by at most 0.002 (e.g.\ design~E: 0.251
against 0.249). The rank of $\hat\Sigma$ equals $\dim\beta$ in every replication of
every $2\times2$ design ($\dim\beta$ from 3 to 401); the $3\times3$ fit on all three
columns has rank 12 of 18 by the exact adding-up identity.

\subsection{New Zealand benchmark (ground truth)}

\texttt{benchmarks/run\_ei\_datasets\_benchmark.R}; summary in
\texttt{benchmarks/results/ei\_datasets\_nz\_summary.md}. Units are polling places;
rows the party vote (Labour, National, other); columns the electorate vote; the
truth is the official party $\times$ candidate cross-tabulation, 2002--2020.
\begin{itemize}[nosep]
\item \emph{Per electorate} (492 tables, about 64 polling places each; 4302 cells):
M0's CLR and projection intervals for the actual rate covered every cell, at about
DD width (median width ratio 1.00). Goodman (L6) covered 0.939, with 2.2\% of
non-empty intervals missing.
\item \emph{Pooling each election's 62--69 electorates} and fitting M1 with
electorates as groups: M1 is rejected (empty estimate in all seven elections; M1
test and pooling test $p\approx0$). There the CLR intervals were rejected (ends
crossing) in 21\% of cells and missed the truth in 8\%; the projection interval
never missed but was as wide as DD.
\end{itemize}
Reading: where the model is wrong the tests flag it, and a non-empty CLR interval is
not reassurance.

\subsection{Key simulation tables for M-SIGN (PHASE5a, PACKAGE\_DESIGN)}

\paragraph{How often signs add information (PACKAGE\_DESIGN \S2.2, [V2]; 2000 random
LCE designs with true signs).} The within sign strictly tightens JKST in 83\% of
designs (median width ratio 0.70); the between sign in 89\% (0.67); CR with
unknown sign in 80\% of the 1060 designs where CR holds (0.69).

\paragraph{CR at the edge (PHASE5a \S3.4; $2\times2$, $n=150$, 600 replications,
conservative intervals for $\tilde D$).}
\begin{center}\small
\begin{tabular}{@{}lcccc@{}}
\toprule
Design & JKST & CR slab & CR union (occupied length) & CR naive (plug-in bracket) \\
\midrule
Truth on the CR edge, $ER^*=0.15$ & 1.000 / 1.106 & \textbf{0.992} / 0.359 & 0.992 / 0.357 & \textbf{0.917} / 0.256 \\
Truth on the edge, $ER^*=0.005$ & 1.000 / 1.223 & 1.000 / 0.242 & 1.000 / 0.242 & 1.000 / 0.147 \\
Interior truth & 1.000 / 1.232 & 1.000 / 0.480 & 1.000 / 0.480 & 1.000 / 0.388 \\
\bottomrule
\end{tabular}
\end{center}
The literal intersection of the fitted bracket with a finished JKST interval
covered 304/600 in the edge design (V13), in line with the at-most-one-half limit;
the analytic limit for CR-naive in the three-point design is at most 0.9335
(R10(e)). The package's gated run of the default CLR branch union in the edge
design gave coverage 0.990 with median width 0.230 (200 replications).

\paragraph{A false CR (EI referee's experiment, V12).} Same generator, with
$m_0=.40-sx$, $m_1=.50+s-sx$, so the aggregate regression is unchanged and CR is
false for $s>0$. At 600 precincts the CR slab covered 398, 110 and 0 of 400 at
$s=0,.08,.25$, and was never empty; JKST alone covered 400 every time. Precision
makes a wrong answer more reproducible, not more credible.

\paragraph{Redundancy (PHASE5a \S2.4, T4).} In $3\times2$ with true-sign rows, a
row was implied in 8\% of cases, cut the set but not the target in 28\%, and
tightened the target in 64\%.

% Section 11: package mapping
\section{Package mapping: \texttt{eiPartialID} 0.3.1}\label{sec:package}

Sources: \path{PHASE4_implementation.md}, \path{package/eiPartialID/NEWS.md},
\path{R/rxc_api.R}, \path{R/rxc_infer.R}, \path{R/rxc_sign.R},
\path{R/rxc_modules.R}, and the vignettes \path{eiPartialID_rxc.Rmd} and
\path{eiPartialID_usage.Rmd}. The original $2\times2$ functions
(\texttt{generateBounds}, \texttt{evaluateBounds}) are unchanged; with
\texttt{range\_offset = 1e-5} and \texttt{vcov = "HC1"} the new code reproduces
them on \texttt{eco::census} to $10^{-10}$.

\begin{longtable}{@{}L{0.30\linewidth}L{0.32\linewidth}L{0.30\linewidth}@{}}
\caption{Functions and options, and the results they implement.}\label{tab:package}\\
\toprule Function / option & What it does & Result \\ \midrule\endfirsthead
\toprule Function / option & What it does & Result \\ \midrule\endhead
\texttt{eiBounds()} & bounds and CIs under one rung & PHASE2a Prop.~7 (LP); PHASE2b Thm~2(b); P3.5 \\
\texttt{model = "M0"} & pooled LCE & PHASE2a Props~2, 7 \\
\texttt{"M2"}, \texttt{"M1"}, \texttt{"M1delta"} & district rungs & PHASE2b Thm~2, Props~D3, D5, D6 \\
\texttt{"L4"}, \texttt{"L5"}, \texttt{"L6"} & no contextual effect for one row; constancy & PHASE2b Prop.~D7 \\
\texttt{range = "hull"} (default) & group-wise sample hull & P3.2(c) (conservative) \\
\texttt{range = "trim"} & dropped units enter through DD & PHASE2a Prop.~7$'$ \\
\texttt{range = "extrapolate"} & pure units or a user matrix (warning) & PHASE2a Prop.~5 (extrapolation; (S3) assumed) \\
\texttt{restrictions = list(zeros, order, sign, slope, linear, external)} & extra LP rows & PHASE2b Thm~2(b); PHASE2c M5(b) for \texttt{order} \\
\texttt{restrictions = list(assoc = ...)} & signs on $\delta_B$ (W1, pair) and $\delta_W$ (W1, W2) & PHASE2c M5(a); PHASE5a R3, R10(a) (pooled M0) \\
\texttt{restrictions = list(cr = ...)} & contextual reinforcement: W1/$2\times2$ slab or branch union & PHASE5a R3(c)(d), R10(b)(c) \\
\texttt{union = "pieces"/"hull"} & print every piece, or the hull & PHASE5a R3(c) \\
\texttt{eiRedundancy()} & face test; added first and added last; already-implied table; CR certificate; weight concentration & PHASE5a R2 \\
\texttt{ci = "clr"} (default) & CLR with simulated candidates (\texttt{ndraw = 100}) & Conjecture P3.5(b3): \emph{algorithmic} \\
\texttt{proj\_model\_*}, \texttt{proj\_actual\_*} & projection intervals, $q=\dim\beta$ (0.3.1) & P3.5(b2)(ii); for $B$, P3.5(c1) \\
\texttt{ci = "conservative"} & simultaneous row inflation & P3.5(a) \\
\texttt{fpc = "bhatia-davis"} (default) / \texttt{"popoviciu"} & residual allowance for $B$ & P3.5(c1) (projection), (c2) (CLR); P3.5(a) with $S_1$ \\
\texttt{jkst = TRUE} & JKST $\mathrm{CI}_x$ with revert-to-DD & JKST Prop.~4 \\
\texttt{fit\$tests} & emptiness test; M1 vs M2; pooling test; wild-bootstrap p-values & P3.6(a)--(c); calibration simulation-supported \\
\texttt{ladder()} & all rungs, most to least plausible & PHASE2b Thm~2(c) \\
\texttt{sensitivity(dial = "eta")}, \texttt{breakdown()} & bounded nonlinearity and breakdown & P3 Prop.~S1; PHASE2c M2(d) (scalar path) \\
\texttt{sensitivity(dial = "slack")} & slack on sign beliefs & PHASE2c M2(d); PHASE5a \S5 \\
\texttt{eiModules()} & registry of restriction types with \texttt{meta\$proof\_status} & PACKAGE\_DESIGN \S1.1 \\
\texttt{eiBenchmark()} & coverage on data with known truth & --- \\
\bottomrule
\end{longtable}

\paragraph{Not implemented.} N1 (sign rows on realised tables), N2, M-LOC, the
M-SHARP series bound ($\ThX$), the finite-sample interval of Theorem~P3.8
(\texttt{ci\_method = "finite"} was deferred because it is uninformative), DiD
wrappers, covariates, integer tables.

\subsection{What \texttt{print()} reports as status (0.3.1)}

From 0.3.1 every fit prints a \texttt{status:} line chosen by
\texttt{.rxc\_ci\_status()} in \path{R/rxc_api.R}:
\begin{itemize}[nosep]
\item \textbf{Pooled M0}: ``\texttt{proj\_*}: asymptotic guarantee (P3.5(b2)(ii)),
pooled model'' (or ``\texttt{ci\_*} (conservative): \ldots P3.5(a)'' with
\texttt{ci = "conservative"}); with sign or CR rows the line adds ``with sign/CR
rows PHASE5a R10, pooled M0''. The rank of $\hat V$ is printed as a diagnostic, and
the number of candidates dropped by the rank guard (a safeguard) if any.
\item \textbf{Grouped rungs without sign rows}: ``asymptotic guarantee only with few,
large groups (P3 regime R1); with many small groups a simulation finding''.
\item \textbf{Grouped rungs with sign or CR rows}: ``\texttt{proj\_*} and
\texttt{ci\_*}: no coverage theorem for sign/CR restrictions under grouped rungs
(R16 open)''.
\end{itemize}
Further labels:
\begin{itemize}[nosep]
\item CLR lines carry ``CLR: algorithmic, coverage supported by simulation, not
proved''; for $B$, ``its Bhatia--Davis bound also uses a CLR interval''; with CR,
\texttt{ci\_*} is ``the union of the CLR branch intervals (algorithmic)''.
\item Every sign or CR restriction prints an \texttt{EXPERIMENTAL module} line with
its \texttt{proof\_status}. Grouped M-SIGN fits print that each sign is one pooled
restriction across groups and that the bound is an ``outer set; sharpness under
grouped rungs not yet proved''.
\item A note when the Bhatia--Davis allowance's interval for $\bar m$ lies within
$[0,.05]$ or $[.95,1]$, recommending \texttt{fpc = "popoviciu"} (P3.5(c1)'s
condition is doubtful there).
\item Notes for near-boundary fits (Slater slack below 0.005), few units per group
(median below 30: CLR coverage rests on simulation), unbounded estimates
(hull not full-dimensional), empty estimates (``a sample diagnostic, not a test''),
and joint rejection (naming the combination and whether the rung alone is
rejected).
\item With restrictions, the rung alone is fitted and printed beside the restricted
result (\texttt{base\_*}); with several restrictions it is labelled their joint
effect.
\end{itemize}

\subsection{Version history relevant to the reviews}

\begin{description}[nosep]
\item[0.2.0] R$\times$C and district rungs, CLR default (PHASE1--4).
\item[0.2.1 work (WS5)] \texttt{cluster} removed; projection intervals reported beside
CLR; pooling test; rank errors naming groups; boundary and small-group notes
(Phases~1--3 referee, round~1).
\item[0.3.0] M-SIGN (experimental); Batch~B response: wording on refutation,
every union piece printed, base intervals, weight concentration, projection rank
guard, zero-variance candidate check, wild-bootstrap calibration of both tests,
\texttt{proj\_actual} from a projection interval.
\item[0.3.1] Round~2: \texttt{proj\_actual} per P3.5(c1) (rung alone); $q=\dim\beta$;
per-fit \texttt{status:} line; the $\bar m$ boundary note; pairwise-CR label naming
both sides; joint-effect label for the base block. No reference estimate changed.
\end{description}

% Section 12: status table
\section{Status of every result}\label{sec:status}

Status labels are those of Section~\ref{sec:purpose}. ``Check'' names the
verification script and, where the repository gives one, the check label with its
class [F] (fails on the earlier statement), [V] (verifies a corrected or new
statement) or [I] (illustration). ``---'' means no numerical check (for example, a
purely measure-theoretic step). Section numbers refer to the source document;
``Report'' gives the result's place in this report. PACKAGE\_DESIGN \S8's proof
register (R1--R18) is mapped in Table~\ref{tab:register}.

{\footnotesize
\setlength{\tabcolsep}{3pt}
\begin{longtable}{@{}L{0.27\linewidth}L{0.14\linewidth}L{0.17\linewidth}L{0.30\linewidth}L{0.07\linewidth}@{}}
\caption{Every numbered result, with status and verification.}\label{tab:status}\\
\toprule Result & Source & Status & Check & Report \\ \midrule\endfirsthead
\toprule Result & Source & Status & Check & Report \\ \midrule\endhead
\multicolumn{5}{@{}l}{\textbf{PHASE1 (foundations; its ``to prove'' items are proved in 2a--2b)}}\\
Lemma 1 (identification decomposition) & P1 \S1.3 & Proved as 2a L1 & 2a & \S\ref{sec:rxc} \\
Prop.\ 2-RC (outer set) & P1 \S1.4 & Proved as 2a L3, P4, P5, P7 & 2a & \S\ref{sec:rxc} \\
Closed form $R=2$; joint vs collapsed corollary & P1 \S1.4 & Proved as 2a P6; corollary corrected (2a Cor.~6.1) & 2a & \S\ref{sec:rxc} \\
Theorem 1 (random-set characterisation) & P1 \S1.5 & Proved as 2a T1 & --- & \S\ref{sec:rxc} \\
Nesting (1.13); equality conjecture & P1 \S1.5 & Nesting proved (2a P1); equality: sufficient condition only (2a P3) & 2a (46/400 strict) & \S\ref{sec:rxc} \\
(D0)--(D5) structural facts & P1 \S1.6 & Proved as 2b D0--D5; (D2) count and (D4) conditioning corrected & 2b & \S\ref{sec:districts} \\
Assumptions A1--A10 & P1 \S2 & Assumptions (A1 turnover open) & --- & \S\ref{sec:assumptions} \\
\multicolumn{5}{@{}l}{\textbf{PHASE2a (R$\times$C under LCE)}}\\
Lemma 0 & 2a \S1 & Proved & --- & \ref{res:p2a-l0} \\
Lemma 1 (a)--(c) & 2a \S2 & Proved & \texttt{verify\_phase2a.py} (rank, symbolic); [supp] [F] & \ref{res:p2a-l1} \\
Remarks 1.0--1.2 & 2a \S2 & Proved (1.1, 1.2); remark (1.0) & 2a (five points never identify $Q$) & \S\ref{sec:rxc} \\
Lemma 2; Remark 2.1 & 2a \S3 & Proved & 2a & \ref{res:p2a-l2} \\
Theorem 1 & 2a \S4 & Proved (cited measure theory) & ``$\Rightarrow$'' 227/227 & \ref{res:p2a-t1} \\
Remark 4.1 & 2a \S4 & Proved & --- & \S\ref{sec:rxc} \\
Remark 4.2 (1) / (2) & 2a \S4 & (1) proved (standard); (2) sketch, not used & 2a (200 integer cases) & \S\ref{sec:rxc} \\
Corollary 1 & 2a \S4 & Proved & --- & \S\ref{sec:rxc} \\
Lemma 3 (a)--(c); Remarks 5.1, 5.2 & 2a \S5 & Proved & 2a (singleton, unbounded LP) [F] & \ref{res:p2a-l3} \\
Proposition 1 (nesting) & 2a \S5 & Proved & 2a (400 designs) & \ref{res:p2a-p1} \\
Proposition 2 (sharp for the regression) & 2a \S5 & Proved & P1 script & \ref{res:p2a-p2} \\
Corollary 2; Proposition 3; Remark 3.1 & 2a \S5 & Proved (P3 sufficient only) & 2a & \S\ref{sec:rxc} \\
Proposition 4 (JKST equivalence) & 2a \S6.1 & Proved & 2a symbolic; $10^{-16}$ & \ref{res:p2a-p4} \\
Proposition 5 ($A=\Delta^{R-1}$) & 2a \S6.2 & Proved & 2a (200 LPs; negative $Q$) & \S\ref{sec:rxc} \\
Lemma 4 & 2a \S6.3 & Proved & 2a (via P6) & \S\ref{sec:rxc} \\
Proposition 6 & 2a \S6.4 & Proved & 2a (300 LPs) & \S\ref{sec:rxc} \\
Corollary 6.1; $R\ge3$ characterisation & 2a \S6.4 & Proved; $R\ge3$ open & 2a (strict cases) & \S\ref{sec:rxc} \\
Proposition 7; value-function regularity & 2a \S6.5 & Proved & 2a ($\gamma$, 100 targets; $V=\mu Q$ [F]) & \ref{res:p2a-p7} \\
Proposition 7$'$ & 2a \S6.6 & Proved & 2a & \S\ref{sec:rxc} \\
Proposition 8 & 2a \S7 & Proved & --- & \S\ref{sec:rxc} \\
Proposition 9 & 2a \S7 & Proved & --- & \ref{res:p2a-p9} \\
\multicolumn{5}{@{}l}{\textbf{PHASE2b (districts)}}\\
Proposition D0; joint DD set & 2b \S1 & Proved & 2b (300 tables) & \ref{res:d0} \\
Lemma D1 (a)--(d); Remark D1.1 & 2b \S2 & Proved & 2b (recovery; stacked rank) & \ref{res:d1} \\
Lemma D2 & 2b \S2 & Proved & --- & \S\ref{sec:districts} \\
Theorem 1$'$ & 2b \S3 & Proved (by reduction to 2a T1) & --- & \ref{res:t1prime} \\
Theorem 2 (a)--(c) & 2b \S3 & Proved & 2b (ladder nesting) & \ref{res:t2} \\
Proposition D3 (+ pure vertices) & 2b \S4 & Proved & 2b (126, 195 LPs) & \ref{res:d3} \\
Proposition D4 (a), (b); Remark D4.1 & 2b \S4 & Proved & 2b D4(b) [F] & \ref{res:d4} \\
Propositions D5, D6, D7; Remark D7.1 & 2b \S4 & Proved & 2b & \S\ref{sec:districts} \\
Proposition D8 (a), (b) & 2b \S5 & Proved & 2b D8b [F] & \ref{res:d8} \\
Remarks 6.1--6.3 & 2b \S5 & Remarks & --- & \S\ref{sec:districts} \\
Strict-gain characterisation $R\ge3$; M1$'$ & 2b \S8 & Open & --- & \S\ref{sec:open} \\
\multicolumn{5}{@{}l}{\textbf{PHASE2c (monotone EI)}}\\
Lemma M1 (a)--(d) & 2c \S1 & Proved & C1 & \ref{res:m1} \\
Proposition M2 (a)--(d) & 2c \S2 & Proved (continuous tables) & C7 [F], C7b [V] & \ref{res:m2} \\
Remark M2.1 ($2\times2$ lattice proposition) & 2c \S2 & Proved ($2\times2$ margins only); otherwise open (R15) & C7 [F], C7a [V] & \ref{res:m21} \\
Proposition M3 (a)--(d); Reading & 2c \S2 & Proved & C2; C8 [F]; C8b [F]; C15 [F] & \ref{res:m3} \\
Proposition M4 & 2c \S2 & Proved & C3; C3b [F] & \S\ref{sec:monotone} \\
Proposition M5 (a), (b); (c) & 2c \S3 & Proved; (c) M0 only & C4 [V] (rebuilt); C9 [F] & \ref{res:m5} \\
Remark M5.1 & 2c \S3 & Proved (bound under P3-S1a) & --- & \S\ref{sec:monotone} \\
Lemma N2.1 & 2c \S4 & Proved & --- & \ref{res:n22} \\
Proposition N2.2 (a)--(e) & 2c \S4 & Proved (finite support) & C5, C5(c$'$) [V], C10 [F], C10b [F] & \ref{res:n22} \\
Conjecture N2.2$'$ & 2c \S4 & Conjecture (sufficiency remark) & C10b & \ref{res:n22prime} \\
Proposition N2.3 (a)--(d) & 2c \S4 & Proved (structural identity); not observable & C6; C11 [F]; C13 [I]; C16 [F] & \ref{res:n23} \\
Lemma N2.4 & 2c \S4 & Proved & C12 [F] & \ref{res:n24} \\
Degree selection (\S6) & 2c \S6 & Prespecify (counterexample proved by simulation) & C12b [F] & \S\ref{sec:monotone} \\
\multicolumn{5}{@{}l}{\textbf{PHASE3 (inference)}}\\
Lemma P3.1 & P3 \S2 & Proved (fixed system) & \texttt{verify\_phase3.py} & \ref{res:p31} \\
Proposition P3.2 (a)--(c) & P3 \S3 & Proved (Slater for (a)) & 3 & \ref{res:p32} \\
Theorem P3.4 (i)--(iii); Remark P3.4.1 & P3 \S4 & Cited structure (fixed system) & 3 (tied constraints; null direction) & \ref{res:p34} \\
Proposition P3.4$'$ & P3 \S4 & Proved & ``P3.4$'$ [F]'' & \ref{res:p34prime} \\
Theorem P3.5(a) & P3 \S5 & Proved & MC & \ref{res:p35} \\
P3.5(b1) & P3 \S5 & Cited (CLR Thm 4, Condition P) & --- & \ref{res:p35} \\
P3.5(b2)(i), (ii) & P3 \S5 & Proved & weak duality (500); null direction [F] & \ref{res:p35} \\
Conjecture P3.5(b3) & P3 \S5 & Conjecture / algorithm; reduction proved & MC (simulation only) & \ref{res:p35} \\
P3.5(c1) / (c2) & P3 \S5 & (c1) proved; (c2) inherits (b3) & ``P3.5(c)'' [V] & \ref{res:p35} \\
Proposition P3.6 (a), (c); (b) & P3 \S5 & Proved; (b) fixed dimension & 2b M1 df & \ref{res:p36} \\
Wild-bootstrap calibration & P3 \S5 & Simulation-supported & \texttt{benchmarks/*\_size.R} & \S\ref{sec:inference} \\
Theorem P3.8 (a)--(d) & P3 \S5.1 & Proved & 3; MC & \ref{res:p38} \\
Proposition P3.7 (a), (b) & P3 \S6 & Proved & --- & \S\ref{sec:inference} \\
Proposition S1; breakdown value & P3 \S7 & Proved (outer bound) & 3 ($\eta=0.02$) & \ref{res:s1} \\
M1/M2 in R3; CLR in R2 & P3 \S0.5 & Conjecture; simulation finding & MC & Table~\ref{tab:regimes} \\
\multicolumn{5}{@{}l}{\textbf{PHASE5a (M-SIGN; pooled M0)}}\\
Theorem R3 (a)--(e); Remarks R3.1, R3.2 & 5a \S1 & Proved (pooled M0) & S1--S5; S6 [V]; S7 [F] & \ref{res:r3} \\
Proposition R2 (a)--(e) & 5a \S2 & Proved & T1--T4; T1b; T1c [F]; T4b [F] & \ref{res:r2} \\
Theorem R10 (a)--(d) & 5a \S3 & By version (Table~\ref{tab:r10}); pooled M0 & V11 [V]; V12 [I]; V13 [V] & \ref{res:r10} \\
Theorem R10(e) & 5a \S3 & Proved (analytic) & V9 [V]; V10 [F] & \ref{res:r10} \\
\multicolumn{5}{@{}l}{\textbf{PACKAGE\_DESIGN}}\\
G2 (JKST + signs) & PD \S2.2 & Proved & [V1], [V2] & \ref{res:g2} \\
P-S (sandwich) & PD \S5.1 & Proved & [V4] & \ref{res:ps} \\
P-R (a), (b), (d); (c) & PD \S5.2 & Proved; (c) checked & [V5], [V7] & \ref{res:pr} \\
P-N (a), (b); Conjecture C1 & PD \S5.3 & Proved (M0); C1 conjecture & [V3] & \ref{res:pn} \\
DiD-S & PD \S7 & Proved (validity; not reviewed) & --- & \S\ref{sec:did} \\
\multicolumn{5}{@{}l}{\textbf{PHASE1b and PHASE2d (proposals; not reviewed)}}\\
1b \S\S A1--A5, B1--B5 (proposed objects, rungs, DiD restrictions) & 1b & Proposal; superseded by 2c, 2d & \texttt{verify\_phase1\_mei\_did.py} & \S\ref{sec:did} \\
Lemma G; Props $\Delta0$, FD1--FD4, $\Delta4$, $\Delta6$, C1, C2, CL, P2, P3; Remark C3 & 2d & Proposal (proofs written, not reviewed) & \texttt{verify\_phase2d.py} (16 PASS) & \S\ref{sec:did} \\
\bottomrule
\end{longtable}
}

{\footnotesize
\begin{longtable}{@{}L{0.06\linewidth}L{0.44\linewidth}L{0.22\linewidth}L{0.20\linewidth}@{}}
\caption{The proof register of PACKAGE\_DESIGN \S8 (R1--R18).}\label{tab:register}\\
\toprule ID & Claim & Status & Where in this report \\ \midrule\endfirsthead
\toprule ID & Claim & Status & Where in this report \\ \midrule\endhead
R1 & C1: $x^*\in$ hull ($2\times2$), so weighted-local $\Rightarrow\delta_W\ge0$ & Conjecture (checked [V3]) & Result~\ref{res:pn} \\
R2 & Face criterion for tightening; implied $\Leftrightarrow$ tightens nothing; G2 closed form & Proved (5a \S2; \texttt{verify\_R2.py}) & Result~\ref{res:r2} \\
R3 & Pooled LCE $\cap$ closed convex restriction is sharp for the model-implied target; CR union and slab & Proved for pooled M0 (grouped: R16) & Result~\ref{res:r3} \\
R4 & Refutability classification for $R\ge3$ & Open & \S\ref{sec:open} \\
R5 & P-R(c): any spanning tree reaches all directions & Proved (star); checked (chain) & Result~\ref{res:pr} \\
R6 & When joint column restrictions in $2\times C$ are strictly tighter & Checked [V6] & \S\ref{sec:open} \\
R7 & Global N2 LP with estimated $\tau$ & Open (Lemma N2.4 proved; N2.2$'$ conjecture) & Results~\ref{res:n22prime}, \ref{res:n24} \\
R8 & B1(M), $\kappa$, CAR-$\delta$ valid and monotone in their dials & Partly (NEW\_AVENUES) & \S\ref{sec:open} \\
R9 & M-SHARP series bound: consistency; never wider than $\ThO$ & Open & \S\ref{sec:open} \\
R10 & Inference with added rows, unions, near-emptiness & By version: (a), (b2) proved; (b1) cited; (b3) algorithm; (e) proved & Result~\ref{res:r10} \\
R11 & DiD: DiD-S sharp within FD; FD4 $\cap\,\Delta4$; levels $\cap$ sign-$\Delta$; union inference & DiD-S validity proved; rest open & \S\ref{sec:did} \\
R12 & P-S (sandwich) & Proved [V4] & Result~\ref{res:ps} \\
R13 & G2 table & Proved [V1, V2] & Result~\ref{res:g2} \\
R14 & Single-index dimension $C-1$ & Conjecture & \S\ref{sec:open} \\
R15 & Integer unit tables (N1, $\Delta4$) & Open ($2\times2$ margins-only proved) & Result~\ref{res:m21} \\
R16 & Grouped rungs: R3$'$, R10$'$ & Open & Results~\ref{res:m5}, \ref{res:r3} \\
R17 & P3.4$'$ triangular-array limit & Proved & Result~\ref{res:p34prime} \\
R18 & P3.5(c): (c1) valid; (c2) with CLR & (c1) proved; (c2) inherits (b3) & Result~\ref{res:p35} \\
\bottomrule
\end{longtable}
}

% Section 13: open items and next steps
\section{What is open, and what comes next}\label{sec:open}

Sources: PHASE3 \S10, PACKAGE\_DESIGN \S\S6, 8, \path{NEXT_BUILD.md},
\path{EXTERNAL_REVIEW_PLAN.md}, \path{MSIGN_REVIEW_LOG.md}.

\subsection{Open theory}

\paragraph{Conjectures.}
\begin{itemize}[nosep]
\item \textbf{Conjecture P3.5(b3)}: validity of the CLR interval with simulated
candidates (the package default). A reduction to the cited full-candidate CLR
theorem is proved; recovery of rarely optimal vertices, the downward bias of $k_p$
over a subset, growing $|\Lambda|$, and uniform validity are open. Proving it
would make the default a theorem.
\item \textbf{Conjecture N2.2$'$}: the sharp N2 set when size is informative.
\item \textbf{Conjecture C1} (R1): $x^*$ lies in the composition hull ($2\times2$).
\item \textbf{R14}: single-index contextual effects have dimension $C-1$.
\item M1/M2 in regime R3 (high-dimensional Gaussian approximation).
\end{itemize}

\paragraph{Open items from the review rounds.}
\begin{itemize}[nosep]
\item \textbf{R15}: integer unit tables beyond the $2\times2$ margins-only case.
\item \textbf{R16}: sharpness and inference for M-SIGN under district rungs. Until
then grouped M-SIGN is experimental and its intervals have no coverage theorem.
\item Uniform (in the design sequence) limit theory for the plug-in endpoint
beyond P3.4$'$'s localisation condition.
\item Boundary cases where Slater fails: rate and limit of the endpoint estimator;
uniform validity near point identification for CLR (Stoye-type).
\item The rate of the hull error, and inference for the population-hull endpoint.
\item Informative finite-sample inference with many small groups (R2): Theorem P3.8
is valid but uninformative.
\item Calibration of the pooling and M1 tests for a growing number of restrictions
(the wild bootstrap is simulation-supported).
\item Inference for the superpopulation target $B^*$; dependence within districts
beyond the district effect; joint relaxation for population turnover (A1);
coverage of the plug-in under misspecification.
\item The characterisation of strict joint-vs-collapsed gains, and of strict M1
gains, for $R\ge3$.
\end{itemize}

\paragraph{Other proof-register items.} R4 (refutability for $R\ge3$), R5 (general
spanning trees), R6 (strict joint-column gains), R7 (global N2 with estimated
$\tau$), R8 (B1(M), $\kappa$, CAR-$\delta$), R9 (M-SHARP consistency), R11 (DiD).

\subsection{Unreviewed material}

PHASE1b, PHASE2d (review batch~C), PACKAGE\_DESIGN beyond the results used here
and NEW\_AVENUES (batch~D), and a dedicated code review of the M-SIGN and WS5
package changes (batch~E). PHASE2c and PHASE5a were reviewed in Batch~B (two
rounds). The original Phases~1--3 have had two rounds from their referee.

\subsection{Build order (\texttt{NEXT\_BUILD.md})}

\begin{enumerate}[nosep]
\item Release hygiene: a pull request to \texttt{master} and a version tag.
\item M-SHARP, the series-based sharper bound (no new assumption; R9).
\item M-LOC in $2\times2$ and $2\times C$ (needs C1), $2\times C$ joint column
signs, $R\times2$ pairwise ordering on spanning trees.
\item DiD wrappers (gated on review batch~C), L1$'$ quadratic contextual effects,
the non-negative-correlation option for the unit-level term.
\item Later: R16, R15, per-group W2, the drop-one display and the L5/L6 sign check,
the N2 global LP, data-dependent modules (E4/E5, B1(M), $\kappa$, contests,
CAR-$\delta$, dispersion), covariate adjustment, informative R2 inference.
\end{enumerate}


\bibliographystyle{plainnat}
\bibliography{references}

\clearpage
\appendix
% Appendix: consolidated revision table
\section{Consolidated revision table}\label{app:revisions}

This appendix maps every referee point to the result it changed and the change
made. It covers both rounds and all three referees. IDs are those of the
repository's ledgers, so each row can be traced:
\begin{itemize}[nosep]
\item \textbf{R-01 to R-37}: the Phases~1--3 referee, round~1 (two reports:
Phases~1--2a, then 1--3). Ledger: \texttt{REVIEW\_RESPONSE\_PLAN.md} \S1;
letter: \texttt{REVIEW\_RESPONSE.md}.
\item \textbf{P-B1 to P-B43}: the Batch~B partial-identification referee, round~1.
\textbf{P-E1 to P-E33}: the Batch~B ecological-inference referee, round~1.
Ledger: \texttt{reviews/B\_PI\_RESPONSE\_PLAN.md} \S2; responses:
\texttt{reviews/B\_*\_response.md}.
\item \textbf{A2-1 to A2-7}: the Phases~1--3 referee, round~2, numbered by the
report's sections~1--7 (\texttt{reviews/A\_phases1to3\_report\_round2.md};
response \texttt{REVIEW\_RESPONSE\_ROUND2.md}).
\item \textbf{I.1 to I.13}: the partial-identification referee, round~2
(\texttt{reviews/ROUND2\_THEORY\_PLAN.md} Part~I;
\texttt{reviews/B\_partial\_identification\_response\_round2.md}).
\item \textbf{E2-1 to E2-5}: the ecological-inference referee, round~2, numbered
by the response's items (corrections 1--4 and the final paragraph;
\texttt{reviews/B\_ecological\_inference\_response\_round2.md}).
\item \textbf{M-1 to M-17}: the internal consistency log of the merged M-SIGN
material (\texttt{MSIGN\_REVIEW\_LOG.md}), listed for completeness; it is not
an external referee.
\end{itemize}
Types: \textbf{E} a claim wrong as stated; \textbf{S} scope; \textbf{P} proof
detail; \textbf{X} exposition; \textbf{C} code. ``Check'' names the numerical
check added, with its [F]/[V]/[I] class where the repository gives one.

{\footnotesize
\setlength{\tabcolsep}{3pt}
\begin{longtable}{@{}L{0.06\linewidth}L{0.03\linewidth}L{0.30\linewidth}L{0.17\linewidth}L{0.30\linewidth}L{0.08\linewidth}@{}}
\toprule
ID & T & Referee point & Result changed & Change made & Check \\
\midrule\endfirsthead
\toprule ID & T & Referee point & Result changed & Change made & Check \\ \midrule\endhead
\multicolumn{6}{@{}l}{\textbf{Phases 1--3 referee, round 1, first report (Phases 1, 2a)}}\\
R-01 & X & $\sum_cH_c=J$ argued from an unassumed open set & P1 (1.5); 2a L1(b) & Derived from A5 (2a Lemma~1(b)). & --- \\
R-02 & S & A5 must be a matrix rank condition; three points not enough for $R=3$ & P1 A5; 2a L1, Rem.~1.1 & A5 in matrix form; Remark~1.1 (six points off a conic). Package rank error. & 2a \\
R-03 & X & Derivative wording in Lemma~2 & 2a Lemma~2 & Coefficient in the free coordinates of Prop.~7. & --- \\
R-04 & E & (S3) does not follow from (S1)--(S2) for every $A$ & 2a (S3), Lemma~3(c), Rem.~5.1 & (S3) automatic only when $A\subseteq\conv(\supp_AP_X)$; singleton counterexample. & 2a (singleton) \\
R-05 & S & $P(X\in A)=1$ used silently & 2a Props 4--7, Cor.~1 & Hypothesis added; new Prop.~7$'$ (units outside $A$; \texttt{range="trim"}). & 2a \\
R-06 & E & Value function claimed jointly concave & 2a Prop.~7 (regularity) & Concave in $Q$, convex in $\gamma$, neither jointly; example $V=\mu Q$. & 2a [F] \\
R-07 & E & ``Nonempty relative interior'' too weak for compactness & 2a Thm~1, Lemma~3(b), Rem.~5.2 & $\operatorname{affdim}(A)=R-1$; A5 implies it; $\ThO$ a polyhedron until bounded. & 2a (unbounded LP) \\
R-08 & P & Sufficiency proof of Thm~1 too brief & 2a Thm~1 & Facts F1--F7; ``all $u$ at once'' paragraph. & --- \\
R-09 & P & Citations from memory & 2a \S0.4 & Citation table; CLR ``Thm~2'' corrected to Thm~4. & --- \\
R-10 & S & $\ThO$ called sharp & 2a Props 2--4, 7 & ``Sharp for the regression function'' only; ``outer bounds'' in package. & --- \\
R-11 & S & Hulls conflated & P1 \S1.7 & Three hulls named. & --- \\
R-12 & S & Sharpness vs integer counts & 2a Rem.~4.2 & Sharp only for continuous latent rates. & 2a (integer vertices) \\
R-13 & X & Prop.~5 nonemptiness & 2a Prop.~5 & Empty iff some $Q_{c,rs}<0$. & 2a \\
R-14 & X & Cor.~6.1(c) vague & 2a Cor.~6.1(c) & Exact sign condition (a negative pure-poll rate). & 2a \\
R-15 & S & Covariates not covered by the proofs & P1 \S1.5--1.6, A3 & $Z$ removed from $\mathcal G$; covariates a planned extension. & --- \\
R-16 & X & Prop.~9 domain; motivation & 2a Prop.~9, Rem.~1.0 & Domain stated; moved forward as Remark~1.0. & --- \\
R-17 & X & What each set needs from the data & 2a \S4.3 & New table. & --- \\
\multicolumn{6}{@{}l}{\textbf{Phases 1--3 referee, round 1, revised report (adds 2b and 3)}}\\
R-18 & X & D0: per-cell vs joint set & 2b Prop.~D0 & Joint set $\mathcal B^*_{DD}$ is not a product. & 2b \\
R-19 & S & ``$R$ distinct compositions'' insufficient under M1 & 2b Lemma~D1(d); P1 A9 & Stacked rank condition (A5-M1); affinely independent points; Rem.~D1.1. & 2b \\
R-20 & P & Thm~1$'$ hypotheses implicit & 2b Thm~1$'$ & H1--H4 listed. & --- \\
R-21 & X & M1(0) status & 2b Thm~2(c), Prop.~D6(d) & Labelled a sensitivity benchmark. & --- \\
R-22 & S & D3 divides by $t$ and $1-t$ & 2b Prop.~D3 & $0<l_d\le u_d<1$ visible; pure-composition vertices treated. & 2b (195 inst.) \\
R-23 & E & D4(b) conditioned on $X$ only & 2b Prop.~D4(b), Rem.~D4.1; P1 (D4) & Condition on $(X,N)$; $D\perp X$ not enough. & 2b [F] \\
R-24 & X & DiD remark deserves a proposition & 2b Prop.~D8 & Promoted; exact moment condition. & --- \\
R-25 & S & Lemma P3.1 scope & P3 Lemma P3.1 & Fixed constraint system only. & --- \\
R-26 & S/P & P3.2 hulls and Slater & P3 Prop.~P3.2 & Non-nested hulls; Slater as interiority; error decomposition. & --- \\
R-27 & E & P3.4(ii) negative mean claimed too widely & P3 Thm~P3.4(ii), Rem.~P3.4.1 & Needs active gradients differing in a direction where the estimator varies; $C-1$ columns. & P3 [F] \\
R-28 & E & P3.5(a) used independence not granted; clustering & P3-S1a, \S0.4; P1 A6 & Independence given design and district effects; no clustering; \texttt{cluster} removed; pooling test P3.6(c). & --- \\
R-29 & E & CLR with simulated candidates claimed as theorem & P3.5(b) & Split: (b1) cited, (b2) proved, (b3) conjecture; package labels CLR algorithmic. & P3 (weak duality) \\
R-30 & E & Overidentification count & P1 (D2); P3.6(b) & $(D^\star-1)(C-1)R(R-1)/2$. & 2b \\
R-31 & E & Asymptotic scope by regime & P3 \S0.5, P3.7, new P3.8 & Regime table; finite-sample Thm~P3.8; CLR in R2 a simulation finding. & P3 \\
R-32 & S & Coverage vs hull & P3.2, P3.5 & Coverage of the true parameter does not depend on the hull. & --- \\
R-33 & S/X & Three targets; ``both linear in $K$'' & P3 \S0.2; P1 \S1.7 & $B$, $\tilde B$, $B^*$; only $\tilde B$ affine in $K$. & --- \\
R-34 & S & S1 sharpness & P3 Prop.~S1 & Labelled an outer bound. & --- \\
R-35 & E & Turnover said to be treated & P1 A1 & Open extension. & --- \\
R-36 & S & A6, A10 alignment & P1 A6, A10 & Aligned with P3-S1a and \S0.5. & --- \\
R-37 & X & Monte Carlo too small & P3 \S8 & 1000 reps, MC SEs; designs E (Slater fails), R2, S0--S2 (shocks). & MC \\
\multicolumn{6}{@{}l}{\textbf{Batch B partial-identification referee, round 1}}\\
P-B1 & S & M2 sharp only for continuous tables & 2c M2, Rem.~M2.1 & Relabelled; 2$\times$2 lattice proposition; integer N1 = R15 (U1). & C7 [F] \\
P-B2 & S & Scalar breakdown needs a path & 2c M2(d) & Componentwise nesting; frontier; scalar path. & C7b [V] \\
P-B3 & E & M3(d) fails at $\gamma=1$ & 2c M3(d) & Cases $\gamma<1$, $\gamma=1$, $D_{ER}=0$. & C8 [F] \\
P-B4 & S & M3(c) witness removable & 2c M3 hypotheses & No unit restrictions, no other sign rows. & C8b [F] \\
P-B5 & E & ``CR never refuted'' & 2c M3, Reading; 5a; guide; package & Names its base system. & C7 \\
P-B6 & --- & M4 established & 2c M4 & No action. & --- \\
P-B7 & E & M5 wrong for district rungs & 2c M5; PD P-N, \S4.2 & Stacked form (a); per-district ordering (b); (c) for M0 only; R16 opened. & C9 [F] \\
P-B8 & E & $p^{-1/2}$ claim for $\delta-\delta^*$ & 2c Rem.~M5.1 & Root sum of squared coefficients, under P3-S1a. & --- \\
P-B9 & S & N2 support ambiguous & 2c N2.2(a)(b) & Finite empirical design law; grids are relaxations. & C10 [F] \\
P-B10 & E & Proposed sharp N2 not sharp & 2c N2.2(a), Rem.~N2.2.1 & $V$-conditional statement as Conjecture N2.2$'$. & C10b [F] \\
P-B11 & E & Derivative not identified on finite support & 2c N2.3 & Structural identity only; N2-local rows dropped (U4(a)). & C11 [F] \\
P-B12 & X & Smooth witness for N2.3(d) & 2c N2.3(d) & Smooth example; reverse direction. & C13 [I] \\
P-B13 & E & Binning argument wrong & 2c Lemma~N2.4 & Person-weighted bin rates (new lemma). & C12 [F] \\
P-B14 & E & CV-selected degree then fixed-degree inference & 2c \S6 & Prespecified degree; $q$ vs $J_X$. & C12b [F] \\
P-B15 & X/S & R3(b) ``LP''; full law vs regression & 5a R3(b), Rem.~R3.1 & Finite LP for polyhedral $C$; example. & S6 [V] \\
P-B16 & E & R3(e) ``never $C$ alone''; ``unrefutable'' & 5a R3(e), Rem.~R3.2 & Conjunction rejected; LCE can exclude CR. & S7 [F] \\
P-B17 & X & Certificate only sufficient & 5a \S1.4 & Square example. & --- \\
P-B18 & E & Face LP omitted target intercept & 5a R2 step~4 & $h^\top k\ge U-\varphi_0$ (code already correct). & T1c [F] \\
P-B19 & X & Positive dual does not certify & 5a R2(b)(c) & Duplicate-row example; witness need not be a reported target. & --- \\
P-B20 & S & R2(d) qualifications & 5a R2(d) & $\gamma<1$; non-constant target. & --- \\
P-B21 & E & ``Redundant given the others'' & 5a R2(e) & Zero incremental gain for this target. & T4b [F] \\
P-B22 & S & Width-gain inference & 5a \S2.5 & Joint treatment required. & --- \\
P-B23 & S & R10 conditions per version & 5a R10 status table & Version table; Condition P stated; zero-variance candidates. & --- \\
P-B24 & P & Residual for contrasts & 5a R10(a) & Contrast envelope; per-cell sum valid. & --- \\
P-B25 & S/P & R10(b)--(d) qualifications & 5a R10(b)--(d) & One simultaneous event; point coverage; triangular arrays. & --- \\
P-B26 & P & R10(e) only by Monte Carlo & 5a R10(e) & Analytic proof (referee's, credited): $\le0.9335$. & V9 [V] \\
P-B27 & E & ``A fortiori'' transfer false & 5a \S3.5 item~4 & Direct $\le\tfrac12$ argument. & V10 [F] \\
P-B28 & X & Lemma M1 hypotheses; W2 sums & 2c Lemma~M1, Def.~1 & Per-part hypotheses; W2 is one sum. & --- \\
P-B29 & X/S & Ladder wording; simplex chart & 2c \S5, N2.2(e) & Same-direction ``both signs''; chart. & C5(c$'$) [V] \\
P-B30 & X/C & Population vs sample emptiness & 5a \S1.6; package & ``Empty at the point estimate''. & pkg test \\
P-B31 & S & EGGHM mixed-unit condition & 2c \S0, M3 & Added. & --- \\
P-B32 & S & EGGHM sharpness is continuous & 2c M2, M3 & Cited with P-B1. & --- \\
P-B33 & X & EGGHM theorem mapping & 1b & Fixed. & --- \\
P-B34 & S & EGGHM local results & 2c N2.3 & With P-B11. & --- \\
P-B35 & C & C4 construction invalid & \texttt{verify\_phase2c.py} & C4 rebuilt. & C4 [V] \\
P-B36 & C/X & S2 labels & \texttt{verify\_R3.py} & Empty / point / positive-width. & --- \\
P-B37 & C & T4 solver status & \texttt{verify\_R2.py} & Status printed. & --- \\
P-B38 & C/X & Union length double-counted & 5a \S3.4 & Occupied length. & V11 [V] \\
P-B39 & E & [supp]: support of the restricted measure & 2a \S0.1, L1, L3(c), Rems 5.1--5.2; P1 A5; 2b D1 & $\supp_AP_X$ throughout. & 2a [supp] [F] \\
P-B40 & E & Projection under singular $\Sigma$ & P3.5(b2)(ii) & Positive definite $\Sigma$ or exact reduced parameterisation; package guard. & P3 [F] \\
P-B41 & P & Residual CLT gap & P3.5(a) step~5 & Chebyshev/Lindeberg subsequence argument (referee's). & P3 [F] \\
P-B42 & X & 1b mapping; pure units & 1b & Fixed. & --- \\
P-B43 & X & Process: internal review missed items & 5a \S4; review plan & Pre-screen must cover scope documents. & --- \\
\multicolumn{6}{@{}l}{\textbf{Batch B ecological-inference referee, round 1}}\\
P-E1 & E & W2 is one pooled sum; W1 and W2 can disagree & 2c Def.~1; guide; vignette & Wording; two-unit example. & C14 [F] \\
P-E2 & E & W1 predictor & vignette; guide & Reworded. & --- \\
P-E3 & X & $\delta_B$ weights & 2c Def.~1 remarks & Weighting sentence; weight diagnostics. & --- \\
P-E4 & X & Sorting; causal reading & 2c remarks; guide & Sentence added. & --- \\
P-E5 & E & Goodman bias not signed & 2c Def.~1 remarks; 1b & $g-B=M^{-1}\Delta$. & --- \\
P-E6 & S/C & Pairwise CR mixes pairings & 5a \S1.4; package & Both terms printed. & pkg test \\
P-E7 & X/C & Pre-specification; pooled restriction & guide; vignette; package & Checklist; base interval printed (U7). & --- \\
P-E8 & X & EGGHM extensions & 2c remarks & W2, pair, CR marked extensions. & --- \\
P-E9 & X & What a sign does not mean & guide & List added. & --- \\
P-E10 & X & ``Exact'' & guide; vignette & Dropped. & --- \\
P-E11 & S & Integer endpoint failure & 2c Rem.~M2.1 & Second example. & C7 \\
P-E12 & E & Equal-rates witness; three events & 2c M3 Reading; 5a & R$\times$C witness; three distinct events. & C8b \\
P-E13 & E & M3(d) at $\gamma=1$ & 2c M3(d) & As P-B3. & C8 \\
P-E14 & X & Slab orientation & vignette & Hull notation. & --- \\
P-E15 & C & Print hides pieces & package & Every piece printed. & pkg test \\
P-E16 & C/X & Estimated set vs CI & package; vignette & Separate lines; interpretation. & --- \\
P-E17 & X/C & CLR CR status; numbers per method & guide; 5a & Method named with each number (U6). & --- \\
P-E18 & X/C & Descriptive values & package; guide & Labels; design-specific miss rates. & --- \\
P-E19 & C/X & Verdict labels; unions & 5a \S2.5; package & Base and order named; hull span. & pkg test \\
P-E20 & C & Report table; L5/L6 signs & package & Base column; L5/L6 check (later). & --- \\
P-E21 & C/X & Union lengths & 5a \S3.4 & As P-B38. & V11 \\
P-E22 & X/C & False-belief experiment & 5a \S3.4; guide & V12 added; slack dial exists (5a \S5 corrected). & V12 [I] \\
P-E23 & E/S & N2 points & 2c \S\S4, 6 & As P-B9 to P-B14. & C10--C12 \\
P-E24 & X & Turnout example & guide; vignette & Added. & --- \\
P-E25 & C/X & Weight concentration & 2c Rem.~M5.1; package & Concentration in \texttt{eiRedundancy()}. & --- \\
P-E26 & X & King's EI comparison & guide & Paragraph added. & --- \\
P-E27 & X & Adding-up; complementary columns & 2c \S5; guide & Sentences added. & --- \\
P-E28 & X & One status key; stale slack line & 5a \S5; 2c header & Fixed. & --- \\
P-E29 & X & 1b CR as two LPs superseded & 1b & Marked superseded. & --- \\
P-E30 & X & 1b W2 wording & 1b & Fixed. & --- \\
P-E31 & X & Guide \S2: requirements per rung & guide & Fixed. & --- \\
P-E32 & X & Aggregate reversal not evidence & 1b & Fixed. & --- \\
P-E33 & X & Belief selection & 5a R10(d) & Sentence added. & --- \\
\multicolumn{6}{@{}l}{\textbf{Phases 1--3 referee, round 2}}\\
A2-1 & --- & Overall assessment & --- & No action asked. & --- \\
A2-2 & E/S & P3.4 with changing sample hulls ($f_n$ counterexample) & P3 Thm~P3.4; new Prop.~P3.4$'$ (R17) & Fixed constraint system; triangular-array limit under localisation. & P3.4$'$ [F] \\
A2-3 & S & Projection interval for $B$ not covered & new P3.5(c) (R18); P3 \S8 & (c1) valid, (c2) algorithmic; MC column rerun; package 0.3.1 \texttt{proj\_actual} per (c1). & P3.5(c) [V] \\
A2-4 & E & D8 ``iff'' false & 2b Prop.~D8 & (a) cancellation; (b) given the identified set; no causal reading. & D8b [F] \\
A2-5 & S & Independence is substantive & P3 \S0.4, after P3.6(c) & Two arguments separated; pooling test does not test it. & --- \\
A2-6 & --- & P3.8 a successful addition & P3.8 & Emphasis kept; R2 open. & --- \\
A2-7 & --- & CLR treatment appropriate & P3 \S10 & ``Uniformly'' added to the (b3) question. & --- \\
\multicolumn{6}{@{}l}{\textbf{Batch B partial-identification referee, round 2}}\\
I.1 & E & Equal-rates witness gives the CR union, not both branches & 2c M3 Reading & Referee's wording; three-row example. & C15 [F] \\
I.2 & S & R10 scope said ``or M1/M2 in R1'' & 5a Thm~R10 & Pooled M0 only. & --- \\
I.3 & S & Numerical rank cutoff is not (ii-b) & P3.5(b2) remark; package & Safeguard only; $q=\dim\beta$ in 0.3.1. & --- \\
I.4 & E & N2.3(d) premise & 2c N2.3(d) & Monotone on a neighbourhood of $x$. & C16 [F] \\
I.5 & X & Lattice step needs both totals & 2c Rem.~M2.1 & $1/N_1+1/N_0$; worked example. & --- \\
I.6 & S & Model-implied bin mean & 2c Lemma~N2.4, \S6 & $\bar\tau_b=E_b\tau$; specification caveat. & --- \\
I.7 & X & ``Every computation is a finite LP'' & 5a \S1.1 & Narrowed to $\ThO\cap C$. & --- \\
I.8 & E & M4 example backwards & 2c M4 & Same sign on complementary columns forces $D=ER$. & C3b \\
I.9 & X & Branch vs image & 2c M3(d) & Distinguished. & --- \\
I.10 & X & Set vs frontier & 2c M2(d) & Boundary is the frontier. & C7b \\
I.11 & --- & N2.2$'$ sufficiency argument & 2c Rem.~N2.2.1 & Credited remark; conjecture stays open. & --- \\
I.12 & X & ``Every check fails on the old statement'' overstated & 2c, 5a headers; scripts & [F]/[V]/[I] classes. & labels \\
I.13 & --- & R2, R3 (pooled), R10 (restricted) accepted & --- & No action. & --- \\
\multicolumn{6}{@{}l}{\textbf{Batch B ecological-inference referee, round 2}}\\
E2-1 & E & W1 example and pairwise CR & guide; vignette; 5a \S1.4; package label & Either part may be negative; pair changes one side only. & pkg tests \\
E2-2 & E & Coverage figure for the procedure described & 5a \S3.5 item~4; guide & Literal intersection 304/600; CR-naive 550/600. & V13 [V] \\
E2-3 & E & False-belief results & guide; vignette & All three counts (398, 110, 0 of 400). & --- \\
E2-4 & S & Projection guarantee scoped to the fit & 5a R10; 2c M5; package \texttt{print()} & Status line per fit; R16 for grouped M-SIGN. & pkg tests \\
E2-5 & X/C & Rung-only comparison; open items & guide; vignette; package & Joint effect of removing all beliefs; ``Not yet available'' list. & pkg test \\
\multicolumn{6}{@{}l}{\textbf{Internal consistency log of merged M-SIGN material (not an external referee)}}\\
M-1 & S & R3 dimension hypothesis & 5a R3, R2 & $\operatorname{affdim}(A)=R-1$. & --- \\
M-2 & S & R10 after the P3.5 split & 5a R10 & Status note by version. & --- \\
M-3 & S & Relaxed-candidate device & 5a \S5 & Justified by weak duality; coverage = (b3). & --- \\
M-4 & X & 1b ``uniformly'' claim & 1b & Dropped. & --- \\
M-5 & S & Clustering in 1b/2d & 1b, 2d & One unit observed twice; county effects in the mean. & --- \\
M-6, M-7 & X & Old theorem numbers; ``proved'' for CLR & PD, 5a, 2c & Renumbered; CLR algorithmic. & --- \\
M-8 & C & \texttt{cluster} argument & package & Removed. & --- \\
M-9--M-13 & --- & Items needing external review & 5a, 2c, 1b, 2d & Reviewed in Batch~B (M-9 to M-12); 1b/2d (M-13) not yet. & --- \\
M-14 & X & Name clash D4 & 1b, 2d, PD & DiD rungs renamed $\Delta0$--$\Delta8$. & --- \\
M-15 & --- & NEW\_AVENUES cluster claims & NEW\_AVENUES & Marked for re-derivation. & --- \\
M-16, M-17 & C & Code review; coverage MC & package & Partly via referees; batch~E code review open. & --- \\
\bottomrule
\end{longtable}
}

% Appendices B-H: the repository's theory documents, reproduced verbatim.
% The files in generated/ are produced by tools/make_appendices.py (pandoc);
% do not edit them by hand.
\clearpage
\section*{About the reproduced appendices}
\addcontentsline{toc}{section}{About the reproduced appendices}
Appendices~\ref{app:phase1}--\ref{app:pkgdesign} reproduce the repository's theory
documents at commit \texttt{b270285}. They were converted mechanically from
Markdown to \LaTeX\ with pandoc (\texttt{tools/make\_appendices.py}); the text,
proofs, justification tags and correction notes are those of the source. Internal
cross-references in them (``2a Lemma~1'', ``P3 \S5'', ``[2c:M3]'') use the
repository's labels, which this report keeps. Section numbers inside each appendix
are the source document's own. References to files are to the repository.
PHASE1b and PHASE2d are not reproduced (proposals, not reviewed; see
Section~\ref{sec:did}).

\begingroup
\small
\setcounter{secnumdepth}{1}

\clearpage
\section{PHASE1: structural foundations and assumptions}\label{app:phase1}
\textit{Source: \texttt{docs/extension/PHASE1\_foundations.md}.}\par
% GENERATED by tools/make_appendices.py from docs/extension/PHASE1_foundations.md
% Do not edit by hand; edit the source and regenerate.
\emph{Extending Jiang, King, Schmaltz \& Tanner (2019), ``Ecological Regression with Partial Identification'' (JKST below), and the \texttt{eiPartialID} package to (1) R×C tables and (2) covariates, especially district fixed effects. This is a draft for review. Results marked \textbf{(to prove)} are stated precisely here and get proved in Phase 2. Nothing in this phase is a derivation.}

\begin{center}\rule{0.5\linewidth}{0.5pt}\end{center}

\hypertarget{pA-what-the-source-materials-actually-do-and-what-that-means-for-sharpness}{%
\subsection{0. What the source materials actually do (and what that means for ``sharpness'')}\label{pA-what-the-source-materials-actually-do-and-what-that-means-for-sharpness}}

I read the paper, the supplementary appendix (\path{appendix.tex}, Appendices A--C) and \path{package/eiPartialID/R/main.R}. Six facts shape everything below.

\begin{enumerate}
\def\labelenumi{\arabic{enumi}.}
\tightlist
\item
  \textbf{Model.} JKST Assumption 1: \((\beta^b_i,\beta^w_i,X_i,N_i)\) are iid, with \(E(\beta^w_i\mid X_i,N_i)=w_0+w_1X_i\) and \(E(\beta^b_i\mid X_i,N_i)=b_0+b_1X_i\). The accounting identity then implies \(E(T_i\mid X_i)=w_0+c_1X_i+d_1X_i^2\). That identifies \(\theta=(w_0,c_1,d_1)\) and leaves one free parameter, \(w_1\).
\item
  \textbf{Prop. 1 (JKST)} bounds \(w_1\) through \(E(L_i\mid X_i)\le w_0+w_1X_i\le E(U_i\mid X_i)\) on a set \(A\). Prop. 2 relaxes this with Jensen's inequality and the quadratic regression. \textbf{The package implements only Prop. 2} with \(A=[\min X_i+10^{-5},\ \max X_i-10^{-5}]\) and \(\lambda=0\), so the package's bound is an \emph{outer} set, not the sharp one.
\item
  \textbf{Inference (Prop. 4, App. A).} \(\mathrm{CI}_x=[\hat B_L-xS_L,\ \hat B_U+xS_U]\cap DD\), with HC1 sandwich variance and the Popoviciu bound \(\operatorname{Var}(\beta)\le 1/4\) for the finite-population term (\(S_1\)). It needs tie-breaking (a unique active constraint) and covers the \emph{parameter} \(B\) with probability at least \(\Phi(x)\), in the spirit of Imbens--Manski.
\item
  \textbf{Selection heuristic.} The package falls back to DD when \(\mathrm{CI}_0\) intersected with DD is empty. App. B justifies this: under the model, \(\mathrm{CI}_0\neq\emptyset\) in the large-\(p\) limit.
\item
  \textbf{Covariates (App. C).} These enter additively in the intercept: \(E(\beta^w\mid X,Z)=\tilde w_0+\tilde w_1X+\tilde w_2Z\), and likewise for \(\beta^b\), with a common slope. Five parameters are identified and only \(\tilde w_1\) is not. The appendix bounds it by restricting \(A\) to a line segment \(Z=a_0+a_1X\) (Prop. 3*). That choice throws away most of what a covariate can contribute. Apart from the \(\log N\) example, it does not analyse \emph{discrete} covariates such as district dummies. None of App. C is in the package.
\item
  \textbf{Verified equivalence.} I rewrote Prop. 2 in the homogeneous form used below (§1.4) and checked numerically that it reproduces \texttt{generateBounds} to \(10^{-16}\) (\path{docs/extension/check_2x2_equivalence.py}). The R×C generalisation therefore contains the existing package as the special case \(R=C=2\).
\end{enumerate}

\textbf{Consequence for the brief.} You asked for a proof that ``the bounds'' are sharp. Under JKST's assumptions the bound the package computes is \textbf{not} sharp. There are three nested sets, \(\Theta_I\subseteq\Theta_X\subseteq\Theta_O\) (§1.5), and Phase 2 will prove sharpness for \(\Theta_I\), prove validity for \(\Theta_O\), and characterise when they coincide. I will not claim sharpness for \(\Theta_O\).

\begin{center}\rule{0.5\linewidth}{0.5pt}\end{center}

\hypertarget{pA-identification-regime}{%
\subsection{1. Identification regime}\label{pA-identification-regime}}

\hypertarget{pA-sampling-scheme}{%
\subsubsection{1.1 Sampling scheme}\label{pA-sampling-scheme}}

There are units (precincts) \(i=1,\dots,p\) nested in districts \(d=1,\dots,D\). \(D_i\in\{1,\dots,D\}\) is unit \(i\)'s district, \(\mathcal I_d=\{i:D_i=d\}\), and \(p_d=|\mathcal I_d|\). Each unit gives
\[
W_i=(N_i,\;X_i,\;T_i,\;Z_i,\;D_i),\qquad X_i\in\Delta^{R-1},\quad T_i\in\Delta^{C-1},
\]
where \(\Delta^{k-1}=\{v\in[0,1]^k:\mathbf 1^\top v=1\}\), \(X_{ir}=N_{ir}/N_i\) is the row margin (e.g.~vote shares in election 1, including abstention), and \(T_{ic}\) is the column margin (e.g.~shares in election 2: abstain, party A, party B, \ldots). \(Z_i\in\mathcal Z\subseteq\mathbb R^{k_z}\) are other covariates, such as \(\log N_i\) or census variables. \emph{(Scope after review: Phases 2a--3 use no covariates other than the district \(D_i\). Covariate adjustment is a planned extension, Phase 2c, because it needs stronger assumptions; see §1.6.)}

\textbf{Latent quantities.} These are the unit cell counts \(N_{irc}\) (the unobserved \(R\times C\) table) and the transition probabilities \(\beta_{irc}=N_{irc}/N_{ir}\), defined when \(N_{ir}>0\). Write \(\beta_i=[\beta_{irc}]\in\mathfrak B:=(\Delta^{C-1})^R\) (each row is a point in the simplex).

\textbf{Accounting identity (definitional for a closed population, see A1):}
\[
T_{ic}=\sum_{r=1}^R X_{ir}\beta_{irc}\quad(c=1,\dots,C),\qquad\text{i.e.}\quad T_i=\beta_i^\top X_i. \tag{1.1}
\]

\hypertarget{pA-estimands}{%
\subsubsection{1.2 Estimands}\label{pA-estimands}}

\begin{itemize}
\tightlist
\item
  \textbf{Finite-population (district- or region-wide) transition rates} for any target set of units \(\mathcal S\) (all units, or \(\mathcal I_d\)):
  \[B_{rc}(\mathcal S)=\frac{\sum_{i\in\mathcal S}N_iX_{ir}\beta_{irc}}{\sum_{i\in\mathcal S}N_iX_{ir}}.\tag{1.2}\]
  The JKST \(B\) is \(B_{\text{black},\text{vote}}(\{1..p\})\).
\item
  \textbf{Superpopulation analogue} (the object that is \emph{identified}; \(B-B^*=O_p(p^{-1/2})\), Phase 3):
  \[B^*_{rc}=\frac{E[N X_r\beta_{rc}]}{E[NX_r]}.\tag{1.3}\]
\item
  \textbf{Transition functionals} for your application (\(R=2\) rows A, B; \(C=3\) columns \(0\) = abstain, A, B): loyalty \(\lambda_r=B_{r,r}\), defection \(\delta_r=B_{r,r'}\) (\(r'\neq r\), \(r'\neq 0\)), demobilisation \(\mu_r=B_{r,0}\). ``Same/different party'' is a row-dependent relabelling of columns, so every quantity of interest is a \textbf{linear functional} \(\varphi=\sum_{rc}a_{rc}B_{rc}\). Since \(\sum_c B_{rc}=1\), the three are linked by \(\lambda_r+\delta_r+\mu_r=1\).
\end{itemize}

\hypertarget{pA-structural-model-linear-contextual-effects-in-homogeneous-form}{%
\subsubsection{1.3 Structural model: linear contextual effects in homogeneous form}\label{pA-structural-model-linear-contextual-effects-in-homogeneous-form}}

Because \(\mathbf 1^\top X=1\), an affine function of \(X\) is a \emph{homogeneous linear} function of \(X\): \(a+\sum_{s<R}g_sX_s=\sum_s(a+g_s)X_s\) with \(g_R:=0\). This removes the intercept/slope bookkeeping and makes the R×C algebra symmetric.

\textbf{Model M0 (pooled LCE, generalising JKST Assumption 1).} For each column \(c\) there is a matrix \(H_c\in\mathbb R^{R\times R}\) such that
\[
E[\beta_{irc}\mid X_i,N_i]=e_r^\top H_cX_i=\sum_{s=1}^R H_{c,rs}X_{is}\qquad\text{a.s. on }\{X_i\in A\}. \tag{1.4}
\]
Rows of \(\beta\) sum to 1, so \(\sum_c e_r^\top H_c x=1=\mathbf 1^\top x\) for \(P_X\)-a.e. \(x\in A\). Under the rank condition A5 this forces
\[\textstyle\sum_{c=1}^C H_c=\mathbf 1\mathbf 1^\top=:J. \tag{1.5}\]
(Proof in 2a Lemma 1(b). With \(d_r^\top=e_r^\top(\sum_cH_c-J)\), the identity gives \((d_r^\top x)(\mathbf 1^\top x)=0\) on \(\operatorname{supp}_AP_X\), which is a quadratic form in \(x\), so A5 gives \(d_r=0\). \emph{Corrected after review:} an earlier version argued from ``an open subset of the simplex'', which was never assumed.)
(2×2 dictionary, with row 1 = black and \(X_1=X\): \(w_0=H_{22}\), \(w_1=H_{21}-H_{22}\), \(b_0=H_{12}\), \(b_1=H_{11}-H_{12}\).)

\textbf{Observable implication.} Taking \(E[\cdot\mid X]\) of (1.1),
\[
\tau_c(x):=E[T_c\mid X=x]=x^\top H_cx=x^\top Q_cx,\qquad Q_c:=\tfrac12(H_c+H_c^\top). \tag{1.6}
\]
Write \(H_c=Q_c+K_c\) with \(K_c:=\tfrac12(H_c-H_c^\top)\) skew-symmetric. A skew part never shows up in a quadratic form, since \(x^\top Kx=0\).

\phantomsection\label{pA-res-lemma-1}\textbf{Lemma 1 (identification decomposition) (to prove).} Assume A5 below (the support of \(X\) is not contained in a quadric). Then:

\begin{enumerate}
\def\labelenumi{(\alph{enumi})}
\tightlist
\item
  \(Q=(Q_1,\dots,Q_C)\) is point-identified from \(\tau(\cdot)\);
\item
  \(K=(K_1,\dots,K_C)\) is \emph{completely} unidentified by \(\tau\) and is constrained only by \(\sum_cK_c=0\) (from 1.5 together with \(\sum_cQ_c=J\));
\item
  so the unidentified parameter lives in the linear space
  \[\mathcal K=\Big\{(K_1,\dots,K_C):K_c^\top=-K_c,\ \textstyle\sum_cK_c=0\Big\},\qquad \dim\mathcal K=(C-1)\,\tfrac{R(R-1)}{2}. \tag{1.7}\]
\end{enumerate}

Check: \(2\times2\Rightarrow1\) (JKST's \(w_1\), with \(K_{1,12}=c_1/2-w_1\)). \textbf{\(2\times3\Rightarrow 2\).} \(3\times3\Rightarrow 6\). Interpretation: for each pair of rows \((r,s)\) and each column \(c\), the data identify the \emph{sum} of ``the effect of \(s\)-share on row \(r\)'' and ``the effect of \(r\)-share on row \(s\)'', but not how that sum splits between them. This is the R×C version of ``is turnout low because blacks vote less, or because whites in black areas vote less?'' (JKST Fig. 1).

\textbf{Estimand as a function of the parameters.} Under A3 (mean independence from \(N\)),
\[
B^*_{rc}(K)=\frac{1}{\mu_r}\Big[\sum_{s}Q_{c,rs}\mu_{rs}+\sum_{s\neq r}K_{c,rs}\mu_{rs}\Big],\qquad \mu_{rs}:=E[NX_rX_s],\ \mu_r:=E[NX_r]. \tag{1.8}
\]
This is \emph{affine in \(K\)}. In the ambient entries \(K_{c,rs}\) (before skew-symmetry and \(\sum_cK_c=0\) are imposed), the coefficient of \(K_{c,rs}\) is \(\mu_{rs}/\mu_r\); for 2×2 this is JKST's \(r\), eq. 18. In free coordinates see 2a Prop. 7. The formula requires \(P(X\in A)=1\) (2a Lemma 2; 2a Prop. 7′ treats units outside \(A\)).

\hypertarget{pA-constraint-sets-the-outer-set-generalised-jkst-prop.-23}{%
\subsubsection{1.4 Constraint sets: the outer set (generalised JKST Prop. 2/3)}\label{pA-constraint-sets-the-outer-set-generalised-jkst-prop.-23}}

Let \(A\subseteq\Delta^{R-1}\) be a compact convex polytope on which (1.4) is assumed, with vertex set \(\operatorname{vert}(A)\).

\[
\boxed{\;\Theta_O(P;A)=\Big\{K\in\mathcal K:\ (Q_c(P)+K_c)\,v\ \ge 0\ \text{(elementwise)}\quad\forall c,\ \forall v\in\operatorname{vert}(A)\Big\}\;} \tag{1.9}
\]

\phantomsection\label{pA-res-proposition-2-RC}\textbf{Proposition 2-RC (to prove).} (i) \(\Theta_O\) is a convex polyhedron containing the true \(K\) under M0; it is bounded (a polytope) when \(A\) is full-dimensional in the simplex, which A5 implies (proved in 2a Lemma 3(b) and Remark 5.2). (ii) For \(R=C=2\) and \(A=[l,u]\) it equals JKST Prop. 2; for \(A=\Delta^{R-1}\), where \(\operatorname{vert}(A)=\{e_s\}\) and the constraint becomes \(H_{c,rs}\ge0\), it equals JKST Prop. 3. (iii) It is exactly the set of \(K\) for which every implied conditional-mean row \(E[\beta_{r\cdot}\mid X=x]\) lies in \(\Delta^{C-1}\) for all \(x\in A\).

\emph{Why it has this form:} the JKST Jensen step \(E[U\mid X]\le\min\{1,E[T\mid X]/(1-X)\}\) generalises to ``the conditional mean of a selection of the DD polytope lies in the DD polytope evaluated at \(E[T\mid X]\)''. Because \(x^\top H_c x=\tau_c(x)\) holds by construction, the only constraint left is nonnegativity, and a linear function on a polytope is minimised at a vertex. (Remark 1.4: this is the form checked numerically against the package.)

\textbf{Bounds on any linear functional} are linear programs:
\[
[\varphi_L,\varphi_U]=\Big[\min_{K\in\Theta_O}\varphi(K),\ \max_{K\in\Theta_O}\varphi(K)\Big]. \tag{1.10}
\]
\textbf{Closed form for \(R=2\) (any \(C\)) (to prove).} Each constraint involves a single \(\kappa_c:=K_{c,12}\), so \(\Theta_O=\big(\prod_c[\ell_c,h_c]\big)\cap\{\sum_c\kappa_c=0\}\). Its projection is
\[\kappa_c\in\Big[\max\{\ell_c,-\textstyle\sum_{c'\neq c}h_{c'}\},\ \min\{h_c,-\sum_{c'\neq c}\ell_{c'}\}\Big].\]
\textbf{Corollary (to prove):} the joint 2×C bound is weakly tighter than running \(C\) separate collapsed 2×2 analyses (``column \(c\) vs the rest''), and strictly tighter whenever the sum-to-zero coupling binds. So the 2×3 transition analysis gains from being done jointly.

\hypertarget{pA-the-sharp-identification-region}{%
\subsubsection{1.5 The sharp identification region}\label{pA-the-sharp-identification-region}}

\textbf{Primitive objects.} Let \((\Omega,\mathcal F,\mathbb Q)\) carry \((W,\beta)\). The observable law is \(P=\mathbb Q\circ W^{-1}\). The conditioning σ-algebra is \(\mathcal G=\sigma(X,N)\) (with \(D\) added in Phase 2b). \emph{(Corrected after review: an earlier version conditioned on \(Z\) as well while the mean \(M_H(X)\) did not depend on \(Z\), so (S2) did not follow from A2--A3 as stated. Covariates are a planned extension, §1.6.)} A \textbf{structure} is a pair \(s=(H,\mathbb Q_{\beta\mid W})\). The model \(\mathcal M(A)\) is the set of laws \(\mathbb Q\) such that

\begin{itemize}
\tightlist
\item
  (S1) \(\beta\in\mathfrak B\) a.s. and \(T=\beta^\top X\) a.s. (A1);
\item
  (S2) \(E_{\mathbb Q}[\beta\mid\mathcal G]=M_H(X):=[e_r^\top H_cX]_{rc}\) a.s. on \(\{X\in A\}\) (A2, A3).
\end{itemize}

\textbf{Definition (sharp identified set).}
\[
\Theta_I(P;A)=\big\{K\in\mathcal K:\ \exists\,\mathbb Q\in\mathcal M(A)\text{ with }\mathbb Q\circ W^{-1}=P\text{ and }H=Q(P)+K\big\}. \tag{1.11}
\]
\(\Theta_I\) is a correspondence \(P\mapsto 2^{\mathcal K}\): it contains exactly the parameters that are observationally equivalent to the data. For the estimand, \(\mathcal B^*_I(P)=\{B^*(K):K\in\Theta_I(P)\}\), which is an interval for a scalar \(\varphi\) because \(\Theta_I\) is convex and \(B^*\) is affine.

\textbf{Random-set characterisation.} Let \(S(x,t)=\{b\in\mathfrak B: b^\top x=t\}\) be the unit-level DD polytope (a transportation polytope in cell-share coordinates \(n_{rc}=x_rb_{rc}\)), with support function \(h_S(u;x,t)=\max_{b\in S(x,t)}\langle u,b\rangle\), an LP that is concave in \(t\).

\phantomsection\label{pA-res-theorem-1}\textbf{Theorem 1 (to prove).} Assume A1--A3 and that \(\mathcal F\) is rich enough to allow auxiliary randomisation. Then
\[
K\in\Theta_I(P;A)\iff \langle u,M_{Q+K}(x)\rangle\le E\big[h_S(u;X,T)\mid X=x,N=n\big]\quad\forall u\in\mathbb S^{RC-1},\ P\text{-a.e. }(x,n),\ x\in A, \tag{1.12}
\]
equivalently \(M_{Q+K}(x)\in\mathbb E[S(X,T)\mid\mathcal G]\), the conditional Aumann (selection) expectation. Tools: Hiai--Umegaki conditional selection expectation, the Kuratowski--Ryll-Nardzewski measurable selection theorem, and Beresteanu, Molchanov \& Molinari (2011). For 2×2 conditioned on \(X\) only, (1.12) reduces \emph{exactly} to JKST Prop. 1, which confirms their ``tightest bound'' claim in that case.

\textbf{Nesting (to prove).}
\[
\Theta_I(P;A)\ \subseteq\ \Theta_X(P;A)\ \subseteq\ \Theta_O(P;A). \tag{1.13}
\]
\(\Theta_X\) is (1.12) conditioning on \(X\) only, i.e.~JKST Prop. 1; the gap \(\Theta_I\subsetneq\Theta_X\) appears when the model conditions on \(N\) and \(E[h_S\mid X,N]\) varies with \(N\). \(\Theta_O\) is (1.9). The second inclusion is the Jensen step. \textbf{Equality condition (conjecture, to prove or refute):} \(\Theta_X=\Theta_O\) if, for a.e. \(x\), \(h_S(u;x,\cdot)\) is affine on the conditional support of \(T\mid X=x\) for every \(u\) that is active at the optimum. In 2×2 this means \(T/(1-X)\) does not straddle 1 (and \(T/X\) does not straddle 0 or 1) within an \(X\)-slice.

\textbf{Support caveat.} (1.8) uses the linear model on all of \(\operatorname{supp}(X)\), so \(A\supseteq\operatorname{supp}(X)\) is required for the formula. If \(A\subsetneq\operatorname{supp}(X)\), the bound for \(B^*\) adds the model-free Duncan--Davis term of the units outside \(A\); this is formalised in 2a Prop. 7′. The package sets \(A=\) the range of the data, which is consistent with this.

\hypertarget{pA-covariates-and-district-fixed-effects}{%
\subsubsection{1.6 Covariates and district fixed effects}\label{pA-covariates-and-district-fixed-effects}}

Write the district-specific coefficient matrices as \(H_c^{(d)}\).

\begin{longtable}[]{@{}
  >{\raggedright\arraybackslash\hspace{0pt}}p{(\columnwidth - 6\tabcolsep) * \real{0.2500}}
  >{\raggedright\arraybackslash\hspace{0pt}}p{(\columnwidth - 6\tabcolsep) * \real{0.2500}}
  >{\raggedright\arraybackslash\hspace{0pt}}p{(\columnwidth - 6\tabcolsep) * \real{0.2500}}
  >{\raggedright\arraybackslash\hspace{0pt}}p{(\columnwidth - 6\tabcolsep) * \real{0.2500}}@{}}
\toprule\noalign{}
\begin{minipage}[b]{\linewidth}\raggedright\hspace{0pt}Model
\end{minipage} & \begin{minipage}[b]{\linewidth}\raggedright\hspace{0pt}Assumption on \(E[\beta_{rc}\mid X,N,D=d]\)
\end{minipage} & \begin{minipage}[b]{\linewidth}\raggedright\hspace{0pt}Unidentified parameter
\end{minipage} & \begin{minipage}[b]{\linewidth}\raggedright\hspace{0pt}Identified set
\end{minipage} \\
\midrule\noalign{}
\endhead
\bottomrule\noalign{}
\endlastfoot
\textbf{M0} pooled & \(e_r^\top H_cX\) (no district) & \(K\in\mathcal K\) & \(\Theta_O(P;A)\) \\
\textbf{M1} district intercepts, common contextual effects & \(e_r^\top H_cX+\zeta_{c,d,r}\), \(\ \sum_c\zeta_{c,d}=0\), \(\ \zeta_{c,1}=0\) & common \(K\in\mathcal K\) & \(\bigcap_d\Theta_O^{(d)}\) \\
\textbf{M1′} slopes depend on district covariates \(V_d\) \emph{(planned, Phase 2c; needs linearity in the covariates and a joint support condition)} & \(e_r^\top(H_c+\textstyle\sum_\ell V_{d\ell}G_{c\ell})X+\zeta_{c,d,r}\) & \(K\) and the skew parts of \(G\) & polyhedron; to derive \\
\textbf{M2} saturated & \(e_r^\top H^{(d)}_cX\), all district-specific & \(K^{(d)}\) for each \(d\) & \(\prod_d\Theta_O(P_d;A_d)\) \\
\end{longtable}

Here \(\Theta_O^{(d)}=\{K:(Q^{(d)}_c+\operatorname{skew}(\zeta_{c,d}\mathbf1^\top)+K_c)v\ge0\ \forall v\in\operatorname{vert}(A_d),\ \forall c\}\), with \(A_d\) the convex hull of the support of \(X\) in district \(d\). JKST App. C is M1 with a continuous \(Z\).

\textbf{Structural facts (to prove in Phase 2b):}

\begin{itemize}
\tightlist
\item
  \textbf{(D0) Without a model, dummies add nothing.} Under A1 alone, the sharp set for \(B_{rc}\) is the DD bound whatever \(Z\) or \(D\) is, because DD is already computed unit by unit and aggregation is additive. \emph{Every gain from district dummies is therefore model-borne}: it comes only from cross-district restrictions such as common \(K\).
\item
  \textbf{(D1) Intercept shifts are point-identified.} Under M1, \(Q_c^{(d)}=Q_c+\tfrac12(\zeta_{c,d}\mathbf1^\top+\mathbf1\zeta_{c,d}^\top)\), so \(\zeta_{c,d}=\operatorname{diag}(Q^{(d)}_c-Q^{(1)}_c)\) is identified. Additive covariates add constraints but no new unidentified dimensions: \(\dim\) of the unidentified part is still \((C-1)R(R-1)/2\), the same as M0. This generalises App. C, where \(\tilde w_2\) and \(\tilde b_2\) are identified.
\item
  \textbf{(D2) Overidentification.} M1 implies \((Q^{(d)}_c-Q^{(1)}_c)_{rs}=\tfrac12(\zeta_{c,d,r}+\zeta_{c,d,s})\) for \(r\ne s\). Because \(\sum_cQ^{(d)}_c=J\) in every district, the restrictions for column \(C\) follow from those for \(c<C\). So there are \(R(R-1)/2\) independent restrictions per \((c<C,\ d\neq1)\), i.e.~\((D-1)(C-1)R(R-1)/2\) in total. This equals the parameter count of M2, \((C-1)R(R+1)/2\) per district, minus that of M1, \((C-1)R\) shifts per extra district: \((D-1)(C-1)\big[R(R+1)/2-R\big]\). \emph{(Corrected after review; the earlier count of \(R(R-1)/2\) per \((c,d)\) for every \(c\) ignored this dependence.)} In 2×2 it is ``the \(X^2\) coefficient is equal across districts''. In addition, \(\bigcap_d\Theta_O^{(d)}=\emptyset\) refutes M1, which generalises the JKST revert-to-DD rule.
\item
  \textbf{(D3) Intersection gain.} \(\Theta^{M1}=\bigcap_d\Theta_O^{(d)}\subseteq\Theta_O^{(d)}\) for every \(d\). For \(R=2\) the bound on \(\kappa_c\) is \([\max_d\ell_{c,d},\ \min_d h_{c,d}]\), so a single district whose intercept puts some row's conditional mean close to 0 or 1 at an extreme of its own \(X\)-range pins down the common contextual effect for \emph{all} districts.
\item
  \textbf{(D4) When dummies matter.} Under M1, \(E[\beta\mid X,N]=M_H(X)+E[\zeta_D\mid X,N]\). So M0, which conditions on \((X,N)\), holds (with a shifted intercept) iff \(E[\zeta_D\mid X,N]\) is affine in \(X\) and does not depend on \(N\), e.g.~when \(D\perp(X,N)\). \textbf{M0 is misspecified when district effects are correlated with \(X\) non-affinely, or with precinct size.} \emph{(Corrected after review: an earlier version conditioned on \(X\) only and claimed \(D\perp X\) suffices; see 2b Prop. D4(b).)} This is an ecological version of omitted-variable bias. If \(\zeta\equiv0\), then \(\Theta^{M1}=\Theta_O(P;\operatorname{conv}\bigcup_dA_d)=\Theta^{M0}\): dummies cost precision but neither help nor hurt identification. Corollary: dummies (i) protect against between-district aggregation bias, and (ii) can tighten bounds through (D3). M0 and M1 are \textbf{non-nested}; M2 is weaker than both.
\item
  \textbf{(D5) M2 sums intervals.} \(\mathcal B^{*,M2}_I=\sum_d\omega_d[\,B^{*}_{d,L},B^{*}_{d,U}]\), a Minkowski sum. It is valid under district-specific contextual effects and can be wider or narrower than M0.
\item
  \textbf{By-product:} M1 and M2 give \textbf{district-level bounds} \(B_{rc}(\mathcal I_d)\) as well, e.g.~transition rates per constituency.
\end{itemize}

\hypertarget{pA-objects-hulls-and-targets-added-after-review}{%
\subsubsection{1.7 Objects, hulls and targets (added after review)}\label{pA-objects-hulls-and-targets-added-after-review}}

The phases use three identified sets, three composition hulls and three targets. Each must be named when a result is stated.

\textbf{Three identified sets} (nested: \(\Theta_I\subseteq\Theta_X\subseteq\Theta_O\), 2a Prop. 1):

\begin{longtable}[]{@{}
  >{\raggedright\arraybackslash\hspace{0pt}}p{(\columnwidth - 6\tabcolsep) * \real{0.2500}}
  >{\raggedright\arraybackslash\hspace{0pt}}p{(\columnwidth - 6\tabcolsep) * \real{0.2500}}
  >{\raggedright\arraybackslash\hspace{0pt}}p{(\columnwidth - 6\tabcolsep) * \real{0.2500}}
  >{\raggedright\arraybackslash\hspace{0pt}}p{(\columnwidth - 6\tabcolsep) * \real{0.2500}}@{}}
\toprule\noalign{}
\begin{minipage}[b]{\linewidth}\raggedright\hspace{0pt}Set
\end{minipage} & \begin{minipage}[b]{\linewidth}\raggedright\hspace{0pt}Information it uses
\end{minipage} & \begin{minipage}[b]{\linewidth}\raggedright\hspace{0pt}Sharp relative to
\end{minipage} & \begin{minipage}[b]{\linewidth}\raggedright\hspace{0pt}How computed
\end{minipage} \\
\midrule\noalign{}
\endhead
\bottomrule\noalign{}
\endlastfoot
\(\Theta_I\) & law of \((N,X,T)\); conditions on \((X,N)\) & the full observable law under the model (2a Thm 1) & conditional support-function inequalities; not a finite LP \\
\(\Theta_X\) & law of \((X,T)\); conditions on \(X\) only & the law of \((X,T)\) (JKST Prop. 1) & the same, conditioning on \(X\) \\
\(\Theta_O\) & the regression \(E[T\mid X]\) plus probability validity & the regression function only (2a Prop. 2) & a finite LP (2a Prop. 7) \\
\end{longtable}

Rule: \(\Theta_O\) is never called ``sharp'' without ``for the regression function''. 2a Prop. 3 gives a \emph{sufficient} condition for \(\Theta_I=\Theta_O\), not a characterisation.

\textbf{Three hulls:}

\begin{itemize}
\tightlist
\item
  \textbf{Population hull} \(A=\operatorname{conv}\operatorname{supp}P_X\) (per district in 2b). This is the default region in all identification results, and it makes (S3) automatic (2a Lemma 3(c)).
\item
  \textbf{Hull of the empirical law} \(\operatorname{conv}\{X_1,\dots,X_n\}\), read as the population hull of \(P_n\). Identification statements with \(P\) replaced by \(P_n\) use it exactly. That is a finite-sample identification analysis, with no estimation step.
\item
  \textbf{Sample hull as an estimate of \(A\).} It is the same set, but used as a plug-in for the population hull. It is not a population assumption. Since \(\hat A\subseteq A\), the plug-in outer set is weakly larger, so estimating the hull only widens the bound (P3 Prop. P3.2(c)). The true parameter satisfies the constraints at the sample-hull vertices, because they lie in \(A\), so coverage of the true parameter does not depend on how well \(\hat A\) approximates \(A\) (to be stated in P3, WS4). What the hull error affects is width, and the distribution of the endpoint estimator, which P3 states for the fixed sample hull.
\end{itemize}

\textbf{Three targets:}

\begin{itemize}
\tightlist
\item
  \(B^*_{rc}=E[NX_r\beta_{rc}]/E[NX_r]\): the superpopulation rate of the identification results.
\item
  \(B_{rc}\): the realised rate for this electorate (1.2), JKST's target.
\item
  \(\tilde B_{rc}=\sum_iw_{ir}m_{rc}(V_i)\) with \(w_{ir}=N_iX_{ir}/\sum_jN_jX_{jr}\): the finite-design conditional-mean rate. Under A6 it is \(E[B\mid\text{design}]\).
\end{itemize}

Only \(\tilde B\) (and \(B^*\)) is affine in \(K\). \(B=\tilde B+E\) adds a residual term that does not depend on \(K\). Inference is conditional on the design and covers \(B\) or \(\tilde B\). Inference for \(B^*\) would also need the sampling variation of \((N_i,X_i,D_i)\), and is not pursued.

\textbf{Continuous latent rates.} Every ``sharp'' statement is relative to the model with real-valued \(\beta\in\mathfrak B\) (2a Remark 4.2), not to the finite set of integer cell counts.

\begin{center}\rule{0.5\linewidth}{0.5pt}\end{center}

\hypertarget{pA-assumptions}{%
\subsection{2. Assumptions}\label{pA-assumptions}}

Each entry gives the statement, its mathematical role in the bounds, what it means for voter transitions, the consequence of a violation, and whether it can be tested.

\textbf{A1. Closed population / accounting identity.} \(T_i=\beta_i^\top X_i\) exactly, with \(N_i\) the same people in both margins.
\emph{Role:} defines \(S(x,t)\), DD and \(\tau(x)=x^\top Q x\); without it nothing is bounded.
\emph{Transitions:} both elections are measured on the \emph{same} electorate (the same register), and rows include previous abstainers and the ineligible if they can enter.
\emph{Violation:} turnover (deaths, moves, new registrants) or boundary changes give \(T_i=\beta_i^\top X_i+\varepsilon_i\). The DD sets shift, \(Q\) absorbs \(E[\varepsilon\mid X]\), and the bounds can exclude the truth. Bounded turnover, \(\|\varepsilon_i\|_1\le\epsilon\), is \textbf{an open extension} (P3 §10). It changes both the unit DD sets and the regression identity \(E[T\mid X]=X^\top QX\), so it cannot be handled by widening the final interval while keeping the estimated \(Q\). \emph{(Corrected after review: an earlier version said Phase 3 treats it.)} Partial remedy: extra ``entered'' row and ``exited'' column categories.
\emph{Testable:} only through implied infeasibility, e.g.~\(T_i\) outside the range compatible with \(X_i\).

\textbf{A2. Linear contextual effects (LCE) on \(A\):} equation (1.4).
\emph{Role:} the \emph{only} identifying restriction; it turns the infinite-dimensional \(E[\beta\mid X]\) into \(H\), and gives (1.6), (1.8) and (1.9).
\emph{Transitions:} A-voters' loyalty in precinct \(i\) is linear in the precinct's previous shares, e.g.~``A-voters in A strongholds are more loyal'' (social context, local organisation).
\emph{Violation:} threshold or tipping effects (party machines mobilise only above 50\%), or U-shapes (competitive precincts mobilise both sides). The quadratic regression is then misspecified, \(\Theta_O\) is computed for the wrong \(Q\), and the bound can miss the truth by \(O(1)\) that does not shrink with \(p\). There are \emph{undetectable} violations: in the R×C analogue of JKST Example 6, quadratic contextual effects cancel in \(\tau\).
\emph{Testable:} partly. A cubic term in the \(T\sim X\) regression rejects A2; emptiness of \(\Theta_O\) or of \(\Theta_O\cap DD\) rejects A2 (App. B generalised); exact cancellations cannot be detected.

\textbf{A3. Mean independence from size:} \(E[\beta\mid X,N]=E[\beta\mid X]\) (with \(D\) added to both sides under M1/M2). Together with A2 this gives \(E[\beta\mid X,N]=M_H(X)\), which is (S2) of 2a. \emph{(Corrected after review: covariates \(Z\) are no longer in the conditioning set.)}
\emph{Role:} makes the finite-population term \(\sum_iN_iX_{ir}e_{irc}/\sum_iN_iX_{ir}\) mean zero, so that (1.8) holds with weights \(N X_r\).
\emph{Transitions:} large (urban) precincts do not differ in loyalty from small ones \emph{given} their partisan mix.
\emph{Violation:} \(B^*\neq\) the regression functional; the bias is \(\operatorname{Cov}(N X_r,\beta_{rc}-M(X))/E[NX_r]\). Remedy: district dummies when size differs by district (2b Prop. D4); \(\log N\) as a covariate is part of the planned covariate extension (App. C).
\emph{Testable:} partly, by adding \(X\cdot\log N\) interactions to the \(T\) regression.

\textbf{A4. Support and extrapolation:} \(\operatorname{supp}(X)\subseteq A\), compact, and A2 holds on all of \(A\).
\emph{Role:} \(\Theta_O\) uses the vertices of \(A\). A larger \(A\) gives tighter bounds (more constraints), but only if A2 is believed there.
\emph{Transitions:} the linear form holds from the least to the most A-dominated precinct observed. For \(R\ge3\), \(A\) is the convex hull of \(X\) in the simplex, and its vertices are the extreme precincts.
\emph{Violation:} extrapolating beyond the data, or outlier precincts at the hull vertices, make the constraints bind spuriously, so bounds are too narrow or empty. The vertices are estimated (App. C, Remark 4*).
\emph{Testable:} sensitivity to trimming \(A\).

\textbf{A5. Rank / richness:} for every symmetric \(R\times R\) matrix \(D\),
\[x^\top Dx=0\ \ \forall x\in\operatorname{supp}_AP_X\ \Longrightarrow\ D=0,\]
i.e.~\(\operatorname{supp}_AP_X\) lies in no quadric of the simplex. Here \(\operatorname{supp}_AP_X:=\operatorname{supp}(P_X|_A)\) is the support of the restricted measure \(E\mapsto P(X\in A\cap E)\). It equals \(\operatorname{supp}P_X\) when \(P(X\in A)=1\), but can be smaller than \(\operatorname{supp}P_X\cap A\) otherwise (\(X\sim U[0,1]\), \(A=\{1/2\}\) gives the empty set). \emph{(Corrected after review, Batch B, Nov 2026: earlier stated on \(\operatorname{supp}P_X\cap A\); see 2a §0.1.)} This is a rank condition on the quadratic regression, not just ``variation in \(X\)''. For \(R=2\) it means at least 3 distinct values of \(X\). For \(R=3\) it needs six support points that do not lie on a common conic: three, four or five distinct compositions are never enough (2a Remark 1.1). A5 also implies that \(A\) is full-dimensional (2a Remark 5.2).
\emph{Role:} point-identifies \(Q\) (Lemma 1a).
\emph{Transitions:} fails, or nearly fails, when shares barely vary, e.g.~stable turnout across precincts for \(R=3\). This is JKST's ``gender'' problem.
\emph{Violation:} \(Q\) is weakly identified, the bound inherits huge variance, and \(\Theta_O\) is very wide or unstable.
\emph{Testable:} yes, through the condition number of the design matrix.

\textbf{A6. Source of randomness and independence (revised after review).} Every precinct and district is observed, so the uncertainty is not from sampling precincts. It is model-based: given the design \(\{(N_i,X_i,D_i)\}\) and the district effects, the precinct rates \(\beta_i\) are random with conditional mean \(m(V_i)\), and the residuals \(e_i=\beta_i-m(V_i)\) are independent across precincts (P3-S1a). Inference is conditional on the design, and the targets (\(B\) and \(\tilde B\), §1.7) refer to this electorate.
\emph{Role:} CLT for \(\hat Q\), and the unit-level allowance \(S_1\) for \(B-\tilde B\).
\emph{Transitions:} precincts in a constituency share candidates, campaigns and shocks. Shocks common to a district are fixed features of this electorate. Under M1/M2 they are the district effects, part of the mean. Under M0 they are not a variance problem but the misspecification described in 2b Prop. D4. No clustered variance is used. A superpopulation target averaging over re-drawn district shocks is not pursued.
\emph{Violation:} dependence \emph{within} a district beyond its district effect, e.g.~neighbourhood spillovers. The residual allowance and OLS variances are then too small. This is a stated limitation.
\emph{Testable:} partly: residual correlation between neighbouring precincts after M1.

\textbf{A7. Regularity for inference} (JKST tie-breaking): a unique active constraint at the optimum, \(\mu_{rs}>0\), and \(\Theta_O\) not a singleton.
\emph{Role:} the LP value function is differentiable at \(Q\), so the delta method and \(\mathrm{CI}_x\) are valid.
\emph{Violation:} at ties the bound is only directionally differentiable; the naive bootstrap and delta method fail, and Phase 3 needs Fang--Santos type inference. This is not a bias in identification.
\emph{Testable:} partly (estimated gaps between constraint values).

\textbf{A8 (M1). Homogeneous contextual effects across districts:} \(K\) (and the symmetric slope parts) are common to all \(d\); only intercepts vary.
\emph{Role:} this restriction is what produces the intersection \(\bigcap_d\Theta_O^{(d)}\), which is the entire gain (D0, D3).
\emph{Transitions:} how a precinct's previous A-share shifts A-voters' loyalty is the same in every constituency, while constituency-level popularity shifts (candidate quality, local swing) are free.
\emph{Violation:} e.g.~a strong incumbent changes the contextual gradient in one district. The intersection can exclude the truth, often with \emph{falsely narrow} bounds, because the max/min across districts picks the misspecified district.
\emph{Testable:} yes, by the overidentification (D2) and emptiness.

\textbf{A9 (M1/M2). Within-district richness.} Under M2, A5 within each district. Under M1, each district's support must be affinely spanning (\(R\) affinely independent compositions), and the stacked design must have full column rank: pooled quadratic monomials plus the district × \(X_r\) shifts for \(d\ge2\). This is less than A5 within every district (2b Remark D1.1). \emph{(Revised after review: ``at least \(R\) distinct values'' is not enough; they must be affinely independent.)}
\emph{Role:} identifies \(\zeta_{c,d}\) (M1) or \(Q^{(d)}\) (M2).
\emph{Transitions:} districts with only a few, similar precincts identify almost nothing.
\emph{Violation:} weak identification per district. Combined with \(\max_d\)/\(\min_d\) this biases the bounds \emph{inward}.
\emph{Testable:} yes.

\textbf{A10. Asymptotic regime (a choice, not a fact):} (R1) \(D\) fixed, \(p_d\to\infty\); (R2) \(D\to\infty\), \(p_d\) bounded; (R3) both grow.
\emph{Role:} under R2 the \(\zeta_d\) are incidental parameters that cannot be estimated consistently, and \(\max_d\hat\ell_{c,d}\) is biased upward. This needs intersection-bounds methods with precision correction (Chernozhukov, Lee \& Rosen 2013) or a hierarchical restriction on \(\zeta\).
\emph{Transitions:} the answer depends on your design (see the questions at the end).
\emph{Coverage of each regime (after review; full table in P3 §0.5):}

\begin{itemize}
\tightlist
\item
  M0: fixed-dimension asymptotics in every regime.
\item
  M1/M2 in R1: fixed-dimension asymptotics.
\item
  M1/M2 in R3: high-dimensional Gaussian approximation under a rate condition.
\item
  Every regime, including R2: a finite-sample confidence interval (P3 Theorem P3.8). It is valid, but uninformative in width in the designs tried (P3 §5.1).
\item
  CLR coverage in R2: a simulation finding only.
\end{itemize}

\textbf{Assumptions explicitly not made:} no distributional form for \(\beta\) (unlike King's truncated bivariate normal), no independence between \(\beta\) and \(X\) (contextual effects are the point), no stationarity over time, and no point-identifying nonlinearity.

\textbf{Optional shape restrictions} (each is a convex constraint on \(K\), so the bounds stay LPs):

\begin{itemize}
\tightlist
\item
  \textbf{Sign or monotone contextual effects}, e.g.~``A-voters are no less loyal where A was stronger'': \(\partial_{x_A}E[\beta_{A,A}\mid x]\ge0\), a half-space in \(K\). The bounds shrink; if the sign is wrong, the truth can be excluded.
\item
  \textbf{Bounded heterogeneity} \(|K^{(d)}-K|\le\delta\) (between M1 and M2): a sensitivity parameter that interpolates between the two (Phase 3).
\item
  \textbf{Bhatia--Davis variance bound} \(\operatorname{Var}(\beta)\le m(1-m)\le\tfrac14\): tightens \(S_1\) (Phase 3); valid without extra assumptions.
\end{itemize}

\begin{center}\rule{0.5\linewidth}{0.5pt}\end{center}

\hypertarget{pA-notation-dictionary}{%
\subsection{3. Notation dictionary}\label{pA-notation-dictionary}}

Also available as a LaTeX \texttt{tabular} in \path{docs/extension/notation.tex}.

\begin{longtable}[]{@{}
  >{\raggedright\arraybackslash\hspace{0pt}}p{(\columnwidth - 2\tabcolsep) * \real{0.5000}}
  >{\raggedright\arraybackslash\hspace{0pt}}p{(\columnwidth - 2\tabcolsep) * \real{0.5000}}@{}}
\toprule\noalign{}
\begin{minipage}[b]{\linewidth}\raggedright\hspace{0pt}Symbol
\end{minipage} & \begin{minipage}[b]{\linewidth}\raggedright\hspace{0pt}Meaning
\end{minipage} \\
\midrule\noalign{}
\endhead
\bottomrule\noalign{}
\endlastfoot
\(i=1,\dots,p\) & unit (precinct) index \\
\(d=1,\dots,D\); \(D_i\); \(\mathcal I_d\); \(p_d\) & district index; unit \(i\)'s district; units in \(d\); their number \\
\(r,s=1,\dots,R\) & row categories (election-1 choice / group) \\
\(c=1,\dots,C\) & column categories (election-2 choice) \\
\(N_i\), \(N_{ir}\), \(N_{irc}\) & unit size; row count; latent cell count \\
\(X_i\in\Delta^{R-1}\) & observed row shares \(N_{ir}/N_i\) \\
\(T_i\in\Delta^{C-1}\) & observed column shares \\
\(Z_i\) & additional covariates (e.g.~\(\log N_i\)) \\
\(\beta_i=[\beta_{irc}]\in\mathfrak B=(\Delta^{C-1})^R\) & latent transition matrix of unit \(i\) \\
\(W_i=(N_i,X_i,T_i,Z_i,D_i)\) & observables of unit \(i\) \\
\(P\) & law of \(W\) (observable); \(\mathbb Q\): law of \((W,\beta)\) \\
\((\Omega,\mathcal F)\) & underlying measurable space \\
\(\mathcal G=\sigma(X,N)\) (or \(\sigma(X,N,D)\); covariates \(Z\) were removed after review, R-15) & conditioning σ-algebra of the contextual model \\
\(\mathcal O_p=\sigma(W_1,\dots,W_p)\) & σ-algebra of the sample \\
\(B_{rc}(\mathcal S)\) & finite-population transition rate over units \(\mathcal S\) (1.2) \\
\(B^*_{rc}\) & superpopulation analogue (1.3) \\
\(\varphi=\sum a_{rc}B_{rc}\) & generic linear target (loyalty \(\lambda_r\), defection \(\delta_r\), demobilisation \(\mu_r\)) \\
\(A\subseteq\Delta^{R-1}\), \(\operatorname{vert}(A)\) & region where LCE is assumed; its vertices (\(R=2\): \(\{l,u\}\)) \\
\(H_c\in\mathbb R^{R\times R}\) & homogeneous LCE coefficients, column \(c\) (1.4) \\
\(Q_c=\operatorname{sym}(H_c)\) & identified part, from \(E[T_c\mid X]=X^\top Q_cX\) \\
\(K_c=\operatorname{skew}(H_c)\); \(\mathcal K\) & unidentified part; its space, \(\dim=(C-1)R(R-1)/2\) \\
\(\kappa_c=K_{c,12}\) & the unidentified scalar per column when \(R=2\) \\
\(M_H(x)=[e_r^\top H_cx]_{rc}\) & implied conditional-mean transition matrix \\
\(\tau_c(x)=E[T_c\mid X=x]\) & regression function (quadratic under A2) \\
\(\mu_r=E[NX_r]\), \(\mu_{rs}=E[NX_rX_s]\) & weighting moments; \(\mu_{rs}/\mu_r\) = sensitivity of \(B^*_{rc}\) to \(K_{c,rs}\) \\
\(S(x,t)\); \(h_S(u;x,t)\) & unit DD polytope; its support function \\
\(\mathbb E[S\mid\mathcal G]\) & conditional Aumann (selection) expectation \\
\(\Theta_I,\ \Theta_X,\ \Theta_O\) & sharp set; JKST-Prop.-1 set; outer (Jensen / vertex) set \\
\(\mathcal B^*_I\) & identified set for the estimand, \(\{B^*(K):K\in\Theta_I\}\) \\
\(\zeta_{c,d}\in\mathbb R^R\) & district-\(d\) intercept shifts (M1), identified \\
\(Q^{(d)}_c\), \(\Theta^{(d)}_O\), \(A_d\) & district-level identified part, constraint set, support \\
\(\mathrm{CI}_x\), \(S_1\), \(S_L\), \(S_U\) & JKST confidence interval and its standard-error terms \\
\(\Phi\) & standard normal CDF \\
\end{longtable}

\begin{center}\rule{0.5\linewidth}{0.5pt}\end{center}

\hypertarget{pA-proposed-refinements-of-the-phase-plan}{%
\subsection{4. Proposed refinements of the phase plan}\label{pA-proposed-refinements-of-the-phase-plan}}

The three-phase structure works, but goal 2 and the code need their own steps. I propose:

\begin{itemize}
\tightlist
\item
  \textbf{Phase 2a: R×C derivations.} Lemma 1; Prop. 2-RC with the vertex reduction and 2×2/JKST equivalence; Theorem 1 (sharpness of \(\Theta_I\) via selection expectations, with both directions of observational equivalence); the nesting (1.13) and the equality condition; the closed form for \(R=2\) and the joint-vs-collapsed corollary; non-emptiness (generalised App. B); the R×C version of JKST Example 6 (undetectable violation).
\item
  \textbf{Phase 2b: covariates and districts.} D0--D5 with proofs; sharpness of \(\bigcap_d\Theta^{(d)}\) under M1; a characterisation of \emph{when the gain is strict} in terms of \(\zeta_d\) and \(A_d\); the misspecification bias of M0 under M1, in closed form; M1′.
\item
  \textbf{Phase 3: inference and sensitivity.} Consistency at rate \(\sqrt p\) (or \(\sqrt{D}\) under clustering) of LP-value estimators; asymptotic distribution at a unique active vertex (normal) and at ties (max/min of Gaussians, directional delta method, Fang \& Santos 2019); generalised \(\mathrm{CI}_x\) (cover-the-parameter, Imbens--Manski/Stoye) vs cover-the-set; the many-districts problem (CLR 2013 intersection bounds); breakdown-frontier sensitivity for A1 (\(\epsilon\); later left open, see A1), A2 (sup-norm deviation \(\delta\) from linearity) and A8 (slope heterogeneity).
\item
  \textbf{Phase 4 (new): R implementation.} New exported functions (suggested names \texttt{eiBoundsRxC()} and \texttt{eiBoundsDistrict()}) that leave the existing API untouched. Closed forms for \(R=2\); a small LP otherwise. The package's dependencies are \texttt{sandwich}, \texttt{MASS} and \texttt{eco}; general \(R\) needs an LP solver (\texttt{lpSolve} is light), which is a dependency decision for you. Cluster-robust \texttt{vcovCL}. Regression test: for \(R=C=2\) the new code must reproduce \texttt{generateBounds} and \texttt{evaluateBounds} to \(10^{-10}\).
\item
  \textbf{Phase 5 (new): validation.} Monte Carlo (correctly specified, \(K\) at a vertex for tightness of coverage, and each assumption violated in turn); ground-truth data (\texttt{eco::census} for 2×2; \texttt{eiPack::senc}, NC party registration by race with known cells, for a true 3×3 check); width ratio and coverage tables as in JKST Table 1.
\end{itemize}

\hypertarget{pA-keeping-the-work-rigorous}{%
\subsection{5. Keeping the work rigorous}\label{pA-keeping-the-work-rigorous}}

\begin{enumerate}
\def\labelenumi{\arabic{enumi}.}
\tightlist
\item
  \textbf{Label every claim}: \emph{Theorem / Proposition} (proved), \emph{Conjecture}, or \emph{Heuristic} (e.g.~the revert-to-DD rule). Never call an outer set ``sharp''.
\item
  \textbf{State the measure-theoretic level explicitly}: conditional statements a.s.; the existence of measurable selections cited, not assumed; auxiliary randomisation stated.
\item
  \textbf{Cite exact theorems}, e.g.~Artstein (1983); Hiai \& Umegaki (1977); Molchanov (2017, \emph{Theory of Random Sets}, 2nd ed.); Beresteanu, Molchanov \& Molinari (2011); Molinari (2020, \emph{Handbook of Econometrics} 7A); Cross \& Manski (2002), the canonical partial-identification treatment of EI with covariates; Chernozhukov, Lee \& Rosen (2013); Imbens \& Manski (2004); Stoye (2009); Fang \& Santos (2019). Use Rosen, Jiang, King \& Tanner (2001) and Greiner \& Quinn (2009) for the R×C EI context. I will check each citation against the source text before relying on it.
\item
  \textbf{Check algebra mechanically}: verify each reparametrisation symbolically (sympy) and numerically, and solve every LP closed form two ways (vertex enumeration and a generic solver), as already done for the 2×2 equivalence.
\item
  \textbf{Nested-model sanity checks}: every R×C or district result must collapse to JKST when \(R=C=2\), \(D=1\).
\item
  \textbf{Adversarial simulation for every assumption}: for each assumption, a DGP that violates it, to confirm the predicted failure mode (direction and order of the error).
\item
  \textbf{Separate identification from inference}: Phase 2 statements concern population objects (\(P\) known). Sampling enters only in Phase 3.
\item
  \textbf{Justification tags}, as you require, will reference specific results (``{[}Lemma 1(a){]}'', ``{[}Hiai--Umegaki 1977, Thm 5.1{]}'') rather than generic labels wherever possible.
\end{enumerate}

\hypertarget{pA-questions-that-change-the-derivations}{%
\subsection{6. Questions that change the derivations}\label{pA-questions-that-change-the-derivations}}

\begin{enumerate}
\def\labelenumi{\arabic{enumi}.}
\tightlist
\item
  \textbf{Rows of the transition table.} Is the electorate restricted to election-1 voters of A and B (true 2×3, \(\dim\mathcal K=2\)), or do previous abstainers and new entrants form rows (3×3 or more, \(\dim\mathcal K=6\))? Are both margins on the same register (A1)?
\item
  \textbf{District design.} Roughly how many districts \(D\), and how many units per district \(p_d\)? Do units within a district vary in \(X\)? This decides between regimes R1 and R2 (A10), and whether Phase 3 needs intersection-bounds corrections.
\item
  \textbf{Targets.} Overall rates only, district-level rates \(B_{rc}(\mathcal I_d)\) too, or both?
\item
  \textbf{Priority.} I recommend implementing \(\Theta_O\) first (closed form / LP, extends the package directly), then a nonparametric estimate of the sharp \(\Theta_I\) as an option. Do you agree?
\item
  \textbf{Dependencies.} May the package take on \texttt{lpSolve}, needed only for \(R\ge3\)?
\end{enumerate}

\begin{center}\rule{0.5\linewidth}{0.5pt}\end{center}

\hypertarget{pA-decisions-after-first-review-2728-oct-2026}{%
\subsection{7. Decisions after first review (27--28 Oct 2026)}\label{pA-decisions-after-first-review-2728-oct-2026}}

\begin{itemize}
\tightlist
\item
  \textbf{General R×C} is the target. The running example is BC transitions, \textbf{3×3} (in-party voters, out-party voters, non-voters in each election), so \(\dim\mathcal K=(3-1)\cdot3=6\); \(X_i\) lives in the 2-simplex, and \(A\) is the convex hull of the observed polls.
\item
  \textbf{Target:} overall (province-wide) rates \(B_{rc}\) only. District-level rates are a by-product, not a goal.
\item
  \textbf{Data design.}

  \begin{itemize}
  \tightlist
  \item
    \emph{Stage 1:} 93 ridings as the units, so \(p=93\) and \(D=1\). Riding dummies are impossible at this stage (one unit per riding). Only M0 applies, and with \(p=93\) and six quadratic coefficients per column, A5 and A7 (precision) are the binding concerns.
  \item
    \emph{Stage 2:} about 10 polls per riding, so \(p\approx930\), \(D=93\), \(p_d\approx10\). This is regime \textbf{R2} (many districts, few units each). The incidental-parameter problem of A10 is first-order here, not a footnote (see the illustration below). Phase 3 must deliver a precision-corrected intersection (Chernozhukov--Lee--Rosen) and/or partial pooling of the riding intercepts.
  \end{itemize}
\item
  \textbf{Priority:} \(\Theta_O\) first, as an extension of the package; the nonparametric \(\Theta_I\) comes later as an option.
\item
  \textbf{Dependencies:} \texttt{lpSolve} is approved.
\end{itemize}

\hypertarget{pA-illustration-what-riding-dummies-do-illustrate_district_dummies.py}{%
\subsubsection{\texorpdfstring{7.1 Illustration: what riding dummies do (\texttt{illustrate\_district\_dummies.py})}{7.1 Illustration: what riding dummies do (illustrate\_district\_dummies.py)}}\label{pA-illustration-what-riding-dummies-do-illustrate_district_dummies.py}}

The script uses a 2×2 setup for transparency: 93 ridings × 10 polls, riding baselines correlated with the riding's mean \(X\), and contextual slopes common to all ridings (M1 true). Results over 300 replications:

\begin{longtable}[]{@{}
  >{\raggedright\arraybackslash\hspace{0pt}}p{(\columnwidth - 4\tabcolsep) * \real{0.3333}}
  >{\raggedright\arraybackslash\hspace{0pt}}p{(\columnwidth - 4\tabcolsep) * \real{0.3333}}
  >{\raggedright\arraybackslash\hspace{0pt}}p{(\columnwidth - 4\tabcolsep) * \real{0.3333}}@{}}
\toprule\noalign{}
\begin{minipage}[b]{\linewidth}\raggedright\hspace{0pt}Bound
\end{minipage} & \begin{minipage}[b]{\linewidth}\raggedright\hspace{0pt}Covers truth
\end{minipage} & \begin{minipage}[b]{\linewidth}\raggedright\hspace{0pt}Median width
\end{minipage} \\
\midrule\noalign{}
\endhead
\bottomrule\noalign{}
\endlastfoot
DD (model-free) & 1.00 & about 0.54 \\
Pooled M0 (current package) & 1.00 & 0.305 \\
Riding M1, \textbf{true} coefficients (identification only) & 1.00 & 0.244 \\
Riding M1, \textbf{estimated} from 10 polls per riding & \textbf{0.72} & 0.172 \\
\end{longtable}

\begin{itemize}
\tightlist
\item
  \textbf{Identification gain (row 3 vs row 2):} about a 20\% narrower width, because each riding constrains the common contextual effect from its own baseline.
\item
  \textbf{Estimation problem (row 4):} taking the max/min over 93 noisy riding-level constraints makes the bounds \emph{too narrow} and they miss the truth 28\% of the time. The plain plug-in intersection must not be used at 10 polls per riding; this is what Phase 3 has to fix.
\item
  With step-shaped riding effects (pooled model misspecified), the pooled bound still covered the truth here (0.997). The robustness argument (D4) is a theoretical guarantee, but in these designs its practical value was small compared with the tightening.
\end{itemize}

\hypertarget{pA-what-h-and-k-mean-and-what-m1-shares-refines-a8}{%
\subsubsection{\texorpdfstring{7.2 What \(H\) and \(K\) mean, and what M1 shares (refines A8)}{7.2 What H and K mean, and what M1 shares (refines A8)}}\label{pA-what-h-and-k-mean-and-what-m1-shares-refines-a8}}

\begin{itemize}
\tightlist
\item
  \(E[\beta_{rc}\mid x]=\sum_sH_{c,rs}x_s\), so \(H_{c,rs}\) is the rate at which row-\(r\) voters go to column \(c\) \textbf{in a (hypothetical) poll made up entirely of row-\(s\) voters}. Diagonal entries are ``own-territory'' rates. Off-diagonal entries are rates of \(r\)-voters in \(s\)-territory, and the contextual effect for row \(r\) is how \(H_{c,r\cdot}\) varies across \(s\).
\item
  \(Q_{c,rs}=\tfrac12(H_{c,rs}+H_{c,sr})\) is identified. \(K_{c,rs}=\tfrac12(H_{c,rs}-H_{c,sr})\), half the difference between ``\(r\)-voters in \(s\)-territory'' and ``\(s\)-voters in \(r\)-territory'', is not. \(K\) is therefore the \textbf{unidentified part of the contextual effect, not the whole of it}.
\item
  \textbf{M1 shares all of \(H\) except an additive riding shift:} \(E[\beta_{rc}\mid x,d]=\zeta_{c,d,r}+\sum_sH_{c,rs}x_s\). Equivalently, there is no interaction between riding and poll composition on the probability scale. The shared symmetric part is testable (D2); the shared \(K\) is not, and it is the substantive leap.
\item
  Additivity on the probability scale is least plausible near 0 or 1 (ceiling and floor compression). A weaker option for Phase 2b is to let \(Q^{(d)}\) be free and share only \(K\), or to allow \(\|K^{(d)}-K\|\le\delta\) as a sensitivity parameter. The free-\(Q^{(d)}\) version needs 6 coefficients per column per riding in 3×3, which 10 polls cannot support.
\end{itemize}

\hypertarget{pA-worked-22-example-m0-vs-m1-example_2x2_ridings.py}{%
\subsubsection{\texorpdfstring{7.3 Worked 2×2 example: M0 vs M1 (\texttt{example\_2x2\_ridings.py})}{7.3 Worked 2×2 example: M0 vs M1 (example\_2x2\_ridings.py)}}\label{pA-worked-22-example-m0-vs-m1-example_2x2_ridings.py}}

Rows: \(P\) = voted P last time (share \(x\)), \(O\) = everyone else. Outcome: votes P now. \emph{Loyalty} is the rate for P-voters; \emph{crossover} is the rate for O-voters. The example is noise-free, so it shows identification only.

\begin{longtable}[]{@{}
  >{\raggedright\arraybackslash\hspace{0pt}}p{(\columnwidth - 4\tabcolsep) * \real{0.3333}}
  >{\raggedright\arraybackslash\hspace{0pt}}p{(\columnwidth - 4\tabcolsep) * \real{0.3333}}
  >{\raggedright\arraybackslash\hspace{0pt}}p{(\columnwidth - 4\tabcolsep) * \real{0.3333}}@{}}
\toprule\noalign{}
\begin{minipage}[b]{\linewidth}\raggedright\hspace{0pt}\end{minipage} & \begin{minipage}[b]{\linewidth}\raggedright\hspace{0pt}M0: one model for the province
\end{minipage} & \begin{minipage}[b]{\linewidth}\raggedright\hspace{0pt}M1: riding dummies
\end{minipage} \\
\midrule\noalign{}
\endhead
\bottomrule\noalign{}
\endlastfoot
Loyalty in a poll & \(a_P+g_Px\) & \(a_{P,d}+g_Px\) \\
Crossover in a poll & \(a_O+g_Ox\) & \(a_{O,d}+g_Ox\) \\
Same everywhere & all four numbers & the slopes \(g_P,g_O\) only \\
Riding matters & only through its polls' \(x\) & shifts each group's baseline \\
Data reveal & \(T=a_O+(a_P-a_O+g_O)x+(g_P-g_O)x^2\), pooled & the same, riding by riding, with common \(x^2\) term \\
Unknown & \(g_O\) (one number) & \(g_O\) (one number, shared) \\
Testable & cubic term, empty bound & plus: \(x^2\) term equal across ridings \\
Not testable & how the curvature splits into \(g_P,g_O\) & that the split is the same in every riding \\
\end{longtable}

In words, M1 says that \textbf{the gap between two ridings' loyalty is the same in every kind of poll} (parallel lines), and likewise for crossover.

Truth: \(g_P=g_O=0.10\). Riding 1 (P stronghold, polls \(x\in[0.5,0.9]\)): loyalty \(0.89+0.1x\), crossover \(0.20+0.1x\). Riding 2 (P weak, \(x\in[0.1,0.5]\)): loyalty \(0.45+0.1x\), crossover \(0.02+0.1x\).

\begin{itemize}
\tightlist
\item
  Riding 1's polls reveal \(T=0.20+0.79x\). Loyalty at \(x=0.5\) cannot exceed 1, so \(g\ge-0.02\). Riding 2's polls reveal \(T=0.02+0.53x\). Loyalty at \(x=0.1\) cannot fall below 0, so \(g\le0.611\).
\item
  \textbf{M1:} \(g\in[-0.020,0.611]\), with riding 1 setting the lower limit and riding 2 the upper. This gives \(B\in[0.618,0.866]\) (truth \(0.819\); DD \([0.453,0.879]\)).
\item
  \textbf{M0} fits one curve through both ridings: \(T=-0.095+0.936x+0.305x^2\). It reads the stronghold's higher baselines as a contextual effect, and the implied bound on the slope is empty (\([1.160,0.969]\)). The package's rule then reverts to DD. Here M0's failure was \emph{detectable}; it need not always be.
\end{itemize}


\clearpage
\section{PHASE2a: R$\times$C bounds, derivations and proofs}\label{app:phase2a}
\textit{Source: \texttt{docs/extension/PHASE2a\_rxc.md}.}\par
% GENERATED by tools/make_appendices.py from docs/extension/PHASE2a_rxc.md
% Do not edit by hand; edit the source and regenerate.
\emph{Builds on \texttt{PHASE1\_foundations.md} (cited as P1). Every proof step ends with a justification tag in brackets; the tag legend is in §0.3. Every closed form and every inclusion claimed here is also checked numerically or symbolically in \texttt{verify\_phase2a.py} (output in \texttt{verify\_phase2a\_output.txt}; 32/32 PASS). §10 lists the corrections this phase makes to P1.}

\begin{center}\rule{0.5\linewidth}{0.5pt}\end{center}

\hypertarget{pB-setting-conventions-and-tags}{%
\subsection{0. Setting, conventions and tags}\label{pB-setting-conventions-and-tags}}

\hypertarget{pB-objects}{%
\subsubsection{0.1 Objects}\label{pB-objects}}

\begin{itemize}
\tightlist
\item
  \(\mathcal W=\mathbb N\times\Delta^{R-1}\times\Delta^{C-1}\) is the space of one unit's observables \(W=(N,X,T)\); it is Polish. \(V:=(X,N)\) and \(\mathcal G:=\sigma(V)\). Districts come in Phase 2b. Other covariates are a planned extension (P1 §1.6); here the mean model depends on \(X\) only.
\item
  \(\mathfrak B=(\Delta^{C-1})^R\subset\mathbb R^{R\times C}\), with \(\langle u,b\rangle=\sum_{rc}u_{rc}b_{rc}\) and the Frobenius norm. Every \(b\in\mathfrak B\) satisfies \(\|b\|\le\sqrt R\).
\item
  \(J=\mathbf 1\mathbf 1^\top\). \(\operatorname{Skew}_R=\{K\in\mathbb R^{R\times R}:K^\top=-K\}\).
\item
  \(\mathcal K=\{(K_1,\dots,K_C)\in\operatorname{Skew}_R^C:\sum_cK_c=0\}\).
\item
  For \(H=(H_1,\dots,H_C)\), \(M_H(x)\in\mathbb R^{R\times C}\) has entries \(M_H(x)_{rc}=e_r^\top H_cx\).
\item
  \(A\subseteq\Delta^{R-1}\) is compact and convex. The default is \(A=\operatorname{conv}(\operatorname{supp}P_X)\), the \textbf{population} hull. It is not the sample hull: P1 §1.7 separates the population hull, the hull of the empirical law, and the sample hull used as an estimate. A \emph{polytope} \(A\) has a finite vertex set \(\operatorname{vert}(A)\), and that is always the case for the sample hull.
\item
  \textbf{Support on \(A\).} \(\operatorname{supp}_AP_X:=\operatorname{supp}(P_X|_A)\) is the support of the restricted measure \(E\mapsto P(X\in A\cap E)\): the smallest closed set \(F\) with \(P(X\in A\setminus F)=0\). It is contained in \(A\cap\operatorname{supp}P_X\) and can be strictly smaller: for \(X\sim U[0,1]\) (as a share) and \(A=\{1/2\}\), \(\operatorname{supp}P_X\cap A=\{1/2\}\) but \(\operatorname{supp}_AP_X=\emptyset\). When \(P(X\in A)=1\) (the default \(A\), and every result that uses Lemma 2), \(P_X|_A=P_X\) and \(\operatorname{supp}_AP_X=\operatorname{supp}P_X\). \emph{(Corrected after review, Batch B, Nov 2026: earlier versions used \(\operatorname{supp}P_X\cap A\) throughout. Continuous functions that agree \(P_X\)-a.e. on \(A\) need not agree on that set, as the example shows; they agree on \(\operatorname{supp}_AP_X\). A5, the tag {[}supp{]}, Lemma 1, Lemma 3(c), Remarks 1.1, 5.1--5.2, Props 7′ and 9 now use \(\operatorname{supp}_AP_X\). Nothing changes when \(P(X\in A)=1\). Check: \texttt{verify\_phase2a.py}, ``{[}supp{]}''.)}
\item
  \textbf{Standing integrability:} \(E[N]<\infty\), and \(\mu_r:=E[NX_r]>0\) for all \(r\). \(\mu_{rs}:=E[NX_rX_s]\).
\end{itemize}

\hypertarget{pB-model-identified-sets-these-refine-p1-1.5}{%
\subsubsection{0.2 Model, identified sets (these refine P1 §1.5)}\label{pB-model-identified-sets-these-refine-p1-1.5}}

\textbf{Model \(\mathcal M(H;A)\):} the set of laws \(\mathbb Q\) of \((W,\beta)\) on \(\mathcal W\times\mathfrak B\) such that

\begin{itemize}
\tightlist
\item
  (S1) \(\beta\in\mathfrak B\) and \(T=\beta^\top X\), \(\mathbb Q\)-a.s. \emph{(A1, accounting identity)};
\item
  (S2) \(E_{\mathbb Q}[\beta\mid V]=M_H(X)\) \(\mathbb Q\)-a.s. on \(\{X\in A\}\) \emph{(A2 + A3, linear contextual effects, no dependence on \(N\))};
\item
  (S3) \(M_H(x)\in\mathfrak B\) for every \(x\in A\) \emph{(A4: the linear model describes valid probabilities on all of \(A\). Lemma 3(c) shows this follows from (S1)--(S2) when \(A\subseteq\operatorname{conv}(\operatorname{supp}_AP_X)\), in particular for the default \(A=\operatorname{conv}\operatorname{supp}P_X\). It is a genuine extra assumption whenever some part of \(A\) lies outside the convex hull of the support points }at which (S2) is imposed\emph{: when \(A\) extrapolates beyond the data, as in JKST Prop. 3, and also when \(A\) is a subset of the hull that the data barely touch (Remark 5.1))}.
\end{itemize}

\textbf{Identified part.} \(Q(P)=(Q_1,\dots,Q_C)\) denotes the unique symmetric matrices with \(E_P[T_c\mid X=x]=x^\top Q_cx\) for \(P_X\)-a.e. \(x\in A\), when they exist (Lemma 1). If they do not exist, set \(\Theta_I(P;A):=\emptyset\).

\textbf{Sharp identified set:}
\[\Theta_I(P;A)=\{K\in\mathcal K:\ \exists\,\mathbb Q\in\mathcal M(Q(P)+K;A),\ \mathbb Q\circ W^{-1}=P\}.\]

\textbf{\(\Theta_X(P;A)\)} is the same with (S2) replaced by the weaker \(E_{\mathbb Q}[\beta\mid X]=M_H(X)\); this is JKST Prop. 1.

\textbf{Outer set:}
\[\Theta_O(P;A)=\{K\in\mathcal K:\ (Q_c(P)+K_c)x\ge0\ \text{elementwise}\ \ \forall c,\ \forall x\in A\}.\]

\textbf{Unit Duncan--Davis polytope and its support function:}
\[S(x,t)=\{b\in\mathfrak B:b^\top x=t\},\qquad h_S(u;x,t)=\max_{b\in S(x,t)}\langle u,b\rangle.\]
\(P_{T\mid V}\) is a regular conditional distribution of \(T\) given \(V\) under \(P\), and
\[\psi(u;v):=\int h_S(u;x,t)\,P_{T\mid V=v}(dt).\]

\textbf{Estimand:} \(B^*_{rc}=E[NX_r\beta_{rc}]/\mu_r\), and \(\varphi=\sum_{rc}a_{rc}B^*_{rc}\).

\hypertarget{pB-tag-legend}{%
\subsubsection{0.3 Tag legend}\label{pB-tag-legend}}

Theorem numbers marked † are cited from memory; §0.4 records which have been checked, and against what. The mathematical content of each is standard.

\begin{longtable}[]{@{}
  >{\raggedright\arraybackslash}p{(\columnwidth - 2\tabcolsep) * \real{0.5000}}
  >{\raggedright\arraybackslash}p{(\columnwidth - 2\tabcolsep) * \real{0.5000}}@{}}
\toprule\noalign{}
\begin{minipage}[b]{\linewidth}\raggedright
Tag
\end{minipage} & \begin{minipage}[b]{\linewidth}\raggedright
Meaning / source
\end{minipage} \\
\midrule\noalign{}
\endhead
\bottomrule\noalign{}
\endlastfoot
{[}Def{]} & the definition just introduced or referenced \\
{[}S1{]}, {[}S2{]}, {[}S3{]}, {[}A1{]}--{[}A5{]} & model conditions §0.2, assumptions P1 §2 \\
{[}lin{]} & linearity (of expectation, inner product, matrix products) \\
{[}alg{]} & elementary algebra, shown in the same line \\
{[}CE-lin{]}, {[}CE-mono{]}, {[}CE-pull{]}, {[}CE-tower{]} & conditional expectation: linearity, monotonicity, ``taking out what is known'', tower property \\
{[}CE-cvx{]} & if \(Y\in F\) a.s. with \(F\) closed convex and \(Y\) integrable, then \(E[Y\mid\mathcal G]\in F\) a.s. (conditional Jensen / separation) \\
{[}RCD{]} & existence of regular conditional distributions on Polish spaces, and disintegration: \(E[f(V,T)\mid V]=\int f(V,t)P_{T\mid V}(dt)\) a.s. for bounded jointly measurable \(f\) (Kallenberg 2002, Thms 6.3--6.4†) \\
{[}supp{]} & if continuous \(f,g\) agree \(P_X\)-a.e. on \(A\), they agree on \(\operatorname{supp}_AP_X\) (§0.1). \emph{Proof:} if \(f(x_0)\ne g(x_0)\) at \(x_0\in\operatorname{supp}_AP_X\), then \(f\ne g\) on an open \(U\ni x_0\), and \(P(X\in A\cap U)>0\) by the definition of the support, a contradiction. The weaker set \(\operatorname{supp}P_X\cap A\) does not work (§0.1 example). \emph{(Corrected after review, Batch B.)} \\
{[}Frob{]} & \(\lvert\langle u,b\rangle\rvert\le\lVert u\rVert\,\lVert b\rVert\) (Cauchy--Schwarz) \\
{[}Mink{]} & a polytope is the convex hull of its vertices (Minkowski / Krein--Milman for polytopes) \\
{[}Rock{]} & closed convex sets are determined by their support functions (Rockafellar 1970, Cor. 13.1.1†) \\
{[}Berge{]} & Berge maximum theorem: max of a continuous function over a compact-valued u.h.c. correspondence is u.s.c. \\
{[}MMT{]} & measurable maximum theorem: a measurable argmax selector exists (Aliprantis \& Border 2006, Thm 18.19†) \\
{[}DP{]} & Dunford--Pettis: uniformly integrable subsets of \(L^1\) are relatively weakly compact \\
{[}Mazur{]} & convex norm-closed subsets of a normed space are weakly closed \\
{[}HB{]} & Hahn--Banach strict separation of a point from a closed convex set, with \((L^1)^*=L^\infty\) for a probability space \\
{[}wk-cont{]} & bounded linear operators are weak-to-weak continuous \\
{[}Riesz{]} & \(L^1\) convergence implies a.s. convergence along a subsequence \\
{[}FH{]} & Fréchet--Hoeffding bounds for a cell of a table with given margins \\
{[}poly{]} & a polynomial in \(k\) variables vanishing on a nonempty open set is the zero polynomial \\
{[}IVT{]} & intermediate value theorem / the image of a convex set under a linear map is convex \\
{[}exch{]} & exchange argument for linear objectives \\
{[}L0{]}--{[}L4{]}, {[}T1{]}, {[}P1{]}--{[}P9{]} & results of this document \\
\end{longtable}

\hypertarget{pB-citation-check-after-review}{%
\subsubsection{0.4 Citation check (after review)}\label{pB-citation-check-after-review}}

Status:
* \textbf{checked (text)}: the theorem number was read in the source text, in the version given.
* \textbf{standard}: a textbook result whose number was checked from memory only; confirm against the printed edition before submission.
* \textbf{to replace}: the number cited earlier was wrong.

\begin{longtable}[]{@{}
  >{\raggedright\arraybackslash}p{(\columnwidth - 6\tabcolsep) * \real{0.2500}}
  >{\raggedright\arraybackslash}p{(\columnwidth - 6\tabcolsep) * \real{0.2500}}
  >{\raggedright\arraybackslash}p{(\columnwidth - 6\tabcolsep) * \real{0.2500}}
  >{\raggedright\arraybackslash}p{(\columnwidth - 6\tabcolsep) * \real{0.2500}}@{}}
\toprule\noalign{}
\begin{minipage}[b]{\linewidth}\raggedright
Tag
\end{minipage} & \begin{minipage}[b]{\linewidth}\raggedright
Statement used
\end{minipage} & \begin{minipage}[b]{\linewidth}\raggedright
Reference and location
\end{minipage} & \begin{minipage}[b]{\linewidth}\raggedright
Status
\end{minipage} \\
\midrule\noalign{}
\endhead
\bottomrule\noalign{}
\endlastfoot
{[}RCD{]} & regular conditional distributions on Polish spaces; disintegration & Kallenberg (2002), \emph{Foundations of Modern Probability}, 2nd ed., Thms 6.3 and 6.4 & standard \\
{[}Rock{]} & closed convex sets are ordered by their support functions & Rockafellar (1970), \emph{Convex Analysis}, Cor. 13.1.1 & standard \\
{[}Berge{]} & Berge maximum theorem & Aliprantis \& Border (2006), \emph{Infinite Dimensional Analysis}, 3rd ed., Thm 17.31 & standard \\
{[}MMT{]} & measurable maximum theorem (Carathéodory objective, weakly measurable compact-valued correspondence) & Aliprantis \& Border (2006), 3rd ed., Thm 18.19 & standard \\
{[}DP{]} & Dunford--Pettis: uniformly integrable ⇒ relatively weakly compact in \(L^1\) & Dunford \& Schwartz (1958), \emph{Linear Operators I}, Cor. IV.8.11; or Bogachev (2007), \emph{Measure Theory}, Thm 4.7.18 & standard \\
{[}Mazur{]} & convex norm-closed sets are weakly closed & Brezis (2011), \emph{Functional Analysis, Sobolev Spaces and PDEs}, Thm 3.7 & standard \\
{[}HB{]} & strict separation of a point from a closed convex set; \((L^1)^*=L^\infty\) on a probability space & Rudin (1991), \emph{Functional Analysis}, 2nd ed., Thm 3.4(b); Folland (1999), \emph{Real Analysis}, 2nd ed., Thm 6.15 & standard \\
{[}wk-cont{]} & norm-continuous linear maps are weak-to-weak continuous & Brezis (2011), Thm 3.10 & standard \\
{[}Riesz{]} & \(L^p\) convergence implies a.s. convergence along a subsequence & Brezis (2011), Thm 4.9 & standard \\
{[}Gale{]} (2b) & supply--demand feasibility of transportation problems & Gale (1957), \emph{Pacific J. Math.} 7:1073--1082, Thm 1 & standard \\
{[}dDM{]} (P3) & delta method for Hadamard directionally differentiable maps & Fang \& Santos (2019), \emph{Rev.~Econ. Stud.} 86:377--412, Thm 2.1; Shapiro (1991), \emph{Ann. Oper. Res.} 30:169--186 & checked (text): arXiv:1404.3763v2 (2016), Thm 2.1 \\
Fang--Santos bootstrap (P3) & standard bootstrap consistent iff the map is fully differentiable tangentially to the support of the Gaussian limit; non-Gaussian limit ⇒ inconsistent & Fang \& Santos (2019), Thm 3.1 and Cor. 3.1 & checked (text): arXiv v2, Thm 3.1, Cor. 3.1; confirm in the published version \\
{[}CLR{]} (P3) & analytic/simulation intersection-bound inference with a parametrically estimated bounding function over a compact (here finite) index set & Chernozhukov, Lee \& Rosen (2013), \emph{Econometrica} 81:667--737: Condition P, Lemma 4, \textbf{Thm 4} (via Thms 1--2); finite index set: Example 1 & \textbf{to replace}: P3 cited ``Thm 2''; the finite-index parametric case is Thm 4. Checked (text): arXiv:0907.3503v5 (2013) \\
{[}Lyap{]} (P3) & Lyapunov CLT for independent, non-identically distributed summands & Billingsley (1995), \emph{Probability and Measure}, 3rd ed., Thm 27.3 & standard \\
{[}BD{]} (P3) & Bhatia--Davis variance bound & Bhatia \& Davis (2000), \emph{Amer. Math. Monthly} 107:353--357 & standard \\
Stoye (P3) & uniform coverage of the parameter near point identification & Stoye (2009), \emph{Econometrica} 77:1299--1315 & standard \\
\end{longtable}

CLR's Condition P(iv) requires that the gradient of every candidate bounding function be bounded away from zero, and that the limiting covariance be positive definite. In our setting that means using \(C-1\) columns (P3 Remark P3.4.1), a full-rank design, and \(g_\lambda\ne0\) for every candidate.

\begin{center}\rule{0.5\linewidth}{0.5pt}\end{center}

\hypertarget{pB-the-unit-duncandavis-polytope}{%
\subsection{1. The unit Duncan--Davis polytope}\label{pB-the-unit-duncandavis-polytope}}

\textbf{Lemma 0.} For all \(x\in\Delta^{R-1}\) and \(t\in\Delta^{C-1}\):
(a) \(S(x,t)\neq\emptyset\);
(b) \(S(x,t)\) is compact and convex;
(c) \((x,t)\mapsto S(x,t)\) has closed graph, is upper hemicontinuous and Borel measurable;
(d) \(h_S(u;x,t)\) is sublinear in \(u\), \(\sqrt R\)-Lipschitz in \(u\), concave in \(t\), and Borel in \((u,x,t)\);
(e) \(S(x,t)=\{b:\langle u,b\rangle\le h_S(u;x,t)\ \forall u\}\).

\emph{Proof.}
1. Put \(b_{rc}:=t_c\). Each row is \(t\in\Delta^{C-1}\), so \(b\in\mathfrak B\). {[}Def{]}
2. \((b^\top x)_c=\sum_rx_rt_c=t_c\), using \(\sum_r x_r=1\). {[}alg{]} So (a) holds. {[}1--2{]}
3. \(\mathfrak B\) is a finite product of simplices, hence compact and convex. {[}Def{]}
4. \(\{b:b^\top x=t\}\) is an affine subspace, hence closed and convex. {[}lin{]}
5. \(S(x,t)\) is the intersection of 3 and 4, hence compact and convex. That is (b). {[}3--4{]}
6. The graph \(\{(x,t,b):b\in\mathfrak B,\ b^\top x=t\}\) is the zero set of a continuous map intersected with a compact set, hence closed. {[}alg{]}
7. A closed-graph correspondence into the compact \(\mathfrak B\) is upper hemicontinuous with compact values. Such correspondences are Borel measurable. That is (c). {[}6; Aliprantis \& Border 2006, Ch. 17--18†{]}
8. \(h_S\) is a maximum of the linear functions \(u\mapsto\langle u,b\rangle\), hence convex and positively homogeneous in \(u\), i.e.~sublinear. {[}Def; lin{]}
9. \(\lvert h_S(u)-h_S(u')\rvert\le\max_{b\in\mathfrak B}\lvert\langle u-u',b\rangle\rvert\le\sqrt R\,\lVert u-u'\rVert\). {[}Frob; \(\lVert b\rVert\le\sqrt R\){]}
10. If \(b\in S(x,t)\), \(b'\in S(x,t')\) and \(\lambda\in[0,1]\), then \(\lambda b+(1-\lambda)b'\in\mathfrak B\) and \((\lambda b+(1-\lambda)b')^\top x=\lambda t+(1-\lambda)t'\). {[}3; lin{]}
11. Taking \(b,b'\) to be maximisers, \(h_S(u;x,\lambda t+(1-\lambda)t')\ge\lambda h_S(u;x,t)+(1-\lambda)h_S(u;x,t')\): concavity in \(t\). {[}10; Def{]}
12. \(h_S\) is u.s.c. in \((u,x,t)\), hence Borel. That completes (d). {[}Berge; 7{]}
13. (e) holds because \(S(x,t)\) is closed and convex. {[}Rock; 5{]} ∎

\begin{center}\rule{0.5\linewidth}{0.5pt}\end{center}

\hypertarget{pB-what-the-data-identify}{%
\subsection{2. What the data identify}\label{pB-what-the-data-identify}}

\textbf{Lemma 1 (identification decomposition).} Let \(\mathbb Q\in\mathcal M(H;A)\) with \(\mathbb Q\circ W^{-1}=P\), and assume {[}A5{]}: for symmetric \(D\), \(x^\top Dx=0\) for all \(x\in\operatorname{supp}_AP_X\) implies \(D=0\). Then
(a) \(E[T_c\mid X=x]=x^\top Q_cx\) with \(Q_c=\tfrac12(H_c+H_c^\top)\), and \(Q_c\) is unique;
(b) \(\sum_cH_c=J\), so \(\sum_cQ_c=J\) and \(\sum_cK_c=0\) where \(K_c=\tfrac12(H_c-H_c^\top)\);
(c) for every \(K'\in\mathcal K\), \(Q+K'\) gives the same regression function; \(\dim\mathcal K=(C-1)R(R-1)/2\).

\emph{Proof of (a).}
1. \(T_c=\sum_rX_r\beta_{rc}\) a.s. {[}S1{]}
2. \(E[T_c\mid X]=\sum_rX_rE[\beta_{rc}\mid X]\) a.s. {[}CE-lin; CE-pull, since \(X\) is \(\sigma(X)\)-measurable{]}
3. \(E[\beta_{rc}\mid X]=E\big[E[\beta_{rc}\mid V]\mid X\big]=e_r^\top H_cX\) a.s. on \(\{X\in A\}\). {[}CE-tower; S2; \(M_H(X)\) is \(\sigma(X)\)-measurable{]}
4. Hence \(E[T_c\mid X]=\sum_{r,s}X_rH_{c,rs}X_s=X^\top H_cX\) on \(\{X\in A\}\). {[}2--3; alg{]}
5. Since \(X^\top K_cX\) is a scalar, \(X^\top K_cX=(X^\top K_cX)^\top=X^\top K_c^\top X=-X^\top K_cX\), so \(X^\top K_cX=0\). {[}alg; \(K_c\) skew{]}
6. Therefore \(X^\top H_cX=X^\top Q_cX\). {[}4--5; \(H_c=Q_c+K_c\){]}
7. If symmetric \(Q_c,Q'_c\) both satisfy 6 \(P_X\)-a.e. on \(A\), then \(x^\top(Q_c-Q'_c)x=0\) on \(\operatorname{supp}_AP_X\). {[}supp; both sides continuous{]}
8. \(Q_c-Q'_c\) is symmetric, so \(Q_c=Q'_c\). {[}A5{]} ∎(a)

\emph{Proof of (b).}
1. \(\sum_c\beta_{rc}=1\) a.s. {[}S1; \(\beta\in\mathfrak B\){]}
2. \(\sum_ce_r^\top H_cx=1\) for \(P_X\)-a.e. \(x\in A\). {[}CE-lin on 1; S2; CE-tower{]}
3. Let \(d_r^\top:=e_r^\top(\sum_cH_c-J)\). Then \(d_r^\top x=1-\mathbf 1^\top x=0\) on \(\operatorname{supp}_AP_X\). {[}2; \(\mathbf 1^\top x=1\); supp{]}
4. \(x^\top\tfrac12(d_r\mathbf 1^\top+\mathbf 1d_r^\top)x=(d_r^\top x)(\mathbf 1^\top x)=0\) on the same set. {[}alg; 3{]}
5. So \(d_r\mathbf 1^\top+\mathbf 1d_r^\top=0\). {[}A5; the matrix is symmetric{]}
6. Its \((s,s)\) entry is \(2d_{r,s}=0\) for every \(s\), so \(d_r=0\). {[}alg{]}
7. Hence \(\sum_cH_c=J\). Taking symmetric and skew parts, \(\sum_cQ_c=J\) and \(\sum_cK_c=0\). {[}6; \(J\) symmetric; lin{]} ∎(b)

\emph{Proof of (c).}
1. For \(K'\in\mathcal K\): \(x^\top(Q_c+K'_c)x=x^\top Q_cx\). {[}(a) step 5{]}
2. \((K_1,\dots,K_{C-1})\mapsto(K_1,\dots,K_{C-1},-\sum_{c<C}K_c)\) is a linear bijection from \(\operatorname{Skew}_R^{C-1}\) onto \(\mathcal K\). {[}Def; lin{]}
3. \(\dim\operatorname{Skew}_R=R(R-1)/2\), so \(\dim\mathcal K=(C-1)R(R-1)/2\). {[}2{]} ∎

\textbf{Remark 1.0 (why A2 is a substantive restriction).} Lemma 1(c) shows that the data cannot see the skew part \(K\). Proposition 9 below goes further: whole families of \emph{nonlinear} departures from A2 leave the distribution of the data unchanged while moving \(B^*\). So the linear contextual-effects assumption cannot be validated by goodness of fit. It is a structural assumption, and \(\mathcal K\) is only the linear part of what the data cannot see.

\textbf{Remark 1.1 (A5 in practice).} For \(R=2\), A5 holds iff \(\operatorname{supp}_AP_X\) contains at least three distinct points.
\emph{Proof.}
1. On \(x=(t,1-t)\), \(x^\top Dx\) is a polynomial of degree at most 2 in \(t\). {[}alg{]}
2. If it vanishes at three distinct \(t\), it is the zero polynomial. {[}poly{]}
3. If \(D_{11}t^2+2D_{12}t(1-t)+D_{22}(1-t)^2\equiv0\), then \(t=0\) gives \(D_{22}=0\), \(t=1\) gives \(D_{11}=0\), and then \(D_{12}=0\). {[}alg{]}
4. With only two support points \(t_1,t_2\), the nonzero form \(D\) given by \((t-t_1)(t-t_2)\) in homogenised form vanishes on both. {[}alg{]} ∎

For general \(R\), A5 says that the design of the quadratic regression on the polls has full rank: all \(R(R+1)/2\) coefficients of a symmetric \(D\) must be pinned down. For \(R=3\) that needs at least six support points that do not lie on a common conic. Any five points lie on some conic, so three, four or five distinct compositions are never enough. It is checkable from the condition number of the regressor matrix (Phase 4).

\textbf{Remark 1.2 (sum-to-one in estimation).} Fit each column by OLS on the \(R(R+1)/2\) quadratic monomials of \(X\) (equivalently the homogeneous form). Then \(\sum_c\hat Q_c=J\) exactly.
\emph{Proof.}
1. \(\sum_cT_{ic}=1=(\mathbf 1^\top X_i)^2=X_i^\top JX_i\) for every \(i\). {[}Def; alg{]}
2. OLS is linear in the response, so \(\sum_c\hat Q_c\) is the OLS fit of the response \(1\). {[}lin{]}
3. By 1 that response is fitted exactly by \(J\), and the fit is unique when the design has full rank (sample A5). {[}1; A5{]} ∎

So the estimated columns are automatically coherent; no constrained regression is needed.

\begin{center}\rule{0.5\linewidth}{0.5pt}\end{center}

\hypertarget{pB-the-estimand-in-terms-of-k}{%
\subsection{\texorpdfstring{3. The estimand in terms of \(K\)}{3. The estimand in terms of K}}\label{pB-the-estimand-in-terms-of-k}}

\textbf{Lemma 2.} Under (S1)--(S2) with \(P(X\in A)=1\):
\[B^*_{rc}=\frac1{\mu_r}\Big[\sum_sQ_{c,rs}\mu_{rs}+\sum_{s\neq r}K_{c,rs}\mu_{rs}\Big].\]
In the ambient coordinates (the entries \(K_{c,rs}\), before skew-symmetry and \(\sum_cK_c=0\) are imposed), the coefficient of \(K_{c,rs}\) is \(\mu_{rs}/\mu_r\ge0\). The entries are not free: in the free coordinates \(k_{c,rs}\) (\(c<C\), \(r<s\)) of Prop. 7, one coordinate enters four cells, with the coefficients given there.

\emph{Proof.}
1. \(0\le NX_r\beta_{rc}\le N\), so \(E[NX_r\beta_{rc}]\) is finite. {[}S1; standing integrability{]}
2. \(E[NX_r\beta_{rc}]=E\big[NX_rE[\beta_{rc}\mid V]\big]\). {[}CE-tower; CE-pull, since \(NX_r\) is \(V\)-measurable{]}
3. \(=E[NX_r\,e_r^\top H_cX]\). {[}S2; \(P(X\in A)=1\){]}
4. \(=\sum_sH_{c,rs}E[NX_rX_s]=\sum_sH_{c,rs}\mu_{rs}\). {[}lin{]}
5. \(H_{c,rs}=Q_{c,rs}+K_{c,rs}\) and \(K_{c,rr}=0\). {[}Def; skew{]}
6. Divide by \(\mu_r>0\). {[}standing integrability{]}
7. The coefficients are read off from the affine form. \(\mu_{rs}\ge0\) because \(N,X_r,X_s\ge0\). {[}alg{]} ∎

\textbf{2×2 check.} \(B^*=[Q_{11}\mu_{11}+(Q_{12}+K)\mu_{12}]/\mu_1\). With the dictionary of Prop. 4 below and \(\mu_{11}+\mu_{12}=\mu_1\), this is JKST eq. (16) in the population: \(w_0+c_1+d_1\mu_{11}/\mu_1-w_1\mu_{12}/\mu_1\).

\textbf{Remark 2.1 (\(A\subsetneq\operatorname{supp}\)).} If \(P(X\notin A)>0\), then
\[B^*_{rc}=\mu_r^{-1}\Big\{E\big[NX_rM_{H,rc}(X)1\{X\in A\}\big]+E\big[NX_r\beta_{rc}1\{X\notin A\}\big]\Big\},\]
and the second term is bounded only by the Duncan--Davis terms of the polls outside \(A\). Prop. 7′ (§6.6) turns this into bounds, and Phase 2b lets \(A\) vary by district. Every result below that uses Lemma 2 assumes \(P(X\in A)=1\) unless it says otherwise. The default \(A=\operatorname{conv}\operatorname{supp}P_X\) satisfies it automatically.

\begin{center}\rule{0.5\linewidth}{0.5pt}\end{center}

\hypertarget{pB-the-sharp-identified-set}{%
\subsection{4. The sharp identified set}\label{pB-the-sharp-identified-set}}

\textbf{Theorem 1 (characterisation of \(\Theta_I\)).} Assume (A1), (A5), standing integrability, and \(H=Q(P)+K\) with \(K\in\mathcal K\). The following are equivalent:
(i) \(K\in\Theta_I(P;A)\);
(ii) (S3) holds, and there is a \(P_V\)-null set \(\mathcal N\) such that for all \(v=(x,n)\notin\mathcal N\) with \(x\in A\),
\[\langle u,M_H(x)\rangle\le\psi(u;v)\quad\text{for all }u\in\mathbb R^{R\times C}.\tag{T1}\]
Moreover, \(\Theta_I(P;A)\) is convex and closed. It is compact when \(\operatorname{affdim}(A)=R-1\), i.e.~\(A\) has nonempty interior relative to \(\Delta^{R-1}\). Under (A5) this always holds (Remark 5.2), so under the theorem's own assumptions \(\Theta_I(P;A)\) is compact.

\emph{Interpretation:} the implied mean transition matrix at composition \(x\) must be an average of selections from the polls' Duncan--Davis polytopes among polls with that composition and size, i.e.~\(M_H(x)\in\mathbb E[S(X,T)\mid V=v]\). Directions \(u\) and \(-u\) along the equality constraints force \(x^\top M_H(x)=E[T\mid V=v]\), so (T1) also requires \(E[T\mid V]=E[T\mid X]\), an observable implication of A3.

\textbf{Why (T1) must hold for all \(u\) at once.} The inequality for one direction \(u\) is a single half-space restriction on \(M_H(x)\). The conclusion \(M_H(x)\in\mathbb E[S(X,T)\mid V=v]\) needs every half-space at the same \(v\), because a closed convex set is the intersection of all its supporting half-spaces {[}Rock{]}. Each fixed \(u\) gives its own exceptional null set \(\mathcal N_u\), and there are uncountably many \(u\), so the union of the \(\mathcal N_u\) need not be null. Steps 9--11 below avoid this: they take the union over a countable dense set and extend by continuity in \(u\). Sufficiency step 15 then uses (T1) at the random direction \(u=g(V(w))\), which is legitimate only because one null set serves all \(u\).

\emph{Proof of (i)⇒(ii).}
1. Take \(\mathbb Q\in\mathcal M(H;A)\) with \(\mathbb Q\circ W^{-1}=P\). (S3) holds by membership. {[}Def \(\Theta_I\); S3{]}
2. \(\beta\in S(X,T)\) \(\mathbb Q\)-a.s. {[}S1; Def \(S\){]}
3. For each fixed \(u\), \(\langle u,\beta\rangle\le h_S(u;X,T)\) a.s. {[}2; Def \(h_S\){]}
4. Both sides are bounded by \(\sqrt R\lVert u\rVert\), hence integrable. {[}Frob; L0(d){]}
5. \(E[\langle u,\beta\rangle\mid V]\le E[h_S(u;X,T)\mid V]\) a.s. {[}3--4; CE-mono{]}
6. The left side equals \(\langle u,E[\beta\mid V]\rangle=\langle u,M_H(X)\rangle\) a.s. on \(\{X\in A\}\). {[}CE-lin; S2{]}
7. The right side equals \(\psi(u;V)\) a.s. The RCD of \(T\) given \(V\) under \(\mathbb Q\) is a version of \(P_{T\mid V}\), because \(\mathbb Q\circ W^{-1}=P\). {[}RCD; L0(d){]}
8. So for each \(u\) there is a null set \(\mathcal N_u\) outside which (T1) holds at that \(u\). {[}5--7{]}
9. Let \(U_0\) be a countable dense subset of \(\mathbb R^{R\times C}\) and \(\mathcal N=\bigcup_{u\in U_0}\mathcal N_u\). Then \(\mathcal N\) is null. {[}countable union of null sets{]}
10. For \(v\notin\mathcal N\), the left side of (T1) is linear in \(u\), and \(\lvert\psi(u;v)-\psi(u';v)\rvert\le\sqrt R\lVert u-u'\rVert\). {[}lin; L0(d) integrated{]}
11. So (T1) extends from \(U_0\) to all \(u\). {[}9--10; density and continuity{]} ∎

\textbf{Functional-analytic facts used in (ii)⇒(i).} Write \(L^1:=L^1(\mathcal W,\mathcal B(\mathcal W),P;\mathbb R^{R\times C})\) with the Frobenius norm. Since \(\mathbb R^{R\times C}\) is finite-dimensional, every statement below can be read coordinate by coordinate for scalar \(L^1\).
* \textbf{F1.} \(\mathcal S\subset L^1\). Each \(\xi\in\mathcal S\) is Borel with \(\lVert\xi\rVert\le\sqrt R\) a.s., hence integrable.
* \textbf{F2.} A family bounded in \(L^\infty\) on a probability space is uniformly integrable: \(\sup_\xi E[\lVert\xi\rVert1\{\lVert\xi\rVert>m\}]=0\) for \(m>\sqrt R\).
* \textbf{F3.} A convex, norm-closed subset of a normed space is weakly closed {[}Mazur{]}.
* \textbf{F4.} A uniformly integrable subset of \(L^1\) is relatively weakly compact {[}DP{]}. With F3, \(\mathcal S\) is weakly compact.
* \textbf{F5.} \(\Pi=E[\cdot\mid V]\) is linear with \(\lVert\Pi\xi\rVert_1\le\lVert\xi\rVert_1\). A norm-continuous linear map between normed spaces is weak-to-weak continuous {[}wk-cont{]}, so \(\Pi(\mathcal S)\) is weakly compact.
* \textbf{F6.} For a bounded \(\mathcal G\)-measurable direction \(g\), \(\sup_{\xi\in\mathcal S}E\langle g,\xi\rangle\) is attained by a Borel selector of \(w\mapsto\arg\max_{b\in S(x,t)}\langle g(v),b\rangle\). This needs: \(S\) nonempty and compact-valued with a measurable graph (L0); the objective is Carathéodory (Borel in \(w\), continuous in \(b\)); and {[}MMT{]}.
* \textbf{F7.} Separation in \(L^1(\mathcal G)\) uses its dual \(L^\infty(\mathcal G)\) {[}HB{]}. A point outside a weakly closed convex set is strictly separated from it by a bounded \(\mathcal G\)-measurable \(g\).

\emph{Proof of (ii)⇒(i).} Work on the canonical space \((\mathcal W,\mathcal B(\mathcal W),P)\). Let \(\mathcal S=\{\xi:\mathcal W\to\mathbb R^{R\times C}\text{ Borel}:\ \xi(w)\in S(x,t)\ P\text{-a.s.}\}\), viewed in \(L^1\).

\begin{enumerate}
\def\labelenumi{\arabic{enumi}.}
\tightlist
\item
  \(\xi^0(w)_{rc}:=t_c\) is Borel, and \(\xi^0\in\mathcal S\), so \(\mathcal S\neq\emptyset\). {[}L0(a) proof{]}
\item
  \(\mathcal S\) is convex. {[}L0(b){]}
\item
  \(\lVert\xi\rVert\le\sqrt R\) for every \(\xi\in\mathcal S\), so \(\mathcal S\) is bounded in \(L^\infty\) and uniformly integrable. {[}\(\lVert b\rVert\le\sqrt R\){]}
\item
  If \(\xi_n\to\xi\) in \(L^1\) with \(\xi_n\in\mathcal S\), a subsequence converges a.s. Since \(S(x,t)\) is closed, \(\xi\in S\) a.s., so \(\mathcal S\) is norm-closed. {[}Riesz; L0(b){]}
\item
  \(\mathcal S\) is weakly compact in \(L^1\). {[}3; DP; 2, 4; Mazur{]}
\item
  \(\Pi\xi:=E[\xi\mid V]\) is a bounded linear operator \(L^1\to L^1(\mathcal G)\), hence weak-to-weak continuous. {[}CE contraction; wk-cont{]}
\item
  \(\mathcal C:=\Pi(\mathcal S)\) is convex and weakly compact, hence weakly closed, in \(L^1(\mathcal G)\). {[}5--6; lin{]}
\item
  Put \(\eta:=M_H(X)1\{X\in A\}+\Pi(\xi^0)1\{X\notin A\}\). \(\eta\) is \(\mathcal G\)-measurable and bounded, because \(M_H\) is linear and \(A\) is compact. {[}Def; alg{]}
\item
  Suppose, for contradiction, that \(\eta\notin\mathcal C\). {[}---{]}
\item
  Then there exist \(g\in L^\infty(\mathcal G;\mathbb R^{R\times C})\) and \(\varepsilon>0\) with \(E\langle g,\eta\rangle\ge\sup_{\zeta\in\mathcal C}E\langle g,\zeta\rangle+\varepsilon\). {[}HB on \((L^1(\mathcal G),\text{weak})\); 7{]}
\item
  For \(\xi\in\mathcal S\): \(E\langle g,\Pi\xi\rangle=E\big[E[\langle g,\xi\rangle\mid V]\big]=E\langle g,\xi\rangle\). {[}CE-pull, \(g\) is \(\mathcal G\)-measurable; CE-tower{]}
\item
  \(\sup_{\xi\in\mathcal S}E\langle g,\xi\rangle\le E[h_S(g(V);X,T)]\), pointwise. {[}Def \(h_S\){]}
\item
  A Borel selector \(\xi^\star(w)\in\arg\max_{b\in S(x,t)}\langle g(v),b\rangle\) exists, belongs to \(\mathcal S\) and attains the bound in 12. So \(\sup_{\mathcal S}E\langle g,\xi\rangle=E[h_S(g(V);X,T)]\). {[}MMT; L0(c); the objective is Carathéodory{]}
\item
  \(E[h_S(g(V);X,T)]=E[\psi(g(V);V)]\). {[}RCD applied to \(f(v,t)=h_S(g(v);x,t)\), bounded and Borel by L0(d); CE-tower{]}
\item
  On \(\{X\in A\}\setminus\{V\in\mathcal N\}\): \(\langle g(V),\eta\rangle=\langle g(V),M_H(X)\rangle\le\psi(g(V);V)\). {[}8; (T1) with \(u=g(V(w))\), valid for all \(u\) at once{]}
\item
  On \(\{X\notin A\}\): \(\langle g(V),\eta\rangle=E[\langle g(V),\xi^0\rangle\mid V]\le E[h_S(g(V);X,T)\mid V]=\psi(g(V);V)\) a.s. {[}8; CE-pull; \(\xi^0\in S\); CE-mono; RCD{]}
\item
  Integrating 15--16: \(E\langle g,\eta\rangle\le E\,\psi(g(V);V)=\sup_{\zeta\in\mathcal C}E\langle g,\zeta\rangle\), which contradicts 10. {[}15--16; CE-mono; 11, 13, 14{]}
\item
  Hence there is \(\xi^\star\in\mathcal S\) with \(E[\xi^\star\mid V]=\eta\) a.s. {[}9--17; Def \(\mathcal C\){]}
\item
  Let \(\beta:=\xi^\star(W)\) and \(\mathbb Q:=\) the law of \((W,\beta)\) under \(P\). Then \(\mathbb Q\circ W^{-1}=P\). {[}Def{]}
\item
  \(\beta\in\mathfrak B\) and \(\beta^\top X=T\) a.s., which is (S1). {[}18; Def \(\mathcal S\), \(S\){]}
\item
  \(E_{\mathbb Q}[\beta\mid V]=\eta=M_H(X)\) a.s. on \(\{X\in A\}\), which is (S2). {[}18; 8{]}
\item
  (S3) holds by hypothesis (ii). So \(\mathbb Q\in\mathcal M(H;A)\) and \(K\in\Theta_I(P;A)\). {[}20--21; Def \(\Theta_I\){]} ∎
\end{enumerate}

\emph{Proof of convexity, closedness, compactness.}
1. \(K\mapsto M_{Q+K}(x)\) is affine, so the left side of (T1) is affine in \(K\), and (S3) is a family of affine inequalities. {[}lin{]}
2. If \(K,K'\) satisfy (ii) outside null sets \(\mathcal N,\mathcal N'\), then every convex combination satisfies (ii) outside \(\mathcal N\cup\mathcal N'\). {[}1{]}
3. If \(K_n\to K\) with \(K_n\) satisfying (ii) outside \(\mathcal N_n\), then \(K\) satisfies (ii) outside \(\bigcup_n\mathcal N_n\), which is null. {[}continuity in \(K\); countable union{]}
4. \(\Theta_I\subseteq\Theta_O\), and \(\Theta_O\) is bounded when \(\operatorname{affdim}(A)=R-1\), which (A5) implies. {[}P1 below; L3(b); Remark 5.2{]} ∎

\textbf{Remark 4.1 (no auxiliary randomisation).} P1 §1.5 said the construction needs extra randomisation. It does not: \(S(x,t)\) is convex, so the selection can be chosen as a function of \(W\) (step 18).

\textbf{Remark 4.2 (sharpness is relative to continuous latent rates).} The model takes \(\beta\in\mathfrak B\) with real entries. With integer cell counts, \(\beta_{irc}=N_{irc}/N_{ir}\) lies on a grid, so the exact finite-population feasible set is a subset of \(S(X,T)\). Theorem 1 asserts sharpness for the continuous model of (S1)--(S3), not for that integer set. Two observations suggest the distinction rarely matters, but they are not used anywhere below:
1. In count coordinates \(n_{rc}=Nx_rb_{rc}\), \(S(x,t)\) is a transportation polytope with integer margins \(Nx_r\) and \(Nt_c\). Such polytopes have integer vertices (total unimodularity), so every vertex of \(S(x,t)\) is an integer table, and the Duncan--Davis bounds are attained by integer tables.
2. \emph{(Sketch, not verified.)} If the construction may randomise, choosing an integer vertex of \(S(X,T)\) with probabilities equal to measurable barycentric weights of \(\xi^\star(W)\), the resulting integer-valued \(\beta\) has \(E[\beta\mid W]=\xi^\star(W)\). Then \(E[\beta\mid V]=\eta\), so the sharp set would be unchanged. This needs a measurable choice of the weights and gives up Remark 4.1's selection without randomisation.

\textbf{Corollary 1 (sharp bounds for any linear target).} Assume Theorem 1 holds, \(P(X\in A)=1\) and \(\Theta_I(P;A)\neq\emptyset\). Then the sharp identified set for \(\varphi=\sum a_{rc}B^*_{rc}\) is the interval
\[\Big[\min_{K\in\Theta_I}\varphi(K),\ \max_{K\in\Theta_I}\varphi(K)\Big],\]
where \(\varphi(K)\) is given by Lemma 2.
\emph{Proof.}
1. For every \(\mathbb Q\) consistent with \(P\) (i.e.~in \(\mathcal M(Q(P)+K;A)\) with \(\mathbb Q_W=P\)), \(\varphi(B^*(\mathbb Q))=\varphi(K)\). {[}L2; lin{]}
2. So the identified set is the image \(\varphi(\Theta_I)\), and every point of it is attained by some consistent \(\mathbb Q\). {[}1; T1, (ii)⇒(i){]}
3. \(\varphi\) is affine and \(\Theta_I\) is convex and compact, so the image is a compact interval with the stated endpoints. {[}lin; T1; IVT{]} ∎

\hypertarget{pB-what-each-set-needs-from-the-data}{%
\subsubsection{4.3 What each set needs from the data}\label{pB-what-each-set-needs-from-the-data}}

\begin{longtable}[]{@{}
  >{\raggedright\arraybackslash}p{(\columnwidth - 4\tabcolsep) * \real{0.3333}}
  >{\raggedright\arraybackslash}p{(\columnwidth - 4\tabcolsep) * \real{0.3333}}
  >{\raggedright\arraybackslash}p{(\columnwidth - 4\tabcolsep) * \real{0.3333}}@{}}
\toprule\noalign{}
\begin{minipage}[b]{\linewidth}\raggedright
Set
\end{minipage} & \begin{minipage}[b]{\linewidth}\raggedright
Information used
\end{minipage} & \begin{minipage}[b]{\linewidth}\raggedright
Computation
\end{minipage} \\
\midrule\noalign{}
\endhead
\bottomrule\noalign{}
\endlastfoot
\(\Theta_I\) (Theorem 1) & the conditional law of \(T\) given \((X,N)\), through \(\psi(u;v)=E[h_S(u;X,T)\mid V=v]\) & infinitely many conditional constraints; not a finite LP in general \\
\(\Theta_X\) (JKST Prop. 1) & the conditional law of \(T\) given \(X\) & the same, conditioning on \(X\) only \\
\(\Theta_O\) (§5, Prop. 7) & the regression \(E[T\mid X]\), i.e.~\(Q\), plus probability validity & a finite LP \\
\end{longtable}

Theorem 1 is the identification result. Prop. 7 is the computable outer bound that extends the JKST implementation. In data, \((X_i,N_i)\) is typically unique for every precinct, so \(P_{T\mid V}\) cannot be estimated without smoothing or pooling over nearby compositions and sizes. Estimating \(\Theta_I\) or \(\Theta_X\) therefore needs extra structure: the series-based \(\Theta_X\) is planned for the package's second release (Phase 4). Results stated for \(\Theta_O\) are called sharp only relative to the regression function (Prop. 2).

\begin{center}\rule{0.5\linewidth}{0.5pt}\end{center}

\hypertarget{pB-the-outer-set-reduction-nesting-and-when-it-is-sharp}{%
\subsection{5. The outer set: reduction, nesting, and when it is sharp}\label{pB-the-outer-set-reduction-nesting-and-when-it-is-sharp}}

\textbf{Lemma 3.}
(a) If \(A\) is a polytope and \(\ell\) is linear, then \(\ell\ge0\) on \(A\) iff \(\ell(v)\ge0\) for all \(v\in\operatorname{vert}(A)\). Hence the constraint ``\(\forall x\in A\)'' in \(\Theta_O\) can be replaced by ``\(\forall v\in\operatorname{vert}(A)\)''.
(b) If \(\operatorname{affdim}(A)=R-1\) (equivalently, \(A\) has nonempty interior relative to \(\Delta^{R-1}\)), then \(\Theta_O(P;A)\) is bounded, hence a polytope. Without this condition \(\Theta_O\) is a closed convex polyhedron, possibly unbounded (Remark 5.2).
(c) Under (S1)--(S2), if \(A\subseteq\operatorname{conv}(\operatorname{supp}_AP_X)\), then (S3) holds automatically. In particular this is so for \(A=\operatorname{conv}(\operatorname{supp}P_X)\), since then \(P(X\in A)=1\) and \(\operatorname{supp}_AP_X=\operatorname{supp}P_X\).

\emph{Proof.}
1. (a) ``⇒'' is immediate since \(\operatorname{vert}(A)\subseteq A\). ``⇐'': write \(x=\sum_k\lambda_kv_k\) with \(\lambda\ge0\), \(\sum\lambda_k=1\). {[}Mink{]}
2. Then \(\ell(x)=\sum_k\lambda_k\ell(v_k)\ge0\). {[}lin{]} ∎(a)
3. (b) For \(K\in\Theta_O\) and \(x\in A\), every row of \(M_{Q+K}(x)\) is nonnegative and sums to \(e_r^\top Jx=1\), so it lies in \(\Delta^{C-1}\). {[}Def \(\Theta_O\); L1(b){]}
4. Hence \(\lvert e_r^\top K_cx\rvert\le1+\lvert e_r^\top Q_cx\rvert\) for all \(x\in A\). {[}3; alg{]}
5. Since \(\operatorname{affdim}(A)=R-1\), \(A\) contains \(R\) affinely independent points \(x^{(1)},\dots,x^{(R)}\). These are linearly independent because \(0\notin\{\mathbf 1^\top x=1\}\). {[}alg{]}
6. \(K_c=[\,e_r^\top K_cx^{(k)}\,]_{r,k}\,[x^{(1)}\cdots x^{(R)}]^{-1}\) is therefore bounded uniformly over \(\Theta_O\). {[}4--5; lin{]} ∎(b)
7. (c) \(E[\beta\mid V]\in\mathfrak B\) a.s., since \(\mathfrak B\) is closed and convex. {[}CE-cvx{]}
8. (S2) holds on \(\{X\in A\}\), so \(M_H(x)\in\mathfrak B\) for \(P_X\)-a.e. \(x\in A\), and by continuity at every \(x\in\operatorname{supp}_AP_X\). {[}S2; supp{]}
9. \(M_H\) is linear in \(x\) and \(\mathfrak B\) is convex, so \(M_H(x)\in\mathfrak B\) on \(\operatorname{conv}(\operatorname{supp}_AP_X)\), which contains \(A\) by hypothesis. {[}lin; L1(b); Mink{]} ∎(c)

\textbf{Remark 5.1 (S3 is not automatic for every \(A\) inside the hull).} The hypothesis of (c) cannot be weakened to \(A\subseteq\operatorname{conv}\operatorname{supp}P_X\). Take \(A=\{x_0\}\) with \(x_0\) in the interior of the hull but \(P(X=x_0)=0\). Then (S2) restricts nothing, and \(M_H(x_0)\) need not be a valid transition matrix. Here \(\operatorname{supp}_AP_X=\emptyset\), so the hypothesis of (c) fails, as it should. \emph{(Corrected after review, Batch B: with the earlier wording \(\operatorname{supp}P_X\cap A\), this set was \(\{x_0\}\) whenever \(x_0\in\operatorname{supp}P_X\), so (c)'s hypothesis appeared to hold and (c) contradicted this remark. The restricted support removes the tension.)} A numerical instance is in \texttt{verify\_phase2a.py}. The same happens, less starkly, whenever \(A\) reaches into a part of the hull where (S2) is imposed on no support points.

\textbf{Remark 5.2 (dimension of \(A\)).} ``Nonempty relative interior'' in the earlier draft was too weak: a singleton has nonempty relative interior in its own affine hull. For \(R=3\), \(C=2\), \(A=\{x_0\}\) with \(x_0=(\tfrac13,\tfrac13,\tfrac13)\), \(Q_1=Q_2=\tfrac12J\), and \(K_1=-K_2=t\begin{psmallmatrix}0&1&-1\\-1&0&1\\1&-1&0\end{psmallmatrix}\): since \(K_1x_0=0\), every \(t\in\mathbb R\) is in \(\Theta_O\). Under (A5), however, \(\operatorname{supp}_AP_X\) lies in no hyperplane of the simplex. (If \(a^\top x=b\) there, for a proper hyperplane, so that \(a-b\mathbf 1\ne0\), then \(\mathbf 1^\top x=1\) gives \(x^\top Dx=\big((a-b\mathbf 1)^\top x\big)(\mathbf 1^\top x)=0\) with \(D=\tfrac12[(a-b\mathbf 1)\mathbf 1^\top+\mathbf 1(a-b\mathbf 1)^\top]\ne0\), contradicting (A5).) So \(A\supseteq\operatorname{supp}_AP_X\) (\(A\) is closed) is full-dimensional, and the condition of (b) holds automatically.

\textbf{Proposition 1 (nesting).} \(\Theta_I(P;A)\subseteq\Theta_X(P;A)\subseteq\Theta_O(P;A)\).
\emph{Proof.}
1. Let \(\mathbb Q\in\mathcal M(H;A)\) (conditioning on \(V\)). Then \(E[\beta\mid X]=E[E[\beta\mid V]\mid X]=M_H(X)\) on \(\{X\in A\}\), so \(\mathbb Q\) belongs to the \(X\)-only model. {[}CE-tower; S2{]}
2. Every \(K\) admitted by the stronger model is admitted by the weaker one, so \(\Theta_I\subseteq\Theta_X\). {[}1; Def{]}
3. Let \(K\in\Theta_X\) with witness \(\mathbb Q\). Then \(M_H(X)=E[\beta\mid X]\in\mathfrak B\) a.s. on \(\{X\in A\}\). {[}S2 (\(X\)-version); CE-cvx{]}
4. (S3) gives \(M_H(x)\in\mathfrak B\) on all of \(A\); in particular \((Q_c+K_c)x\ge0\) for all \(x\in A\). {[}S3; Def \(\mathfrak B\){]}
5. So \(K\in\Theta_O\). {[}4; Def \(\Theta_O\){]} ∎

The JKST ``Jensen'' step (their proof of Prop. 2) is this inclusion, \(\Theta_X\subseteq\Theta_O\). In R×C it needs no Jensen inequality: conditional expectations preserve the convex set \(\mathfrak B\), and the accounting identity in means (\(x^\top M_H(x)=\tau(x)\)) holds by Lemma 1.

\textbf{Proposition 2 (\(\Theta_O\) is sharp for the regression function).} Let \(\mathcal P(P)\) be the set of laws \(P'\) of \((N,X,T')\) with \(P'_{(N,X)}=P_{(N,X)}\) and \(E_{P'}[T'\mid V]=E_P[T\mid X]\) a.s. Assume \(P(X\in A)=1\). Then
\[\Theta_O(P;A)=\bigcup_{P'\in\mathcal P(P)}\Theta_I(P';A).\]
\emph{Proof.}
1. ⊇: For \(P'\in\mathcal P(P)\), \(Q(P')=Q(P)\), because the regression function and \(\operatorname{supp}P_X\) are the same. {[}L1(a); Def \(\mathcal P\){]}
2. \(\Theta_O\) depends on \(P\) only through \(Q(P)\) and \(A\), so \(\Theta_O(P')=\Theta_O(P)\). {[}Def \(\Theta_O\); 1{]}
3. \(\Theta_I(P')\subseteq\Theta_O(P')=\Theta_O(P)\). {[}P1; 2{]}
4. ⊆: Let \(K\in\Theta_O(P;A)\), \(H=Q(P)+K\). Put \(\beta:=M_H(X)\) and \(T':=\beta^\top X\). {[}Def{]}
5. \(\beta\in\mathfrak B\) a.s.: it is nonnegative on \(A\) and its rows sum to 1. {[}Def \(\Theta_O\); L1(b); \(P(X\in A)=1\){]}
6. (S1) holds by construction of \(T'\). {[}4--5{]}
7. \(E[\beta\mid V]=M_H(X)\), because \(\beta\) is \(V\)-measurable. This is (S2). {[}CE-pull{]}
8. (S3) holds by the definition of \(\Theta_O\). {[}Def{]}
9. \(T'_c=X^\top H_cX=X^\top Q_cX=E_P[T_c\mid X]\), so the law \(P'\) of \((N,X,T')\) is in \(\mathcal P(P)\). {[}L1(a) step 5; Def \(Q(P)\){]}
10. The law of \((N,X,T',\beta)\) is in \(\mathcal M(H;A)\) with \(W\)-marginal \(P'\), so \(K\in\Theta_I(P';A)\). {[}6--9; Def \(\Theta_I\){]} ∎

\emph{Meaning:} \(\Theta_O\) is exactly what can be learned when only the regression of \(T\) on \(X\) is used. It is also the sharp set for any dataset in which polls with the same composition have identical outcomes.

\textbf{Corollary 2.} If \(T=E[T\mid X]\) a.s., then \(\Theta_I(P;A)=\Theta_O(P;A)\).
\emph{Proof.}
1. \(\Theta_I(P)\subseteq\Theta_O(P)\). {[}P1{]}
2. In the proof of P2, \(T'=E_P[T\mid X]=T\) a.s., so \(P'=P\). {[}P2 step 9; hypothesis{]}
3. Hence every \(K\in\Theta_O(P)\) lies in \(\Theta_I(P)\). {[}P2 step 10{]} ∎

\textbf{Proposition 3 (sufficient condition for \(\Theta_I=\Theta_O\)).} Assume \(P(X\in A)=1\) and \(E[T\mid V]=E[T\mid X]=:\tau(X)\) a.s. Suppose that for \(P_V\)-a.e. \(v=(x,n)\) and every \(u\), the map \(t\mapsto h_S(u;x,t)\) is affine on \(\operatorname{conv}\operatorname{supp}P_{T\mid V=v}\). Then \(\Theta_I(P;A)=\Theta_O(P;A)\).
\emph{Proof.}
1. Let \(K\in\Theta_O\) and \(H=Q+K\). Fix such a \(v\). {[}Def{]}
2. \(\psi(u;v)=\int h_S(u;x,t)P_{T\mid V=v}(dt)=h_S\big(u;x,\int t\,P_{T\mid V=v}(dt)\big)\): the integral of an affine function equals the function at the mean, and the mean lies in the closed convex hull of the support. {[}hypothesis; lin{]}
3. \(=h_S(u;x,\tau(x))\). {[}RCD; hypothesis \(E[T\mid V]=\tau(X)\){]}
4. \(M_H(x)\in\mathfrak B\) and \(x^\top M_H(x)=(x^\top H_cx)_c=\tau(x)\), so \(M_H(x)\in S(x,\tau(x))\). {[}Def \(\Theta_O\); L1(a){]}
5. So \(\langle u,M_H(x)\rangle\le h_S(u;x,\tau(x))=\psi(u;v)\) for all \(u\). {[}Def \(h_S\); 2--3{]}
6. Condition (ii) of Theorem 1 holds, with (S3) from the definition of \(\Theta_O\). So \(K\in\Theta_I\). {[}T1{]}
7. The reverse inclusion is Proposition 1. {[}P1{]} ∎

\textbf{Remark 3.1 (2×2 reading of Prop. 3).} In 2×2, \(S(x,t)\) is the segment with \(\beta^w\in[L,U]\), where \(L=\max\{0,(t_1-x)/(1-x)\}\) and \(U=\min\{1,t_1/(1-x)\}\).
1. The segment's direction depends on \(x\) only. So for fixed \(u\) the maximising endpoint is the same for every \(t\), and \(h_S(u;x,\cdot)\) is an affine image of \(L\) or of \(U\). {[}alg{]}
2. \(L\) and \(U\) are piecewise affine in \(t_1\), with kinks at \(t_1=x\) and \(t_1=1-x\) respectively. {[}alg{]}
3. So Prop. 3 applies when, among polls with the same composition, the vote share never lies on both sides of \(x\) or of \(1-x\). {[}1--2{]}

In words: \(\Theta_O\) loses nothing unless polls with similar composition are so different that some sit at a Duncan--Davis corner and others do not. Verified numerically: inclusion is strict in 46 of 400 random valid designs, and never reversed.

\begin{center}\rule{0.5\linewidth}{0.5pt}\end{center}

\hypertarget{pB-closed-forms}{%
\subsection{6. Closed forms}\label{pB-closed-forms}}

\hypertarget{pB-equivalence-with-jkst-22}{%
\subsubsection{6.1 Equivalence with JKST (2×2)}\label{pB-equivalence-with-jkst-22}}

\textbf{Proposition 4.} Let \(R=C=2\), with row 1 = black, \(X=X_1\), and \(A=\{(t,1-t):t\in[l,u]\}\), \(0<l<u<1\). Write \(E[T\mid X]=w_0+c_1X+d_1X^2\) and \(K:=K_{1,12}\). Parts (a)--(c) concern \(\Theta_O\) only. Part (d), and any use of it for \(B^*\), also needs \(P(l\le X\le u)=1\) (Lemma 2). Then
(a) \(Q_{1,22}=w_0\), \(Q_{1,12}=w_0+c_1/2\), \(Q_{1,11}=w_0+c_1+d_1\), and \(K=c_1/2-w_1\);
(b) \(\Theta_O\), written in terms of \(w_1\), is exactly JKST Prop. 2;
(c) with \(A=\Delta^1\) it is exactly JKST Prop. 3;
(d) \(B^*(K)\) equals JKST (16).

The resulting bounds are bounds over the \textbf{outer} set, as are JKST's. They are sharp relative to the regression function (Prop. 2), not relative to the full law of the data (Theorem 1).

\emph{Derivation (a).}
1. \(X^\top Q_1X=Q_{11}t^2+2Q_{12}t(1-t)+Q_{22}(1-t)^2\) with \(t=X\). {[}Def{]}
2. \(=Q_{22}+(2Q_{12}-2Q_{22})t+(Q_{11}-2Q_{12}+Q_{22})t^2\). {[}alg{]}
3. Matching coefficients with \(w_0+c_1t+d_1t^2\): \(Q_{22}=w_0\); \(Q_{12}=w_0+c_1/2\); \(Q_{11}=d_1+2Q_{12}-Q_{22}=w_0+c_1+d_1\). {[}L1(a) uniqueness; alg{]}
4. White row: \(E[\beta^w\mid X]=H_{21}t+H_{22}(1-t)=H_{22}+(H_{21}-H_{22})t\), so \(w_0=H_{22}\) and \(w_1=H_{21}-H_{22}\). {[}Def{]}
5. \(H_{21}=Q_{12}-K\) and \(H_{22}=Q_{22}\), so \(w_1=w_0+c_1/2-K-w_0\), i.e.~\(K=c_1/2-w_1\). {[}skew; 3--4{]}

\emph{Derivation (b).} Column 2 is \(H_2=J-H_1\) (L1(b)). At a vertex \(v=(t,1-t)\):
1. Row 2, column 1: \(H_{21}t+H_{22}(1-t)=w_0+w_1t\ge0\iff w_1\ge-w_0/t\). {[}4 above; alg{]}
2. Row 2, column 2: \(1-(w_0+w_1t)\ge0\iff w_1\le(1-w_0)/t\). {[}L1(b); alg{]}
3. Row 1, column 1: \(H_{11}=Q_{11}\) and \(H_{12}=Q_{12}+K=w_0+c_1-w_1\). {[}(a); skew{]}
4. So \(H_{11}t+H_{12}(1-t)=(w_0+c_1+d_1)t+(w_0+c_1-w_1)(1-t)=w_0+c_1+d_1t-w_1(1-t)\). {[}alg{]}
5. \(\ge0\iff w_1\le(w_0+c_1+d_1t)/(1-t)\). {[}alg; \(t<1\){]}
6. Row 1, column 2: \(1-[\,4\,]\ge0\iff w_1\ge(w_0+c_1+d_1t-1)/(1-t)\). {[}L1(b); alg{]}
7. JKST write the last bound as \((w_0+c_1+d_1-1)/(1-x)-d_1=(w_0+c_1+d_1x-1)/(1-x)\), which is the same. {[}alg{]}
8. Taking \(\max\) / \(\min\) over \(t\in\{l,u\}\) gives their \(w_l,w_u\). {[}L3(a); 1--7{]} ∎(b)

This is also verified numerically against \texttt{generateBounds}: agreement to \(10^{-16}\) (P1 §0).

\emph{Derivation (c).} With \(\operatorname{vert}(A)=\{(1,0),(0,1)\}\):
1. At \((1,0)\): rows 1 and 2 give \(Q_{11}\in[0,1]\), which does not involve \(K\), and \(w_0+w_1\in[0,1]\). {[}Def; 1--2 of (b) at \(t=1\){]}
2. At \((0,1)\): \(H_{22}=w_0\in[0,1]\), which does not involve \(K\), and \(H_{12}=w_0+c_1-w_1\in[0,1]\). {[}Def; (b)3{]}
3. Together: \(\max\{-w_0,\,c_1+w_0-1\}\le w_1\le\min\{1-w_0,\,c_1+w_0\}\), which is JKST Prop. 3. {[}alg{]} ∎(c)

The same result also follows from Prop. 5 below: verified symbolically in \texttt{verify\_phase2a.py}.

\emph{Derivation (d).}
1. \(\mu_1B^*=Q_{11}\mu_{11}+(Q_{12}+K)\mu_{12}=(w_0+c_1+d_1)\mu_{11}+(w_0+c_1-w_1)\mu_{12}\). {[}L2; (a){]}
2. \(=(w_0+c_1)(\mu_{11}+\mu_{12})+d_1\mu_{11}-w_1\mu_{12}\). {[}alg{]}
3. \(\mu_{11}+\mu_{12}=E[NX(X+1-X)]=\mu_1\). {[}lin{]}
4. So \(B^*=w_0+c_1+d_1\mu_{11}/\mu_1-w_1\mu_{12}/\mu_1\), the population version of JKST (16). {[}1--3{]} ∎

\hypertarget{pB-a-the-whole-simplex-general-rc}{%
\subsubsection{\texorpdfstring{6.2 \(A\) = the whole simplex (general R×C)}{6.2 A = the whole simplex (general R×C)}}\label{pB-a-the-whole-simplex-general-rc}}

\textbf{Proposition 5.} With \(A=\Delta^{R-1}\):
\[\Theta_O=\Big\{K\in\mathcal K:\ \lvert K_{c,rs}\rvert\le Q_{c,rs}\ \ \forall c,\ r<s\Big\},\]
which also requires \(Q_{c,rr}\ge0\) (a testable condition not involving \(K\)). The set separates by row pair \((r,s)\), and for each pair and column
\[K_{c,rs}\in[-m_{c,rs},\,m_{c,rs}],\qquad m_{c,rs}:=\min\{Q_{c,rs},\,1-Q_{c,rs}\}.\]
Hence
\[B^*_{rc}\in\frac1{\mu_r}\Big[\sum_sQ_{c,rs}\mu_{rs}\ \mp\ \sum_{s\neq r}m_{c,rs}\mu_{rs}\Big],\]
with the minus sign giving the lower bound. These formulas presuppose \(\Theta_O\neq\emptyset\). By steps 1--3 below, \(\Theta_O\neq\emptyset\) iff \(Q_{c,rs}\ge0\) for all \(c,r,s\). Then \(\sum_cQ_{c,rs}=1\) gives \(Q_{c,rs}\le1\) and \(0\le m_{c,rs}\le\tfrac12\). If some \(Q_{c,rs}<0\), the set is empty and the model with \(A=\Delta^{R-1}\) is refuted. \(P(X\in A)=1\) holds automatically here.

\emph{Derivation.}
1. \(\operatorname{vert}(\Delta^{R-1})=\{e_1,\dots,e_R\}\), so \(\Theta_O=\{K:(Q_c+K_c)e_s\ge0\ \forall c,s\}=\{K:Q_{c,rs}+K_{c,rs}\ge0\ \forall c,r,s\}\). {[}L3(a); alg{]}
2. For \(r=s\): \(K_{c,rr}=0\), so the condition is \(Q_{c,rr}\ge0\). {[}skew{]}
3. For \(r\ne s\), the pair of conditions \((r,s)\) and \((s,r)\) reads \(Q_{c,rs}+K_{c,rs}\ge0\) and \(Q_{c,rs}-K_{c,rs}\ge0\), i.e.~\(\lvert K_{c,rs}\rvert\le Q_{c,rs}\). {[}\(Q\) symmetric; \(K\) skew{]}
4. The only other restriction is \(\sum_cK_{c,rs}=0\), which involves one pair at a time, so the set is a product over pairs. {[}Def \(\mathcal K\); 3{]}
5. For a fixed pair, apply Lemma 4(b) with \(\ell_c=-Q_{c,rs}\), \(h_c=Q_{c,rs}\): the coordinate range is \([\max\{-Q_{c,rs},-\sum_{c'\ne c}Q_{c',rs}\},\ \min\{Q_{c,rs},\sum_{c'\ne c}Q_{c',rs}\}]\). {[}L4(b){]}
6. \(\sum_{c'}Q_{c',rs}=J_{rs}=1\), so \(\sum_{c'\ne c}Q_{c',rs}=1-Q_{c,rs}\), and the range is \([-m_{c,rs},m_{c,rs}]\). {[}L1(b); alg{]}
7. \(B^*_{rc}\) depends on \(\{K_{c,rs}\}_{s\ne r}\), which belong to different pairs and so vary independently, each with coefficient \(\mu_{rs}/\mu_r\ge0\). {[}L2; 4{]}
8. Its range is therefore attained coordinate-wise at \(\mp m_{c,rs}\). {[}6--7; lin{]} ∎

\emph{Meaning:} \(H_{c,rs}\) and \(H_{c,sr}\) (``\(r\)-voters in pure-\(s\) polls'' and ``\(s\)-voters in pure-\(r\) polls'') are two probabilities whose average \(Q_{c,rs}\) is identified. Splitting a known average of two numbers in \([0,1]\) leaves exactly \(\pm\min(Q,1-Q)\) of freedom. This holds only if the linear model extrapolates validly to pure polls (A4 with \(A=\Delta\)). Verified against an LP solver in 200 random tables.

\hypertarget{pB-box-hyperplane-used-by-props-56}{%
\subsubsection{6.3 Box ∩ hyperplane (used by Props 5--6)}\label{pB-box-hyperplane-used-by-props-56}}

\textbf{Lemma 4.} Let \(\ell\le h\) in \(\mathbb R^C\) and \(\mathcal P=\{\kappa:\ell\le\kappa\le h,\ \sum_c\kappa_c=0\}\).
(a) \(\mathcal P\neq\emptyset\iff\sum\ell_c\le0\le\sum h_c\).
(b) \(\{\kappa_c:\kappa\in\mathcal P\}=\big[\max\{\ell_c,-\sum_{c'\ne c}h_{c'}\},\ \min\{h_c,-\sum_{c'\ne c}\ell_{c'}\}\big]\).
(c) \(\max_{\mathcal P}g^\top\kappa\) is attained by the greedy rule: start from \(\kappa=\ell\), and raise coordinates in decreasing order of \(g_c\), each up to \(h_c\), until \(\sum\kappa_c=0\).

\emph{Proof.}
1. The box \([\ell,h]\) is convex, so its image under \(\kappa\mapsto\sum\kappa_c\) is an interval containing \(\sum\ell\) and \(\sum h\), and contained in \([\sum\ell,\sum h]\) by monotonicity. So it equals \([\sum\ell,\sum h]\). {[}IVT; lin{]}
2. (a) follows: \(0\) is attained iff \(0\in[\sum\ell,\sum h]\). {[}1{]}
3. (b): \(\kappa_c=k\) extends to a point of \(\mathcal P\) iff \(\ell_c\le k\le h_c\) and the box in the other coordinates contains a point with sum \(-k\). {[}Def{]}
4. By 1 applied to the other coordinates, that second condition is \(\sum_{c'\ne c}\ell_{c'}\le-k\le\sum_{c'\ne c}h_{c'}\). Intersecting gives (b). {[}1; alg{]}
5. (c): substitute \(\kappa=\ell+y\) with \(0\le y\le h-\ell\) and \(\sum y=s:=-\sum\ell\ge0\). The objective becomes \(g^\top\ell+g^\top y\). {[}alg; (a){]}
6. If \(y\) is optimal, \(g_a>g_b\), \(y_a<h_a-\ell_a\) and \(y_b>0\), moving \(\delta>0\) from \(b\) to \(a\) stays feasible and raises the objective by \(\delta(g_a-g_b)>0\), a contradiction. So the optimum fills coordinates in decreasing order of \(g\), which is the greedy rule. {[}exch{]} ∎

\hypertarget{pB-two-rows-any-number-of-columns-alu}{%
\subsubsection{\texorpdfstring{6.4 Two rows, any number of columns, \(A=[l,u]\)}{6.4 Two rows, any number of columns, A={[}l,u{]}}}\label{pB-two-rows-any-number-of-columns-alu}}

\textbf{Proposition 6.} Let \(R=2\) and \(\operatorname{vert}(A)=\{(t_k,1-t_k)\}_{k=1,2}\) with \(0<t_1=l<t_2=u<1\). Write \(\kappa_c:=K_{c,12}\), so that \(K_c=\kappa_c\begin{psmallmatrix}0&1\\-1&0\end{psmallmatrix}\). Define
\[\lambda_c(t)=-Q_{c,12}-Q_{c,11}\tfrac{t}{1-t},\qquad\rho_c(t)=Q_{c,12}+Q_{c,22}\tfrac{1-t}{t},\]
\[\ell_c=\max_k\lambda_c(t_k),\qquad h_c=\min_k\rho_c(t_k).\]
Then \(\Theta_O=\{\kappa:\ell\le\kappa\le h,\ \sum_c\kappa_c=0\}\). The projected interval for each \(\kappa_c\) is given by Lemma 4(b). If \(P(X\in A)=1\), the bounds are
* \(B^*_{1c}=[Q_{c,11}\mu_{11}+(Q_{c,12}+\kappa_c)\mu_{12}]/\mu_1\), attained at the endpoints of \(\kappa_c\)'s projected interval (increasing in \(\kappa_c\));
* \(B^*_{2c}=[(Q_{c,12}-\kappa_c)\mu_{12}+Q_{c,22}\mu_{22}]/\mu_2\) (decreasing in \(\kappa_c\));
* any \(\varphi=\varphi_0+\sum_cg_c\kappa_c\), with \(g_c=\mu_{12}(a_{1c}/\mu_1-a_{2c}/\mu_2)\), by the greedy rule of Lemma 4(c).

\emph{Derivation.}
1. Row 1 at \(v=(t,1-t)\): \(e_1^\top(Q_c+K_c)v=Q_{c,11}t+(Q_{c,12}+\kappa_c)(1-t)\ge0\). {[}Def{]}
2. \(\iff\kappa_c\ge-[Q_{c,11}t+Q_{c,12}(1-t)]/(1-t)=\lambda_c(t)\). {[}alg; \(1-t>0\){]}
3. Row 2: \((Q_{c,12}-\kappa_c)t+Q_{c,22}(1-t)\ge0\iff\kappa_c\le\rho_c(t)\). {[}alg; \(t>0\){]}
4. Over both vertices: \(\ell_c\le\kappa_c\le h_c\). Adding \(\sum_c\kappa_c=0\) gives the stated set. {[}L3(a); Def \(\mathcal K\){]}
5. The \(B^*\) formulas are Lemma 2 with \(R=2\). Their monotonicity follows from \(\mu_{12}\ge0\). {[}L2{]}
6. \(\varphi\) is affine in \(\kappa\) with the stated \(g\), and Lemma 4(c) applies. {[}L2; lin; L4(c){]} ∎

Verified against an LP solver in 300 random cases, for both the projection and the greedy rule.

\textbf{Corollary 6.1 (joint vs collapsed; corrects P1 §1.4).} The collapsed analysis of column \(c\) treats the table as two columns \((c,\neg c)\) with \(Q_{\neg c}=J-Q_c\). For \(R=2\):
(a) the joint interval for \(\kappa_c\) is contained in the collapsed one;
(b) the two are equal whenever, for each \(r\in\{1,2\}\), the numbers \(Q_{c',rr}\) (\(c'\ne c\)) are all of one sign. That holds in particular when the extrapolated pure-poll rates \(H_{c',rr}=Q_{c',rr}\) are nonnegative;
(c) strict tightening therefore requires that the common-sign condition of (b) fail: for some row \(r\), the numbers \(Q_{c',rr}\) (\(c'\ne c\)) take both signs. In particular at least one extrapolated pure-poll rate \(H_{c',rr}=Q_{c',rr}\) is negative, so the linear model is not valid on the whole simplex.

For \(R\ge3\), rows are coupled through the shared \(K_{c,rs}\), and the joint bound can be strictly tighter even when every extrapolation is valid. A characterisation for \(R\ge3\) is left open.

\emph{Proof.}
1. \(\lambda_{\neg c}(t)=-(1-Q_{c,12})-(1-Q_{c,11})\tfrac t{1-t}=\sum_{c'\ne c}\lambda_{c'}(t)\), because \(\sum_{c'}Q_{c'}=J\). {[}Def; L1(b){]}
2. The collapsed upper end is \(\min\{h_c,-\ell_{\neg c}\}\), and \(-\ell_{\neg c}=-\max_k\sum_{c'\ne c}\lambda_{c'}(t_k)\). {[}L4(b) with two columns; 1{]}
3. \(\max_k\sum_{c'}\lambda_{c'}(t_k)\le\sum_{c'}\max_k\lambda_{c'}(t_k)=\sum_{c'\ne c}\ell_{c'}\). {[}the max of a sum is at most the sum of the maxima{]}
4. So the collapsed upper end is at least \(\min\{h_c,-\sum_{c'\ne c}\ell_{c'}\}\), the joint upper end. The lower ends compare the same way, using \(\rho\). That proves (a). {[}2--3; L4(b){]}
5. \(t\mapsto t/(1-t)\) is increasing, so \(\lambda_{c'}\) is monotone with direction set by the sign of \(Q_{c',11}\). Likewise \((1-t)/t\) is decreasing and \(\rho_{c'}\) is monotone with direction set by the sign of \(Q_{c',22}\). {[}alg{]}
6. If the signs agree across \(c'\ne c\), all \(\lambda_{c'}\) are maximised at the same vertex, so equality holds in 3. The same argument applies to \(\rho\). That proves (b). {[}5{]}
7. (c) is the contrapositive of (b). A set of numbers that is not of one sign contains a negative one, and \(H_{c',rr}=Q_{c',rr}\). {[}skew{]} ∎

Numerical evidence (\texttt{verify\_phase2a.py}):
* 2×3 with nonnegative pure-poll rates: never strict, as (b) predicts.
* 2×3 with the model valid only on \(A\): strict in 46 cases.
* 3×3 with all rates valid: strict in 20 cases.
* Never looser in any case.

Joint fitting remains necessary anyway for targets that combine columns (e.g.~loyalty minus defection) and for \(R\ge3\).

\hypertarget{pB-general-rc-the-linear-program}{%
\subsubsection{6.5 General R×C: the linear program}\label{pB-general-rc-the-linear-program}}

\textbf{Proposition 7.} Assume \(P(X\in A)=1\) (needed for the target; Prop. 7′ drops it). Parametrise \(\mathcal K\) by \(k=(k_{c,rs})_{c<C,\ r<s}\in\mathbb R^d\), \(d=(C-1)R(R-1)/2\), with \(K_c=\sum_{r<s}k_{c,rs}(e_re_s^\top-e_se_r^\top)\) for \(c<C\) and \(K_C=-\sum_{c<C}K_c\). For a polytope \(A\):
\[\Theta_O=\{k:Gk\le g\},\]
with one row for each \((v,r,c)\in\operatorname{vert}(A)\times[R]\times[C]\):
\[-\big(e_r^\top K_c(k)v\big)\le(Q_cv)_r,\qquad e_r^\top K_c(k)v=\sum_{s>r}k_{c,rs}v_s-\sum_{s<r}k_{c,sr}v_s\ \ (c<C),\quad e_r^\top K_C(k)v=-\sum_{c<C}e_r^\top K_c(k)v.\]
\(\Theta_O\) is a polyhedron. It is bounded (a polytope) when \(\operatorname{affdim}(A)=R-1\) (Lemma 3(b)). This is a finite LP for the \textbf{outer} set. It is not an LP representation of the sharp set \(\Theta_I\) of Theorem 1, which in general needs infinitely many conditional constraints (§4.3).

For a linear target, \(\varphi(k)=\varphi_0+\gamma^\top k\) with
\[\varphi_0=\sum_{rc}\frac{a_{rc}}{\mu_r}\sum_sQ_{c,rs}\mu_{rs},\]
and, for \(c<C\) and \(r<s\),
\[\gamma_{c,rs}=\mu_{rs}\Big[\frac{a_{rc}-a_{rC}}{\mu_r}-\frac{a_{sc}-a_{sC}}{\mu_s}\Big].\]
This follows from Lemma 2: \(k_{c,rs}\) enters \(B^*_{rc}\) with \(+\mu_{rs}/\mu_r\) and \(B^*_{sc}\) with \(-\mu_{rs}/\mu_s\) (since \(K_{c,sr}=-k_{c,rs}\)), and enters column \(C\) with the opposite signs (since \(K_C=-\sum_{c<C}K_c\)). The bounds are \(\varphi_0+\min/\max\{\gamma^\top k:Gk\le g\}\). With \(A\) the sample hull of \(\{X_i\}\), imposing the rows at every observed \(X_i\) is equivalent to imposing them at the hull vertices.

\emph{Proof.}
1. The parametrisation is a bijection onto \(\mathcal K\). {[}L1(c) step 2{]}
2. \((K_cv)_r=\sum_sK_{c,rs}v_s\), with \(K_{c,rs}=k_{c,rs}\) for \(r<s\), \(K_{c,rs}=-k_{c,sr}\) for \(r>s\), and \(K_{c,rr}=0\). {[}Def; skew{]}
3. Each constraint \((Q_c+K_c)v\ge0\) becomes one row of \(Gk\le g\). {[}L3(a); 2{]}
4. \(\varphi\) is affine in \(K\) by Lemma 2, hence affine in \(k\) by 1. {[}L2; lin{]}
5. The vertices of \(\operatorname{conv}\{X_i\}\) are among the \(X_i\), and every \(X_i\in A\). So constraints at all \(X_i\) imply constraints at the vertices, and conversely by Lemma 3(a). {[}Mink; L3(a){]} ∎

\textbf{Regularity of the value function.} Write \(V(Q,\mu)=\max\{\gamma(\mu)^\top k:Gk\le g(Q)\}\). Only \(g\) depends on \(Q\), and only \(\gamma\) depends on \(\mu\) (through ratios such as \(\mu_{rs}/\mu_r\)). Three statements must be kept apart:
* For fixed \(\mu\), \(V(\cdot,\mu)\) is concave and piecewise linear in \(Q\), on the set where the LP is feasible and bounded (right-hand-side parametric LP).
* For fixed \(Q\), \(V(Q,\cdot)\) is convex and piecewise linear in \(\gamma\): it is the support function of the polyhedron \(\{Gk\le g(Q)\}\). As a function of \(\mu\) it is a composition with the nonlinear map \(\mu\mapsto\gamma(\mu)\).
* Jointly in \((Q,\mu)\), \(V\) is in general neither concave nor piecewise linear. Example: \(\max\{\mu k:-Q\le k\le Q\}=\mu Q\) for \(\mu,Q>0\), which is bilinear (checked in \texttt{verify\_phase2a.py}).

\(V\) is differentiable at \((Q,\mu)\) when the optimal basis is unique and nondegenerate, because there \(V\) coincides locally with a smooth function of both arguments. This generalises JKST's ``tie-breaking'' condition. Phase 3 conditions on the design, so \(\gamma\) and \(G\) are fixed. Its Lemma P3.1 is the first bullet, applied to the estimated coefficients.

\hypertarget{pB-units-outside-a}{%
\subsubsection{\texorpdfstring{6.6 Units outside \(A\)}{6.6 Units outside A}}\label{pB-units-outside-a}}

\textbf{Proposition 7′.} Let \(A\) be a polytope with \(P(X\in A)>0\), (S1) everywhere, (S2)--(S3) on \(A\), and (A5) on \(\operatorname{supp}_AP_X\), so that \(Q\) is identified from \(\{X\in A\}\). Put \(\mu^A_{rs}:=E[NX_rX_s1\{X\in A\}]\), and keep \(\mu_r=E[NX_r]\) over all units. For a linear target \(\varphi=\sum_{rc}a_{rc}B^*_{rc}\), let \(\varphi^A_0\) and \(\gamma^A\) be the \(\varphi_0\) and \(\gamma\) of Prop. 7 with \(\mu_{rs}\) replaced by \(\mu^A_{rs}\). Let \(\tilde a(x)_{rc}:=a_{rc}x_r/\mu_r\) and
\[U_{\text{out}}=E\big[N\,h_S(\tilde a(X);X,T)\,1\{X\notin A\}\big],\qquad L_{\text{out}}=-E\big[N\,h_S(-\tilde a(X);X,T)\,1\{X\notin A\}\big].\]
Then
\[\varphi\in\Big[\varphi^A_0+\min_{k\in\Theta_O}\gamma^{A\top}k+L_{\text{out}},\ \ \varphi^A_0+\max_{k\in\Theta_O}\gamma^{A\top}k+U_{\text{out}}\Big].\]
Each endpoint is attained by a law that has the same regression function on \(A\) and the same law of \(W\) off \(A\). For a single cell, \(U_{\text{out}}=E[N\min\{X_r,T_c\}1\{X\notin A\}]/\mu_r\) and \(L_{\text{out}}=E[N\max\{0,X_r+T_c-1\}1\{X\notin A\}]/\mu_r\), the Duncan--Davis terms of the units outside \(A\). The package option \texttt{range\ =\ "trim"} computes this bound for single cells.

\emph{Proof.}
1. \(\varphi=\sum_{rc}\frac{a_{rc}}{\mu_r}\big\{E[NX_r\beta_{rc}1\{X\in A\}]+E[NX_r\beta_{rc}1\{X\notin A\}]\big\}\). {[}Def; lin{]}
2. In the first term, \(NX_r1\{X\in A\}\) is \(V\)-measurable, so \(E[NX_r\beta_{rc}1\{X\in A\}]=E[NX_r(M_H(X))_{rc}1\{X\in A\}]=\sum_sH_{c,rs}\mu^A_{rs}\). Summed with the weights \(a_{rc}/\mu_r\), this is \(\varphi^A_0+\gamma^{A\top}k\) in the coordinates of Prop. 7. {[}CE-tower; CE-pull; S2 on \(A\); L2 steps 3--5; P7{]}
3. The second part is \(E[N\langle\tilde a(X),\beta\rangle1\{X\notin A\}]\). Since \(\beta\in S(X,T)\), \(-h_S(-\tilde a;X,T)\le\langle\tilde a,\beta\rangle\le h_S(\tilde a;X,T)\), so it lies in \([L_{\text{out}},U_{\text{out}}]\). {[}S1; Def \(h_S\); CE-mono{]}
4. Validity: \(K\in\Theta_O\) by Prop. 1 applied to the law on \(\{X\in A\}\). Adding 2 and 3 gives the interval. {[}P1; 2--3{]}
5. Attainment: no restriction links \(\beta\) off \(A\) to \(K\), so the two terms can be chosen separately. On \(A\), Prop. 2's construction attains any \(K\in\Theta_O\) with the same regression function. Off \(A\), a Borel argmax selector of \(b\mapsto\pm\langle\tilde a(x),b\rangle\) over \(S(x,t)\) attains \(U_{\text{out}}\) or \(L_{\text{out}}\). {[}P2; MMT; L0{]}
6. For a single cell, \(\max_{b\in S(x,t)}x_rb_{rc}=\min\{x_r,t_c\}\) and the minimum is \(\max\{0,x_r+t_c-1\}\). {[}FH{]} ∎

For a target that combines several cells, \(U_{\text{out}}\) can be smaller than the sum of the single-cell Duncan--Davis terms, since \(h_S\) of a sum is at most the sum of the \(h_S\). Summing per-cell terms is valid but can be wider.

\begin{center}\rule{0.5\linewidth}{0.5pt}\end{center}

\hypertarget{pB-refutability-and-undetectable-violations}{%
\subsection{7. Refutability and undetectable violations}\label{pB-refutability-and-undetectable-violations}}

\textbf{Proposition 8 (non-emptiness; generalises JKST App. B).} Suppose \(P\) is generated by some \(\mathbb Q\in\mathcal M(H_0;A)\) with \(P(X\in A)=1\). Then \(K_0\in\Theta_I\subseteq\Theta_O\), and \(B^*_{rc}\in[\varphi_L,\varphi_U]\cap DD^*_{rc}\), where
\[DD^*_{rc}=\Big[\frac{E[N\max\{0,X_r+T_c-1\}]}{\mu_r},\ \frac{E[N\min\{X_r,T_c\}]}{\mu_r}\Big].\]
Hence an empty \(\Theta_O\), or an empty intersection with \(DD^*\), refutes (A1)--(A5) jointly.
\emph{Proof.}
1. \(K_0\in\Theta_I\) by definition, and \(\Theta_I\subseteq\Theta_O\). {[}Def; P1{]}
2. \(B^*_{rc}=\varphi(K_0)\in[\varphi_L,\varphi_U]\). {[}L2; Cor 1; 1{]}
3. Write \(n_{rc}=X_r\beta_{rc}\). Then \(\sum_cn_{rc}=X_r\), \(\sum_rn_{rc}=T_c\) and \(n\ge0\). {[}S1{]}
4. \(n_{rc}\le\min\{X_r,T_c\}\), and \(n_{rc}=X_r-\sum_{c'\ne c}n_{rc'}\ge X_r-\sum_{c'\ne c}T_{c'}=X_r+T_c-1\). {[}3; FH{]}
5. Multiply by \(N\ge0\), take expectations, and divide by \(\mu_r\). {[}CE-mono; standing integrability{]} ∎

In samples an empty set can also arise from sampling noise; Phase 3 turns this into a test.

\textbf{Proposition 9 (what the data can never rule out).} Let \(\Delta:\Delta^{R-1}\to\mathbb R^{R\times C}\) be Borel with (i) \(x^\top\Delta(x)=0\), (ii) \(\Delta(x)\mathbf 1=0\), and (iii) \(M_H(x)+\Delta(x)\in\mathfrak B\), all three for every \(x\in\operatorname{supp}P_X\). Only the support matters: the values of \(\Delta\) off \(\operatorname{supp}P_X\) never enter. Then the structure \(\beta=M_H(X)+\Delta(X)\) and the structure \(\beta'=M_H(X)\) generate the same \(P\), but
\[B^*_{rc}(\beta)-B^*_{rc}(\beta')=\frac{E[NX_r\Delta_{rc}(X)]}{\mu_r}.\]
Under (A5) on \(\operatorname{supp}_AP_X\) (with \(A\supseteq\operatorname{supp}P_X\)), the linear \(\Delta\)'s satisfying (i)--(ii) on the support are exactly \(M_{K'}\) with \(K'\in\mathcal K\), i.e.~the unidentified part. A nonlinear \(\Delta\) violates A2 yet leaves no trace in the data.

\emph{Proof.}
1. \(T=\beta^\top X=M_H(X)^\top X+\Delta(X)^\top X=\tau(X)+0\), and likewise \(T'=\tau(X)\). With the same law of \((N,X)\), both give the same \(P\). {[}S1; (i); L1(a){]}
2. Both are valid structures. {[}(ii)--(iii); \(M_H\) has unit row sums by L1(b){]}
3. The difference in \(B^*\) follows by the computation of Lemma 2. {[}L2 steps 1--4{]}
4. If \(\Delta=M_D\) is linear, then (i) gives \(x^\top D_cx=0\) on the support, so \(\operatorname{sym}D_c=0\) by A5. (ii) gives \((\sum_cD_c)x=0\) on the support, which by A5 lies in no hyperplane (Remark 5.2) and so spans \(\mathbb R^R\); hence \(\sum_cD_c=0\) and \(D\in\mathcal K\). Conversely every \(M_{K'}\) with \(K'\in\mathcal K\) satisfies (i)--(ii). {[}L1(a) step 5; A5; L1(b){]} ∎

\emph{Example (\(R=2\)).} Take \(\Delta_{1c}=\delta_c(x)x_2\) and \(\Delta_{2c}=-\delta_c(x)x_1\) with \(\sum_c\delta_c=0\). A constant \(\delta_c\) is a shift of \(\kappa_c\), invisible but inside the model. \(\delta_c(x)=b_cx_1\) instead gives a \emph{quadratic} contextual effect for row 1 that is exactly cancelled in row 2. It shifts \(B^*_{1c}\) by \(b_cE[NX_1^2X_2]/\mu_1\) and can push the truth outside \(\Theta_O\); JKST Example 6 is the case \(C=2\). This is the precise sense in which A2 is an untestable leap. Phase 3 bounds its consequences through a sensitivity parameter \(\sup_x\lVert\Delta(x)\rVert\).

\begin{center}\rule{0.5\linewidth}{0.5pt}\end{center}

\hypertarget{pB-the-bc-33-case}{%
\subsection{8. The BC 3×3 case}\label{pB-the-bc-33-case}}

Rows and columns are both ordered \((P,O,N)\): in-party, out-party, non-voter, both last time and this time. There are \(d=6\) unknowns:
\[k_{c,PO},\ k_{c,PN},\ k_{c,ON}\quad\text{for }c\in\{P,O\},\]
with the abstain column fixed by \(K_N=-K_P-K_O\).

\emph{Reading one unknown.} \(k_{P,PO}=\tfrac12(H_{P,PO}-H_{P,OP})\) is half the gap between two hypothetical rates:
* the P-loyalty of P-voters in a poll otherwise all out-party; and
* the rate at which out-party voters switch to P in a poll otherwise all in-party.

The data identify the average of the two rates, not the gap.

\emph{One constraint} (poll \(x=(x_P,x_O,x_N)\), row \(P\), column \(P\)):
\[Q_{P,PP}x_P+Q_{P,PO}x_O+Q_{P,PN}x_N+k_{P,PO}x_O+k_{P,PN}x_N\ \ge\ 0.\]
There are 9 such rows per poll vertex.

\emph{Loyalty:}
\[B^*_{PP}=\frac{Q_{P,PP}\mu_{PP}+(Q_{P,PO}+k_{P,PO})\mu_{PO}+(Q_{P,PN}+k_{P,PN})\mu_{PN}}{\mu_P}.\]

\textbf{Illustration} (noise-free, identification only; arbitrary linear DGP; 2000 polls; \(A\) = their hull; \texttt{verify\_phase2a.py}):

\begin{longtable}[]{@{}
  >{\raggedright\arraybackslash}p{(\columnwidth - 8\tabcolsep) * \real{0.2000}}
  >{\raggedright\arraybackslash}p{(\columnwidth - 8\tabcolsep) * \real{0.2000}}
  >{\raggedright\arraybackslash}p{(\columnwidth - 8\tabcolsep) * \real{0.2000}}
  >{\raggedright\arraybackslash}p{(\columnwidth - 8\tabcolsep) * \real{0.2000}}
  >{\raggedright\arraybackslash}p{(\columnwidth - 8\tabcolsep) * \real{0.2000}}@{}}
\toprule\noalign{}
\begin{minipage}[b]{\linewidth}\raggedright
Target
\end{minipage} & \begin{minipage}[b]{\linewidth}\raggedright
Truth
\end{minipage} & \begin{minipage}[b]{\linewidth}\raggedright
\(\Theta_O\) bound
\end{minipage} & \begin{minipage}[b]{\linewidth}\raggedright
DD bound
\end{minipage} & \begin{minipage}[b]{\linewidth}\raggedright
Width ratio
\end{minipage} \\
\midrule\noalign{}
\endhead
\bottomrule\noalign{}
\endlastfoot
Loyalty P→P & 0.848 & {[}0.396, 0.932{]} & {[}0.131, 0.973{]} & 0.64 \\
Defection P→O & 0.086 & {[}0.022, 0.361{]} & {[}0.000, 0.628{]} & 0.54 \\
Demobilisation P→abstain & 0.066 & {[}0.022, 0.301{]} & {[}0.000, 0.532{]} & 0.52 \\
Mobilisation N→P & 0.169 & {[}0.048, 0.590{]} & {[}0.000, 0.896{]} & 0.60 \\
\end{longtable}

All four contain the truth, and all are narrower than DD. They are still wide, because six numbers are unidentified. Narrowing them is what Phase 2b (riding dummies) and the optional shape restrictions (P1 §2) are for.

\begin{center}\rule{0.5\linewidth}{0.5pt}\end{center}

\hypertarget{pB-verification-log}{%
\subsection{9. Verification log}\label{pB-verification-log}}

\texttt{python3\ -I\ verify\_phase2a.py} (31 PASS, 0 FAIL; full output in \texttt{verify\_phase2a\_output.txt}) checks:
* Lemma 1: the dimension \((C-1)R(R-1)/2\) by rank computation for 2×2, 2×3, 3×3 and 4×3; \(x^\top Kx=0\) and uniqueness of \(Q\) symbolically for \(R=3\).
* Prop. 4: the JKST Prop. 3 constraints symbolically; Prop. 2 numerically (P1 script).
* Prop. 5 against an LP solver (200 tables).
* Prop. 6 projection and greedy rule against an LP solver (300 cases).
* Cor. 6.1: inclusion, the equality case, and strict cases.
* Prop. 1: \(\Theta_X\subseteq\Theta_O\) in 400 random valid 2×2 designs, strict in 46.
* Theorem 1 ``⇒'': a known feasible structure lies in \(\Theta_X\) (227/227).
* Prop. 7: the closed-form \(\gamma\) against direct computation (100 targets).
* The 3×3 example: the truth lies inside the bounds.
* Corrections after review (each check would fail on the earlier statement):
* Remark 5.1: with \(A\) a singleton inside the hull, (S1)--(S2) hold but \(M_H(x_0)\notin\mathfrak B\);
* Remark 5.2: with \(A\) a singleton, \(\Theta_O\) is unbounded, and the LP is unbounded;
* Prop. 7: \(V(Q,\mu)=\mu Q\) violates joint concavity (midpoint test).
* Scope and proof additions after review:
* Remark 1.1: five compositions never identify \(Q\) for \(R=3\); six generic ones do;
* Remark 4.2(1): LP optima over integer-margin transportation polytopes are integer tables (200 cases);
* Prop. 5: a negative \(Q_{c,rs}\) empties \(\Theta_O\);
* Prop. 7′: with \(A\) a strict subset of the support, the truth lies in the bound, which lies inside DD.

Not verified by computer: the measure-theoretic steps of Theorem 1, which rest on the cited theorems (status in §0.4), and Remark 4.2(2), which is a sketch.

\hypertarget{pB-corrections-to-phase-1}{%
\subsection{10. Corrections to Phase 1}\label{pB-corrections-to-phase-1}}

\begin{enumerate}
\def\labelenumi{\arabic{enumi}.}
\tightlist
\item
  \textbf{Joint vs collapsed (P1 §1.4).} For \(R=2\), the joint bound is strictly tighter only if some extrapolated pure-poll rate is negative (Cor. 6.1(c)). The general claim of strict gain is withdrawn.
\item
  \textbf{A4.} The model now includes (S3). It is implied when \(A\subseteq\operatorname{conv}(\operatorname{supp}_AP_X)\), e.g.~\(A\) = the population hull (Lemma 3(c)). It is a real assumption otherwise: when extrapolating, as in JKST Prop. 3 and our Prop. 5, and for regions \(A\) that the support does not span (Remark 5.1).
\item
  \textbf{Prop. 2-RC(i).} \(\Theta_O\) is a polyhedron. It is a polytope (bounded) when \(\operatorname{affdim}(A)=R-1\), which A5 implies (Lemma 3(b), Remark 5.2).
\item
  \textbf{Theorem 1} needs no auxiliary randomisation (Remark 4.1).
\item
  \textbf{New result (Prop. 2):} \(\Theta_O\) is exactly the sharp set for an analyst who uses only the regression function.
\end{enumerate}

\hypertarget{pB-open-items-carried-forward}{%
\subsection{11. Open items carried forward}\label{pB-open-items-carried-forward}}

\begin{itemize}
\tightlist
\item
  \textbf{2b:} riding dummies (M1/M2), including \(A\subsetneq\operatorname{supp}\) per district (Remark 2.1), and the 3×3 counterpart of P1 §7.3.
\item
  \textbf{3:} inference for LP value functions (Prop. 7) with non-unique bases; the finite-population term \(B-B^*\); the many-ridings intersection; sensitivity through \(\sup\lVert\Delta\rVert\) (Prop. 9).
\item
  \textbf{Open question:} characterise strict joint-vs-collapsed gains for \(R\ge3\).
\end{itemize}

\begin{center}\rule{0.5\linewidth}{0.5pt}\end{center}

\hypertarget{pB-relation-to-cho-manski-2008-and-cross-manski-2002}{%
\subsection{12. Relation to Cho \& Manski (2008) and Cross \& Manski (2002)}\label{pB-relation-to-cho-manski-2008-and-cross-manski-2002}}

JKST's conclusion suggests extending to R×C tables ``perhaps generalizing Cho and Manski 2008''. That chapter surveys EI and summarises Cross \& Manski (2002, \emph{Econometrica} 70:357--368). The correspondence with this document:

\begin{longtable}[]{@{}
  >{\raggedright\arraybackslash}p{(\columnwidth - 2\tabcolsep) * \real{0.5000}}
  >{\raggedright\arraybackslash}p{(\columnwidth - 2\tabcolsep) * \real{0.5000}}@{}}
\toprule\noalign{}
\begin{minipage}[b]{\linewidth}\raggedright
Cross--Manski / Cho--Manski
\end{minipage} & \begin{minipage}[b]{\linewidth}\raggedright
Here
\end{minipage} \\
\midrule\noalign{}
\endhead
\bottomrule\noalign{}
\endlastfoot
outcome \(Y\), group \(Z\), precinct \(X\) & column \(c\), row \(r\), unit \(i\) \\
data: \(P(Y\mid X)\), \(P(Z\mid X)\) & \(T_i\), \(X_i\) \\
target: \(P(Y\mid X,Z)\) & \(\beta_i\); aggregated, \(B\) \\
``data alone'': feasible \(P(Y\mid X=x,Z)\) = solutions of the finite-mixture problem \(\sum_z\eta_zP(Z=z\mid x)=P(Y\mid x)\) & the unit DD polytope \(S(x,t)\) (Lemma 0) \\
that region is convex with finitely many extreme points & \(S(x,t)\) is a transportation polytope \\
region across precincts is a Cartesian product (no cross-precinct restriction) & DD for \(B\) is a weighted sum of unit polytopes; P1 (D0) \\
identifying assumption: \(V\) (part of the precinct identifier) as an instrument, \(E(Y\mid V,W,Z)=E(Y\mid W,Z)\); Goodman's constancy is the case with no \(W\) & the special case of LCE with \textbf{zero contextual effects}, \(H_c=\alpha_c\mathbf 1^\top\). Then \(K\) is pinned down and \(B\) is point-identified \\
sharp region \(H^*_w\) under that assumption; outer region \(C^*_w\) from iterated expectations only, with \(H^*_w\subseteq C^*_w\) & \(\Theta_I\subseteq\Theta_O\); \(\Theta_O\) is exactly what iterated expectations of the accounting identity deliver (Props 1--2) \\
rank of \(\Pi\) = \(\lvert Z\rvert\) gives point identification & A5 (rank of the quadratic design) \\
an empty region refutes the assumption & Prop. 8 \\
``layer assumptions in decreasing order of plausibility'' & DD ⊇ LCE bounds ⊇ Goodman point (and P1 §6, M2 ⊇ M1 ⊇ within-riding constancy) \\
\end{longtable}

\textbf{What is new relative to Cross--Manski:}
1. \textbf{The assumption is weaker than constancy.} Rates may vary linearly with composition. This leaves the \((C-1)R(R-1)/2\)-dimensional \(K\) unidentified, so the result is a set rather than a point. JKST's 2×2 bounds and our R×C bounds sit \emph{between} their data-alone region and their Goodman point.
2. \textbf{It is a mean restriction on a superpopulation.} Cross--Manski treat each precinct's population as fixed, so constancy means \(\beta_i=\beta_j\) exactly. LCE restricts only \(E[\beta\mid X]\), and each poll may deviate. This is why the bounds are for aggregates only, why individual polls stay at DD, and why the finite-population gap \(B-B^*\) must be handled in Phase 3.
3. \textbf{Composition is continuous.} Their instrument \(V\) is a discrete precinct attribute with several precincts per cell. Here the ``instrument'' is the composition \(X\) itself, entering through the linear model.

\textbf{Bearing on Phase 2b.} Their \(W\)-conditioned assumption (constancy across precincts \emph{within} cells of \(W\)) is, with \(W\) = riding, ``Goodman within each riding''. It is a stronger layer than M1, which lets rates vary linearly with composition inside each riding. It is a natural extra layer to report, and its point estimate offers a check: does it fall inside the M1 bounds?


\clearpage
\section{PHASE2b: units nested in groups, and the assumption ladder}\label{app:phase2b}
\textit{Source: \texttt{docs/extension/PHASE2b\_districts.md}.}\par
% GENERATED by tools/make_appendices.py from docs/extension/PHASE2b_districts.md
% Do not edit by hand; edit the source and regenerate.
\emph{Builds on \texttt{PHASE1\_foundations.md} (P1) and \texttt{PHASE2a\_rxc.md} (2a). Every proof step ends with a justification tag; tags are those of 2a §0.3 plus the new ones in §0.3. All closed forms, inclusions and examples are checked in \texttt{verify\_phase2b.py} (output: \texttt{verify\_phase2b\_output.txt}, 22/22 PASS).}

\emph{``District'' means any partition of the units into groups: ridings, counties within states, schools within boards, or, as Remark 6.2 shows, election periods. The results do not use anything specific to voter transitions.}

\begin{center}\rule{0.5\linewidth}{0.5pt}\end{center}

\hypertarget{pC-setting}{%
\subsection{0. Setting}\label{pC-setting}}

\hypertarget{pC-data-and-notation}{%
\subsubsection{0.1 Data and notation}\label{pC-data-and-notation}}

\begin{itemize}
\tightlist
\item
  Observables per unit: \(W=(N,X,T,D)\) with \(D\in\{1,\dots,D^\star\}\) the group, finite, and \(\pi_d=P(D=d)>0\). \(V:=(X,N,D)\), \(\mathcal G=\sigma(V)\).
\item
  \(A_d\) is a compact convex set with \(P(X\in A_d\mid D=d)=1\); the default is \(A_d=\operatorname{conv}\operatorname{supp}P_{X\mid D=d}\), the district's own composition range. In samples it is a polytope with vertices among the district's polls.
\item
  Moments: \(\mu^{(d)}_{rs}=E[NX_rX_s1\{D=d\}]\), \(\mu^{(d)}_r=E[NX_r1\{D=d\}]\), and \(\mu_{r}=\sum_d\mu^{(d)}_r>0\), \(\mu_{rs}=\sum_d\mu^{(d)}_{rs}\).
\item
  For \(Z\in\mathbb R^{R}\), \(Z\mathbf 1^\top\) is the \(R\times R\) matrix whose row \(r\) is constant and equal to \(Z_r\). Since \(\mathbf 1^\top x=1\) on the simplex, \((Z\mathbf 1^\top)x=Z\): an additive shift written in homogeneous form.
\end{itemize}

\hypertarget{pC-the-district-model-and-the-rungs}{%
\subsubsection{0.2 The district model and the rungs}\label{pC-the-district-model-and-the-rungs}}

\textbf{District model \(\mathcal M(\mathbf H;\mathbf A)\)}, with \(\mathbf H=(H^{(1)},\dots,H^{(D^\star)})\) and each \(H^{(d)}=(H^{(d)}_1,\dots,H^{(d)}_C)\): laws \(\mathbb Q\) of \((W,\beta)\) with

\begin{itemize}
\tightlist
\item
  (S1) \(\beta\in\mathfrak B\) and \(T=\beta^\top X\) a.s.;
\item
  (S2-d) \(E[\beta\mid V]=M_{H^{(D)}}(X)\) a.s.;
\item
  (S3-d) \(M_{H^{(d)}}(x)\in\mathfrak B\) for all \(x\in A_d\), all \(d\).
\end{itemize}

Write \(Q^{(d)}=\operatorname{sym}H^{(d)}\) and \(K^{(d)}=\operatorname{skew}H^{(d)}\) column by column. Let \(k^{(d)}\in\mathbb R^{d_0}\) be the free coordinates of \(K^{(d)}\) (2a Prop. 7), \(d_0=(C-1)R(R-1)/2\), and let \(c_d\) denote the coordinates of \(\operatorname{skew}(\zeta_{\cdot,d}\mathbf 1^\top)\).

\textbf{Rungs} (restrictions on \(\mathbf H\); decreasing plausibility; P1 §7, \texttt{ASSUMPTION\_MENU.md}):

\begin{longtable}[]{@{}
  >{\raggedright\arraybackslash}p{(\columnwidth - 2\tabcolsep) * \real{0.5000}}
  >{\raggedright\arraybackslash}p{(\columnwidth - 2\tabcolsep) * \real{0.5000}}@{}}
\toprule\noalign{}
\begin{minipage}[b]{\linewidth}\raggedright
Rung
\end{minipage} & \begin{minipage}[b]{\linewidth}\raggedright
Restriction \(\mathcal R_\rho\) on \(\mathbf H\)
\end{minipage} \\
\midrule\noalign{}
\endhead
\bottomrule\noalign{}
\endlastfoot
\textbf{M2} & none: every district has its own linear contextual effects \\
\textbf{M1(\(\delta\))} & \(\exists\bar k\): \(\lVert k^{(d)}-\bar k-c_d\rVert_\infty\le\delta\) for all \(d\), where \(\zeta_{c,d}:=\operatorname{diag}(Q^{(d)}_c-Q^{(1)}_c)\) \\
\textbf{M1} & \(H^{(d)}_c=H_c+\zeta_{c,d}\mathbf 1^\top\) with \(\sum_c\zeta_{c,d}=0\) and \(\zeta_{c,1}=0\): shared contextual effects, district baselines \\
\textbf{L4(\(r\))} & M1 and, for the designated row \(r\), \(H_{c,rs}=H_{c,rr}\) for all \(s,c\): row \(r\) has no contextual effect \\
\textbf{L5} & \(H^{(d)}_c=\alpha_{c,d}\mathbf 1^\top\): constancy within each district \\
\textbf{L6} & L5 with \(\alpha_{c,d}\) the same for all \(d\) (Goodman) \\
\emph{M0} & \emph{(not nested with the others)} \(H^{(d)}=H\) for all \(d\) with \(A=\operatorname{conv}\bigcup_dA_d\): the pooled model of 2a \\
\end{longtable}

The district \(\delta\)-ball in M1(\(\delta\)) is centred on the M1 prediction \(\bar k+c_d\). \(K\) entries are half-gaps between two hypothetical rates (2a §6.2), so \(\delta\) is measured in probability units: \(\delta=0.05\) lets the gap between ``\(r\)-voters in pure-\(s\) polls'' and ``\(s\)-voters in pure-\(r\) polls'' differ across districts by up to 10 points.

\textbf{Identified sets.} For rung \(\rho\), \(\Theta^\rho_I(P)\) is the set of \((k^{(1)},\dots,k^{(D^\star)})\) (and \(\bar k\) for M1(\(\delta\))) for which some \(\mathbb Q\in\mathcal M(\mathbf H;\mathbf A)\) with \(\mathbf H\in\mathcal R_\rho\), \(Q^{(d)}=Q^{(d)}(P)\), reproduces \(P\). \(\Theta^\rho_O\) replaces the sharp condition by nonnegativity at the vertices of each \(A_d\). \(\mathcal B^\rho\) denotes the corresponding set of values of a linear target \(\varphi\).

\hypertarget{pC-additional-tags}{%
\subsubsection{0.3 Additional tags}\label{pC-additional-tags}}

\begin{longtable}[]{@{}
  >{\raggedright\arraybackslash}p{(\columnwidth - 2\tabcolsep) * \real{0.5000}}
  >{\raggedright\arraybackslash}p{(\columnwidth - 2\tabcolsep) * \real{0.5000}}@{}}
\toprule\noalign{}
\begin{minipage}[b]{\linewidth}\raggedright
Tag
\end{minipage} & \begin{minipage}[b]{\linewidth}\raggedright
Meaning
\end{minipage} \\
\midrule\noalign{}
\endhead
\bottomrule\noalign{}
\endlastfoot
{[}2a:X{]} & result X of \texttt{PHASE2a\_rxc.md} \\
{[}cond-D{]} & conditioning on \(\{D=d\}\), an atom of positive probability: (S1)--(S3-d) restricted to \(\{D=d\}\) are the 2a model for the conditional law \(P(\cdot\mid D=d)\) \\
{[}Gale{]} & Gale's supply--demand theorem (equivalently max-flow/min-cut): a bipartite transportation problem is feasible iff every set of supply nodes demands no more than the capacity of its neighbours \\
{[}sep{]} & minimising a separable sum over a product set equals summing the separate minima \\
{[}nest{]} & if model class \(a\subseteq b\), then \(\Theta^a_I(P)\subseteq\Theta^b_I(P)\) and the identified sets of any functional are nested \\
{[}LP-par{]} & the optimal value of a feasible bounded LP is a continuous, piecewise-linear and convex (min) / concave (max) function of its right-hand side on the region of feasibility \\
\end{longtable}

\begin{center}\rule{0.5\linewidth}{0.5pt}\end{center}

\hypertarget{pC-without-a-model-grouping-adds-nothing}{%
\subsection{1. Without a model, grouping adds nothing}\label{pC-without-a-model-grouping-adds-nothing}}

\textbf{Proposition D0.} Under (S1) alone, the sharp identified set for \(B^*_{rc}\) is
\[DD^*_{rc}=\Big[\tfrac{E[N\max\{0,X_r+T_c-1\}]}{\mu_r},\ \tfrac{E[N\min\{X_r,T_c\}]}{\mu_r}\Big],\]
whatever the law of \(D\) (or of any covariate). Any gain from district dummies therefore comes entirely from a cross-district restriction.

\emph{Proof.}
1. Every \(\mathbb Q\) satisfying (S1) gives \(B^*_{rc}\in DD^*_{rc}\). {[}2a:P8 steps 3--5{]}
2. Fix \((x,t)\) and put \(m=\min\{x_r,t_c\}\). Consider the transportation problem with row margins \(x\), column margins \(t\) and cell \((r,c)\) fixed at \(m\). The residual margins \(x-me_r\) and \(t-me_c\) are nonnegative with equal totals. {[}alg{]}
3. The residual problem has all edges present, so Gale's condition reduces to equal totals, and it is feasible. Hence \(\max_{b\in S(x,t)}x_rb_{rc}=m\). {[}Gale; 2a:L0{]}
4. Put \(m'=\max\{0,x_r+t_c-1\}\) and forbid further mass in cell \((r,c)\). The residual row \(r\) needs \(x_r-m'\le\sum_{c'\ne c}t_{c'}=1-t_c\), which holds because \(m'\ge x_r+t_c-1\). Symmetrically, column \(c\) needs \(t_c-m'\le1-x_r\). {[}alg{]}
5. Any other set of rows is adjacent to every column, so Gale's condition holds for it automatically. So the residual problem is feasible and \(\min_{b\in S(x,t)}x_rb_{rc}=m'\). {[}Gale{]}
6. A Borel selection attaining the maximum (or minimum) of \(x_rb_{rc}\) over \(S(x,t)\) exists. {[}2a tag MMT; 2a:L0(c){]}
7. Integrating \(N\times\) that selection gives the endpoints of \(DD^*_{rc}\). Both are attained by valid structures (S1), with \(\beta\) a function of \(W\). {[}3, 5, 6; lin{]}
8. Neither the constraint \(\beta\in S(X,T)\) nor the objective involves \(D\), so the set does not depend on the law of \(D\). {[}Def{]} ∎

\textbf{Cell intervals versus the joint set.} \(DD^*_{rc}\) is the sharp set for one cell. The joint sharp set for the whole matrix \(B^*\) is
\[\mathcal B^*_{DD}=\Big\{\big[E[NX_r\xi_{rc}]/\mu_r\big]_{rc}:\ \xi\text{ a measurable selection of }S(X,T)\Big\},\]
also free of \(D\). It is convex and compact. It is \textbf{not} the Cartesian product \(\prod_{rc}DD^*_{rc}\): the accounting identity links the cells of each unit, so in general the per-cell endpoints cannot be attained together. For a linear target \(\sum a_{rc}B^*_{rc}\), the sharp interval is \(\big[-E[N\,h_S(-\tilde a(X);X,T)],\ E[N\,h_S(\tilde a(X);X,T)]\big]\) with \(\tilde a(x)_{rc}=a_{rc}x_r/\mu_r\) (the off-\(A\) term of 2a Prop. 7′, applied to every unit), which can be strictly inside the sum of the cell intervals. Verified: per-cell extremes are attained in 300 random tables.

\begin{center}\rule{0.5\linewidth}{0.5pt}\end{center}

\hypertarget{pC-what-each-rung-identifies}{%
\subsection{2. What each rung identifies}\label{pC-what-each-rung-identifies}}

\textbf{Lemma D1.} Assume (A5-d): for each \(d\), \(\operatorname{supp}(P_{X\mid D=d}|_{A_d})\) lies in no quadric of the simplex. Here \(\operatorname{supp}(P_{X\mid D=d}|_{A_d})\) is the support of the restricted measure (2a §0.1); since \(P(X\in A_d\mid D=d)=1\) (§0.1), it equals \(\operatorname{supp}P_{X\mid D=d}\). \emph{(Corrected after review, Batch B, Nov 2026: earlier written \(\operatorname{supp}P_{X\mid D=d}\cap A_d\), which differs when \(A_d\) carries less than full mass; see 2a §0.1.)}
(a) Under any rung, \(Q^{(d)}\) is identified from \(E[T\mid X,D=d]=X^\top Q^{(d)}X\), \(\sum_cH^{(d)}_c=J\), and each \(K^{(d)}\) is unidentified with \(K^{(d)}\in\mathcal K\).
(b) Under M1: \(\zeta_{c,d}=\operatorname{diag}(Q^{(d)}_c-Q^{(1)}_c)\); the off-diagonal entries satisfy the \textbf{overidentifying restrictions}
\[(Q^{(d)}_c-Q^{(1)}_c)_{rs}=\tfrac12(\zeta_{c,d,r}+\zeta_{c,d,s})\quad(r\ne s);\]
and \(K^{(d)}=K+\operatorname{skew}(\zeta_{\cdot,d}\mathbf 1^\top)\) with \(K=K^{(1)}\) common. The unidentified dimension is \(d_0\), the same as the pooled model.
(c) Under L5: \(\alpha_{c,d}=\operatorname{diag}Q^{(d)}_c\), \(K^{(d)}=\operatorname{skew}(\alpha_{\cdot,d}\mathbf 1^\top)\) is identified, and \(Q^{(d)}_c=\operatorname{sym}(\alpha_{c,d}\mathbf 1^\top)\) is testable.

\begin{enumerate}
\def\labelenumi{(\alph{enumi})}
\setcounter{enumi}{3}
\tightlist
\item
  \textbf{Under M1, less richness suffices (revised after review).} Replace (A5-d) by the stacked rank condition (A5-M1): for symmetric \(D\) and vectors \(z_2,\dots,z_{D^\star}\),
  \[x^\top Dx+z_d^\top x=0\ \ \forall x\in\operatorname{supp}(P_{X\mid D=d}|_{A_d}),\ \forall d\quad(z_1:=0)\ \Longrightarrow\ D=0,\ z_d=0.\]
  This is the population version of ``the design with the pooled quadratic monomials and the district × \(X_r\) shifts (\(d\ge2\)) has full column rank''. Then \(Q\), the \(\zeta_{c,d}\) and the conclusions of (b) hold. A sufficient condition: A5 holds within one district, and every other district's support contains \(R\) affinely independent points. Having \(R\) \emph{distinct} compositions is not enough: they must be affinely independent.
\end{enumerate}

\emph{Proof of (a).}
1. On \(\{D=d\}\), (S1)--(S3-d) are the 2a model with \(H=H^{(d)}\) and \(A=A_d\) for the law \(P(\cdot\mid D=d)\). {[}cond-D{]}
2. Apply 2a Lemma 1 to that conditional law. {[}2a:L1; A5-d{]} ∎(a)

\emph{Proof of (d).}
1. Under M1, \(E[T_c\mid X=x,D=d]=x^\top Q_cx+x^\top(\zeta_{c,d}\mathbf 1^\top)x=x^\top Q_cx+\zeta_{c,d}^\top x\) on the district's support. {[}2a:L1(a) steps 1--6 within \(d\); Def M1; \(\mathbf 1^\top x=1\){]}
2. If two parameter sets give the same regression in every district, their differences \((D,z_d)\) satisfy the premise of (A5-M1), so they coincide. {[}1; A5-M1{]}
3. Sufficiency: if A5 holds in district \(d_0\), then for \(d_0=1\) the premise gives \(x^\top Dx=0\) there, so \(D=0\). For \(d_0\ne1\), \(x^\top Dx+z_{d_0}^\top x=x^\top(D+\operatorname{sym}(z_{d_0}\mathbf 1^\top))x=0\) gives \(D=-\operatorname{sym}(z_{d_0}\mathbf 1^\top)\), and district 1's equation then forces \(z_{d_0}^\top x=0\) on an affinely spanning set, so \(z_{d_0}=0\) and \(D=0\). Then \(z_d^\top x=0\) on each district's affinely spanning support. Such a support spans \(\mathbb R^R\) linearly, because it lies in \(\{\mathbf 1^\top x=1\}\), which misses the origin. So \(z_d=0\). {[}alg; A5{]}
4. (b)'s formulas follow as before, from the identified \(Q\) and \(\zeta\). ∎(d)

\emph{Proof of (b).}
1. \(H^{(d)}_c=H_c+\zeta_{c,d}\mathbf 1^\top\). {[}Def M1{]}
2. \(\operatorname{sym}(\zeta\mathbf 1^\top)_{rs}=\tfrac12(\zeta_r+\zeta_s)\) and \(\operatorname{skew}(\zeta\mathbf 1^\top)_{rs}=\tfrac12(\zeta_r-\zeta_s)\). {[}alg{]}
3. So \(Q^{(d)}_c=Q_c+\operatorname{sym}(\zeta_{c,d}\mathbf 1^\top)\) and \(K^{(d)}_c=K_c+\operatorname{skew}(\zeta_{c,d}\mathbf 1^\top)\), with \(Q=Q^{(1)}\) and \(K=K^{(1)}\) because \(\zeta_{\cdot,1}=0\). {[}1--2; lin{]}
4. Taking diagonals in 3 gives \(\operatorname{diag}(Q^{(d)}_c-Q^{(1)}_c)=\zeta_{c,d}\), identified by (a). {[}2--3{]}
5. The off-diagonal entries of 3 give the overidentifying restrictions. {[}2--3{]}
6. The only unidentified object left is \(K\); \(\sum_c\zeta_{c,d}=0\) keeps \(K^{(d)}\in\mathcal K\). {[}3; (a){]} ∎(b)

\emph{Proof of (c).}
1. \(H^{(d)}_c=\alpha_{c,d}\mathbf 1^\top\), so \(Q^{(d)}_c=\operatorname{sym}(\alpha_{c,d}\mathbf 1^\top)\), whose diagonal is \(\alpha_{c,d}\). {[}Def L5; (b)2{]}
2. Hence \(K^{(d)}_c=\operatorname{skew}(\alpha_{c,d}\mathbf 1^\top)\) is a known function of identified quantities. The off-diagonal entries of 1 are testable restrictions. {[}1; (a){]} ∎(c)

\textbf{Remark D1.1 (why (d) matters in practice).} Given the common \(Q\), district \(d\)'s regression is \(E[T_c\mid X,d]=X^\top Q_cX+\zeta_{c,d}^\top X\), which is linear in \(\zeta_{c,d}\). So under M1 (Lemma D1(d)), \(\zeta_{c,d}\) needs only \(R\) affinely independent compositions in district \(d\), rather than A5-d.~This is identification, not precision: with about 10 polls in a riding, \(\hat\zeta_d\) is identified but noisy (Phase 3, regime R2). For example, 2 distinct polls per district in 2×2, provided \(Q\) is identified elsewhere. This matters with about 10 polls per riding.

Verified on the 3×3 example: \(\zeta\) recovered to \(10^{-15}\); the overidentifying restrictions and the \(K^{(d)}\) relation hold to \(10^{-15}\).

\textbf{Lemma D2 (estimand).} Under (S1)--(S2-d) with \(P(X\in A_D)=1\):
\[B^*_{rc}=\frac1{\mu_r}\sum_d\Big[\sum_sQ^{(d)}_{c,rs}\mu^{(d)}_{rs}+\sum_{s\ne r}K^{(d)}_{c,rs}\mu^{(d)}_{rs}\Big].\]
Under M1, its sensitivity to the common \(K_{c,rs}\) is \(\mu_{rs}/\mu_r\), exactly as in the pooled model.
\emph{Proof.}
1. \(E[NX_r\beta_{rc}]=\sum_dE[NX_r\beta_{rc}1\{D=d\}]\). {[}lin; \(D\) finite{]}
2. \(=\sum_dE[NX_r1\{D=d\}\,e_r^\top H^{(d)}_cX]\). {[}CE-tower; CE-pull (\(NX_r1\{D=d\}\) is \(V\)-measurable); S2-d{]}
3. \(=\sum_d\sum_sH^{(d)}_{c,rs}\mu^{(d)}_{rs}\), then split \(H=Q+K\). {[}lin; skew{]}
4. Under M1, \(K^{(d)}=K+\operatorname{skew}(\zeta\mathbf 1^\top)\), and \(\sum_d\mu^{(d)}_{rs}=\mu_{rs}\). {[}D1(b); Def{]} ∎

\begin{center}\rule{0.5\linewidth}{0.5pt}\end{center}

\hypertarget{pC-sharp-and-outer-sets-for-each-rung}{%
\subsection{3. Sharp and outer sets for each rung}\label{pC-sharp-and-outer-sets-for-each-rung}}

\textbf{Theorem 1′ (sharpness for any affine mean model).} Let \(\theta\) range over a closed subset \(\Theta_0\) of \(\mathbb R^p\). Let \(\mathcal V_A=\{v=(x,n,d):x\in A_d\}\), and assume:
* (H1) for each \(\theta\), \(v\mapsto m_\theta(v)\in\mathbb R^{R\times C}\) is Borel on \(\mathcal V_A\);
* (H2) \(\sup_{v\in\mathcal V_A}\lVert m_\theta(v)\rVert<\infty\) for each \(\theta\);
* (H3) \(\theta\mapsto m_\theta(v)\) is affine for each \(v\), and continuous uniformly in \(v\) on bounded sets of \(\theta\);
* (H4) the selection set \(\mathcal S\) of 2a Theorem 1 is unchanged: \(S(x,t)\) is nonempty, compact and convex, with a measurable graph (2a Lemma 0), and is not affected by \(\theta\).

Define the model by (S1), \(E[\beta\mid V]=m_\theta(V)\) a.s. on \(\{V\in\mathcal V_A\}\), and \(m_\theta(v)\in\mathfrak B\) for all \(v\in\mathcal V_A\). Then \(\theta\in\Theta_I\) iff the latter holds and, for \(P_V\)-a.e. \(v\) and all \(u\),
\[\langle u,m_\theta(v)\rangle\le\psi(u;v).\]
\emph{Proof.} Replace \(M_H(X)\) by \(m_\theta(V)\) in the proof of 2a Theorem 1. (H4) leaves facts F1--F7 there unchanged. That proof uses only boundedness (H2) and \(V\)-measurability (H1) of \(M_H(X)\) (steps 6, 8, 15, 16 and 21 of ``(ii)⇒(i)''; steps 6--7 of ``(i)⇒(ii)''), and affinity and continuity in the parameter (H3; the convexity and closedness part). Each step therefore carries over with the same tag. For the rungs of Theorem 2, \(m_\theta(x,n,d)=M_{Q^{(d)}+K^{(d)}(\theta)}(x)\) is linear in \(x\), so (H1)--(H2) hold on compact \(A_d\), and (H3) holds because \(K^{(d)}(\theta)\) is linear in \(\theta\). {[}2a:T1{]} ∎

\textbf{Theorem 2 (the rungs).} Assume (A1), (A5-d) and \(P(X\in A_D)=1\).
(a) \textbf{Sharp set.} For each rung \(\rho\),
\[\Theta^\rho_I=\Big\{\theta\in\mathcal R_\rho:\ \text{for each }d,\ \langle u,M_{Q^{(d)}+K^{(d)}(\theta)}(x)\rangle\le\psi(u;x,n,d)\ \ \forall u,\ P_V\text{-a.e. }(x,n)\text{ given }D=d\Big\}.\]
That is the rung restriction intersected with each district's own sharp condition (2a Theorem 1 within district \(d\)).
(b) \textbf{Outer set.}
\[\Theta^\rho_O=\Big\{\theta\in\mathcal R_\rho:(Q^{(d)}_c+K^{(d)}_c(\theta))v\ge0\ \ \forall d,c,\ v\in\operatorname{vert}(A_d)\Big\}\supseteq\Theta^\rho_I.\]
It is a polyhedron, and every target bound is an LP. For M1 the variables are just the common \(k\in\mathbb R^{d_0}\) and
\[\Theta^{M1}_O=\bigcap_d\Theta^{(d)}_O,\qquad\Theta^{(d)}_O=\{k:(Q^{(d)}_c+\operatorname{skew}(\zeta_{c,d}\mathbf 1^\top)+K_c(k))v\ge0\ \ \forall c,\ v\in\operatorname{vert}A_d\}.\]
(c) \textbf{Ladder.} For \(0\le\delta\le\delta'\):
\[\mathcal B^{L6}\subseteq\mathcal B^{L5}\subseteq\mathcal B^{L4(r)}\subseteq\mathcal B^{M1}\subseteq\mathcal B^{M1(0)}\subseteq\mathcal B^{M1(\delta)}\subseteq\mathcal B^{M1(\delta')}\subseteq\mathcal B^{M2}\subseteq DD^*,\]
for both the sharp and the outer versions, whenever the smaller rung's model is non-empty.

\emph{Proof of (a).}
1. Each rung is an affine mean model in the sense of Theorem 1′, with \(\theta=(k^{(d)})_d\) (plus \(\bar k\)) restricted to the polyhedron \(\mathcal R_\rho\), and \(m_\theta(x,n,d)=M_{Q^{(d)}+K^{(d)}(\theta)}(x)\). {[}Def; D1(a); lin{]}
2. Apply Theorem 1′. The condition holds \(P_V\)-a.e. iff it holds a.e. given each atom \(D=d\). {[}T1′; cond-D{]} ∎(a)

\emph{Proof of (b).}
1. Within district \(d\), sharp ⊆ outer by 2a Prop. 1 applied to the conditional law, with vertex reduction by 2a Lemma 3(a). {[}cond-D; 2a:P1; 2a:L3{]}
2. Intersecting over \(d\) and with \(\mathcal R_\rho\) preserves the inclusion. {[}1{]}
3. All constraints are linear in \(\theta\), and \(\mathcal R_\rho\) is polyhedral for every rung (an \(\ell_\infty\) ball is a polyhedron). The target is affine in \(\theta\) by Lemma D2. So the bounds are LPs. {[}D2; Def{]}
4. For M1, \(K^{(d)}(k)=K(k)+\operatorname{skew}(\zeta\mathbf 1^\top)\) with \(\zeta\) identified, so the constraints of district \(d\) involve only the common \(k\). {[}D1(b){]} ∎(b)

\emph{Proof of (c).}
1. \textbf{M2 ⊆ DD:} the M2 model class is contained in the (S1)-only class. {[}Def; nest; D0{]}
2. \textbf{M1(\(\delta\)) ⊆ M1(\(\delta'\)) ⊆ M2:} the \(\delta\)-ball is contained in the \(\delta'\)-ball, and M2 has no restriction. {[}Def; nest{]}
3. \textbf{M1 ⊆ M1(0):} under M1, \(k^{(d)}=k(K)+c_d\) by D1(b), so \(\bar k=k(K)\) satisfies the M1(0) condition. {[}D1(b); nest{]}
4. \textbf{L4(\(r\)) ⊆ M1:} L4 is M1 plus a further restriction. {[}Def; nest{]}
5. \textbf{L5 ⊆ L4(\(r\)):} under L5 every row is constant in \(x\). Writing \(\alpha_{c,d}=\alpha_{c,1}+(\alpha_{c,d}-\alpha_{c,1})\) shows L5 is M1 with \(H_c=\alpha_{c,1}\mathbf 1^\top\) and \(\zeta_{c,d}=\alpha_{c,d}-\alpha_{c,1}\), where \(\sum_c\zeta_{c,d}=\mathbf 1-\mathbf 1=0\); and row \(r\) is constant. {[}Def; alg{]}
6. \textbf{L6 ⊆ L5:} L6 is L5 with \(\alpha\) common. {[}Def; nest{]}
7. \textbf{Outer versions:} each step adds constraints to, or restricts the variables of, the same polyhedral system, so the feasible sets and hence the LP ranges are nested. {[}(b); Def{]} ∎

Verified on the 3×3 example (§6): the chain DD ⊇ M2 ⊇ M1(.10) ⊇ M1(.05) ⊇ M1(.02) ⊇ M1 holds for all four targets, and the truth lies in every rung.

\begin{center}\rule{0.5\linewidth}{0.5pt}\end{center}

\hypertarget{pC-closed-forms-and-properties}{%
\subsection{4. Closed forms and properties}\label{pC-closed-forms-and-properties}}

\textbf{Proposition D3 (M1 with two rows).} Let \(R=2\), \(\operatorname{vert}(A_d)=\{(t,1-t):t\in\{l_d,u_d\}\}\) with \(0<l_d\le u_d<1\) (needed: the formula divides by \(t\) and \(1-t\); pure-composition vertices are treated below), and let \(\lambda_c,\rho_c\) be defined from the reference \(Q=Q^{(1)}\) as in 2a Prop. 6. Under M1:
\[\ell_c=\max_d\max_{t\in\{l_d,u_d\}}\Big[\lambda_c(t)-\frac{\zeta_{c,d,1}}{1-t}\Big],\qquad h_c=\min_d\min_{t\in\{l_d,u_d\}}\Big[\rho_c(t)+\frac{\zeta_{c,d,2}}{t}\Big],\]
and \(\Theta^{M1}_O=\{\kappa:\ell\le\kappa\le h,\ \sum_c\kappa_c=0\}\), projected by 2a Lemma 4.

\emph{Derivation.}
1. Row 1 in district \(d\) at \(v=(t,1-t)\): \(e_1^\top(H_c+\zeta_{c,d}\mathbf 1^\top)v=Q_{c,11}t+(Q_{c,12}+\kappa_c)(1-t)+\zeta_{c,d,1}\ge0\). {[}Def M1; \(\mathbf 1^\top v=1\){]}
2. \(\iff\kappa_c\ge-\frac{Q_{c,11}t+Q_{c,12}(1-t)+\zeta_{c,d,1}}{1-t}=\lambda_c(t)-\frac{\zeta_{c,d,1}}{1-t}\). {[}alg; 2a:P6{]}
3. Row 2: \((Q_{c,12}-\kappa_c)t+Q_{c,22}(1-t)+\zeta_{c,d,2}\ge0\iff\kappa_c\le\rho_c(t)+\frac{\zeta_{c,d,2}}{t}\). {[}alg; 2a:P6{]}
4. Intersecting over the vertices and the districts gives the max/min. {[}D3 statement; Thm 2(b){]}
5. Adding \(\sum_c\kappa_c=0\) gives a box intersected with a hyperplane. {[}Def \(\mathcal K\); 2a:L4{]} ∎

Verified against the LP in 126 random column instances.

\textbf{Pure-composition vertices.} If \(u_d=1\) (a poll with only group 1), the row-1 constraint at \(t=1\) reads \(Q_{c,11}+\zeta_{c,d,1}\ge0\). It does not involve \(\kappa_c\): it is a testable restriction, and if it holds it adds no bound on \(\kappa_c\). The term \(t=1\) is therefore dropped from the maximum defining \(\ell_c\). The row-2 constraint at \(t=1\) is \(\kappa_c\le Q_{c,12}+\zeta_{c,d,2}=\rho_c(1)+\zeta_{c,d,2}\), as the formula gives. Symmetrically, if \(l_d=0\), the row-2 constraint at \(t=0\) becomes the testable \(Q_{c,22}+\zeta_{c,d,2}\ge0\) and is dropped from the minimum defining \(h_c\), while row 1 gives \(\kappa_c\ge\lambda_c(0)-\zeta_{c,d,1}=-Q_{c,12}-\zeta_{c,d,1}\). With \(\zeta=0\) the same rule applies to 2a Prop. 6.
\emph{Derivation.} Set \(t=1\) (or \(t=0\)) in steps 1 and 3 above before dividing; the coefficient of \(\kappa_c\) vanishes in exactly one of the two rows. {[}alg{]}

\textbf{Reading the formula.} A district raises the lower bound when its group-1 baseline for column \(c\) is low (\(\zeta_{c,d,1}<0\)) and it contains polls that are mostly group 1 (\(t\) near 1, where \(1/(1-t)\) is large). In such polls group 1's rate is nearly observed, the shared contextual effect carries that information to every district, and so \textbf{one riding with lopsided polls and an extreme baseline can tighten the province-wide bound.} The upper bound behaves symmetrically, through group 2.

\textbf{Proposition D4 (when dummies matter).}
(a) If \(\zeta\equiv0\), then \(\Theta^{M1}_O=\Theta^{M0}_O(\operatorname{conv}\bigcup_dA_d)\): dummies neither help nor hurt identification.
(b) Under M1 with \(P(X\in A_D)=1\), write \(Z_d:=[\zeta_{c,d,r}]_{rc}\). Here M0 means the pooled model of 2a, a restriction on \(E[\beta\mid X,N]\), which does not involve \(D\). Then
\[E[\beta\mid X,N]=M_H(X)+E[Z_D\mid X,N],\]
and M0 holds (for some \(H'\)) iff there is a \(G\) with \(E[Z_D\mid X,N]=M_G(X)\) a.s.; that is, iff \(E[Z_D\mid X,N]\) is a.s. equal to a function of \(X\) alone that is affine on \(\operatorname{supp}P_X\). In particular \(D\perp(X,N)\) suffices, and so does \(D\perp N\mid X\) together with \(E[Z_D\mid X]\) affine. \textbf{\(D\perp X\) alone does not suffice} (Remark D4.1). When M0 holds this way, its (S3) on \(A=\operatorname{conv}\operatorname{supp}P_X\) is automatic by 2a Lemma 3(c). For any other pooled region \(A\) it remains a separate condition.

\emph{Proof of (a).}
1. \(\zeta\equiv0\) gives \(Q^{(d)}=Q\) for all \(d\), so the pooled regression function is also \(X^\top QX\). {[}D1(b); CE-tower{]}
2. The M1 constraints then require \((Q_c+K_c)v\ge0\) for \(v\in\bigcup_d\operatorname{vert}(A_d)\). By linearity that is equivalent to the constraint on \(\operatorname{conv}\bigcup_dA_d\), i.e.~the M0 outer set. {[}Thm 2(b); 2a:L3(a){]} ∎(a)

\emph{Proof of (b).}
1. \(E[\beta\mid X,N]=E\big[E[\beta\mid X,N,D]\mid X,N\big]=E[M_H(X)+Z_D\mid X,N]=M_H(X)+E[Z_D\mid X,N]\). {[}CE-tower, \(\sigma(X,N)\subseteq\sigma(V)\); S2-d with M1; 0.1, \((Z\mathbf 1^\top)x=Z\); CE-pull, \(M_H(X)\) is \(\sigma(X,N)\)-measurable{]}
2. On the simplex, an affine function of \(x\) equals a homogeneous linear one: \(a+b^\top x=(a\mathbf 1+b)^\top x\). Every linear map \(\mathbb R^R\to\mathbb R^{R\times C}\) is \(x\mapsto M_G(x)\) for some \(G\). {[}alg; \(\mathbf 1^\top x=1\){]}
3. ``If'': if \(E[Z_D\mid X,N]=M_G(X)\) a.s., then \(E[\beta\mid X,N]=M_{H+G}(X)\), which is (S2) of 2a with \(H'=H+G\). {[}1--2{]}
4. ``Only if'': if M0 holds with \(H'\), then \(E[Z_D\mid X,N]=M_{H'}(X)-M_H(X)=M_{H'-H}(X)\) a.s., a function of \(X\) only, linear in \(X\). {[}1; 2a Def of (S2){]}
5. If \(D\perp(X,N)\), then \(E[Z_D\mid X,N]=E[Z_D]\), a constant, hence affine in \(X\). If \(D\perp N\mid X\), then \(E[Z_D\mid X,N]=E[Z_D\mid X]\). {[}indep.{]}
6. (S3) on \(\operatorname{conv}\operatorname{supp}P_X\): under M0 with \(H'\), (S1)--(S2) hold, so 2a Lemma 3(c) applies. {[}2a:L3(c){]} ∎(b)

\textbf{Remark D4.1 (\(D\perp X\) is not enough; corrected after review).} An earlier version conditioned on \(X\) only and claimed \(D\perp X\) suffices. The size-weighted estimand of 2a Lemma 2 uses \(E[\beta\mid X,N]\), not \(E[\beta\mid X]\). Counterexample: two districts with the same distribution of \(X\), district 2 has larger precincts, and \(\zeta_{\cdot,2}\neq0\). Then \(E[Z_D\mid X,N]=Z_2\,P(D=2\mid N)\) varies with \(N\), so (S2) of M0 fails. The \(X\)-only pooled model of \(\Theta_X\) still holds, since \(E[Z_D\mid X]=Z_2\pi_2\) is constant. But the identity \(E[NX_r\beta_{rc}]=E[NX_rM_{H'}(X)_{rc}]\) behind Lemma 2 fails whenever size is correlated with the district shift, so pooled bounds on \(B^*\) can be biased. The violation is generally visible in the data: \(E[T\mid X,N]\) then depends on \(N\) unless \(x^\top Z_2=0\), which (T1) of 2a Theorem 1 rules out. It is checked numerically in \texttt{verify\_phase2b.py}.

\emph{Meaning:} for the conditional-mean model, dummies are needed exactly when the expected district baseline, given composition and size, either varies non-affinely with composition (e.g.~stronghold districts with different baselines) or varies with precinct size. The pooled \emph{outer} bounds additionally need probability validity on the pooled region, which is automatic for the population hull and an assumption otherwise. When M0 fails, the pooled quadratic fit is a projection with no bounding guarantee. In the 2×2 example (P1 §7.3) the pooled outer set is empty. In the 3×3 example it happens to cover the truth, which is not guaranteed.

\textbf{Proposition D5 (M2 adds district intervals).}
\[\mathcal B^{M2}_O=\sum_d\big[\varphi^{(d)}_L,\varphi^{(d)}_U\big],\]
where \(\varphi^{(d)}\) is district \(d\)'s term in Lemma D2, bounded over \(\Theta_O(P_d;A_d)\) (2a Prop. 7 within the district). The same holds for the sharp version.
\emph{Proof.}
1. Under M2, the feasible set is the product \(\prod_d\Theta^{(d)}\). {[}Thm 2; Def M2{]}
2. \(\varphi=\sum_d\varphi^{(d)}(k^{(d)})\) is separable. {[}D2{]}
3. So the min (max) of the sum is the sum of the mins (maxes). {[}sep{]} ∎

Verified numerically. As a consequence, M2 widths add across districts. This is why M2 is usually little tighter than DD unless districts have wide, informative composition ranges.

\textbf{Proposition D6 (bounded deviation interpolates).}
(a) \(\delta\mapsto\mathcal B^{M1(\delta)}_O\) is nested (shrinks as \(\delta\) falls).
(b) For \(\delta\ge\bar\delta:=\max_d\max_{k\in\Theta^{(d)}_O}\lVert k-c_d\rVert_\infty\) (finite), \(\mathcal B^{M1(\delta)}_O=\mathcal B^{M2}_O\).
(c) The endpoints are continuous, piecewise-linear functions of \(\delta\) wherever the set is non-empty.
(d) \(\mathcal B^{M1(0)}_O=\mathcal B^{M1}_O\) computed with the same \(Q^{(d)}\) and \(\zeta\).
\emph{Proof.}
1. (a) follows from Theorem 2(c), step 2. {[}Thm 2(c){]}
2. (b): each \(\Theta^{(d)}_O\) is bounded, so \(\bar\delta<\infty\). Take \(\bar k=0\); every point of \(\prod_d\Theta^{(d)}_O\) then satisfies the \(\delta\)-condition, so the restriction is slack. {[}2a:L3(b); Def{]}
3. (c): \(\delta\) enters only the right-hand side of the LP, linearly. {[}LP-par{]}
4. (d): at \(\delta=0\) the M1(\(\delta\)) rows become \(k^{(d)}=\bar k+c_d\), which are the M1 equalities. The objective is identical. {[}Def; D2; Thm 2(b){]} ∎

Note on (d): M1(0) does not impose the testable restrictions of D1(b). So it stays defined even when the data reject M1's symmetric part; it is ``M1 without its testable implication''. It is therefore a \textbf{sensitivity benchmark}, not a substantive rung on a par with M1: it separates a failure of M1's testable implications from uncertainty about the unidentified contextual effects. \emph{(Emphasis added after review.)}

\textbf{Proposition D7 (the most restrictive rungs).}
(a) \textbf{L5} point-identifies the target:
\[B^*_{rc}=\sum_d\frac{\mu^{(d)}_r}{\mu_r}\,\alpha_{c,d,r},\qquad\alpha_{c,d}=\operatorname{diag}Q^{(d)}_c.\]
Equivalently, \(\alpha_{\cdot,d}\) is the Goodman regression of \(T\) on \(X\) without intercept, within district \(d\). If L5 holds, this point lies in every larger rung's bound.
(b) \textbf{L4(\(r\))} pins \(K_{c,rs}=Q_{c,rr}-Q_{c,rs}\) (reference district) for all \(s\ne r\) and all \(c\), i.e.~\((C-1)(R-1)\) of the \(d_0\) coordinates. For \(R=2\) the target is point-identified.
\emph{Proof of (a).}
1. Under L5, \(E[\beta_{rc}\mid V]=\alpha_{c,D,r}\), so \(B^*_{rc}=\sum_dE[NX_r1\{D=d\}]\alpha_{c,d,r}/\mu_r\). {[}CE-tower; Def L5{]}
2. \(\alpha_{c,d}=\operatorname{diag}Q^{(d)}_c\). {[}D1(c){]}
3. Within district \(d\), \(E[T_c\mid X,d]=\sum_rX_r\alpha_{c,d,r}\). {[}2a:L1(a) proof steps 1--2; Def L5{]}
4. The point lies in every larger rung's bound. {[}Thm 2(c){]} ∎(a)

\emph{Proof of (b).}
1. \(H_{c,rs}=H_{c,rr}\), and \(H_{c,rr}=Q_{c,rr}\) because \(K_{c,rr}=0\). {[}Def L4; skew{]}
2. \(H_{c,rs}=Q_{c,rs}+K_{c,rs}\), hence \(K_{c,rs}=Q_{c,rr}-Q_{c,rs}\). {[}Def; alg{]}
3. The pairs containing \(r\) number \(R-1\), in each of \(C-1\) free columns. {[}alg{]}
4. For \(R=2\) that is all \(d_0=C-1\) coordinates. {[}alg{]} ∎(b)

Verified: under a DGP with zero contextual effects, the L5 point equals the truth and lies inside the M1 bound for all four 3×3 targets.

\textbf{Remark D7.1 (L5 as a diagnostic).} When contextual effects are not zero, the L5 ``point'' is a misspecified estimate. In the 3×3 example it is 0.870 for loyalty (truth 0.836): inside M1 but wrong. Reporting it next to the bounds shows how much the zero-contextual-effect assumption is doing.

\begin{center}\rule{0.5\linewidth}{0.5pt}\end{center}

\hypertarget{pC-remarks-on-scope}{%
\subsection{5. Remarks on scope}\label{pC-remarks-on-scope}}

\textbf{Remark 6.1 (estimation regimes).} Everything above is population identification. With about 10 polls per riding (BC stage 2), each \(Q^{(d)}\) and \(\zeta_d\) is noisy. The plug-in intersection in Theorem 2(b) then takes a max/min over many noisy constraints and is biased inward (P1 §7.1: 72\% coverage). Remark D1.1 (fewer parameters per district under M1) and M1(\(\delta\)) with \(\delta>0\) (which softens the max) both help. A valid procedure needs Phase 3: precision-corrected intersection bounds (Chernozhukov, Lee \& Rosen 2013) and inference for LP values with non-unique solutions. \emph{(Revised: an earlier version also listed cluster-robust variances. Inference is conditional on the design and the district effects, with no clustered variance; P3 §0.4.)}

\textbf{Remark 6.2 (periods as districts: the DiD link).} Take ``districts'' to be election periods: the same counties observed twice, \(D\in\{\text{pre},\text{post}\}\).
* \textbf{M1} then says the contextual effects are the same in both periods, with only baselines shifting, which is exactly the ``contextual effects unchanged'' assumption in \texttt{ASSUMPTION\_MENU.md} item 18.
* Lemma D1(b) identifies the baseline shifts \(\zeta\), and Proposition D8 below gives a sufficient moment condition (by cancellation) and the necessary and sufficient condition given the identified set, under which the change \(B^*_\text{post}-B^*_\text{pre}\) is point-identified. That recovers the owner's earlier result as a special case of M1.
* Theorem 2(b) adds that the \emph{levels} in both periods are bounded by the intersection of the two periods' sets.
* \textbf{M1(\(\delta\))} is the bounded-change version (item 19).
* The overidentifying restrictions of D1(b) are the testable part, ``equal quadratic coefficients across periods'' (menu note to item 18).

\textbf{Proposition D8 (the DiD under M1 with periods as districts; promoted from Remark 6.2 after review).} Let \(D\in\{\text{pre},\text{post}\}\), assume M1, (S1)--(S2-d), \(P(X\in A_D)=1\) and (A5-M1). For period \(t\) write \(\mu^{[t]}_{rs}=E[NX_rX_s\mid D=t]\), \(\mu^{[t]}_r=E[NX_r\mid D=t]\) and \(B^*_{rc,t}=E[NX_r\beta_{rc}\mid D=t]/\mu^{[t]}_r\). Then
\[B^*_{rc,\text{post}}-B^*_{rc,\text{pre}}=\text{(identified)}+\sum_{s\ne r}K_{c,rs}\Big(\frac{\mu^{[\text{post}]}_{rs}}{\mu^{[\text{post}]}_r}-\frac{\mu^{[\text{pre}]}_{rs}}{\mu^{[\text{pre}]}_r}\Big).\]
Write \(\Delta q_{rs}:=\mu^{[\text{post}]}_{rs}/\mu^{[\text{post}]}_r-\mu^{[\text{pre}]}_{rs}/\mu^{[\text{pre}]}_r\) and \(L(K):=\sum_{s\ne r}K_{c,rs}\Delta q_{rs}\), and let \(\Theta\) be the identified set for \(K\).
(a) \textbf{Identification by cancellation.} The change does not depend on \(K\) for every \(K\in\mathcal K\) iff \(\Delta q_{rs}=0\) for all \(s\ne r\). This holds in particular when the size-and-composition moments are equal across the periods, \(\mu^{[\text{post}]}=\mu^{[\text{pre}]}\). It is a condition on the \(NX_rX_s\) and \(NX_r\) moments, not merely on similar-looking distributions of vote shares.
(b) \textbf{Identification given the feasible set.} The change is point-identified iff \(\max_{K\in\Theta}L(K)-\min_{K\in\Theta}L(K)=0\). This is necessary and sufficient given \(\Theta\). It holds under (a), and also when \(\Theta\) fixes \(L\) although \(\Delta q\ne0\). Example: 2×2 with \(Q_{1,12}=0\) and LCE on the whole simplex; probability validity at the vertices forces \(K_{1,12}=0\), so \(\Theta=\{0\}\) and the change is point-identified whatever \(\Delta q\) is.

\emph{(Corrected after review, round 2, Nov 2026: the proposition said the change is point-identified, ``whatever the identified set for \(K\)'', \textbf{iff} \(\Delta q=0\). As a statement about identification that is false (the example above, due to the Phases 1--3 referee); \(\Delta q=0\) is sufficient, and necessary only for invariance to every \(K\in\mathcal K\). Check: \texttt{verify\_phase2b.py}, ``D8b''.)}
\emph{Proof.}
1. Within period \(t\), Lemma D2 applied to the conditional law gives \(B^*_{rc,t}=\big[\sum_sQ^{(t)}_{c,rs}\mu^{[t]}_{rs}+\sum_{s\ne r}K^{(t)}_{c,rs}\mu^{[t]}_{rs}\big]/\mu^{[t]}_r\). {[}D2; cond-D{]}
2. Under M1, \(Q^{(t)}\) and \(\operatorname{skew}(\zeta_{\cdot,t}\mathbf 1^\top)\) are identified, and \(K^{(t)}=K+\operatorname{skew}(\zeta_{\cdot,t}\mathbf 1^\top)\) with \(K\) common. {[}D1(b), D1(d){]}
3. Subtracting the two periods, every term except \(\sum_{s\ne r}K_{c,rs}(\cdot)\) is identified. {[}1--2; lin{]}
4. (a): \(L\) is a linear form in the free entries of \(K\in\mathcal K\) with coefficients \(\Delta q_{rs}\); it vanishes on all of \(\mathcal K\) iff every coefficient is zero. {[}3; alg; the entries \(K_{c,rs}\), \(s\ne r\), are free in \(\mathcal K\) for \(c<C\){]}
5. (b): by 3 the change equals an identified constant plus \(L(K)\), and its identified set is \(L(\Theta)\) shifted by that constant. A set of reals is a point iff its max and min agree. {[}3; Def identified set{]} ∎
\emph{Reading.} The change in ecological transition rates between the periods is point-identified when the contextual effects are unchanged across periods (M1) and the electorate's size-weighted composition moments are the same in both periods (a), or when the identified set happens to fix \(L\) (b). Otherwise the change is only bounded; the bound is the projection of the common identified set for \(K\) (Theorem 2(b)) through the coefficients above. This is a statement about the difference of two ecological rates. A causal (difference-in-differences) reading needs further structure (2d). \emph{(Reading corrected after review, round 2.)}

\textbf{Remark 6.3 (other groupings).} Nothing above uses the nesting being geographic. Groups can be states (with counties as units), demographic strata of units, or time.

\begin{center}\rule{0.5\linewidth}{0.5pt}\end{center}

\hypertarget{pC-worked-examples-noise-free-identification-only}{%
\subsection{6. Worked examples (noise-free; identification only)}\label{pC-worked-examples-noise-free-identification-only}}

\textbf{2×2 (P1 §7.3, two ridings).} Province-wide P-voter loyalty, truth 0.819.

\begin{longtable}[]{@{}ll@{}}
\toprule\noalign{}
Rung & Bound \\
\midrule\noalign{}
\endhead
\bottomrule\noalign{}
\endlastfoot
DD & {[}0.453, 0.879{]} \\
M0 pooled & empty (reverts to DD) \\
M2 & {[}0.563, 0.870{]} \\
M1 & {[}0.618, 0.866{]} \\
L5 point & 0.858 (misspecified: slopes are 0.10) \\
\end{longtable}

\textbf{3×3 (six districts with different composition ranges and baselines; common \(H\) from 2a §8).} Each bound is {[}lower, upper{]}.

\begin{longtable}[]{@{}
  >{\raggedright\arraybackslash}p{(\columnwidth - 18\tabcolsep) * \real{0.1000}}
  >{\raggedright\arraybackslash}p{(\columnwidth - 18\tabcolsep) * \real{0.1000}}
  >{\raggedright\arraybackslash}p{(\columnwidth - 18\tabcolsep) * \real{0.1000}}
  >{\raggedright\arraybackslash}p{(\columnwidth - 18\tabcolsep) * \real{0.1000}}
  >{\raggedright\arraybackslash}p{(\columnwidth - 18\tabcolsep) * \real{0.1000}}
  >{\raggedright\arraybackslash}p{(\columnwidth - 18\tabcolsep) * \real{0.1000}}
  >{\raggedright\arraybackslash}p{(\columnwidth - 18\tabcolsep) * \real{0.1000}}
  >{\raggedright\arraybackslash}p{(\columnwidth - 18\tabcolsep) * \real{0.1000}}
  >{\raggedright\arraybackslash}p{(\columnwidth - 18\tabcolsep) * \real{0.1000}}
  >{\raggedright\arraybackslash}p{(\columnwidth - 18\tabcolsep) * \real{0.1000}}@{}}
\toprule\noalign{}
\begin{minipage}[b]{\linewidth}\raggedright
Target
\end{minipage} & \begin{minipage}[b]{\linewidth}\raggedright
Truth
\end{minipage} & \begin{minipage}[b]{\linewidth}\raggedright
DD
\end{minipage} & \begin{minipage}[b]{\linewidth}\raggedright
M0 pooled*
\end{minipage} & \begin{minipage}[b]{\linewidth}\raggedright
M2
\end{minipage} & \begin{minipage}[b]{\linewidth}\raggedright
M1(.10)
\end{minipage} & \begin{minipage}[b]{\linewidth}\raggedright
M1(.05)
\end{minipage} & \begin{minipage}[b]{\linewidth}\raggedright
M1(.02)
\end{minipage} & \begin{minipage}[b]{\linewidth}\raggedright
M1
\end{minipage} & \begin{minipage}[b]{\linewidth}\raggedright
L5 point
\end{minipage} \\
\midrule\noalign{}
\endhead
\bottomrule\noalign{}
\endlastfoot
Loyalty P→P & 0.836 & {[}0.167, 0.953{]} & {[}0.423, 0.898{]} & {[}0.391, 0.916{]} & {[}0.392, 0.916{]} & {[}0.397, 0.914{]} & {[}0.412, 0.899{]} & {[}0.430, 0.880{]} & 0.870 \\
Defection P→O & 0.095 & {[}0.000, 0.610{]} & {[}0.026, 0.331{]} & {[}0.024, 0.404{]} & {[}0.024, 0.403{]} & {[}0.024, 0.390{]} & {[}0.024, 0.367{]} & {[}0.042, 0.339{]} & 0.061 \\
Demobilisation P→abstain & 0.068 & {[}0.000, 0.551{]} & {[}0.029, 0.296{]} & {[}0.024, 0.350{]} & {[}0.024, 0.339{]} & {[}0.024, 0.323{]} & {[}0.030, 0.309{]} & {[}0.044, 0.295{]} & 0.069 \\
Mobilisation N→P & 0.135 & {[}0.001, 0.685{]} & {[}0.038, 0.410{]} & {[}0.045, 0.470{]} & {[}0.045, 0.457{]} & {[}0.051, 0.444{]} & {[}0.070, 0.428{]} & {[}0.089, 0.413{]} & 0.122 \\
\end{longtable}

*M0 is misspecified here (district baselines vary with composition). It happens to cover the truth and to be narrower than M2, but Prop. D4 gives no guarantee for it.

What the table shows:
* The ladder shrinks monotonically and always contains the truth.
* Moving from M2 to M1 narrows loyalty from width 0.525 to 0.450.
* δ traces a smooth path between the two rungs.
* The baseline shifts here are small (SD 0.04). The gain grows with more heterogeneous baselines and more lopsided polls (Prop. D3 reading).

\begin{center}\rule{0.5\linewidth}{0.5pt}\end{center}

\hypertarget{pC-verification-log}{%
\subsection{7. Verification log}\label{pC-verification-log}}

\texttt{python3\ -I\ verify\_phase2b.py} (21 PASS, 0 FAIL) checks:
* Lemma D1: recovery of \(\zeta\), the overidentifying restrictions and the \(K^{(d)}\) relation.
* Theorem 2: ladder nesting for four targets; the truth in every valid rung.
* Prop. D6: monotone in \(\delta\); equals M2 for large \(\delta\).
* Prop. D5: M2 equals the sum of district bounds.
* Prop. D4(a): \(\zeta=0\) makes dummies equal to pooled.
* Prop. D7: the L5 point under a true L5.
* Prop. D3: closed form against the LP (126 instances).
* Prop. D0: per-cell Fréchet extremes attained (300 tables).
* The 2×2 example reproduces P1 §7.3.
* Corrections after review (each check fails on the earlier statement):
* Prop. D4(b): with \(D\perp X\) but \(D\) correlated with \(N\), \(E[\beta\mid X,N]\) depends on \(N\) while \(E[\beta\mid X]\) is affine, and the pooled size-weighted rate is biased. With \(D\perp(X,N)\) it is exact.
* P1 (D2) and P3 Prop. P3.6(b): the rank gap between the M2 and M1 designs is \((D^\star-1)R(R-1)/2\) per column, and column \(C\)'s restrictions are implied. That gives \((D^\star-1)(C-1)R(R-1)/2\) in total (three table shapes).
* Scope additions after review:
* Lemma D1(d): the stacked M1 design has full rank with A5 in one district and three affinely independent points in the others, and is rank-deficient with three distinct but collinear points;
* Prop. D3 with pure-composition vertices: the modified closed form matches the LP (195 instances).

Not machine-checked: the measure-theoretic content of Theorem 1′, which rests on 2a Theorem 1.

\hypertarget{pC-open-items}{%
\subsection{8. Open items}\label{pC-open-items}}

\begin{itemize}
\tightlist
\item
  \textbf{Phase 3:} noisy-intersection inference (Remark 6.1); cluster dependence; the finite-population gap; a test of D1(b)'s overidentifying restrictions; sensitivity in \(\delta\) (the breakdown value of a conclusion).
\item
  \textbf{Not yet done:} M1′ (contextual effects depending on district covariates) is a polyhedral restriction like the others and fits Theorem 2 directly. It is left until a covariate set is chosen.
\item
  \textbf{Strict-gain characterisation for \(R\ge3\):} when \(\Theta^{M1}\) is strictly smaller than the best single district's set, beyond the \(R=2\) reading of Prop. D3.
\end{itemize}


\clearpage
\section{PHASE2c: monotone ecological inference in R$\times$C}\label{app:phase2c}
\textit{Source: \texttt{docs/extension/PHASE2c\_monotone.md}.}\par
% GENERATED by tools/make_appendices.py from docs/extension/PHASE2c_monotone.md
% Do not edit by hand; edit the source and regenerate.
\emph{Builds on P1, 2a, 2b, P3 and the Phase-1 proposal \texttt{PHASE1b\_monotone\_did.md} (cited as 1b). It follows Elzayn, Goldin, Guage, Ho \& Morton, ``Monotone Ecological Inference'' (EGGHM; arXiv v1 numbering, see 1b §0). Every proof step ends with a justification tag: those of 2a §0.3, 2b §0.3 and P3 §0.3, plus §0.3 below. Checks: \texttt{verify\_phase2c.py} (output \texttt{verify\_phase2c\_output.txt}, 29/29 PASS). Owner decisions applied (1b §D):}
* \emph{finite-population bounds for the sign rungs, \textbf{sharp for continuous unit tables given the margins} (corrected after review, Batch B: earlier ``exact''; see Remark M2.1 for integer counts, decision U1(a));}
* \emph{one-vs-rest (W1) as the default, with W2 as an option;}
* \emph{contextual reinforcement available as an option;}
* \emph{API: planned here as \texttt{monotone=} plus \texttt{model=}. The shipped M-SIGN module uses \texttt{restrictions\ =\ list(assoc\ =\ ...,\ cr\ =\ ...)} instead (5a §5). N1, N2 and local monotonicity are not implemented.}

\emph{Batch B review (Nov 2026). Two external referee reports found errors in M2 (integer tables), M3(c)--(d), M5(c) (district rungs), N2.2 and N2.3, and §6 (binning, degree selection). All are corrected below and marked; the ledger is \texttt{reviews/B\_PI\_RESPONSE\_PLAN.md} §2 (rows P-B1 to P-B14, P-B28, P-B29, P-B31, P-B32, P-B34, P-B35; P-E1, P-E3, P-E8, P-E11 to P-E13, P-E23, P-E27, P-E28). Each correction has a check in \texttt{verify\_phase2c.py}, tagged by what it shows: \textbf{{[}F{]}} fails on an earlier statement and passes on the corrected one; \textbf{{[}V{]}} verifies a corrected or new statement (nothing earlier fails); \textbf{{[}I{]}} illustrates or strengthens a statement that was already true. }(Corrected after review, round 2: this note earlier said every check fails on the old statement, which overstated the evidence.)* Status labels follow \texttt{REVIEW\_RESPONSE\_PLAN.md}.*

\emph{Round 2 (Nov 2026): the partial-identification referee's second report (\texttt{reviews/B\_partial\_identification\_report\_round2.md}) asked for four corrections (the equal-rates witness, the premise of N2.3(d), and wording in M2(d), M3(d), M4 and Lemma N2.4); they are marked ``round 2'' below.}

\begin{center}\rule{0.5\linewidth}{0.5pt}\end{center}

\hypertarget{pD-setting}{%
\subsection{0. Setting}\label{pD-setting}}

\hypertarget{pD-finite-population-the-n1-rungs}{%
\subsubsection{0.1 Finite population (the N1 rungs)}\label{pD-finite-population-the-n1-rungs}}

\begin{itemize}
\tightlist
\item
  The units \(i=1,\dots,p\) are fixed (the design is conditioned on, as in P3 (P3-S2)). Unit weights are \(p_i=N_i/\sum_jN_j\).
\item
  \(X_i\in\Delta^{R-1}\), \(T_i\in\Delta^{C-1}\), and the realised rates satisfy \(\beta_i\in S(X_i,T_i)\) (2a §0.2).
\item
  \(E_p[f]:=\sum_ip_if_i\), and \(\mu_r:=E_p[X_r]>0\).
\item
  Row weights \(w^r_i:=p_iX_{ir}/\mu_r\). \(\operatorname{Cov}_w\) is the covariance under weights \(w\).
\item
  \textbf{Target} (P3's actual rate): \(B_{rc}=E_p[X_r\beta_{rc}]/\mu_r=\sum_iw^r_i\beta_{irc}\).
\item
  \textbf{Continuous-rate reading.} The realised rates \(\beta_i\) are real numbers: the cell masses \(N_iX_{ir}\beta_{irc}\) need not be whole numbers of people. All N1 results below are for this reading. With counted ballots (integer cell counts) attainability and refutation can change (Remark M2.1). \emph{(Added after review, Batch B.)}
\end{itemize}

\hypertarget{pD-superpopulation-the-k-rows-and-n2-rungs}{%
\subsubsection{0.2 Superpopulation (the K-rows and N2 rungs)}\label{pD-superpopulation-the-k-rows-and-n2-rungs}}

\begin{itemize}
\tightlist
\item
  \(m(v)=E[\beta\mid V=v]\), and P3's model-implied rate is \(\tilde B_{rc}=\sum_iw^r_im_{rc}(V_i)\).
\item
  A superscript \(^*\) on an association means it is computed with \(m\) in place of \(\beta\).
\end{itemize}

\hypertarget{pD-additional-tags}{%
\subsubsection{0.3 Additional tags}\label{pD-additional-tags}}

\begin{longtable}[]{@{}
  >{\raggedright\arraybackslash}p{(\columnwidth - 2\tabcolsep) * \real{0.5000}}
  >{\raggedright\arraybackslash}p{(\columnwidth - 2\tabcolsep) * \real{0.5000}}@{}}
\toprule\noalign{}
\begin{minipage}[b]{\linewidth}\raggedright
Tag
\end{minipage} & \begin{minipage}[b]{\linewidth}\raggedright
Meaning
\end{minipage} \\
\midrule\noalign{}
\endhead
\bottomrule\noalign{}
\endlastfoot
{[}EGGHM:x{]} & result x of EGGHM, arXiv v1 (2504.14752). Their identities hold for any bounded outcome; checked again in 1b V1 \\
{[}Cheb{]} & Chebyshev's association inequality: for nondecreasing \(f,g:\mathbb R\to\mathbb R\) and any probability measure on \(\mathbb R\), \(\operatorname{Cov}(f(Z),g(Z))\ge0\) (e.g.~Boucheron, Lugosi \& Massart 2013, Thm 2.14†) \\
{[}ext{]} & a function that is monotone on a finite subset of \(\mathbb R\) extends to a monotone function on \(\mathbb R\) (piecewise-constant extension) \\
{[}conn{]} & a continuous image of a connected set is connected; a connected subset of \(\mathbb R\) is an interval \\
{[}CE-cvx{]}, {[}CE-lin{]}, {[}CE-pull{]} & as in 2a §0.3 \\
\end{longtable}

\begin{center}\rule{0.5\linewidth}{0.5pt}\end{center}

\hypertarget{pD-the-associations-in-rc-and-their-identities}{%
\subsection{1. The associations in R×C and their identities}\label{pD-the-associations-in-rc-and-their-identities}}

Fix a column \(c\) and drop it from the notation.

\textbf{Definition 1.}
1. \textbf{Between, pairwise:} \(\delta_B(r,s):=E_p[X_rX_s(\beta_r-\beta_s)]\) for \(r\ne s\). It is antisymmetric.
2. \textbf{Between, one-vs-rest:} \(\delta^r_B:=\sum_{s\ne r}\delta_B(r,s)\).
3. \textbf{Within, W2 (within each original group, pooled into one sum):} \(\Delta:=\sum_s\mu_s\operatorname{Cov}_{w^s}(X,\beta_s)\in\mathbb R^R\). A sign on \(\Delta_r\) signs \textbf{one sum} of row-specific covariances; the summands need not share its sign. \emph{(Corrected after review, Batch B: earlier labelled ``each group separately''.)}
4. \textbf{Within, W1 (one-vs-rest; EGGHM §4 end).} Let \(\beta^{\neg r}_i:=\sum_{s\ne r}X_{is}\beta_{is}/(1-X_{ir})\), the rate of everyone outside row \(r\) (any value if \(X_{ir}=1\), where its weight is 0), and \(w^{\neg r}_i:=p_i(1-X_{ir})/(1-\mu_r)\). Then
\[\delta^r_W:=\mu_r\operatorname{Cov}_{w^r}(\beta_r,X_r)+(1-\mu_r)\operatorname{Cov}_{w^{\neg r}}(\beta^{\neg r},X_r).\]

\textbf{Remarks on Definition 1} \emph{(added after review, Batch B)}.
* \textbf{Weights.} \(\delta^r_B\) is the A-minus-rest gap averaged with weights \(p_iX_{ir}(1-X_{ir})\): pure units get zero weight, and a large mixed unit can dominate. It is not an ordering in every unit, nor against each rival. The pairwise \(\delta_B(r,s)\) uses \(p_iX_{ir}X_{is}\).
* \textbf{W1 versus W2.} In W1 the predictor in \textbf{both} covariances is \(X_r\) (row \(r\)'s share), and the complement's rate is that of a changing mix of the other rows; either covariance can be negative when their sum is positive. W2 conditions on the original rows. The two can have \textbf{opposite signs} while every group's rate rises. Example (EI referee): rows A, B, C, two units of equal size, shares \((.2,.7,.1)\) and \((.6,.05,.35)\), rates \((.40,.80,.10)\) and \((.42,.82,.12)\): every rate rises by 2 points, yet \(\delta^A_W=-0.02633\) and \(\Delta_A=+0.00110\) (check C14). A direction for each named group needs separate restrictions.
* \textbf{Goodman bias.} \(g-B=M^{-1}\Delta\) (Lemma M1(a)): signing \(\Delta_r\) does not in general sign the bias of row \(r\)'s Goodman coefficient.
* \textbf{Scope relative to EGGHM.} The one-vs-rest collapse W1 is EGGHM §4's suggestion. W2, the pairwise weighting, and CR with those pairings are extensions with their own proofs here (M1, M4); they are not EGGHM results. Covariances also reflect sorting and other between-unit differences; a causal reading needs a separate design.

\textbf{Estimators} (functions of the data only):
* Goodman's no-intercept regression: \(g:=M^{-1}E_p[XT_c]\), with \(M:=E_p[XX^\top]\).
* The neighbourhood model: \(NM_r:=E_p[X_rT_c]/\mu_r\).
* The one-vs-rest regression with intercept: \(ER^r:=\operatorname{Cov}_p(T_c,X_r)/\operatorname{Var}_p(X_r)\).
* The one-vs-rest gap: \(D^r:=B_r-B_{\neg r}\), with \(B_{\neg r}:=\sum_{s\ne r}\mu_sB_s/(1-\mu_r)\).

\textbf{Lemma M1 (identities; accounting identity only).} Assume {[}A1{]}. Part (a)'s formula for \(g\) needs \(M\) invertible (the \(X_i\) span \(\mathbb R^R\)); \(\mathbf 1^\top\Delta=0\), (b) and (d) need nothing more; (c) needs \(\operatorname{Var}_p(X_r)>0\) for \(ER^r\) and \(0<\mu_r<1\). Where \(X_{ir}=1\), give the absent rest group any bounded value (its weight is 0). \emph{(Corrected after review, Batch B: earlier ``Assume \(M\) is invertible'' for every part.)}
(a) \(\mathbf 1^\top\Delta=0\) and \(g-B_{\cdot c}=M^{-1}\Delta\).
(b) \(NM_r-B_r=-\delta^r_B/\mu_r\).
(c) The collapsed population \((X_r,\beta_r,\beta^{\neg r})\) is a valid two-group population for column \(c\). Hence
\[ER^r-D^r=\frac{\delta^r_W}{\operatorname{Var}_p(X_r)},\qquad (NM_r-NM_{\neg r})-D^r=-\frac{\delta^r_B}{\mu_r(1-\mu_r)},\]
where \(NM_{\neg r}=E_p[(1-X_r)T_c]/(1-\mu_r)\).
(d) \(B\), \(\delta_B(r,s)\), \(\delta^r_B\), \(\Delta\) and \(\delta^r_W\) are affine functions of \(\beta=(\beta_i)_i\), with coefficients that depend only on the design.

\emph{Proof of (a).}
1. \(T_{ic}=\sum_sX_{is}\beta_{is}\). {[}A1{]}
2. \(E_p[XT_c]=\sum_sE_p[XX_s\beta_s]\). {[}1; lin{]}
3. \(E_p[XX_s\beta_s]=\mu_sE_{w^s}[X\beta_s]=\mu_s\big(E_{w^s}[X]B_s+\operatorname{Cov}_{w^s}(X,\beta_s)\big)\). {[}Def \(w^s\); Def Cov; \(B_s=E_{w^s}\beta_s\){]}
4. \(\mu_sE_{w^s}[X]=E_p[XX_s]=Me_s\). {[}Def{]}
5. So \(E_p[XT_c]=MB_{\cdot c}+\Delta\), and \(g=B_{\cdot c}+M^{-1}\Delta\). {[}2--4; \(M\) invertible{]}
6. \(\mathbf 1^\top\Delta=\sum_s\mu_s\operatorname{Cov}_{w^s}(\mathbf 1^\top X,\beta_s)=0\), since \(\mathbf 1^\top X=1\) is constant. {[}lin; simplex{]} ∎(a)

\emph{Proof of (b).}
1. \(\mu_rNM_r=E_p[X_rT_c]=E_p[X_r^2\beta_r]+\sum_{s\ne r}E_p[X_rX_s\beta_s]\). {[}A1; lin{]}
2. \(\mu_rB_r=E_p[X_r\beta_r]=E_p[X_r^2\beta_r]+\sum_{s\ne r}E_p[X_rX_s\beta_r]\), using \(X_r=1-\sum_{s\ne r}X_s\). {[}Def; simplex{]}
3. Subtract: \(\mu_r(NM_r-B_r)=\sum_{s\ne r}E_p[X_rX_s(\beta_s-\beta_r)]=-\delta^r_B\). {[}1--2; Def 1{]} ∎(b)

\emph{Proof of (c).}
1. \(\beta^{\neg r}_i\) is a convex combination of \(\beta_{is}\in[0,1]\) with weights \(X_{is}/(1-X_{ir})\), so it lies in \([0,1]\). {[}Def; simplex{]}
2. \(T_{ic}=X_{ir}\beta_{ir}+(1-X_{ir})\beta^{\neg r}_i\). {[}A1; Def{]}
3. \(E_p[(1-X_r)\beta^{\neg r}]/(1-\mu_r)=\sum_{s\ne r}\mu_sB_s/(1-\mu_r)=B_{\neg r}\). {[}Def; lin{]}
4. So \((X_r,\beta_r,\beta^{\neg r})\) is a two-group population in EGGHM's sense, with group means \(B_r\) and \(B_{\neg r}\). {[}1--3{]}
5. Its between association is \(E_p[X_r(1-X_r)(\beta_r-\beta^{\neg r})]=E_p[X_r((1-X_r)\beta_r-\sum_{s\ne r}X_s\beta_s)]=\delta^r_B\). Its within association is \(\delta^r_W\) by definition. {[}Def; simplex{]}
6. Apply {[}EGGHM:P4{]} and {[}EGGHM:P6{]} to it. Here the individual-level variance is \(\operatorname{Var}(X)=\mu_r(1-\mu_r)\) and the between-unit variance is \(\operatorname{Var}_p(X_r)\). {[}4--5{]} ∎(c)

\emph{Proof of (d).} \(B\) and \(\delta_B\) are linear in \(\beta\) with design weights. {[}Def{]} \(\operatorname{Cov}_w(\beta_s,X_r)=\sum_iw_i\beta_{is}(X_{ir}-E_wX_r)\) is linear in \(\beta\). {[}Def Cov{]} In \(\delta^r_W\), \((1-\mu_r)w^{\neg r}_i\beta^{\neg r}_i=p_i\sum_{s\ne r}X_{is}\beta_{is}\), so the division by \(1-X_{ir}\) cancels. {[}Def{]} ∎

Checked: C1 (R = C = 3, 100 populations, max error \(6\times10^{-15}\)).

\begin{center}\rule{0.5\linewidth}{0.5pt}\end{center}

\hypertarget{pD-bounds-under-sign-restrictions-rung-n1-sharp-for-continuous-tables}{%
\subsection{2. Bounds under sign restrictions (rung N1): sharp for continuous tables}\label{pD-bounds-under-sign-restrictions-rung-n1-sharp-for-continuous-tables}}

\textbf{Definition 2.}
* A \emph{sign restriction} is \(g_k(\beta)\ge-\varepsilon_k\) with \(g_k=\pm\) one of the associations of Definition 1, for any column, and slack \(\varepsilon_k\ge0\) (default 0).
* Optional \emph{unit-level} restrictions \(\beta_i\in U_i\), with \(U_i\) polyhedral: unit-level ordering, structural zeros (MENU A1, B5).
* The \textbf{N1 feasible set} is
\[F:=\{\beta:\ \beta_i\in S(X_i,T_i)\cap U_i\ \forall i;\ g_k(\beta)\ge-\varepsilon_k\ \forall k\}.\]
* \textbf{Contextual reinforcement} for a chosen pair \((g_B,g_W)\) is \(F_{CR}:=F_+\cup F_-\), where \(F_\pm:=F\cap\{\pm g_B\ge0,\pm g_W\ge0\}\).

\textbf{Proposition M2 (N1 is sharp for continuous unit tables).} Let \(\varphi=\sum a_{rc}B_{rc}\), and read the unit tables as nonnegative \textbf{real} cell masses (§0.1).
(a) The set of values of \(\varphi\) attained by continuous unit tables consistent with the margins and the restrictions is \(\varphi(F)\). If \(F\neq\emptyset\) it equals \([\min_F\varphi,\max_F\varphi]\): \textbf{sharp for continuous unit tables, conditional on the observed margins}.
(b) \(F=\emptyset\) iff the restrictions contradict the margins in this continuous problem.
(c) Under CR the set is \(\varphi(F_+)\cup\varphi(F_-)\), a union of at most two intervals.
(d) \(\varepsilon\mapsto\varphi(F(\varepsilon))\) is nested \textbf{componentwise} (\(\varepsilon\le\varepsilon'\) coordinatewise implies \(\varphi(F(\varepsilon))\subseteq\varphi(F(\varepsilon'))\)), with continuous piecewise-linear endpoints on the region where \(F(\varepsilon)\ne\emptyset\). With several slacks, the set of \(\varepsilon\) that preserves a given conclusion (a sublevel set of a piecewise-linear function) is in general \textbf{not a product of intervals} (check C7b); its boundary is the \textbf{breakdown frontier}. \emph{(Corrected after review, round 2, Nov 2026: the set itself was called the frontier.)} A scalar breakdown value needs a specified conclusion, a rule for equality, and a scalar path \(\varepsilon=t\,v\), \(t\ge0\); along such a path it is well defined. (The package's \texttt{breakdown()} uses one scalar slack per module, i.e.~a stated path.)

\emph{(Corrected after review, Batch B, Nov 2026: (a) earlier said ``attained by unit tables \ldots{} sharp in the finite-population sense'', without the continuous reading; (b) earlier said ``contradict the data''; (d) earlier said ``the breakdown value of a conclusion is well defined'' for vector slacks. Checks C7, C7b.)}

\emph{Proof.}
1. A table \(n_{irc}=N_iX_{ir}\beta_{irc}\) of nonnegative real masses reproduces the margins iff \(\beta_i\in S(X_i,T_i)\). {[}A1; Def \(S\){]}
2. It satisfies the restrictions iff \(\beta\in F\). {[}Def 2; L M1(d){]}
3. \(F\) is a bounded polyhedron, since each \(S(X_i,T_i)\) is compact (2a L0(b)) and the restrictions are linear. \(\varphi\) is affine in \(\beta\). {[}L M1(d){]}
4. So \(\varphi(F)\) is a compact interval with both endpoints attained, which gives (a). (b) is the definition. {[}conn; 3{]}
5. (c): \(F_+\) and \(F_-\) are each of the form in (a). {[}Def 2; 4{]}
6. (d): \(\varepsilon\) enters only the right-hand side, so the sets are nested componentwise. Right-hand-side LP duality gives continuous polyhedral (piecewise-linear) endpoint functions on the region where \(F(\varepsilon)\neq\emptyset\). A conclusion such as ``upper endpoint \(\le c\)'' or ``width \(\le c\)'' holds on a sublevel set of such a function, which need not be a product set. {[}LP-par; nest{]} ∎

\textbf{Remark M2.1 (integer counts; added after review, Batch B).} Step 1 gives real cell masses, not whole numbers of people. With counted ballots, M2(a) and (b) can fail.
* \emph{Referee example (partial identification).} Three units of size 4 with group-1 counts 1, 2, 3 and positive outcomes 0, 1, 0. Each group has six people, so the integer gaps are exactly \(\{-1/6,1/6\}\). By symmetry \(D_{ER}=D_{NM}=0\) and \(\gamma=1/6\): the continuous CR set is \(\{0\}\) (half of the positive outcome in each group), while \textbf{no integer table satisfies CR}.
* \emph{Referee example (ecological inference).} Two four-voter precincts with 1 and 2 A voters and one P vote each. The integer gaps are \(-2/5\), \(2/15\), \(2/3\); CR requires \(D=0\), so integer CR is empty. Under the positive between sign the continuous lower endpoint is 0, the integer one \(2/15\): an \textbf{endpoint} moves, not only emptiness.
* \textbf{Proposition (2×2, margins only, sign rows on one column).} Let \(N_1,N_0\) be the group totals and \(N=N_1+N_0\). The integer gap set is the Duncan--Davis interval \([D^-,D^+]\) intersected with the lattice \(\{-Y/N_0+j(1/N_1+1/N_0)\}\), whose step is \(1/(N\mu(1-\mu))\); both endpoints are attained. \emph{Proof.} \(D=Y_1(1/N_1+1/N_0)-Y/N_0\), where \(Y_1\) is the number of group-1 positives and \(Y\) the total. Each unit's group-1 count ranges over an integer interval, so their sum takes every integer between its extremes, which are the DD extremes. ∎ Each sign row of M3 is a half-line in \(D\) (M3(a)), so the integer set under sign rows is the continuous set intersected with the lattice. Endpoints therefore differ by less than one step, and integer emptiness can occur only when the continuous set is shorter than one step. A numerical value needs both group totals: the step is \(1/N_1+1/N_0\); for example, \(N=10^6\) people with \(\mu=.3\) (\(N_1=3\times10^5\), \(N_0=7\times10^5\)) gives \(4.8\times10^{-6}\). At such sizes the difference matters only for tiny tables or near-ties. \emph{(Corrected after review, round 2, Nov 2026: earlier quoted a step for ``370 units of 500--4,000 people'', which does not fix the group totals.)} Checks C7, C7a.
* \textbf{Open (proof item R15).} With unit restrictions \(U_i\), joint columns, or \(R\ge3\), the integer--continuous gap is not bounded by this argument. Integer N1 would need integer cell variables, discrete target sets and separate integer feasibility checks per branch (decision U1: not built; continuous N1 is labelled as such).
* EGGHM's sharpness results range over compatible joint distributions on the continuous mean-probability model; they do not license whole-ballot interpolation either.

\textbf{Proposition M3 (two rows: closed forms, refutability, CR).} Let \(R=2\), write \(X=X_1\), and fix a column. Write \(D=B_1-B_2\), \(\gamma=\operatorname{Var}_p(X)/(\mu(1-\mu))\), and let \([D^-,D^+]\) be the MOB interval (DD for \(D\)). Assume \(\operatorname{Var}_p(X)>0\), the continuous reading (§0.1), \textbf{no unit-level restrictions} (\(U_i\) the whole simplex) and no sign rows other than those named. \emph{(Corrected after review, Batch B: the last two hypotheses were missing; the proof of (c) uses them. Check C8b.)}
(a) \(\delta_W\ge0\iff D\le D_{ER}\), and \(\delta_B\ge0\iff D\ge D_{NM}\).
(b) \(\{\delta_W\ge0\}\) gives \([D^-,\min\{D^+,D_{ER}\}]\). It is \textbf{empty iff \(D_{ER}<D^-\)}. Symmetrically, \(\{\delta_W\le0\}\) is empty iff \(D_{ER}>D^+\).
(c) \(D_{NM}\in[D^-,D^+]\) always. So in this continuous, model-free problem, sign restrictions on \(\delta_B\) are never refuted: they give \([D_{NM},D^+]\) or \([D^-,D_{NM}]\). Unit restrictions or other sign rows can remove the witness of step 4 (C8b). With integer counts a single weak between sign still has an integer witness on its side of \(D_{NM}\), since the integer extremes are the DD extremes (Remark M2.1).
(d) \(D_{NM}=\gamma D_{ER}\) with \(\gamma\in(0,1]\). Under CR:
* if \(\gamma<1\) (some unit with positive weight is mixed) and \(D_{ER}\neq0\), exactly one of \(F_\pm\) is non-empty;
* if \(\gamma=1\) (every unit with positive weight is pure), \(D\) is point-identified, \(D^-=D^+=D_{ER}=D_{NM}\), and both branches are nonempty with the same gap image \(\{D_{ER}\}\);
* if \(D_{ER}=0\), both branches are nonempty with gap image \(\{0\}\). A branch is a set of latent tables; the singleton is its \textbf{image} under the gap, and corresponds to many latent tables. \emph{(Corrected after review, round 2, Nov 2026: branches and their images were not always distinguished.)}

The identified set is
\[[D_{NM},\min\{D^+,D_{ER}\}]\ \text{if }D_{ER}\ge0,\qquad[\max\{D^-,D_{ER}\},D_{NM}]\ \text{if }D_{ER}<0.\]
In this continuous, model-free problem CR is never refuted, and under it \(\operatorname{sign}D=\operatorname{sign}D_{ER}\) (for \(D_{ER}\neq0\)). \emph{(Corrected after review, Batch B: (d) earlier said ``exactly one of \(F_\pm\) is non-empty when \(D_{ER}\neq0\)'', which fails at \(\gamma=1\); the proof noted the case and the statement dropped it. ``CR is never refuted by the data'' now names its base system. Checks C8, C7.)}

\emph{Proof.}
1. \(\delta_W=(D_{ER}-D)\operatorname{Var}_p(X)\) and \(\delta_B=(D-D_{NM})\mu(1-\mu)\). {[}L M1(c) with \(R=2\), where \(\neg1=2\){]} This gives (a). {[}\(\operatorname{Var}_p(X)>0\){]}
2. \(D\) is affine in \(\beta\) over the product of the unit polytopes. Its range is \([D^-,D^+]\), with every value attained. {[}L M1(d); M2(a); EGGHM:P3{]}
3. (b): \(F=\{\beta:D(\beta)\le D_{ER}\}\). Its image is \([D^-,D^+]\cap(-\infty,D_{ER}]\), which is empty iff \(D_{ER}<D^-\). {[}1; 2; conn{]}
4. (c): put \(\beta_{i1}=\beta_{i2}=T_i\) in every unit. This lies in \(S(X_i,T_i)\) (and in \(U_i\) only because \(U_i\) is the whole simplex), gives \(\delta_B=0\), and hence \(D=D_{NM}\). {[}Def \(S\); hypothesis; 1{]}
5. (d): \(D_{NM}=\gamma D_{ER}\) {[}EGGHM:L1{]}. \(\gamma>0\) by \(\operatorname{Var}_p(X)>0\). \(\gamma\le1\) because \(\mu(1-\mu)-\operatorname{Var}_p(X)=E_p[X(1-X)]\ge0\), with equality iff every unit with \(p_i>0\) is pure. {[}alg{]}
6. Suppose \(D_{ER}>0\). Then \(F_-\) requires \(D\ge D_{ER}\) and \(D\le D_{NM}=\gamma D_{ER}\le D_{ER}\). So \(D=D_{ER}=D_{NM}\), which forces \(\gamma=1\). Hence \(F_-=\emptyset\) when \(\gamma<1\). When \(\gamma=1\) every positive-weight unit is pure, each unit's rate for its only group equals \(T_i\), so \(D\) is point-identified and equals \(D_{ER}=D_{NM}\), and both branches are nonempty with gap image \(\{D_{ER}\}\). {[}5; 1{]}
7. \(F_+=\{D_{NM}\le D\le D_{ER}\}\cap[D^-,D^+]=[D_{NM},\min\{D^+,D_{ER}\}]\). This is non-empty because \(D^-\le D_{NM}\le D_{ER}\) and \(D_{NM}\le D^+\). {[}1; 4; 5{]}
8. Then \(D\ge D_{NM}>0\). The case \(D_{ER}<0\) is symmetric, and \(D_{ER}=0\) gives \(D_{NM}=0=D\) on both branches. {[}5--7{]} ∎

Checked: C2 (300 populations; LP = closed forms, \(\max\) error \(6\times10^{-15}\); emptiness iff \(D_{ER}<D^-\); one CR branch when \(\gamma<1\); sign agreement in all 152 populations satisfying realised CR); C8 (\(\gamma=1\): both branches); C8b (a unit ordering removes the witness); C7 (integer counts).

\textbf{Reading.} In two rows, in the \textbf{continuous, model-free} problem with no unit restrictions, only signs on the \emph{within} association can be contradicted by the margins; a between sign and CR are always feasible there. That non-refutability is specific to this base system:
* with a linear rung (LCE), a between sign and CR can be \textbf{jointly} rejected even when LCE alone fits (5a Remark R3.2);
* with integer counts, CR can be rejected outright (Remark M2.1);
* known cells, structural zeros or unit orderings can remove the equal-rates witness (C8b).

The equal-rates table (\(\beta_{ir}=T_i\) for every row, all columns, in R×C too) sets every pairwise between association to zero. It therefore satisfies every weak between sign and, for each specified pairing, the CR \textbf{union}, in the continuous problem without unit restrictions (EI referee). Which branch contains it depends on its within association; both do only if that is also zero. Example: equal weights on the compositions \((.2,.5,.3)\), \((.5,.2,.3)\), \((.8,.1,.1)\) with every group's rate equal to \(X_1\), so \(T=X_1\). For group 1 against the rest, \(ER=1\), \(\gamma=.24\), and the witness has \(D=NM=.24\), \(\delta_B=0\) and \(\delta_W=\operatorname{Var}_p(X_1)(ER-D)=.0456>0\): it lies in the positive branch only (check C15). \emph{(Corrected after review, round 2, Nov 2026: earlier ``a weak \(\delta_B\) sign and both weak CR branches hold''; the counterexample is the partial-identification referee's.)} Where a belief survives, its credibility still rests on substantive argument (or auxiliary individual data, as in EGGHM §6). An empty plug-in set, an empty population set and a rejection by a calibrated test are three different events. \emph{(Corrected after review, Batch B: earlier ``between signs and CR can never be rejected from aggregates''.)}

\textbf{Proposition M4 (joint restrictions).} Imposing sign restrictions on several columns or rows jointly gives a subset of \(F\) for any single restriction. So joint bounds are weakly tighter.
\emph{Proof.} Adding constraints shrinks the feasible set. {[}Def 2; nest{]} ∎

Checked (C3): in 2×3, signing the within association of two columns jointly narrowed the bound for column 1 strictly in 75 of 139 random populations (max width gain 0.27). This is why the package solves one LP rather than applying EGGHM's 2×2 closed forms column by column. The referee's two-column example makes the point by hand: two 2×2 units at \(x=.25,.75\) with complementary columns, so \(\delta_{W,2}=-\delta_{W,1}\). Imposing the \textbf{same} weak within sign on both columns forces both associations to zero and so \(D=ER\), shrinking \([-.5,0]\) to \(\{0\}\), while each single restriction is feasible. Opposite signs on complementary columns duplicate one restriction (check C3b). \emph{(Corrected after review, round 2, Nov 2026: earlier ``opposite within signs'', which has it backwards.)}

\begin{center}\rule{0.5\linewidth}{0.5pt}\end{center}

\hypertarget{pD-sign-restrictions-under-linear-contextual-effects-k-rows}{%
\subsection{3. Sign restrictions under linear contextual effects (K-rows)}\label{pD-sign-restrictions-under-linear-contextual-effects-k-rows}}

\textbf{Proposition M5.} Under (S2) for an LCE rung, with the conditional means written for the fixed design as
\[m_i=m_{\rm id}(V_i)+L_ik,\]
where \(m_{\rm id}\) is identified and \(k\) stacks the free coordinates of that rung: under M0 the coordinates of \(K\) (2a Prop. 7), \(m_i=M_{Q+K}(X_i)\); under M2 the district-specific coordinates \(K^{(d)}\); under M1 the common \(K\), with the district intercept shifts in the identified part (2b Lemma D1(b)); under M1(δ) possibly also a common centre that enters the association rows with zero coefficient.
(a) \textbf{Every rung.} \(\delta^*_B(r,s)\), \(\delta^{r*}_B\), \(\Delta^*\) and \(\delta^{r*}_W\) are affine in \(k\). So each sign restriction is one LP row added to \(Gk\le g\), and CR is a union of two LPs.
(b) \textbf{Ordering, per district.} If \(m_{rc}(x,d)\ge m_{sc}(x,d)\) for every district \(d\) and every \(x\) in that district's support (or \(A_d\)), then \(\delta^*_B(r,s)\ge0\), pooled over districts.
(c) \textbf{Slopes and the within sign, \(R=2\).}
* \emph{M0 (common lines).} \(\delta^*_W=a(w_1+d_1)+b\,w_1\) with \(a=\mu\operatorname{Var}_{w^1}(X)\ge0\) and \(b=(1-\mu)\operatorname{Var}_{w^2}(X)\ge0\). Hence both contextual slopes \(\ge0\) ⇒ \(\delta^*_W\ge0\).
* \emph{District rungs.} By the law of total covariance under each row weight,
\[\delta^*_W=\sum_{s}\mu_s\Big(\underbrace{E_{w^s}\big[\operatorname{Cov}_{w^s}(X,m_s\mid D)\big]}_{\text{within districts}}+\underbrace{\operatorname{Cov}_{w^s}\big(E_{w^s}[X\mid D],E_{w^s}[m_s\mid D]\big)}_{\text{between districts}}\Big).\]
Nonnegative district slopes make the within-district term \(\ge0\) (as for M0, district by district). The between-district term has no sign: under M1 it contains \(\sum_s\mu_s\operatorname{Cov}_{w^s}(X,\zeta_{s,D})\), the covariance of the district intercept shifts with composition. So the shortcut ``positive slopes ⇒ \(\delta^*_W\ge0\)'' holds for M0, and within each district (district-conditional weights), but \textbf{not} for the pooled association under M1/M1(δ)/M2.
* \emph{Example (referee).} Six equally weighted units, \(X=.1,.2,.3\) in district 1 and \(.7,.8,.9\) in district 2; both rows' mean \(.8+.1x\) in district 1 and \(.2+.1x\) in district 2 (the other column the complement). This satisfies M1 with common strictly positive slopes and valid rates. Within each district \(\delta^*_W=+0.00064\), but the pooled \(\delta^*_W=-0.04927\): the district intercept falls as \(X\) rises. (Check C9; the plan's check I2 also shows that the 2×2 local rows hold everywhere here, so the chain ``signed slopes ⊂ weighted-local ⊂ \(\delta_W\ge0\)'' also fails for pooled M1.)

\emph{(Corrected after review, Batch B, Nov 2026. M5 was stated for every LCE rung with the single formula \(m=M_{Q+K}(X)\), which proves only the pooled case. (a) now uses the stacked parameterisation of each rung; (b) needs district-specific ordering; (c) was false for the district rungs, where the proof treated the district intercept as constant. By decision U2(a), the pooled statements are the proved ones; sharpness and inference for M-SIGN under district rungs (R3′, R10′) are proof item \textbf{R16}, not proved. For confidence intervals this means: pooled M0 is covered by P3.5 and R10; a district rung \textbf{without} sign rows is covered by P3's fixed-district results in regime R1 only; a district rung \textbf{with} sign rows has no coverage theorem yet. }(Scope sentence added after review, round 2, EI referee.)* The W1 slab identity \(\delta^*_W=\operatorname{Var}_p(X)(ER^*-D^*)\) still holds with the district-fitted \(\tau\), so the CR slab rows remain valid rows under every rung.)*

\emph{Proof.}
1. For each rung, \(m_i=m_{\rm id}(V_i)+L_ik\) with \(L_i\) fixed by the design. {[}2a L1; 2b Lemma D1; Def of the rung{]}
2. Each association is linear in \((m_i)_i\) with design coefficients. {[}L M1(d) with \(m\) for \(\beta\){]} Composing with 1, it is affine in \(k\). (a) follows. {[}Def 2{]}
3. (b): \(\delta^*_B(r,s)=\sum_ip_iX_{ir}X_{is}\big(m_r(X_i,D_i)-m_s(X_i,D_i)\big)\), and every term is \(\ge0\) by hypothesis, since \(p_iX_{ir}X_{is}\ge0\). {[}Def; hypothesis{]}
4. (c), M0: \(\operatorname{Cov}_{w^1}(b_0+b_1X,X)=b_1\operatorname{Var}_{w^1}(X)\), and likewise for row 2. Then \(b_1=w_1+d_1\) (P1 §1.3 dictionary). {[}Def Cov; alg{]}
5. (c), district rungs: apply the law of total covariance under \(w^s\), conditioning on \(D\). Within district \(d\), \(m_s(\cdot,d)\) is linear with slope \(b_{s,d}\), so the conditional covariance is \(b_{s,d}\operatorname{Var}_{w^s}(X\mid D=d)\ge0\) when \(b_{s,d}\ge0\). {[}total covariance; 4{]} The example is computed in C9. ∎

Checked: C4 (affinity in \(k\) for R = C = 3; ordering ⇒ \(\delta^*_B\ge0\), with valid transition matrices; \emph{rebuilt after review, Batch B: the old construction changed one column without rebalancing the others}), C9 (district reversal and the decomposition), and 1b V3 (closed form in 2×2).

\textbf{Remark M5.1 (which δ).}
* N1 uses the realised associations of the actual population (owner decision 1).
* K-rows use \(\delta^*\), the behavioural version.
* Under (S2), \(\delta-\delta^*=\sum_i\sum_rc_{ir}e_{ir}\) is a design-weighted sum of unit residuals, with mean zero given the design. Under P3-S1a (independence given the design) its standard deviation is at most \(\tfrac12\{\sum_i(\sum_r\lvert c_{ir}\rvert)^2\}^{1/2}\), the root sum of squared design coefficients; it is small only when no unit keeps a non-negligible coefficient. (S2) alone, a mean condition, gives no rate. \emph{(Corrected after review, Batch B: earlier ``order \(p^{-1/2}\)'', which needs these sampling and weight conditions.)}
* Intersecting N1 with an LCE rung therefore mixes the two readings. The package will report them separately.

\begin{center}\rule{0.5\linewidth}{0.5pt}\end{center}

\hypertarget{pD-monotone-nonparametric-contextual-effects-rung-n2-two-rows}{%
\subsection{4. Monotone nonparametric contextual effects (rung N2; two rows)}\label{pD-monotone-nonparametric-contextual-effects-rung-n2-two-rows}}

\textbf{Support convention} \emph{(added after review, Batch B)}. N2 is defined on the \textbf{finite empirical design law}: \(\operatorname{supp}P_X=\{x_1<\dots<x_{J_X}\}\) is the set of distinct observed compositions, each with positive mass. \(J_X\) denotes the support size; polynomial degrees in §6 are written \(q\). On a continuous population support the observed compositions are only a grid, and the results below do not carry over (N2.2(b)).

\textbf{Model N2(σ).}
* (S1) holds, and \(E[\beta\mid V]=m(X)\) for a Borel \(m:[0,1]\to\mathfrak B\) (A3). Here \(X\) is the row-1 share.
* For designated pairs \((r,c)\) with directions \(\sigma_{rc}\in\{+1,-1\}\), \(\sigma_{rc}m_{rc}\) is nondecreasing on \(\operatorname{supp}P_X\).
* No linearity is assumed. \(\tau(x):=E[T\mid X=x]\) is identified at each support point.

\textbf{Lemma N2.1.} \(m(x)\in S(x,\tau(x))\) for \(P_X\)-a.e. \(x\), i.e.~at every support point.
\emph{Proof.}
1. \(m(X)=E[\beta\mid X]\in\mathfrak B\) a.s. {[}CE-tower with A3; CE-cvx{]}
2. \(m(X)^\top X=E[\beta^\top X\mid X]=E[T\mid X]=\tau(X)\) a.s. {[}CE-pull; CE-lin; A1{]} ∎

\textbf{Outer set and LP.} With the support points \(x_1<\dots<x_{J_X}\),
\[\mathcal M_O:=\{(m(x_j))_j:\ m(x_j)\in S(x_j,\tau(x_j))\ \forall j;\ \sigma_{rc}\big(m_{rc}(x_{j+1})-m_{rc}(x_j)\big)\ge0\ \forall j,(r,c)\}.\]
Bounds on \(\tilde B_{rc}=\sum_iw^r_im_{rc}(X_i)\) are the LP \([\min,\max]_{\mathcal M_O}\), which has \(2CJ_X\) variables.

\textbf{Proposition N2.2.} Under the support convention:
(a) Every \(m\) allowed by N2(σ) restricts to a point of \(\mathcal M_O\). So the LP bounds are an \textbf{outer} set for \(\tilde B\). They are not claimed sharp. Replacing \(S(x,\tau(x))\) by \(\mathbb E[S(X,T)\mid X=x]\) is \textbf{not} the sharp version when size is informative, because the model conditions on \(V=(X,N)\) (Remark N2.2.1).
(b) Conversely, every point of \(\mathcal M_O\) extends to a monotone Borel \(m\) on \([0,1]\) satisfying Lemma N2.1 on the support: a common piecewise-constant extension of the whole matrix preserves every row sum and every designated direction. So, \textbf{on the complete finite support}, restricting to support points loses nothing. On a strict subset of the support (a coarser grid) this fails: extendibility subject to the \textbf{other} accounting equations can fail, and the grid LP is only a relaxation (Remark N2.2.2).
(c) \textbf{Nesting with LCE.} If \(\tau\) is the quadratic of 2a Lemma 1, then \(\{M_{Q+K}:K\in\Theta_O\cap\{\sigma\text{-signed slopes}\}\}\subseteq\mathcal M_O\). Hence \(\mathcal B(\text{M0}\cap\text{signed slopes})\subseteq\mathcal B(\text{N2})\), for both outer sets.
(d) \textbf{Nesting with N1} (superpopulation reading). If \(m_{1c}\) and \(m_{2c}\) are both nondecreasing, then \(\delta^*_W\ge0\) for column \(c\). The two-row shape restriction is stronger than a pooled covariance sign: a sign on \(\delta^*_W\) permits reversals.
(e) For \(R\ge3\) a direction ``in \(x_1\)'' needs a \textbf{simplex chart}: which share adjusts when \(x_1\) changes. With the chart \((x_1,x_3)\) (so \(x_2\) adjusts), \(m_{rc}(x)=x_3\) is constant, hence weakly nondecreasing, in \(x_1\); yet under a symmetric Dirichlet design \(E[X_3\mid X_1]=(1-X_1)/2\), and tilting by \(X_1\) or \(1-X_1\) preserves that conditional mean, so the covariance with \(X_1\) is negative under \textbf{both} W1 weights. So (d) has no general R×C analogue. (A globally increasing \emph{univariate} function of \(X_1\) still obeys the covariance inequality.)

\emph{Proof.}
1. (a): N2(σ) ⇒ Lemma N2.1 at each \(x_j\), and monotonicity holds on the support. {[}L N2.1; Def{]} The target depends on \(m\) only at the support points. {[}Def \(\tilde B\){]}
2. (b): extend the whole matrix \(m\) as a step function, constant between support points. {[}ext{]} The extension is Borel, each row stays in the simplex, each designated entry stays monotone, and its values at the support points are unchanged, so Lemma N2.1 still holds there. {[}Def{]}
3. (c): for \(K\) in the outer set, \(M_{Q+K}(x)\ge0\) on \(A\supseteq\operatorname{supp}\) {[}2a L3(a){]} and its rows sum to 1 {[}2a L1(b){]}, so it lies in \(\mathfrak B\). Also \(x^\top M_{Q+K}(x)=x^\top Qx=\tau(x)\) {[}2a L1(a){]}, so \(M_{Q+K}(x)\in S(x,\tau(x))\) at each atom. A linear function with a slope of sign σ is σ-monotone. {[}Def{]}
4. (d): \(\delta^*_W=\mu\operatorname{Cov}_{w^1}(m_1(X),X)+(1-\mu)\operatorname{Cov}_{w^2}(m_2(X),X)\), with both covariances \(\ge0\). {[}Def 1 with \(m\); Cheb applied under \(w^1\) and under \(w^2\){]}
5. (e): the example. {[}C5(c), C5(c′){]} ∎

\emph{(Corrected after review, Batch B, Nov 2026: ``support points (the distinct observed compositions)'' was ambiguous between the finite empirical support and a grid on a larger support, and (b) is false on a grid; (a)'s proposed sharp version was not sharp; (e) lacked a chart and the weighted covariances. Checks C10, C10b, C5(c′).)}

\textbf{Remark N2.2.1 (size conditioning; Conjecture).} Fix \(X=1/2\) and give two sizes equal probability. At the first size \(T\) is 0 or 1 equiprobably, at the second \(T=1/2\) surely. At the first size both row rates must equal \(T\), so the model forces \(m(1/2)=(.5,.5)\). Conditioning only on \(X\) (and the \(S(x,\tau)\) LP) admits \((.75,.25)\), by selecting \((1,0)\) at the second size (check C10b). \textbf{Conjecture N2.2′.} The sharp N2 set is \(\{m:\ m(X)\in\mathbb E[S(X,T)\mid V]\ P_V\text{-a.s., plus the shape restrictions}\}\), i.e.~every conditional support-function inequality holds for the same \(m(x)\) at almost every size given \(x\). The proof route is 2a Theorem 1's selection construction applied to each admissible measurable \(m\); it is not written. Conditioning on \(X\) alone suffices if size carries no information relevant to the selection sets, or if the model conditions only on \(X\). The same caution applies to M-SHARP (PD step 1b): its series bound conditions on \(X\) only, so it is ``tighter than \(\Theta_O\)'', not sharp, when size is informative.

\emph{Remark (sufficiency under the selection-expectation reading; due to the partial-identification referee, round 2).} If \(\mathbb E[S(X,T)\mid V]\) denotes the selection (Aumann) expectation of the random set, the characterisation of Conjecture N2.2′ is \textbf{sufficient} on the finite support: for each positive-mass \((x,n)\) take a measurable selection of \(S(x,T)\) whose conditional mean is the common candidate \(m(x)\), and glue the selections across the finitely many design atoms; accounting holds a.s., and the shape restrictions already hold for the common \(m\). This does not settle the conjecture as stated (we have not written out necessity in general), and it does not imply that a procedure using only \(E[T\mid X]\) attains the set; the example above shows it does not when size is informative. Conjecture N2.2′ stays open.

\textbf{Remark N2.2.2 (grids).} True support \(\{1/4,1/2,3/4\}\), row-1 means \((.4,.4,.4)\), row-2 means \((8/15,.7,.8)\), both weakly increasing, \(\tau=(.5,.55,.5)\). On a grid keeping only the extremes, assigning both rows \(.5\) at both extremes passes the LP. Any monotone extension then assigns \(.5\) to both rows at the middle and cannot give \(.55\) (check C10). For binned or gridded data use Lemma N2.4.

Checked:
* C5(a): N2 contains the truth under a monotone nonlinear DGP.
* C5(b): the nesting (c) held in 47 random cases.
* C5(c), C5(c′): the R = 3 example, with negative covariance unweighted and under both W1 weights.
* C10, C10b: the grid and size examples.
* 1b V5: Chebyshev.

\textbf{Proposition N2.3 (structural local identity; EGGHM Props 12--13 in this setting).} Let \(m_{1c},m_{2c}\) be differentiable at \(x\) and write \(\tau_{m,c}(z)=z\,m_{1c}(z)+(1-z)m_{2c}(z)\) for the \textbf{structural} aggregate. Write \(D_c(x)=m_{1c}(x)-m_{2c}(x)\) and \(\delta_{W,c}(x)=x\,m_{1c}'(x)+(1-x)m_{2c}'(x)\). Then:
(a) \(\tau_{m,c}'(x)=D_c(x)+\delta_{W,c}(x)\).
(b) \(\delta_{W,c}(x)\ge0\iff D_c(x)\le\tau'_{m,c}(x)\).
(c) Local CR, \(D_c(x)\delta_{W,c}(x)\ge0\), holds iff \(D_c(x)\) lies between 0 and \(\tau'_{m,c}(x)\) (including a zero derivative).
(d) If \(m_{1c}\) and \(m_{2c}\) are nondecreasing \textbf{on a neighbourhood of \(x\)}, the local sign at \(x\) is \(\ge0\). Monotonicity on the finite support (N2) is not enough: both rows equal to \(.5+h\) with \(h=(x-\tfrac14)(x-\tfrac12)(x-\tfrac34)\) are constant on \(\{1/4,1/2,3/4\}\), yet have derivative \(-1/16\) at \(1/2\) (check C16). \emph{(Corrected after review, round 2, Nov 2026: earlier ``Both rows nondecreasing ⇒ local sign \(\ge0\)'', which read with N2's support-only monotonicity does not follow.)} The local sign everywhere does \textbf{not} imply \(\delta^*_W\ge0\), even for smooth \(m\): with \(X\) uniform on \([1/5,2/5]\), \(m_1=x-x^2/2\) and \(m_0=1-x^2/2\), both are probabilities, \(xm_1'+(1-x)m_0'=0\) everywhere, and \(\delta^*_W=-1/47250\) (\(-1/5250\) for the finite design with equal mass at the two endpoints). The reverse non-implication also holds: complement both functions, shrink them slightly into the interior, and subtract a small positive multiple of \(x\) from each; the aggregate within association stays positive by continuity and the local one becomes negative.

\textbf{Not an observable restriction under the support convention.} On a finite support \(\tau'\) is \textbf{not identified}: adding \(h(x)=\prod_j(x-x_j)\) to both rows leaves the law on the support unchanged and moves the derivative at every node. Example: support \(\{1/4,1/2,3/4\}\), \(m_1=.7\), \(m_0=.3\), \(h=(x-\tfrac14)(x-\tfrac12)(x-\tfrac34)\); both models are valid probabilities on \([0,1]\) and induce the same law, but \(\tau'(1/2)\) is \(.4\) in one and \(.3375\) in the other. The first model satisfies weak local CR with gap \(.4\); an LP using the other compatible derivative caps the gap at \(.3375\) and excludes the truth (check C11). An observable version needs an \textbf{identified} derivative: an interior point of an interval support with compatible smooth versions, or a maintained identified regression family, and an estimated derivative adds sampling error. At a pure composition EGGHM's local between association is \(x(1-x)D(x)=0\); replacing its sign by the sign of \(D(x)\) changes the belief.

\textbf{Decision U4(a):} the N2-local rows are \textbf{dropped} from the build; N2.3 is kept as the structural calculus identity above. \emph{(Corrected after review, Batch B, Nov 2026: N2.3 was stated as an observable restriction ``linear in \(m(x)\) given \(\tau'\)'', with an estimated \(\tau'\) at the support points; and (d)'s witness was the kinked 1b V6. Checks C11, C13; C6 checks the identity.)}

\emph{Proof.}
1. Differentiate \(\tau_{m,c}(z)=z\,m_{1c}(z)+(1-z)m_{2c}(z)\). {[}product rule{]} This gives (a), and (b) follows. {[}alg{]}
2. (c), case \(\tau'\ge0\): if \(D\ge0\) and \(\delta_W\ge0\), then \(0\le D=\tau'-\delta_W\le\tau'\). If \(D\le0\) and \(\delta_W\le0\), then \(D=\tau'-\delta_W\ge\tau'\ge0\ge D\), so \(D=0\). Conversely, \(D\in[0,\tau']\) gives \(D\ge0\) and \(\delta_W=\tau'-D\ge0\). The case \(\tau'\le0\) is symmetric. {[}a; alg{]}
3. (d): \(x,1-x\ge0\); the example is computed exactly in C13. {[}Def{]} ∎

\textbf{Lemma N2.4 (person-weighted bins; added after review, Batch B, following the referee's repair).} Let \(b=1,\dots,K\) be ordered, nonoverlapping bins of compositions, \(P(b)>0\) their person-weighted masses and \(E_b\) the person-weighted (\(p\)-weighted) expectation within bin \(b\). Put
\[\bar x_b=E_bX,\qquad\bar m_{1b}=E_b[Xm_1(X)]/\bar x_b,\qquad\bar m_{2b}=E_b[(1-X)m_2(X)]/(1-\bar x_b),\qquad\bar\tau_b=E_b[\tau(X)],\]
the \textbf{model-implied} bin mean. It equals the population person-weighted bin mean of \(T\); in a sample, the observed person-weighted bin mean of \(T\) estimates it with residual error, which the inference must account for. \emph{(Corrected after review, round 2, Nov 2026: \(\bar\tau_b\) was defined as \(E_bT\) without this distinction.)}
with a compatible zero-weight extension for a group absent from a bin (\(\bar x_b\in\{0,1\}\): any value that keeps the bin rates monotone). Under N2(σ):
(i) \(\bar m_b\in S(\bar x_b,\bar\tau_b)\);
(ii) \(\tilde B_{1c}=\sum_bP(b)\bar x_b\bar m_{1bc}/\mu\), and likewise for row 2: the row targets are linear in the bin rates with known weights;
(iii) if \(\sigma_{rc}m_{rc}\) is nondecreasing, so is \(\sigma_{rc}\bar m_{rbc}\) across the ordered bins.
Hence the binned LP \(\{\bar m_b\in S(\bar x_b,\bar\tau_b),\ \text{adjacent order}\}\) with the targets of (ii) is an \textbf{outer relaxation} of the global N2 bounds. It does not cover local derivative rows.
\emph{Proof.} (i) \(\bar x_b\bar m_{1b}+(1-\bar x_b)\bar m_{2b}=E_b[Xm_1+(1-X)m_2]=E_b\tau=\bar\tau_b\) {[}L N2.1{]}, each \(\bar m_{rb}\) is a weighted average of simplex points, so lies in the simplex. (ii) \(E_p[Xm_1]=\sum_bP(b)E_b[Xm_1]\). {[}lin{]} (iii) \(\bar m_{1b}\) is an average of \(m_1\) over bin \(b\) under the tilted weights \(\propto pX\), so lies in \([\min_bm_1,\max_bm_1]\); ordered nonoverlapping bins and monotone \(m_1\) give \(\max_bm_1\le\min_{b+1}m_1\). Row 2 likewise with weights \(\propto p(1-X)\). ∎
The ordinary bin average does \textbf{not} work: with equal bin weights at \(x=.25,.75\), \(m_1=x\), \(m_2=0\), \(E_bT=E[X^2]=.3125\), while the vector of ordinary bin-average rates evaluated at \(E[X]=.5\) gives \(.25\); the person-weighted rate is \(.625\) and reproduces \(.3125\) (check C12). The earlier argument (``averaging Lemma N2.1 within a bin uses {[}CE-cvx{]}'') applied conditional convexity to the varying sets \(S(x,\tau(x))\), which does not prove the step.

\begin{center}\rule{0.5\linewidth}{0.5pt}\end{center}

\hypertarget{pD-the-ladder-2-rows-with-status}{%
\subsection{5. The ladder (2 rows), with status}\label{pD-the-ladder-2-rows-with-status}}

All N1 statements are for continuous tables (Remark M2.1).
\[DD\ \supseteq\ \text{N1}(\delta_B\text{ sign})\ ,\ \text{N1(CR)}\ ,\ \text{N1}(\delta_W\text{ sign})\ \supseteq\ \text{N1(both signs, same direction)}\qquad\text{[M2; nest]}\]
``Both signs'' means both weakly positive or both weakly negative; then N1(both signs) ⊆ N1(CR). Two opposite signs generally do not imply CR.
\[\mathcal B(\text{M0}\cap\text{signed slopes})\ \subseteq\ \mathcal B(\text{N2})\quad(\text{outer, }\tau\text{ quadratic, finite support})\qquad\text{[N2.2(c)]}\]
\[\text{N2(both rows }\uparrow)\ \Rightarrow\ \delta^*_W\ge0\qquad\text{[N2.2(d)]};\qquad \text{mean-level order in each district}\ \Rightarrow\ \delta^*_B\ge0\qquad\text{[M5(b)]};\qquad\text{M0 signed slopes}\ \Rightarrow\ \delta^*_W\ge0\qquad\text{[M5(c), M0 only]}.\]
* N2 admits step functions, where a local derivative is undefined; the structural local identity N2.3 applies only to a differentiable subclass and is not an observable rung (U4(a)).
* Compare a realised association (\(\delta\), N1) with a model-implied one (\(\delta^*\)) only after fixing the same population of tables or means (Remark M5.1).
* For \(R\ge3\), any monotonicity statement needs a simplex chart (N2.2(e)).

\textbf{Not nested:}
* M0 with N1 or N2;
* the local sign everywhere and the aggregate \(\delta^*_W\) sign, in either direction (N2.3(d));
* for \(R\ge3\), N2 with N1(\(\delta_W\)) (N2.2(e));
* pooled \(\delta^*_W\) and signed slopes under the district rungs (M5(c)).

\emph{(Corrected after review, Batch B, Nov 2026: the line ``N2(both rows ↑) ⊆ N2-local'' is removed with the N2-local rung; ``both signs'' now requires the same direction; the M5 cells are scoped to M0.)}

\begin{center}\rule{0.5\linewidth}{0.5pt}\end{center}

\hypertarget{pD-what-phase-3-must-supply}{%
\subsection{6. What Phase 3 must supply}\label{pD-what-phase-3-must-supply}}

\begin{itemize}
\tightlist
\item
  \textbf{N1} needs no sampling inference (owner decision 1). Phase 3 adds the slack dials ε and the breakdown frontier, with a scalar breakdown value along a stated path (M2(d)). Labels: ``sharp for continuous unit tables given the margins''.
\item
  \textbf{K-rows:} the existing P3 machinery, with one extra row each (pooled M0; grouped rungs: R16).
\item
  \textbf{N2 (global only; N2-local dropped, U4(a))} needs \(\tau\) at the support points.

  \begin{itemize}
  \tightlist
  \item
    A weighted polynomial of \textbf{prespecified} degree \(q\) (separate notation from the support size \(J_X\)). Degree \(q\) has \(q+1\) coefficients and is identified only from at least \(q+1\) distinct support points; \(q=J_X\) is never identified. With \(q=2\), N2.2(c) holds exactly.
  \item
    The constraints are then linear in the coefficients, and P3.5 applies \emph{conditional on the degree} ((a) and (b2) proved; CLR as implemented algorithmic, (b3)): valid if \(\tau\) is a polynomial of degree \(q\). Otherwise add an approximation-bias allowance, or label the bound as conditional on the specification. A prespecified degree removes the selection problem below, not the specification problem.
  \item
    \textbf{Do not select the degree by cross-validation and then use fixed-degree inference.} In a balanced two-level design with true slope zero, choosing between intercept-only and intercept-plus-slope by leave-one-out error selects the slope when \(Z^2>2\), and the nominal 95\% slope interval then covers with conditional probability \(1-.05/P(\lvert Z\rvert>\sqrt2)\approx0.68\) (check C12b). Use a prespecified degree, sample splitting, or a selection-aware procedure, and state any approximation bias.
  \item
    With many support points, group compositions into bins and use the \textbf{person-weighted} bin rates of Lemma N2.4; \(\bar\tau_b\) is the model-implied bin mean, estimated by the person-weighted bin mean of \(T\) (with residual error to account for). This is an outer relaxation of global N2.
  \item
    \emph{(Corrected after review, Batch B, Nov 2026: earlier ``a weighted polynomial of degree \(J\), selected by cross-validation'' (with \(J\) also the support size) and ``binning keeps the outer property: averaging Lemma N2.1 within a bin uses {[}CE-cvx{]}''. Checks C12, C12b.)}
  \end{itemize}
\item
  \textbf{Implementation sizes.} N1 has \(p\,R\,C\) variables. That is about 370×6 for Wisconsin in 2×3, which \texttt{lpSolve} handles; the equalities will be eliminated to \(p(R-1)(C-1)\) free variables.
\end{itemize}


\clearpage
\section{PHASE3: estimation, inference and sensitivity}\label{app:phase3}
\textit{Source: \texttt{docs/extension/PHASE3\_inference.md}.}\par
% GENERATED by tools/make_appendices.py from docs/extension/PHASE3_inference.md
% Do not edit by hand; edit the source and regenerate.
\emph{Builds on P1, 2a and 2b. Proof steps carry justification tags (2a §0.3, 2b §0.3, plus §0.3 below). Theorems marked ``(cited)'' apply a published result whose conditions are verified here but whose proof is not repeated. Checks: \path{verify_phase3.py} (17/17 PASS; \path{verify_phase3_output.txt}). Monte Carlo: \path{phase3_montecarlo.py}, 1000 replications per design, with Monte Carlo standard errors (\path{phase3_montecarlo_output.txt}; cross-check against the package in \path{phase3_montecarlo_crosscheck.py}).}

\begin{center}\rule{0.5\linewidth}{0.5pt}\end{center}

\hypertarget{pE-setting}{%
\subsection{0. Setting}\label{pE-setting}}

\hypertarget{pE-sampling}{%
\subsubsection{0.1 Sampling}\label{pE-sampling}}

\begin{itemize}
\tightlist
\item
  There are \(n\) units with observables \(W_i=(N_i,X_i,T_i,D_i)\).
\item
  \textbf{(P3-S1a)} Given the design and the district effects, the unit rates \(\beta_i\) are independent across units, with \(E[\beta_i\mid\text{design}]=m(V_i)\). Equivalently, the residuals \(e_i=\beta_i-m(V_i)\) are independent with mean zero. No dependence is allowed within a district beyond the district effect, and no clustered variance is used. §0.4 explains why. \emph{(Revised after review: the earlier P3-S1 allowed arbitrary within-district dependence for M0 while Theorem P3.5's residual allowance assumed independence.)}
\item
  \textbf{(P3-S2)} The design is conditioned on: inference is conditional on \(\{(N_i,X_i,D_i)\}\), so the vertices of the sample hulls \(\hat A_d\) are fixed.
\item
  \textbf{Regimes} (P1 A10): (R1) \(D^\star\) fixed, \(n_d\to\infty\); (R2) \(D^\star\to\infty\), \(n_d\) bounded; (R3) both grow. Which result holds where is in §0.5.
\end{itemize}

\hypertarget{pE-targets}{%
\subsubsection{0.2 Targets}\label{pE-targets}}

There are three related objects (P1 §1.7):

\begin{itemize}
\tightlist
\item
  \textbf{Superpopulation rate} \(B^*_{rc}=E[NX_r\beta_{rc}]/E[NX_r]\), the object of Phases 2a--2b.
\item
  \textbf{Finite-population rate} (JKST's target): \(B_{rc}=\sum_iN_iX_{ir}\beta_{irc}\big/\sum_iN_iX_{ir}\), the realised rate for this electorate.
\item
  \textbf{Finite-design conditional-mean rate:} \(\tilde B_{rc}=\sum_iw_{ir}\,m_{rc}(V_i)\), with \(w_{ir}=N_iX_{ir}/\sum_jN_jX_{jr}\) and \(m=E[\beta\mid V]\). This is the conditional expectation of \(B\) given the observed design (under P3-S1). It is a sample-weighted analogue of \(B^*\), not equal to it in general. Under suitable design-sampling conditions it estimates \(B^*\), but inference for \(B^*\) would also need the variation of the design, which is not pursued here.
\end{itemize}

All inference below is conditional on the design and covers \(B\) or \(\tilde B\).

They differ by the residual term
\[E_{rc}=B_{rc}-\tilde B_{rc}=\sum_iw_{ir}e_{irc},\qquad e=\beta-m.\]
\(\tilde B\) is the right target when the question is about behaviour rather than the particular electorate. Only \(\tilde B\) is affine in \(K\), with the \emph{sample} weights \(w\), so \(\mu\) is never estimated. \(B\) is \(\tilde B\) plus the residual term \(E\), which does not depend on \(K\) and is handled separately (§5). \emph{(Corrected after review: an earlier version said both are linear in \(K\).)}

\hypertarget{pE-additional-tags}{%
\subsubsection{0.3 Additional tags}\label{pE-additional-tags}}

\begin{longtable}[]{@{}
  >{\raggedright\arraybackslash\hspace{0pt}}p{(\columnwidth - 2\tabcolsep) * \real{0.5000}}
  >{\raggedright\arraybackslash\hspace{0pt}}p{(\columnwidth - 2\tabcolsep) * \real{0.5000}}@{}}
\toprule\noalign{}
\begin{minipage}[b]{\linewidth}\raggedright\hspace{0pt}Tag
\end{minipage} & \begin{minipage}[b]{\linewidth}\raggedright\hspace{0pt}Meaning
\end{minipage} \\
\midrule\noalign{}
\endhead
\bottomrule\noalign{}
\endlastfoot
{[}LPdual{]} & LP strong duality: a feasible LP with bounded value has an optimal dual solution with equal value; the optimum of a pointed polyhedron is attained at a vertex; a polyhedron has finitely many vertices \\
{[}Slater-b{]} & under strict feasibility, the optimal dual vectors are bounded: \(\lVert\lambda\rVert_1\le(\text{value}-\text{value at the strictly feasible point})/\varepsilon\) \\
{[}Hausdorff{]} & Hausdorff/Kuratowski convergence of compact convex sets implies convergence of support functions \\
{[}dDM{]} & delta method for Hadamard directionally differentiable maps (Shapiro 1991; Fang \& Santos 2019, Thm 2.1†) \\
{[}minLin{]} & the directional derivative of \(\beta\mapsto\min_{j\in J}g_j^\top\beta\) (\(J\) finite) at \(\beta_0\), in direction \(h\), is \(\min_{j\in J^*(\beta_0)}g_j^\top h\), with \(J^*\) the active set \\
{[}CLR{]} & Chernozhukov, Lee \& Rosen (2013, \emph{Econometrica} 81:667--737): parametric bounding functions (Condition P, Lemma 4, Thm 4, via Thms 1--2); finite index set as in their Example 1. \emph{Corrected after review: earlier cited as Thm 2; see 2a §0.4.} \\
{[}Pop{]} & Popoviciu: \(\operatorname{Var}(Y)\le(b-a)^2/4\) for \(Y\in[a,b]\); {[}BD{]} Bhatia--Davis: \(\operatorname{Var}(Y)\le(M-\mu)(\mu-m)\) \\
{[}Lyap{]} & Lindeberg--Feller CLT for independent summands \(Y_i\) with \(\lvert Y_i\rvert\le b_i\), applied relative to the \textbf{actual} standard deviation \(s_n\) of the sum: it holds when \(\max_ib_i/s_n\to0\). \emph{(Corrected after review, Batch B, Nov 2026: earlier stated as a CLT under \(\max_iw_i^2/\sum_iw_i^2\to0\) alone, which does not give a variance-standardised CLT when, e.g., only one summand has positive variance. P3.5(a) step 5 now combines it with Chebyshev; see there.)} \\
{[}union{]} & \(P(A\cup B)\le P(A)+P(B)\) \\
{[}OLS{]} & consistency and asymptotic normality of OLS for a fixed-dimensional parameter, with heteroskedasticity-robust (HC) variance, independent units, a full-rank design and bounded outcomes \\
{[}BR{]} & the multivariate CLT holds uniformly over all convex sets of \(\mathbb R^p\) (Bhattacharya \& Rao 1976, \emph{Normal Approximation and Asymptotic Expansions}, Thm 17.5, standard) \\
{[}Hoeff{]} & Hoeffding's inequality: for independent \(Y_i\in[a_i,b_i]\), \(P\big(\sum_i(Y_i-EY_i)\ge t\big)\le\exp\{-2t^2/\sum_i(b_i-a_i)^2\}\) (Hoeffding 1963, \emph{JASA} 58:13--30, Thm 2) \\
{[}CCK{]} & Gaussian approximation for maxima of high-dimensional sums (Chernozhukov, Chetverikov \& Kato 2017, \emph{Ann. Probab.} 45:2309--2352) \\
\end{longtable}

\hypertarget{pE-where-the-randomness-comes-from-added-after-review}{%
\subsubsection{0.4 Where the randomness comes from (added after review)}\label{pE-where-the-randomness-comes-from-added-after-review}}

Abadie, Athey, Imbens and Wooldridge (2020, \emph{Econometrica}; 2023, \emph{QJE}) distinguish two sources of uncertainty:

\begin{itemize}
\tightlist
\item
  \textbf{Sampling-based:} the units are a sample from a larger population.
\item
  \textbf{Design-based:} some assignment is random.
\end{itemize}

They show that clustered standard errors are justified by clustered sampling or clustered assignment, not by the mere fact that residuals are correlated. EI has neither. Every precinct and riding of the electorate is observed, and nothing is assigned. In a purely design-based view, nothing is random at all: \(B\) is a fixed unknown, and the only valid bound is Duncan--Davis.

The contextual-effects assumption is a statement about a conditional expectation, \(E[\beta\mid V]=M_H(X)\). That only has content if the precinct rates are treated as random given the precinct's characteristics. The uncertainty here is therefore \textbf{model-based, conditional on the design}: the \(\beta_i\) are random given \(V_i\), and the design is held fixed (P3-S2). Abadie, Imbens and Zheng (2014) give the parallel fixed-regressor framework for regression.

Within this framework:

\begin{itemize}
\tightlist
\item
  \textbf{District shocks are fixed features of this electorate.} The targets \(B\) and \(\tilde B\) refer to this electorate, so a shock common to a riding (candidate, campaign, local swing) is not re-drawn.

  \begin{itemize}
  \tightlist
  \item
    Under M1/M2 it is a district effect, part of the mean \(m(V_i)\).
  \item
    Under M0 a district shock is not a variance problem but a misspecification. Conditional on the design, M0's mean model holds for the observed units only if the district shifts are zero (Prop. P3.6(c)). Clustered standard errors would not repair it.
  \end{itemize}
\item
  \textbf{What remains assumed} (P3-S1a): no dependence within a district beyond its district effect. Neighbourhood spillovers between adjacent precincts would violate it. This is a stated limitation. A partial check is residual correlation between neighbouring precincts after fitting M1.
\item
  \textbf{Independence is a substantive assumption} \emph{(added after review, round 2, Nov 2026)}. Two arguments must be kept apart.

  \begin{enumerate}
  \def\labelenumi{\arabic{enumi}.}
  \tightlist
  \item
    Conditioning on observed district effects and \textbf{assuming} that the remaining residuals are independent (P3-S1a) is a coherent model.
  \item
    ``Every unit is observed, so dependence need not be considered'' does \textbf{not} follow. Observing every unit removes sampling uncertainty about which units exist; it says nothing about dependence among their residuals.
  \end{enumerate}

  Within-district shocks can make \(\operatorname{Cov}(e_i,e_j\mid\text{design})\ne0\) with no district mean shift. That violates P3-S1a without violating the mean model, so the pooling test P3.6(c), which tests mean shifts, is \textbf{not} a test of within-district dependence. P3-S1a must therefore be argued for each application; residual correlation within districts and between neighbours is a diagnostic, not a test that can confirm it. We do not reintroduce clustered standard errors, for the reasons above. Under dependence, the variance and residual-allowance statements of P3.5 and P3.8 are not covered.
\item
  \textbf{A superpopulation target that averages over re-drawn district shocks} (``behaviour in a typical riding'') would need cluster methods and many districts. It is not pursued.
\end{itemize}

\hypertarget{pE-which-result-holds-in-which-regime}{%
\subsubsection{0.5 Which result holds in which regime}\label{pE-which-result-holds-in-which-regime}}

\begin{longtable}[]{@{}
  >{\raggedright\arraybackslash\hspace{0pt}}p{(\columnwidth - 6\tabcolsep) * \real{0.2500}}
  >{\raggedright\arraybackslash\hspace{0pt}}p{(\columnwidth - 6\tabcolsep) * \real{0.2500}}
  >{\raggedright\arraybackslash\hspace{0pt}}p{(\columnwidth - 6\tabcolsep) * \real{0.2500}}
  >{\raggedright\arraybackslash\hspace{0pt}}p{(\columnwidth - 6\tabcolsep) * \real{0.2500}}@{}}
\toprule\noalign{}
\begin{minipage}[b]{\linewidth}\raggedright\hspace{0pt}Model
\end{minipage} & \begin{minipage}[b]{\linewidth}\raggedright\hspace{0pt}Regime
\end{minipage} & \begin{minipage}[b]{\linewidth}\raggedright\hspace{0pt}Formal result
\end{minipage} & \begin{minipage}[b]{\linewidth}\raggedright\hspace{0pt}Status
\end{minipage} \\
\midrule\noalign{}
\endhead
\bottomrule\noalign{}
\endlastfoot
any rung & any (R1, R2, R3), any \(n\) & Theorem P3.8: finite-sample CI; finite-sample emptiness test & Theorem; currently uninformative in width (§5.1) \\
M0 & any (\(\dim\beta\) fixed) & P3.2, P3.4, P3.5(a), (b1), (b2) & Theorem (asymptotic, \(n\to\infty\)) \\
M1, M2 & R1 (\(D^\star\) fixed) & P3.2, P3.4, P3.5(a), (b1), (b2) & Theorem (asymptotic) \\
M1, M2 & R3 & the same, with a high-dimensional Gaussian approximation {[}CCK{]} under a rate condition on \(D^\star\) versus \(n_d\) & Conjecture \\
M1 & R2 & coverage of the CLR interval & Simulation finding (§8.2: 0.995--1 with up to 200 districts of 10 polls) \\
any & any & the CLR interval as implemented (candidates by simulation) & Conjecture P3.5(b3) \\
\end{longtable}

Under M0 the number of constraints grows with the number of hull vertices, but the regression has fixed dimension \(p\). Every constraint term is a linear functional of one \(p\)-dimensional asymptotically normal vector. The relevant events, \(\{\max_j(\cdot)\le k\}\), are convex sets in \(\mathbb R^p\), and the multivariate CLT holds uniformly over convex sets {[}BR{]}. So the growth in constraints does not matter. Under M1/M2 in R2 or R3, \(\dim\beta\) grows with \(D^\star\), and that argument no longer applies.

\begin{center}\rule{0.5\linewidth}{0.5pt}\end{center}

\hypertarget{pE-estimators}{%
\subsection{1. Estimators}\label{pE-estimators}}

\begin{enumerate}
\def\labelenumi{\arabic{enumi}.}
\tightlist
\item
  \textbf{Regression.} Fit each column \(T_c\) by OLS on the homogeneous quadratic monomials \(X_rX_s\) (\(r\le s\)) for M0. For M1, also include the district shifts \(1\{D=d\}X_r\), \(d\ge2\); this imposes the overidentifying restrictions of 2b Lemma D1(b). For M2, fit one regression per district. The coefficients are \(\hat\beta\), and the estimated columns are coherent automatically (2a Remark 1.2). Use an HC3 covariance \(\hat\Sigma\) (P3-S1a; no clustering, §0.4). Only \(C-1\) columns are needed; column \(C\) follows by adding up (Remark P3.4.1).
\item
  \textbf{Constraint terms.} For each district \(d\), vertex \(v\in\operatorname{vert}(\hat A_d)\), row \(r\) and column \(c\), the constraint is \(s_j(\beta)+a_j^\top k\ge0\), with \(s_j\) linear in \(\beta\) and \(a_j\) fixed (2a Prop. 7, 2b Thm 2(b)).
\item
  \textbf{Target.} \(\tilde B(k)=c(\beta)+\gamma^\top k\), with \(c\) linear in \(\beta\) and \(\gamma\) computed from the sample weights (2a Lemma 2, 2b Lemma D2).
\item
  \textbf{Plug-in bound.} \([\hat\varphi_L,\hat\varphi_U]=c(\hat\beta)+[\min,\max]\{\gamma^\top k:s(\hat\beta)+Ak\ge0\}\).
\end{enumerate}

\begin{center}\rule{0.5\linewidth}{0.5pt}\end{center}

\hypertarget{pE-structure-of-the-bound}{%
\subsection{2. Structure of the bound}\label{pE-structure-of-the-bound}}

\phantomsection\label{pE-res-lemma-P3.1}\textbf{Lemma P3.1 (the endpoint is a minimum of finitely many linear functions).} Suppose the LP is feasible with bounded value at \(\beta\). Then
\[\varphi_U(\beta)=c(\beta)+\min_{\lambda\in\Lambda}\lambda^\top s(\beta),\qquad\Lambda=\operatorname{vert}\{\lambda\ge0:A^\top\lambda=-\gamma\},\]
a finite set. \(\varphi_U\) is therefore concave and piecewise linear in \(\beta\), and likewise \(\varphi_L(\beta)=c(\beta)-\min_{\lambda\in\Lambda'}\lambda^\top s(\beta)\), with \(\Lambda'=\operatorname{vert}\{\lambda\ge0:A^\top\lambda=\gamma\}\).

\emph{Proof.}

\begin{enumerate}
\def\labelenumi{\arabic{enumi}.}
\tightlist
\item
  \(\max\{\gamma^\top k:s+Ak\ge0\}\) has Lagrangian \(\gamma^\top k+\lambda^\top(s+Ak)\). Stationarity in the free \(k\) requires \(\gamma+A^\top\lambda=0\), and the dual is \(\min\{\lambda^\top s:\lambda\ge0,\ A^\top\lambda=-\gamma\}\). {[}LPdual{]}
\item
  Primal feasibility and boundedness give equal optimal values. {[}LPdual; 2a:L3(b){]}
\item
  The dual feasible set lies in \(\{\lambda\ge0\}\), so it contains no line and is pointed. A linear objective bounded below on it attains its minimum at a vertex, and there are finitely many vertices. {[}LPdual{]}
\item
  The dual feasible set does not depend on \(\beta\) (\(A\) and \(\gamma\) are fixed given the design). So \(\varphi_U(\beta)-c(\beta)=\min_{\lambda\in\Lambda}\lambda^\top s(\beta)\) for every \(\beta\) at which the primal is feasible and bounded. {[}3; Def{]}
\item
  A minimum of finitely many linear functions is concave and piecewise linear. {[}alg{]} The lower endpoint follows by applying 1--5 to \(-\gamma\). ∎
\end{enumerate}

Checked: concavity and piecewise linearity along random directions (\path{verify_phase3.py}).

\emph{Scope.} The lemma is about variation in the regression coefficients \(\beta\) for a \textbf{fixed} constraint matrix \(A\) and objective \(\gamma\). Conditioning on the design fixes both, because the target uses the sample weights. It is the first case of 2a Prop. 7's regularity paragraph. It says nothing about joint variation in the identified coefficients, the population moments and the constraint set, where the value function is in general neither concave nor piecewise linear.

\textbf{Reading.} JKST's Proposition 4 is this lemma with \(|\Lambda|=4\) (their \(g_{lj},g_{uj}\)). Its ``tie-breaking'' assumption says the minimising \(\lambda\) is unique. In R×C, and especially with many districts, \(\Lambda\) is large and near-ties are the norm. That is why the inference below treats the endpoint as an intersection bound.

\begin{center}\rule{0.5\linewidth}{0.5pt}\end{center}

\hypertarget{pE-consistency-and-rate}{%
\subsection{3. Consistency and rate}\label{pE-consistency-and-rate}}

\phantomsection\label{pE-res-proposition-P3.2}\textbf{Proposition P3.2.} Assume (P3-S1--S2), {[}A1{]}--{[}A5{]} for the rung used, \(\sqrt n(\hat\beta-\beta_0)=O_p(1)\), and \textbf{(Slater)}: there are \(\tilde k\) and \(\varepsilon>0\) with \(s_j(\beta_0)+a_j^\top\tilde k\ge\varepsilon\) for all constraints defined by the population hulls. Then:

\begin{enumerate}
\def\labelenumi{(\alph{enumi})}
\tightlist
\item
  \(\hat\varphi_U-\varphi_U(\beta_0;\hat A)=O_p(n^{-1/2})\), and likewise for \(\varphi_L\);
\item
  \(\varphi_U(\beta_0;\hat A)\to\varphi_U(\beta_0;A)\) when \(\hat A\to A\) in Hausdorff distance, either along a nested sequence (\(\hat A_n\subseteq\hat A_{n+1}\), as when observations are added), or for any Hausdorff-convergent sequence whose feasible sets \(\Theta(\hat A_n)\) eventually lie in a common bounded set;
\item
  since \(\hat A_d\subseteq A_d\), \(\varphi_U(\beta_0;\hat A)\ge\varphi_U(\beta_0;A)\) and \(\varphi_L(\beta_0;\hat A)\le\varphi_L(\beta_0;A)\): \textbf{estimating the hull only ever widens the bound.}
\end{enumerate}

\emph{Proof.}

\begin{enumerate}
\def\labelenumi{\arabic{enumi}.}
\tightlist
\item
  (c): a vertex of \(\hat A_d\) is an observed composition, hence in \(A_d\). So \(\hat A_d\subseteq A_d\), the constraints over \(\hat A_d\) are a subset of those over \(A_d\), and the feasible set is weakly larger. {[}Def; 2a:L3(a){]}
\item
  (a): by Lemma P3.1, \(\varphi_U(\beta)-\varphi_U(\beta')=c(\beta-\beta')+\big[\min_\lambda\lambda^\top s(\beta)-\min_\lambda\lambda^\top s(\beta')\big]\). {[}P3.1{]}
\item
  The bracket is bounded by \(\max\lVert S^\top\lambda\rVert\,\lVert\beta-\beta'\rVert\), the max taken over optimal duals at \(\beta\) or \(\beta'\). {[}P3.1; \(s\) linear, \(s(\beta)=S\beta\){]}
\item
  For \(\beta\) in a neighbourhood of \(\beta_0\), \(\tilde k\) remains \(\varepsilon/2\)-strictly feasible, so the optimal duals are uniformly bounded. {[}Slater; continuity of \(s\); Slater-b{]}
\item
  Hence \(\lvert\hat\varphi_U-\varphi_U(\beta_0;\hat A)\rvert\le L\lVert\hat\beta-\beta_0\rVert=O_p(n^{-1/2})\), with \(L\) uniform in the hull because the constraint set over \(\hat A\) is a subset of the one over \(A\). {[}2--4; 1{]}
\item
  (b): the sets \(\Theta(\hat A)\) are compact and convex, decrease as \(\hat A\) grows, and contain \(\Theta(A)\). {[}1; 2a:L3(b){]}
\item
  If \(k\in\Theta(\hat A_p)\) for all \(p\) and \(v\in A\), take \(v_p\in\hat A_p\) with \(v_p\to v\). The constraints are continuous in \(v\), so \(k\) satisfies the constraint at \(v\). Hence \(\bigcap_p\Theta(\hat A_p)=\Theta(A)\). {[}Def; continuity{]}
\item
  A decreasing sequence of compact convex sets converges in Hausdorff distance to its intersection, so the support functions, and hence the LP values, converge. {[}Hausdorff{]} ∎(nested)
\item
  Non-nested case: \(\Theta(A)\subseteq\Theta(\hat A_n)\) for every \(n\) by 1. If \(k_n\in\Theta(\hat A_n)\) and \(k_{n_j}\to k\), then for each \(v\in A\) pick \(v_j\in\hat A_{n_j}\) with \(v_j\to v\). Continuity of the constraints in \((k,v)\) gives the constraint at \(v\), so \(k\in\Theta(A)\). The upper Kuratowski limit is therefore \(\Theta(A)\), which is also contained in every term. In a common bounded set this is Hausdorff convergence. {[}1; Def; continuity; Hausdorff{]} ∎(b)
\end{enumerate}

\textbf{Remarks.}

\begin{itemize}
\tightlist
\item
  \textbf{Error decomposition.} \(\hat\varphi_U-\varphi_U(\beta_0;A)=\underbrace{[\hat\varphi_U-\varphi_U(\beta_0;\hat A)]}_{O_p(n^{-1/2})\text{ by (a)}}+\underbrace{[\varphi_U(\beta_0;\hat A)-\varphi_U(\beta_0;A)]}_{\ge0,\ \to0\text{ by (b), rate not shown}}\). The second term is not shown to be \(O_p(n^{-1/2})\), and can be larger. So Theorem P3.4's limit concerns the \textbf{sample-hull} endpoint \(\varphi_U(\beta_0;\hat A)\), and not the population-hull endpoint. Theorem P3.4 itself is for a \textbf{fixed} constraint system; as the sample hull changes with \(n\), the limit is given by Proposition P3.4′ under its localisation condition. \emph{(Corrected after review, round 2, Nov 2026.)}
\item
  \textbf{Slater is a strict-interiority condition.} It requires a \(k\) at which every implied rate at every hull vertex is at least \(\varepsilon>0\). It fails when every feasible \(k\) makes some implied rate exactly 0 at some vertex, e.g.~a true structural zero, which is not exotic in EI. Design E of §8 is an example: the identified set is a single point. (Design C has a rate close to 0 at one vertex: Slater holds, but only just.) The Lipschitz bound in (a) is proved only under Slater; boundary cases are open (§10). The coverage results of §5 do \textbf{not} use Slater: P3.5(a) step 2 uses only that the true \(k_0\) is feasible.
\item
  The hull effect in (c) is a one-sided (conservative) bias. Its size depends on how much mass \(P_X\) puts near the boundary of \(A_d\); it can vanish slowly. Report the bound's sensitivity to trimming the hull (P1 A4).
\item
  The \(O_p(n^{-1/2})\) rate is for the estimation part. The target \(B\) adds \(E=O_p(n^{-1/2})\) as well (§5). Bounds do not shrink to a point: their width converges to the population width, which is \(O(1)\).
\end{itemize}

\begin{center}\rule{0.5\linewidth}{0.5pt}\end{center}

\hypertarget{pE-asymptotic-distribution}{%
\subsection{4. Asymptotic distribution}\label{pE-asymptotic-distribution}}

\phantomsection\label{pE-res-theorem-P3.4}\textbf{Theorem P3.4 (cited structure: {[}dDM{]}; fixed constraint system).} Under Proposition P3.2's assumptions and \(\sqrt n(\hat\beta-\beta_0)\Rightarrow\mathcal Z\sim N(0,\Sigma)\), suppose the \textbf{constraint system is fixed} along the sequence: the vertex set entering the LP does not depend on \(n\), so \(A\), \(\gamma\), \(\Lambda\) and the gradients \(g_\lambda\) are fixed. This holds for a prespecified polytope \(A\) with fixed vertices, or along a design sequence whose sample hull is eventually constant, \(\hat A_n=A\) for \(n\ge n_0\) (for instance, when \(A\) is a polytope whose vertices are atoms of the design law). Then
\[\sqrt n\big(\hat\varphi_U-\varphi_U(\beta_0)\big)\Rightarrow\min_{\lambda\in\Lambda^*}\big(c+S^\top\lambda\big)^\top\mathcal Z,\]
where \(\Lambda^*\subseteq\Lambda\) is the set of optimal dual vertices at \(\beta_0\).

\begin{itemize}
\item
  \begin{enumerate}
  \def\labelenumi{(\roman{enumi})}
  \tightlist
  \item
    If \(\Lambda^*\) is a singleton (JKST's tie-breaking), the limit is \(N(0,g^{*\top}\Sigma g^*)\) with \(g^*=c+S^\top\lambda^*\).
  \end{enumerate}
\item
  \begin{enumerate}
  \def\labelenumi{(\roman{enumi})}
  \setcounter{enumi}{1}
  \tightlist
  \item
    If two active gradients differ in a direction in which the estimator varies, i.e.~\((g_\lambda-g_{\lambda'})^\top\Sigma(g_\lambda-g_{\lambda'})>0\) for some \(\lambda,\lambda'\in\Lambda^*\), the limit is a minimum of correlated normals that are not a.s. equal, and it has strictly negative mean. The plug-in upper endpoint is then biased \emph{downward}, and the lower endpoint upward, both by order \(n^{-1/2}\). That is, the plug-in bound is too narrow. If instead all active gradients agree on the support of \(\mathcal Z\), all the \(g_\lambda^\top\mathcal Z\) are a.s. equal and the limit is normal, as in (i), even though \(\Lambda^*\) is not a singleton. \emph{(Corrected after review: an earlier version claimed a negative mean whenever \(\Lambda^*\) is not a singleton.)}
  \end{enumerate}
\item
  \begin{enumerate}
  \def\labelenumi{(\roman{enumi})}
  \setcounter{enumi}{2}
  \tightlist
  \item
    In case (ii) the standard bootstrap is inconsistent (Fang \& Santos 2019, Cor. 3.1†).
  \end{enumerate}
\end{itemize}

\emph{(Corrected after review, round 2, Nov 2026. The theorem was stated ``with the design and the sample hull \(\hat A\) fixed''. The proof is correct for a fixed set of affine maps, but as \(n\) grows the sample hull, the dual vertices and the active set can change, and gaps between nearly binding candidates can shrink at rate \(n^{-1/2}\). The referee's example: \(f_n(\theta)=\min\{\theta,-\theta+n^{-1/2}\}\), \(\theta_0=0\), \(\sqrt n(\hat\theta-\theta_0)\Rightarrow Z\). The first piece is uniquely active for every \(n\), yet \(\sqrt n(f_n(\hat\theta)-f_n(\theta_0))\Rightarrow\min\{Z,1-Z\}\), which is not normal. The theorem is now stated for a fixed constraint system; changing systems are Proposition P3.4′. Neither affects the coverage results of §5, which do not use P3.4.)}

\emph{Proof.}

\begin{enumerate}
\def\labelenumi{\arabic{enumi}.}
\tightlist
\item
  \(\varphi_U(\beta)=\min_{\lambda\in\Lambda}g_\lambda^\top\beta\) with \(g_\lambda=c+S^\top\lambda\): a finite minimum of linear maps. {[}P3.1{]}
\item
  Such a map is Lipschitz, and its directional derivative is \(h\mapsto\min_{\lambda\in\Lambda^*}g_\lambda^\top h\). A Lipschitz map with directional derivatives in finite dimensions is Hadamard directionally differentiable. {[}minLin{]}
\item
  Apply the directional delta method. {[}dDM; 1--2{]}
\item
  (i): the minimum over a singleton is linear, giving a normal limit. {[}3{]}
\item
  (ii): let \(Y=g_\lambda^\top\mathcal Z\) and \(Y'=g_{\lambda'}^\top\mathcal Z\) with \(\operatorname{Var}(Y-Y')=(g_\lambda-g_{\lambda'})^\top\Sigma(g_\lambda-g_{\lambda'})>0\). Then \(E\min_{\Lambda^*}g^\top\mathcal Z\le E\min(Y,Y')=\tfrac12E(Y+Y')-\tfrac12E\lvert Y-Y'\rvert=-\tfrac12E\lvert Y-Y'\rvert<0\). If \(\operatorname{Var}(g_\lambda^\top\mathcal Z-g_{\lambda'}^\top\mathcal Z)=0\) for all active pairs, the minimum equals each \(g_\lambda^\top\mathcal Z\) a.s., which is normal. {[}3; \(\min(a,b)=\frac{a+b}2-\frac{\lvert a-b\rvert}2\); \(Y-Y'\) is a nondegenerate normal{]}
\item
  (iii): in case (ii) the derivative is nonlinear on the support of \(\mathcal Z\), so the condition for bootstrap consistency fails. {[}Fang \& Santos 2019†{]} ∎
\end{enumerate}

\phantomsection\label{pE-res-remark-P3.4.1}\textbf{Remark P3.4.1 (parametrise \(C-1\) columns).} The shares satisfy \(\sum_cT_c=1\), so the coefficient vectors of the \(C\) column regressions sum to a known constant, and the joint covariance of all \(C\) columns is singular. State \(\hat\beta\), \(\Sigma\) and the gradients \(g_\lambda\) for the first \(C-1\) columns and recover column \(C\) by adding up. Then \(\Sigma\) can be nonsingular. When it is positive definite, the condition in (ii) reduces to \(g_\lambda\neq g_{\lambda'}\) in the reduced parametrisation.

\phantomsection\label{pE-res-proposition-P3.4p}\textbf{Proposition P3.4′ (triangular arrays; added after review, round 2).} Let the design, and hence the constraint system, change with \(n\). Write \(\Lambda_n\) for the dual vertex set of the design-\(n\) LP, with \(\lvert\Lambda_n\rvert\le\bar L\) for all \(n\), gradients \(g_{\lambda,n}=c_n+S_n^\top\lambda\) (bounded), and \textbf{localised slacks}
\[c_{\lambda,n}:=\sqrt n\big(\theta_{\lambda,n}(\beta_0)-\varphi_{U,n}(\beta_0)\big)\ \ge0,\qquad\theta_{\lambda,n}(\beta)=g_{\lambda,n}^\top\beta .\]
\textbf{Localisation condition.} After relabelling the candidates, \(g_{\lambda,n}\to g_\lambda\) and \(c_{\lambda,n}\to c_\lambda\in[0,\infty]\) for each \(\lambda\). Then
\[\sqrt n\big(\hat\varphi_{U,n}-\varphi_{U,n}(\beta_0)\big)\Rightarrow\min_{\lambda:\ c_\lambda<\infty}\big(g_\lambda^\top\mathcal Z+c_\lambda\big).\]
(Without the condition, the same holds along every subsequence on which it holds; the limit can then differ between subsequences.)

\begin{itemize}
\tightlist
\item
  \textbf{Separation condition.} The limit of Theorem P3.4 (a minimum of \(g_\lambda^\top\mathcal Z\) over the active set) obtains when every \(c_\lambda\in\{0,\infty\}\): each candidate is either exactly active at \(\beta_0\) or inactive by more than \(n^{-1/2}\), \(\sqrt n\min_{\lambda\notin\Lambda^*_n}\{\theta_{\lambda,n}(\beta_0)-\varphi_{U,n}(\beta_0)\}\to\infty\), with a stable active set and convergent gradients.
\item
  A candidate inactive by a gap of exact order \(n^{-1/2}\) (\(0<c_\lambda<\infty\)) stays in the limit. The referee's \(f_n\) has \(c=1\) for the second piece and gives \(\min\{Z,1-Z\}\).
\item
  This is a limit for the sample-hull endpoint along the design sequence. It says nothing about the population-hull error (remark after P3.2).
\end{itemize}

\emph{Proof.}

\begin{enumerate}
\def\labelenumi{\arabic{enumi}.}
\tightlist
\item
  \(\varphi_{U,n}(\beta)=\min_{\lambda\in\Lambda_n}\theta_{\lambda,n}(\beta)\) for every \(\beta\) at which the LP is feasible and bounded. {[}P3.1, for the design-\(n\) system{]}
\item
  Hence, exactly, \(\sqrt n(\hat\varphi_{U,n}-\varphi_{U,n}(\beta_0))=\min_\lambda\big(g_{\lambda,n}^\top\sqrt n(\hat\beta-\beta_0)+c_{\lambda,n}\big)\). {[}1; Def \(c_{\lambda,n}\); linearity{]}
\item
  \(\sqrt n(\hat\beta-\beta_0)=O_p(1)\) and the gradients are bounded, so a candidate with \(c_{\lambda,n}\to\infty\) exceeds the active candidates' terms with probability \(\to1\) and drops out of the minimum. {[}OLS; 2{]}
\item
  The remaining finitely many terms converge jointly in distribution to \((g_\lambda^\top\mathcal Z+c_\lambda)_\lambda\), and the minimum is continuous. {[}2--3; continuous mapping{]} ∎
\end{enumerate}

Checked (\path{verify_phase3.py}, ``P3.4′ {[}F{]}''): the referee's \(f_n\) at \(n=10^4\); the limit is \(\min\{Z,1-Z\}\) (mean \(\approx-0.40\), Kolmogorov distance 0.31 from \(N(0,1)\)), not the normal limit the old reading gives.

Checked: with two tied constraints and unit noise, the mean plug-in error is \(-0.564\approx-1/\sqrt\pi\); with a unique active constraint it is 0.007. With two active gradients that differ only in a null direction of \(\Sigma\), the mean is 0 (normal limit), as the corrected (ii) states.

\textbf{Boundary cases.}

\begin{itemize}
\tightlist
\item
  \textbf{Truth at a vertex.} When the true \(K\) lies on the boundary of the identified set, the side of the interval that touches it carries the full error probability, so coverage is close to nominal. When the true \(K\) is interior, coverage tends to 1 (JKST App. A Remark). This is the Imbens--Manski distinction between covering the parameter and covering the set.
\item
  \textbf{Width near zero.} When the bound's width is near zero (point identification: L4 in 2×2, L5, the DiD with unchanged contextual effects), one-sided intervals per endpoint under-cover unless the critical value adapts (Stoye 2009). The two-sided Bonferroni split used below is uniformly valid, at some cost in width.
\end{itemize}

\begin{center}\rule{0.5\linewidth}{0.5pt}\end{center}

\hypertarget{pE-confidence-intervals}{%
\subsection{5. Confidence intervals}\label{pE-confidence-intervals}}

\phantomsection\label{pE-res-theorem-P3.5}\textbf{Theorem P3.5 (validity).} Fix \(\alpha_U,\alpha_L,\alpha_E\in(0,1)\). Write \(\theta_\lambda(\beta)=g_\lambda^\top\beta\) for the candidate upper endpoints (\(\lambda\in\Lambda\)), with standard errors \(\hat\sigma_\lambda\) from \(\hat\Sigma\).

\textbf{(a) Conservative version.} Let \(\hat k\) be the \((1-\alpha_U-\alpha_L)\)-quantile of \(\max_j\) over all constraint terms and the constant of the standardised Gaussian errors (one-sided for the constraints, two-sided for the constant). Put
\[\hat\Theta^+=\{k:s(\hat\beta)+\hat k\hat\sigma_s+Ak\ge0\},\]
\[\resizebox{\linewidth}{!}{$\displaystyle \mathrm{CI}=\Big[c(\hat\beta)-\hat k\hat\sigma_c+\min_{\hat\Theta^+}\gamma^\top k-z_{\alpha_E}S_1,\ \ c(\hat\beta)+\hat k\hat\sigma_c+\max_{\hat\Theta^+}\gamma^\top k+z_{\alpha_E}S_1\Big],$}\]
with \(S_1=\tfrac12\big(\sum_iw_{ir}^2\big)^{1/2}\) and \(z_{\alpha_E}=\Phi^{-1}(1-\alpha_E/2)\). Then
\[\liminf_nP(B_{rc}\in\mathrm{CI})\ge1-\alpha_U-\alpha_L-\alpha_E.\]

\textbf{(b1) CLR with the full candidate set (cited: {[}CLR{]}, Thm 4).} Let \(\Lambda\) be the full set of dual vertices of Lemma P3.1, which is finite for a fixed design. Assume CLR's Condition P: \(\Sigma\) is positive definite (use \(C-1\) columns, Remark P3.4.1, and a full-rank design), and \(g_\lambda\ne0\) for every \(\lambda\in\Lambda\). For the upper end:

\begin{enumerate}
\def\labelenumi{\arabic{enumi}.}
\tightlist
\item
  Select \(\hat V=\{\lambda:\hat\theta_\lambda\le\min_\mu(\hat\theta_\mu+k_p\hat\sigma_\mu)+2k_p\hat\sigma_\lambda\}\), where \(k_p\) is the \(\gamma_n=1-0.1/\log n\) quantile of \(\max_\lambda Z_\lambda\), and \(Z\) is the standardised Gaussian vector with correlation from \(\hat\Sigma\).
\item
  Set \(\hat U=\min_{\lambda\in\hat V}(\hat\theta_\lambda+\hat k\hat\sigma_\lambda)\), with \(\hat k\) the \((1-\alpha_U)\)-quantile of \(\max_{\lambda\in\hat V}Z_\lambda\).
\item
  Do the same for the lower end, and add \(z_{\alpha_E}S_1\) on each side (each side's E-allowance at level \(\alpha_E/2\)).
\end{enumerate}

The same coverage bound holds. With \(\alpha_E=0\) and no \(S_1\) term, either version covers \(\tilde B_{rc}\) with probability at least \(1-\alpha_U-\alpha_L\).

\textbf{(b2) Any candidate set, with a simultaneous critical value (proved here).} Let \(\hat\Lambda\) be any data-dependent set of dual-feasible points (vertices or not; for example collected by simulation), and \(\hat U=\min_{\lambda\in\hat\Lambda}(\hat\theta_\lambda+\hat k\hat\sigma_\lambda)\), with \(\hat U=+\infty\) if \(\hat\Lambda=\emptyset\). Then \(\liminf_nP(\varphi_U(\beta_0)\le\hat U)\ge1-\alpha_U\) in either of two cases:

\begin{itemize}
\item
  \begin{enumerate}
  \def\labelenumi{(\roman{enumi})}
  \tightlist
  \item
    \(\hat k\) is the \((1-\alpha_U)\)-quantile of \(\max_{\lambda\in\bar\Lambda}Z_\lambda\) for a fixed finite set \(\bar\Lambda\supseteq\hat\Lambda\) a.s. (e.g.~all of \(\Lambda\)), with \(\sigma_\lambda>0\) on \(\bar\Lambda\);
  \end{enumerate}
\item
  \begin{enumerate}
  \def\labelenumi{(\roman{enumi})}
  \setcounter{enumi}{1}
  \tightlist
  \item
    projection: \(\hat k=\sqrt{\chi^2_{q,1-\alpha}}\) and \(\hat\sigma_\lambda=(g_\lambda^\top\hat\Sigma g_\lambda/n)^{1/2}\), in either of two settings: (ii-a) \(\Sigma\) is positive definite and \(q=\dim\beta\) (use \(C-1\) columns and a full-rank design, Remark P3.4.1); or (ii-b) a \textbf{reduced full-rank parameterisation}: there are a fixed matrix \(L\) (\(\dim\beta\times q\)) and a vector \(b\) with \(\beta=L\eta+b\) and \(\hat\beta=L\hat\eta+b\) \textbf{exactly} (so every error in a null direction of \(\Sigma\) is identically zero, e.g.~coordinates fixed by known identities), \(\sqrt n(\hat\eta-\eta_0)\Rightarrow N(0,\Sigma_\eta)\) with \(\Sigma_\eta\) positive definite, and \(\hat\Sigma=L\hat\Sigma_\eta L^\top\) with \(\hat\Sigma_\eta\) consistent. This one event covers both ends, any \(\hat\Lambda\), at total level \(\alpha\).
  \end{enumerate}

  \emph{(Corrected after review, Batch B, Nov 2026. The earlier statement took \(q=\operatorname{rank}\Sigma\) for any \(\Sigma\), with a pseudoinverse in the proof. That bounds \(g^\top(\hat\beta-\beta_0)\) only for \(g\) in the range of \(\Sigma\). A singular limit does not control errors of smaller order in its null directions: with \(\hat\beta-\beta_0=(Z/\sqrt n,\,-1/n)\), \(Z\sim N(0,1)\), the limit covariance is \(\operatorname{diag}(1,0)\), the second coordinate gets standard error 0, and the upper bound \(\hat\beta_2+0\) misses \(\beta_{0,2}\) in every sample. Counterexample due to the referee; \path{verify_phase3.py}, ``P3.5(b2)(ii) null direction''.)}
\end{itemize}

\emph{Proof.}

\begin{enumerate}
\def\labelenumi{\arabic{enumi}.}
\tightlist
\item
  Under the model the primal LP at \(\beta_0\) is feasible (\(k_0\) is feasible). For any dual-feasible \(\lambda\) and any primal-feasible \(k\): \(\gamma^\top k=-\lambda^\top Ak\le\lambda^\top s(\beta_0)\). So \(\theta_\lambda(\beta_0)\ge\varphi_U(\beta_0)\) for \textbf{every} dual-feasible \(\lambda\), not only optimal ones. {[}weak duality; \(\lambda\ge0\), \(s+Ak\ge0\){]}
\item
  On \(\mathcal E=\{\theta_\lambda(\beta_0)\le\hat\theta_\lambda+\hat k\hat\sigma_\lambda\ \forall\lambda\in\bar\Lambda\}\): \(\hat U\ge\min_{\lambda\in\hat\Lambda}\theta_\lambda(\beta_0)\ge\varphi_U(\beta_0)\). {[}1; \(\hat\Lambda\subseteq\bar\Lambda\){]}
\item
  (i): \((\theta_\lambda(\beta_0)-\hat\theta_\lambda)/\hat\sigma_\lambda\), \(\lambda\in\bar\Lambda\), is a fixed finite collection of standardised linear functionals of \(\hat\beta\), jointly asymptotically \(N(0,\text{corr})\). So \(P(\mathcal E)\to1-\alpha_U\). {[}OLS; \(\hat\Sigma\) consistent; Def \(\hat k\){]}
\item
  (ii): case (ii-b) contains (ii-a) (take \(L=I\), \(b=0\), \(\eta=\beta\)). Write \(\hat e=\hat\eta-\eta_0\). For every \(g\), \(g^\top(\hat\beta-\beta_0)=(L^\top g)^\top\hat e\) exactly, so by Cauchy--Schwarz in the metric of \(\hat\Sigma_\eta\) (positive definite with probability \(\to1\)),
  \[\resizebox{\linewidth}{!}{$\displaystyle \lvert g^\top(\hat\beta-\beta_0)\rvert\le\big(\hat e^\top\hat\Sigma_\eta^{-1}\hat e\big)^{1/2}\big(g^\top L\hat\Sigma_\eta L^\top g\big)^{1/2}=\big(n\,\hat e^\top\hat\Sigma_\eta^{-1}\hat e\big)^{1/2}\hat\sigma_g .$}\]
  This holds for every \(g\) at once, including \(g\) with \(L^\top g=0\), for which both sides are 0. By continuous mapping, \(n\,\hat e^\top\hat\Sigma_\eta^{-1}\hat e\Rightarrow\chi^2_q\), since \(\Sigma_\eta\) is nonsingular. So \(\mathcal E\) holds for all \(\lambda\) at once, and for both ends, with probability \(\to1-\alpha\). Nonsingularity of \(\Sigma_\eta\) is what the earlier rank-\(q\) pseudoinverse argument lacked: there, a null-direction error of order \(o(n^{-1/2})\) is neither zero nor controlled. {[}alg; OLS; continuous mapping{]} ∎
\end{enumerate}

\emph{Remark (exact versus numerical reduction; added after review, round 2).} Case (ii-b) is claimed only for a documented \textbf{exact} reduction: the adding-up identity that removes column \(C\) (Remark P3.4.1), or coordinates fixed by the design. When the rank of \(\hat\Sigma\) comes from a numerical cutoff, the theorem applies only with \(q=\dim\beta\) (case (ii-a), assuming the true \(\Sigma\) is positive definite). Dropping candidates whose gradients load on a numerically null direction of \(\hat\Sigma\) is then a \textbf{safeguard}: it can only widen the interval (the minimum is over fewer terms), but it neither establishes the identities of (ii-b) nor calibrates a rank-\(q\) critical value. \emph{(An earlier response described that drop rule as case (ii-b); it is not.)}

\begin{enumerate}
\def\labelenumi{(\roman{enumi})}
\setcounter{enumi}{1}
\tightlist
\item
  needs no enumeration of vertices and is valid for any collection rule. It is conservative, increasingly so as \(q\) grows (many districts). (i) is sharper, but needs a fixed superset, and enumerating all dual vertices is feasible only in small problems. Checked (\path{verify_phase3.py}): weak duality on 500 random dual-feasible points. In a small LP, the projection bound built from the plug-in dual alone covered in 100\% of 3000 replications, against 94.5\% for a one-candidate \(z\)-bound.
\end{enumerate}

\phantomsection\label{pE-res-conjecture-P3.5-b3}\textbf{Conjecture P3.5(b3) (the implemented CLR procedure).} The package collects \(\hat\Lambda\) as the optimal dual vertices at \(\hat\beta\) and at \texttt{ndraw} Gaussian draws, and computes both \(k_p\) and \(\hat k\) over \(\hat\Lambda\). The conjecture is that it has asymptotic coverage \(\ge1-\alpha_U\) under (b1)'s conditions as \texttt{ndraw} \(\to\infty\).

\begin{itemize}
\tightlist
\item
  \emph{Proved reduction.} Suppose \(k_p\) is computed over \(\Lambda\), and \(\hat\Lambda\) contains every vertex selected by CLR's rule over \(\Lambda\) and a minimiser of \(\hat\theta_\mu+k_p\hat\sigma_\mu\) over \(\Lambda\). Then the threshold, the selected set, \(\hat k\) and \(\hat U\) all coincide with (b1), so (b1)'s coverage carries over. (The threshold is the same because its minimiser is in \(\hat\Lambda\). The selected set over \(\hat\Lambda\) is CLR's selected set intersected with \(\hat\Lambda\), which is CLR's selected set.)
\item
  \emph{Open.} (1) Recovery: vertices that are optimal only for rare Gaussian directions, degenerate optimal faces, and singular \(\Sigma\). (2) Computing \(k_p\) over \(\hat\Lambda\subseteq\Lambda\) can only lower it, so this error goes in the anti-conservative direction. (3) \(\lvert\Lambda\rvert\) grows with the number of districts, outside CLR's fixed-index setting.
\end{itemize}

Until this is settled, the CLR-type interval as implemented is an \emph{Algorithm} with simulation support (§8). The intervals of (a) and (b2) carry the formal guarantee.

\emph{Proof of (a).}

\begin{enumerate}
\def\labelenumi{\arabic{enumi}.}
\tightlist
\item
  Let \(\mathcal E\) be the event that \(s_j(\beta_0)\le s_j(\hat\beta)+\hat k\hat\sigma_{s_j}\) for all \(j\) and \(\lvert c(\beta_0)-c(\hat\beta)\rvert\le\hat k\hat\sigma_c\). Then \(P(\mathcal E)\to\ge1-\alpha_U-\alpha_L\). {[}OLS; linear maps of an asymptotically normal \(\hat\beta\), uniformly over the convex events of §0.5 (BR); Def of \(\hat k\); \(\hat\Sigma\) consistent. For M1/M2 this needs \(\dim\beta\) fixed, i.e.~R1{]}
\item
  On \(\mathcal E\), the true \(k_0\) satisfies \(s(\beta_0)+Ak_0\ge0\), hence \(s(\hat\beta)+\hat k\hat\sigma_s+Ak_0\ge0\), i.e.~\(k_0\in\hat\Theta^+\). {[}2b:Thm 2 / 2a:P1, the true value is feasible; \(\mathcal E\){]}
\item
  On \(\mathcal E\), \(\tilde B=c(\beta_0)+\gamma^\top k_0\) lies in \([c(\hat\beta)-\hat k\hat\sigma_c+\min_{\hat\Theta^+}\gamma^\top k,\ c(\hat\beta)+\hat k\hat\sigma_c+\max_{\hat\Theta^+}\gamma^\top k]\). {[}2; \(\mathcal E\); 2b:D2 with sample weights{]}
\item
  \(E=\sum_iw_ie_i\) with \(e_i\) independent given the design, mean 0, and \(\operatorname{Var}(e_i\mid V_i)\le\tfrac14\). So \(\operatorname{Var}(E\mid\text{design})\le S_1^2\). {[}P3-S1; CE; Pop{]}
\item
  \(\limsup_nP(\lvert E\rvert>z_{\alpha_E}S_1)\le\alpha_E\), provided \(\max_iw_i^2/\sum_iw_i^2\to0\). Let \(s_n=\operatorname{sd}(E\mid\text{design})\le S_1\) (step 4). It suffices that every subsequence has a further subsequence along which the \(\limsup\) is at most \(\alpha_E\). Take one along which \(s_n/S_1\to\rho\in[0,1]\).

  \begin{itemize}
  \tightlist
  \item
    If \(\rho=0\): by Chebyshev, \(P(\lvert E\rvert>z_{\alpha_E}S_1)\le s_n^2/(z_{\alpha_E}^2S_1^2)\to0\).
  \item
    If \(\rho>0\): the summands \(w_ie_i\) are independent with \(\lvert w_ie_i\rvert\le\lvert w_i\rvert\), and \(\max_i\lvert w_i\rvert/s_n\le2\max_i\lvert w_i\rvert/\big(\rho'(\textstyle\sum_iw_i^2)^{1/2}\big)\to0\) for any \(0<\rho'<\rho\) and \(n\) large. So Lindeberg's condition holds relative to \(s_n\), \(E/s_n\Rightarrow N(0,1)\), and \(P(\lvert E\rvert>z_{\alpha_E}S_1)\le P(\lvert E\rvert/s_n>z_{\alpha_E})\to\alpha_E\).
  \end{itemize}

  {[}Lyap; Chebyshev; 4{]} \emph{(Corrected after review, Batch B, Nov 2026: the earlier step applied a CLT under the weight condition alone, which fails when, e.g., only one summand has positive variance. The subsequence argument, Chebyshev where \(s_n/S_1\to0\) and Lindeberg where it is bounded away from 0, is the referee's; the conclusion and the allowance \(z_{\alpha_E}S_1\) are unchanged.)}
\item
  If \(B\notin\mathrm{CI}\), then either \(\tilde B\) lies outside the interval in 3, or \(\lvert E\rvert>z_{\alpha_E}S_1\). {[}\(B=\tilde B+E\); triangle inequality{]}
\item
  So \(P(B\notin\mathrm{CI})\le P(\mathcal E^c)+P(\lvert E\rvert>z_{\alpha_E}S_1)\to\le\alpha_U+\alpha_L+\alpha_E\). {[}union; 1, 3, 5, 6{]} ∎
\end{enumerate}

\emph{(b1)} replaces step 1 by CLR's coverage of each intersection bound, which needs the finite index set \(\Lambda\) (Lemma P3.1) and their Condition P {[}CLR{]}. Steps 4--7 are unchanged. \emph{(b2)} is proved above. The simulation-based collection of candidates is the subject of Conjecture P3.5(b3); it is not claimed as a theorem.

\textbf{Remarks.}

\begin{itemize}
\tightlist
\item
  \textbf{Better variance bound for the residual term.} \(\operatorname{Var}(e_{irc}\mid V)\le m(1-m)\) {[}BD{]}, with \(m=E[\beta_{irc}\mid V]\), which is bounded through the identified set. A data-dependent version needs its own coverage argument: Proposition P3.5(c) below.
\end{itemize}

\phantomsection\label{pE-res-proposition-P3.5-c}\textbf{Proposition P3.5(c) (variance-adaptive residual allowance; added after review, round 2, Nov 2026).} Fix a cell \((r,c)\) with weights \(w_i=w_{ir}\), put \(\pi_i=w_i^2/\sum_jw_j^2\) and \(\bar m=\sum_i\pi_im_{rc}(V_i)\), the \(\pi\)-weighted model-implied rate (a linear target in \((\beta,k)\)). By concavity of \(m(1-m)\),
\[s_n^2:=\operatorname{Var}(E\mid\text{design})\le\sum_iw_i^2m_i(1-m_i)\le\Big(\sum_iw_i^2\Big)\bar m(1-\bar m)=:\bar s^2 .\]
Let \(\mathrm{CI}_v\) be a confidence interval for \(\bar m\) at level \(1-\alpha_v\), and \(\hat S_{BD}=\big(\sum_iw_i^2\max_{m\in\mathrm{CI}_v\cap[0,1]}m(1-m)\big)^{1/2}\).

\begin{itemize}
\tightlist
\item
  \textbf{(c1) Valid version.} Suppose \(\mathrm{CI}_v\) is a P3.5(a) interval or a P3.5(b2)(ii) projection interval for \(\bar m\), computed on the \textbf{rung's own constraint system} (no added sign or CR rows), and suppose \(\max_iw_i^2/\bar s^2\to0\) (Lindeberg relative to \(\bar s\); it holds, for instance, when \(\bar m\) stays bounded away from 0 and 1). Replacing \(z_{\alpha_E}S_1\) by \(z_{\alpha_E}\hat S_{BD}\) in P3.5(a), or in the projection interval for \(B\), gives
  \[\liminf_nP(B_{rc}\in\mathrm{CI})\ge1-\alpha_{\rm est}-\alpha_E-\alpha_v ,\]
  with \(\alpha_{\rm est}\) the level spent on the estimation part. The Popoviciu version (\(S_1\), no \(\alpha_v\)) needs only \(\max_iw_i^2/\sum_iw_i^2\to0\) (step 5 of (a)).
\item
  \textbf{(c2) Adaptive version.} If \(\mathrm{CI}_v\) is the CLR interval (Conjecture P3.5(b3)), the same bound holds conditional on \(\mathrm{CI}_v\) covering \(\bar m\); the result inherits the algorithmic status of (b3).
\end{itemize}

\emph{Proof of (c1).}

\begin{enumerate}
\def\labelenumi{\arabic{enumi}.}
\tightlist
\item
  \(\bar m\) is a linear target under the rung (2a Prop. 7 with weights \(\pi\)), and the rung alone holds at the truth, so \(\mathcal E_v=\{\bar m\in\mathrm{CI}_v\}\) has \(\liminf P(\mathcal E_v)\ge1-\alpha_v\). {[}P3.5(a) or (b2)(ii){]}
\item
  On \(\mathcal E_v\), \(\hat S_{BD}\ge\bar s\ge s_n\). {[}concavity above; Def \(\hat S_{BD}\){]}
\item
  So \(\{\lvert E\rvert>z_{\alpha_E}\hat S_{BD}\}\subseteq\{\lvert E\rvert>z_{\alpha_E}\bar s\}\cup\mathcal E_v^c\), with \(\bar s\) deterministic given the design. {[}2{]}
\item
  \(\limsup P(\lvert E\rvert>z_{\alpha_E}\bar s)\le\alpha_E\): along a subsequence with \(s_n/\bar s\to0\) use Chebyshev; otherwise \(\max_i\lvert w_i\rvert/s_n\to0\) by the Lindeberg condition relative to \(\bar s\), and Lindeberg's CLT applies relative to \(s_n\), as in step 5 of (a) with \(\bar s\) in place of \(S_1\). {[}Chebyshev; Lyap{]}
\item
  Combine with the estimation event of (a) or (b2)(ii) by the union bound. {[}union; 1--4{]} ∎
\end{enumerate}

\emph{Why the condition matters.} If \(\mathrm{CI}_v\) fails, BD is not conservative: with true \(\bar m=.5\) a wrong interval inside \([0,.1]\) gives a variance bound at most \(.09\) against the true \(.25\). And if \(\bar s/S_1\to0\) (e.g.~\(\bar m\to0\)), Lindeberg relative to \(\bar s\) can fail, so step 4 needs its condition. Check: \path{verify_phase3.py}, ``P3.5(c)''. \emph{(The referee asked that the provably valid projection interval for \(B\) be separated from the variance-adaptive one. In the package (0.3.x), the projection column for \(B\) uses (c1) with \(\mathrm{CI}_v\) the projection interval on the rung alone, and prints a note when \(\mathrm{CI}_v\) comes near 0 or 1; the CLR \texttt{ci\_actual} uses (c2) and is labelled algorithmic. \(\tilde B\) is unaffected.)}

\begin{itemize}
\tightlist
\item
  \textbf{Optional assumption for a much smaller residual term.} If the poll-level deviations of different groups are not negatively correlated, then \(\operatorname{Var}(X_{ir}e_{irc})\le\operatorname{Var}(T_{ic}\mid V_i)\), which is estimable. In simulation this would remove most of the residual allowance. It is an added assumption and belongs on the menu, not in the default.
\item
  \textbf{JKST's \(\mathrm{CI}_x\)} is version (b) with a single candidate per side (tie-breaking), no selection step, and \(S_L=S_1+\hat\sigma_{\lambda^*}\), which combines the two error sources by adding standard deviations instead of by the union bound.
\end{itemize}

\phantomsection\label{pE-res-proposition-P3.6}\textbf{Proposition P3.6 (tests).}

\begin{enumerate}
\def\labelenumi{(\alph{enumi})}
\tightlist
\item
  \textbf{Emptiness test.} Reject the model (rung, foundations) at level \(\alpha_U+\alpha_L\) if \(\hat\Theta^+\) is empty.
\item
  \textbf{Test of M1's testable part.} Run a Wald/F test of the M1 regression against M2 (the overidentifying restrictions of 2b D1(b)), with \((D^\star-1)(C-1)R(R-1)/2\) restrictions. In 2×2: equal curvature across districts.
\item
  \textbf{Test of M0 against M1 (pooling).} Conditional on the design, suppose each district's observed compositions contain \(R\) affinely independent points. Under M1, M0's mean model holds for the observed units iff all district shifts vanish, \(\zeta_{c,d}=0\). So pooling is tested by a Wald/F test of \(\zeta=0\) in the M1 regression, with \((D^\star-1)(C-1)R\) restrictions. \emph{(Added after review. 2b Prop. D4(b) gives a weaker condition, \(E[Z_D\mid X,N]\) affine in \(X\). That condition averages over the district labels, so it describes the superpopulation pooled over re-drawn districts, which §0.4 does not target. Conditional on the observed labels, only \(\zeta=0\) makes M0 correct.)} (c) tests mean shifts only. It has no power against within-district residual dependence without mean shifts, which violates P3-S1a (§0.4). \emph{(Sentence added after review, round 2.)}
  \emph{Proof of (a).}
\end{enumerate}

\begin{enumerate}
\def\labelenumi{\arabic{enumi}.}
\tightlist
\item
  Under the model, \(k_0\in\hat\Theta^+\) on \(\mathcal E\). {[}P3.5(a) step 2{]}
\item
  So \(P(\hat\Theta^+=\emptyset)\le P(\mathcal E^c)\to\le\alpha_U+\alpha_L\). {[}1{]}
\end{enumerate}

\begin{enumerate}
\def\labelenumi{(\alph{enumi})}
\setcounter{enumi}{1}
\tightlist
\item
  is the standard Wald test {[}OLS{]}. Its \(\chi^2\) calibration is established for fixed dimension (R1). Under R2/R3 the number of restrictions grows with \(D^\star\). The F-form with HC3, or a wild bootstrap, is then a practical choice with simulation support only.
\end{enumerate}

\textbf{Calibration in the package (0.3.0; added after the Batch B review and the WS6 Monte Carlo, Nov 2026).} In moderate dimension the HC3 Wald statistic against \(\chi^2\) over-rejects: for the pooling test (c) in design S0 of §8.3 (38 restrictions), 0.078 at nominal 5\%, and 0.072 for the F form. The package therefore calibrates both (b) and (c) by a \textbf{restricted wild bootstrap}:

\begin{itemize}
\tightlist
\item
  fit the null model (M0 for (c); M1 for (b));
\item
  draw \(Y^*=\hat Y_0+\hat e_0v\) with Rademacher \(v_i\), one per unit, shared across the \(C-1\) equations;
\item
  refit the larger model and recompute the same HC Wald statistic; \(p=(1+\#\{W^*\ge W\})/(B+1)\), with \(B=399\) by default.
\end{itemize}

The \(\chi^2\) and F p-values are reported as well. Rejection rates at nominal 5\% in design S0 (MC standard errors in parentheses):

\begin{longtable}[]{@{}
  >{\raggedright\arraybackslash\hspace{0pt}}p{(\columnwidth - 10\tabcolsep) * \real{0.1667}}
  >{\raggedright\arraybackslash\hspace{0pt}}p{(\columnwidth - 10\tabcolsep) * \real{0.1667}}
  >{\raggedright\arraybackslash\hspace{0pt}}p{(\columnwidth - 10\tabcolsep) * \real{0.1667}}
  >{\raggedright\arraybackslash\hspace{0pt}}p{(\columnwidth - 10\tabcolsep) * \real{0.1667}}
  >{\raggedright\arraybackslash\hspace{0pt}}p{(\columnwidth - 10\tabcolsep) * \real{0.1667}}
  >{\raggedright\arraybackslash\hspace{0pt}}p{(\columnwidth - 10\tabcolsep) * \real{0.1667}}@{}}
\toprule\noalign{}
\begin{minipage}[b]{\linewidth}\raggedright\hspace{0pt}Test
\end{minipage} & \begin{minipage}[b]{\linewidth}\raggedright\hspace{0pt}Restrictions
\end{minipage} & \begin{minipage}[b]{\linewidth}\raggedright\hspace{0pt}Replications
\end{minipage} & \begin{minipage}[b]{\linewidth}\raggedright\hspace{0pt}\(\chi^2\)
\end{minipage} & \begin{minipage}[b]{\linewidth}\raggedright\hspace{0pt}F
\end{minipage} & \begin{minipage}[b]{\linewidth}\raggedright\hspace{0pt}Wild bootstrap
\end{minipage} \\
\midrule\noalign{}
\endhead
\bottomrule\noalign{}
\endlastfoot
(c) pooling, M0 vs M1 & 38 & 2000 (\(B=399\)) & 0.078 (0.006) & 0.072 (0.006) & \textbf{0.043} (0.005) \\
(b) M1 vs M2 & 19 & 1000 (\(B=199\)) & 0.057 (0.007) & 0.055 (0.007) & \textbf{0.054} (0.007) \\
\end{longtable}

Scripts: \path{benchmarks/pooling_test_size.R}, \path{benchmarks/m1_test_size.R} (outputs in \path{benchmarks/results/}). The bootstrap is a practical calibration supported by simulation; its validity as the number of restrictions grows is not proved.

\emph{Proof of (c).}

\begin{enumerate}
\def\labelenumi{\arabic{enumi}.}
\tightlist
\item
  Under M1 and P3-S1a, \(E[\beta_i\mid\text{design}]=M_H(X_i)+Z_{D_i}\), with \(Z_d=[\zeta_{c,d,r}]_{rc}\). {[}Def M1; 2b §0.1{]}
\item
  M0 holds for the observed units iff there is \(H'\) with \(M_{H'}(X_i)=M_H(X_i)+Z_{D_i}\) for all \(i\), i.e.~\(G_cX_i=\zeta_{c,D_i}\) with \(G=H'-H\). {[}Def M0{]}
\item
  Within district \(d\), \((G_c-\zeta_{c,d}\mathbf 1^\top)X_i=0\) for every \(i\) in \(d\) (using \(\mathbf 1^\top X_i=1\)). The \(X_i\) include \(R\) affinely independent points, which are linearly independent because they lie in \(\{\mathbf 1^\top x=1\}\). So \(G_c=\zeta_{c,d}\mathbf 1^\top\). {[}alg{]}
\item
  This holds for every \(d\), and \(\zeta_{c,1}=0\), so \(G_c=0\) and every \(\zeta_{c,d}=0\). Conversely, \(\zeta=0\) gives M0 with \(H'=H\). {[}3; Def M1{]}
\item
  The \(\zeta\) are coefficients of the M1 regression, identified by 2b Lemma D1(d). The Wald test of \(\zeta=0\) has \((D^\star-1)(C-1)R\) restrictions for the \(C-1\) free columns. {[}D1(d); OLS{]} ∎
\end{enumerate}

\hypertarget{pE-a-finite-sample-interval-valid-in-every-regime}{%
\subsubsection{5.1 A finite-sample interval, valid in every regime}\label{pE-a-finite-sample-interval-valid-in-every-regime}}

\phantomsection\label{pE-res-theorem-P3.8}\textbf{Theorem P3.8 (finite-sample CI and emptiness test).} Fix a rung and condition on the design. Assume P3-S1a, that the rung's mean model holds, and that its regression design \(Z\) has full column rank. Fit the \(C-1\) free columns by OLS, \(\hat\beta_c=(Z^\top Z)^{-1}Z^\top T_{\cdot c}\).

\emph{Hoeffding radius of a linear functional.} Every linear functional \(f(\beta)=f_0+\sum_{c<C}b_c^\top\beta_c\) satisfies \(f(\hat\beta)=f_0+\sum_i\sum_{c<C}w_{ic}(f)T_{ic}\), with \(w_{ic}(f)=b_c^\top(Z^\top Z)^{-1}z_i\). Put \(w_{iC}(f):=0\), \(\rho_i(f):=\max_cw_{ic}(f)-\min_cw_{ic}(f)\), and
\[h(f;\delta)=\Big(\tfrac12\log(1/\delta)\sum_i\rho_i(f)^2\Big)^{1/2}.\]

\emph{Interval.} Let \(J\) be the number of constraint rows and \(\delta=\alpha_s/(J+2)\). Put \(h_j=h(s_j;\delta)\), \(h_c=h(c;\delta)\), \(t_E=\big(\tfrac12\log(2/\alpha_E)\sum_iw_{ir}^2\big)^{1/2}\), and
\[\hat\Theta^{\rm fin}=\{k:s_j(\hat\beta)+h_j+a_j^\top k\ge0\ \ \forall j\},\]
\[\mathrm{CI}^{\rm fin}=\Big[c(\hat\beta)-h_c+\min_{\hat\Theta^{\rm fin}}\gamma^\top k-t_E,\ \ c(\hat\beta)+h_c+\max_{\hat\Theta^{\rm fin}}\gamma^\top k+t_E\Big].\]
Then, \textbf{for every \(n\)} and in every regime:

\begin{enumerate}
\def\labelenumi{(\alph{enumi})}
\tightlist
\item
  \(P(B_{rc}\in\mathrm{CI}^{\rm fin}\mid\text{design})\ge1-\alpha_s-\alpha_E\);
\item
  without \(t_E\), \(P(\tilde B_{rc}\in\mathrm{CI}^{\rm fin}\mid\text{design})\ge1-\alpha_s\);
\item
  rejecting the rung when \(\hat\Theta^{\rm fin}=\emptyset\) is a test of exact level at most \(\alpha_s\);
\item
  intersecting with Duncan--Davis keeps (a), since \(B\in DD\) always.
\end{enumerate}

\emph{Proof.}

\begin{enumerate}
\def\labelenumi{\arabic{enumi}.}
\tightlist
\item
  \(T_i\in\Delta^{C-1}\), so \(Y_i(f):=\sum_cw_{ic}(f)T_{ic}\) is a convex combination of the \(w_{ic}(f)\). It lies in an interval of length \(\rho_i(f)\). {[}S1; Def{]}
\item
  Under the mean model and full rank, OLS is conditionally unbiased: \(E[f(\hat\beta)\mid\text{design}]=f(\beta_0)\). {[}linearity of OLS; \(E[T_{ic}\mid\text{design}]=z_i^\top\beta_{c,0}\){]}
\item
  The \(Y_i(f)\) are independent given the design, so \(P\big(f(\beta_0)-f(\hat\beta)\ge h(f;\delta)\big)\le\delta\), and the same for the other side. {[}P3-S1a; Hoeff; 1--2{]}
\item
  Let \(\mathcal E\) be the event that \(s_j(\beta_0)\le s_j(\hat\beta)+h_j\) for all \(j\) and \(\lvert c(\beta_0)-c(\hat\beta)\rvert\le h_c\). It is an intersection of \(J+2\) one-sided events, so \(P(\mathcal E^c)\le(J+2)\delta=\alpha_s\). {[}3; union{]}
\item
  On \(\mathcal E\), the true \(k_0\) satisfies \(s(\beta_0)+Ak_0\ge0\), hence \(k_0\in\hat\Theta^{\rm fin}\), and \(\tilde B=c(\beta_0)+\gamma^\top k_0\) lies in the interval without \(t_E\). That proves (b), and (c), since \(\hat\Theta^{\rm fin}\ne\emptyset\) on \(\mathcal E\). {[}P3.5(a) steps 2--3{]}
\item
  \(E=\sum_iw_{ir}e_{irc}\) with independent \(e_{irc}\in[-m_{rc}(V_i),1-m_{rc}(V_i)]\), an interval of length 1. So \(P(\lvert E\rvert\ge t_E)\le2\exp(-2t_E^2/\sum_iw_{ir}^2)=\alpha_E\). {[}P3-S1a; Hoeff{]}
\item
  \(B=\tilde B+E\). A union bound over 4 and 6 gives (a). (d) holds because \(B\in DD\) surely. {[}union; 2b:D0{]} ∎
\end{enumerate}

\emph{Remarks.}

\begin{itemize}
\tightlist
\item
  \textbf{What it needs, and what it does not.}

  \begin{itemize}
  \tightlist
  \item
    Needs: independence given the design (P3-S1a), a correctly specified mean model, and a full-rank design.
  \item
    Does not need: Gaussian approximation, a fixed dimension, Slater, or any regime.
  \end{itemize}

  It is therefore the only formal guarantee for M1 with many small districts (R2), where the asymptotic results are not established.
\item
  \textbf{Width: valid but currently uninformative.} In a 2×2 test design with linear means, Beta-distributed rates and composition ranges of 0.2--0.45 per district, the interval clipped to \([0,1]\) had median width 1.0 in every configuration tried:

  \begin{itemize}
  \tightlist
  \item
    pooled, with 900 and with 5,000 polls;
  \item
    5 districts × 200 polls;
  \item
    10 × 93;
  \item
    30 × 8.
  \end{itemize}

  Two costs combine:

  \begin{enumerate}
  \def\labelenumi{\arabic{enumi}.}
  \tightlist
  \item
    Hoeffding uses the worst-case variance \(\tfrac14\) for each \(T_{ic}\), where the true variance was about 0.01.
  \item
    The separate union bound over all \(J\) constraints and the constant inflates each radius, and each radius is then divided by the distance from the vertex to a pure composition.
  \end{enumerate}

  An \emph{oracle} Gaussian union bound with the true variances still gave widths of 0.76--0.87, and 1.0 at 30 × 8. So most of the cost is the constraint-by-constraint union bound, which the conservative asymptotic interval P3.5(a) shares. Theorem P3.8 settles validity in R2, but it does not yet give useful intervals there. Informative finite-sample inference for R2 is open (§10).
\end{itemize}

Checked (\path{verify_phase3.py}): in the 30 × 8 design, coverage of \(\tilde B\) and \(B\) was 100\% in 300 replications each, and the interval was never empty. In the 1000-replication Monte Carlo (§8.2) it covered in every replication of five M1 designs (10 × 93, 93 × 30, and 30, 93 and 200 districts × 10 polls), with clipped width 1.000 in all of them.

\begin{itemize}
\tightlist
\item
  \textbf{Possible refinements, not proved here.}

  \begin{itemize}
  \tightlist
  \item
    An empirical-Bernstein bound (Maurer \& Pontil 2009) would replace the ranges \(\rho_i\) by estimated variances.
  \item
    A Bhatia--Davis version of \(t_E\), as in §5.
  \item
    Splitting \(\alpha_s\) unevenly across constraints, e.g.~giving less to constraints that cannot bind.
  \end{itemize}
\end{itemize}

\begin{center}\rule{0.5\linewidth}{0.5pt}\end{center}

\hypertarget{pE-many-districts-with-few-units-regime-r2}{%
\subsection{6. Many districts with few units (regime R2)}\label{pE-many-districts-with-few-units-regime-r2}}

\phantomsection\label{pE-res-proposition-P3.7}\textbf{Proposition P3.7.}

\begin{enumerate}
\def\labelenumi{(\alph{enumi})}
\tightlist
\item
  \emph{(Dropping constraints is conservative.)} For any subset \(\mathcal D'\) of districts, \(\Theta^{M1}_O(\mathcal D')\supseteq\Theta^{M1}_O\).
\item
  \emph{(Design-based selection preserves conditional validity.)} Suppose \(\mathcal D'\) is chosen as a function of the design \(\{(N_i,X_i,D_i)\}\) only. Then any procedure that has (finite-sample or asymptotic) coverage when applied with a fixed set of districts keeps that coverage conditionally on the selected \(\mathcal D'\). In particular, Theorem P3.8 holds for any design-selected \(\mathcal D'\). The asymptotic results of §§3--5 transfer only if their own regime conditions (§0.5) hold for the retained districts. Selection does not by itself supply them: retained districts with bounded \(n_d\) are still in R2.
  \emph{Proof.}
\end{enumerate}

\begin{enumerate}
\def\labelenumi{\arabic{enumi}.}
\tightlist
\item
  Fewer constraints give a weakly larger feasible set. {[}Def; 2b:Thm 2(b){]}
\item
  Conditioning on the design fixes \(\mathcal D'\). A procedure's conditional coverage given the design, for the fixed set \(\mathcal D'\), is therefore its coverage given the design and \(\mathcal D'\). {[}P3-S2{]} ∎
\end{enumerate}

\emph{(Revised after review: the earlier statement said every result of §§3--5 holds conditionally on \(\mathcal D'\), which overstated what selection achieves.)}

\textbf{What happens under R2.}

\begin{itemize}
\tightlist
\item
  Each district's shifts \(\hat\zeta_d\) come from about \(n_d\) units, and their sampling error does not shrink.
\item
  The number of candidates grows with \(D^\star\).
\item
  The plug-in is badly too narrow. In the Monte Carlo (§8.2), with 10 polls per district, it covers 31\%, 4\% and 0.5\% at 30, 93 and 200 districts, and is empty 24\%, 46\% and 72\% of the time. This is the Theorem P3.4(ii) bias, compounded across many districts.
\item
  In the Monte Carlo (§8.2), the CLR interval covered 0.995--1 in every R2 design, up to 200 districts. This is a \emph{simulation finding} (§0.5), not a theorem. CLR inflates each candidate in proportion to its own standard error, so noisy districts rarely bind, which is precision weighting built into the method.
\item
  The formal guarantee in R2 is Theorem P3.8, at the cost of width: in §8.2 it is \([0,1]\) in every design. The projection interval is also \([0,1]\) there, because \(q\) grows with \(D^\star\).
\item
  With 10 units per district, CLR's interval is about as wide as Duncan--Davis at 30--93 districts, so the identification gain from district dummies (2b) is consumed by sampling noise. It narrows as districts are added, and is below DD at 200 districts (0.42 for \(\tilde B\) and 0.49 for \(B\), against 0.57).
\end{itemize}

Practical options. Each keeps the identified set valid (P3.7(a)). Their inference is formally covered by Theorem P3.8, and asymptotically only under §0.5's conditions:

\begin{itemize}
\item
  \begin{enumerate}
  \def\labelenumi{(\roman{enumi})}
  \tightlist
  \item
    Keep the constraints of districts with enough units or composition spread (P3.7).
  \end{enumerate}
\item
  \begin{enumerate}
  \def\labelenumi{(\roman{enumi})}
  \setcounter{enumi}{1}
  \tightlist
  \item
    Use M1(\(\delta\)), \(\delta>0\), which bounds how far any one district can pull the shared parameter.
  \end{enumerate}
\item
  \begin{enumerate}
  \def\labelenumi{(\roman{enumi})}
  \setcounter{enumi}{2}
  \tightlist
  \item
    Pool districts into larger groups.
  \end{enumerate}
\item
  \begin{enumerate}
  \def\labelenumi{(\roman{enumi})}
  \setcounter{enumi}{3}
  \tightlist
  \item
    Target \(\tilde B\) to drop the residual term.
  \end{enumerate}
\end{itemize}

A hierarchical restriction on the district shifts (they are draws from a distribution) would change the model and is left open.

\begin{center}\rule{0.5\linewidth}{0.5pt}\end{center}

\hypertarget{pE-sensitivity-analysis}{%
\subsection{7. Sensitivity analysis}\label{pE-sensitivity-analysis}}

\phantomsection\label{pE-res-proposition-S1}\textbf{Proposition S1 (bounded nonlinearity).} Replace (S2) by \(E[\beta\mid V]=M_H(X)+\Delta(X)\), with \(\Delta\) in the undetectable class of 2a Prop. 9 (\(x^\top\Delta(x)=0\), \(\Delta(x)\mathbf 1=0\)) and \(\lVert\Delta(x)\rVert_\infty\le\eta\). Then \(Q\) is still identified as before, and every value of \(\tilde B_{rc}\) consistent with the data lies in
\[\Big[\min_{\Theta_O(\eta)}\varphi-\eta,\ \max_{\Theta_O(\eta)}\varphi+\eta\Big],\qquad\Theta_O(\eta)=\{k:s+Ak\ge-\eta\mathbf 1\}.\]
For a general \(\varphi=\sum a_{rc}B_{rc}\), widen by \(\eta\sum\lvert a_{rc}\rvert\) instead. This is an \textbf{outer} bound under bounded deviations: it is valid, but it is not claimed to be the sharp set of the relaxed model.

\emph{Proof.}

\begin{enumerate}
\def\labelenumi{\arabic{enumi}.}
\tightlist
\item
  \(E[T\mid X]=X^\top Q X+X^\top\Delta(X)=X^\top QX\), so \(Q\) is unchanged. {[}2a:L1(a); \(x^\top\Delta=0\){]}
\item
  \(M_H(x)+\Delta(x)\in\mathfrak B\) implies \(M_H(x)\ge-\Delta(x)\ge-\eta\) elementwise on \(A\). \(M_H\) is linear, so this is equivalent to the same inequality at the vertices: \(s+Ak\ge-\eta\). {[}S3; 2a:L3(a){]}
\item
  \(\tilde B_{rc}=\sum_iw_i[M_H(X_i)+\Delta(X_i)]_{rc}=\varphi(k)+\sum_iw_i\Delta_{rc}(X_i)\), and the last term lies in \([-\eta,\eta]\) because the weights sum to 1. {[}2b:D2; Def{]} ∎
\end{enumerate}

Checked: a nonlinear undetectable violation with \(\eta=0.02\) stays inside the bound, and the bounds are nested in \(\eta\).

\textbf{Breakdown value.} For a conclusion \(\mathcal C\), e.g.~``loyalty \textgreater{} 0.30'', define \(\eta^*(\mathcal C)=\sup\{\eta:\mathcal C\text{ holds for every value in the }\eta\text{-bound}\}\). The bounds are nested in \(\eta\), so the conclusion holds for all \(\eta<\eta^*\) and fails above it, and \(\eta^*\) can be found by bisection. In the 3×3 example, \(\eta^*(\text{loyalty}>0.30)=0.060\): the conclusion survives any undetectable nonlinearity of up to 6 percentage points. Apply the CI version of the bound to account for sampling error.

\textbf{Other dials, same logic} (nested sets, bisection for breakdown):

\begin{itemize}
\tightlist
\item
  \(\delta\) for district heterogeneity (2b Prop. D6);
\item
  \(\delta\) for change between periods (menu item 19);
\item
  population turnover \(\epsilon\) (A1). Turnover requires relaxing the Duncan--Davis sets and the regression identity together. It is left as an open item, with interim advice to add explicit entry/exit rows.
\end{itemize}

\begin{center}\rule{0.5\linewidth}{0.5pt}\end{center}

\hypertarget{pE-monte-carlo}{%
\subsection{8. Monte Carlo}\label{pE-monte-carlo}}

\emph{Rerun for review item R-37 (Nov 2026). It replaces the earlier 200-replication table.}

\path{phase3_montecarlo.py}, \textbf{1000 replications per design} (\path{phase3_montecarlo_output.txt}, summary in \path{phase3_montecarlo_summary.json}). Unit rates are Beta/Dirichlet draws whose conditional means are exactly linear in composition given the district. So M1 always holds, and M0 holds only when the district shifts are zero. Every interval is nominal 95\%:

\begin{itemize}
\tightlist
\item
  \textbf{plug-in}: the estimate of the identified set (not a CI);
\item
  \textbf{conservative} (Thm P3.5(a));
\item
  \textbf{CLR} as implemented (Conjecture P3.5(b3)), for \(B\) with the Bhatia--Davis allowance (the package default) and for \(\tilde B\);
\item
  \textbf{projection} (Prop. P3.5(b2)(ii)): \(k=\sqrt{\chi^2_{q,1-\alpha}}\) with \(q=\operatorname{rank}\hat\Sigma\), applied to the same candidates as CLR. In every 2×2 design here \(\hat\Sigma\) has full rank, so \(q=\dim\beta\) and this is case (ii-a) of the corrected proposition. (Since 0.3.0 the package also drops candidates whose gradients use a numerical null direction of \(\hat V\). \emph{Corrected after review, round 2:} that is a safeguard, not case (ii-b); see the remark after P3.5(b2).) For \(B\): estimation part at level 0.04, plus the Bhatia--Davis allowance with \(\bar m\) bounded by the \textbf{projection} interval at level \(\alpha_V=0.005\) on the rung alone (P3.5(c1); column \texttt{proj\_B}). The earlier runs took \(\bar m\) from a CLR interval; that version is kept as \texttt{proj\_B\_clrbd}, labelled \textbf{not covered} (P3.5(c2));
\item
  \textbf{finite-sample} (Thm P3.8), in the 2×2 M1 designs.
\end{itemize}

Conventions:

\begin{itemize}
\tightlist
\item
  \textbf{Coverage} is the share of 1000 replications whose interval contains the truth. The Monte Carlo standard error, \(\sqrt{p(1-p)/1000}\), is in parentheses. A coverage of \textbf{1} means all 1000 replications covered; the exact one-sided 95\% lower confidence bound is then \(0.05^{1/1000}=0.997\).
\item
  An empty interval (plug-in or conservative set empty, or CLR/projection ends crossing) counts as not covering.
\item
  \textbf{Width} is the median after clipping to \([0,1]\). The Duncan--Davis (DD) width is shown for reference.
\item
  Each cell reads ``coverage / width''.
\end{itemize}

\textbf{Cross-check against the package.} \path{phase3_montecarlo_crosscheck.py} runs one replication of designs A (10 × 93) and S2 through \texttt{eiBounds()} (\path{phase3_montecarlo_crosscheck_output.txt}). The two agree exactly on the plug-in, the projection interval for \(\tilde B\) and the pooling-test statistic. The CLR ends differ by up to 0.01, because the script uses every constraint as a candidate (the full dual-vertex set in 2×2) while the package collects candidates by simulation. The package also intersects \texttt{ci\_actual}/\texttt{proj\_actual} with DD.

\textbf{Which figures correspond to proved results.}

\begin{itemize}
\item
  Conservative and projection intervals: Thm P3.5(a) and Prop. P3.5(b2)(ii) are asymptotic theorems for a fixed-dimensional regression: M0 in every design here, and M1 with few, large districts (A 10 × 93; S 20 × 50). In the R2 rows (93 × 30, and 10 polls per district) \(q\) grows with \(D^\star\), and their coverage there is a simulation finding (§0.5).
\item
  Finite-sample interval: Thm P3.8 holds in every row.
\item
  CLR: Conjecture P3.5(b3) everywhere. All its figures are simulation findings.
\item
  The projection interval for \(B\) (\texttt{proj\_B}) uses the Bhatia--Davis allowance with \(\bar m\) bounded by a projection interval on the rung alone, so it is covered by P3.5(b2)(ii) and P3.5(c1), subject to (c1)'s Lindeberg condition. \emph{(Corrected after review, round 2, Nov 2026: this bullet said the variance bound came from a CLR interval ``as in the package''. Since 0.3.0 the package takes it from the projection interval, and the §8 column has been rerun to match (table below). The CLR-based version, \texttt{proj\_B\_clrbd}, is not covered by any theorem.)}
\item
  \textbf{Rerun (round 2, Nov 2026)} of the projection column for \(B\) with the P3.5(c1) bound, same seeds (no other column uses the changed code, and every other column reproduces the earlier run exactly; \path{phase3_montecarlo_output.txt}):

  \begin{longtable}[]{@{}
    >{\raggedright\arraybackslash\hspace{0pt}}p{(\columnwidth - 4\tabcolsep) * \real{0.3333}}
    >{\raggedright\arraybackslash\hspace{0pt}}p{(\columnwidth - 4\tabcolsep) * \real{0.3333}}
    >{\raggedright\arraybackslash\hspace{0pt}}p{(\columnwidth - 4\tabcolsep) * \real{0.3333}}@{}}
  \toprule\noalign{}
  \begin{minipage}[b]{\linewidth}\raggedright\hspace{0pt}  Design
  \end{minipage} & \begin{minipage}[b]{\linewidth}\raggedright\hspace{0pt}  \texttt{proj\_B} (P3.5(c1), \(\bar m\) from the projection interval): coverage / median width
  \end{minipage} & \begin{minipage}[b]{\linewidth}\raggedright\hspace{0pt}  \texttt{proj\_B\_clrbd} (\(\bar m\) from CLR; earlier column, \textbf{not covered})
  \end{minipage} \\
  \midrule\noalign{}
  \endhead
  \bottomrule\noalign{}
  \endlastfoot
  A 10 × 93 (M1) & 1.000 / 0.662 & 1.000 / 0.662 \\
  A 93 × 30 (M1) & 1.000 / 0.964 & 1.000 / 0.964 \\
  R2 30 × 10, 93 × 10, 200 × 10 (M1) & 1.000 / 1.000 each & 1.000 / 1.000 each \\
  C boundary (M0) & 1.000 / 0.059 & 1.000 / 0.059 \\
  E Slater (M0) & 1.000 / 0.251 & 1.000 / 0.249 \\
  S0, S1 (M0; M1) & 1.000 / 0.436, 0.444; 1.000 / 0.808, 0.773 & identical \\
  S2 non-affine (M0, misspecified) & 0.206 (SE 0.013) / 0.038, empty 0.478 & identical \\
  S2 (M1) & 1.000 / 0.611 & 1.000 / 0.611 \\
  B 3×3, four targets (M0) & 1.000 / 0.679, 0.445, 0.378, 0.661 & 1.000 / 0.679, 0.445, 0.377, 0.661 \\
  \end{longtable}

  Coverage is unchanged in every design, and widths change by at most 0.002 (design E and one design-B target). In these designs the projection interval for \(\bar m\) is only slightly wider than CLR's, and the Bhatia--Davis maximum \(\max m(1-m)\) over it is flat near the plausible rates. So the provably valid version costs essentially nothing here. The S2/M0 row is misspecification, not a failure of the theorem.
\end{itemize}

\textbf{Rank of \(\hat\Sigma\) in the projection interval} (\path{phase3_montecarlo_rank.py}, \path{phase3_montecarlo_rank_output.txt}; same seeds, 1000 replications per design). These runs use \(q=\operatorname{rank}\hat\Sigma\), which equals \(\dim\beta\) in every 2×2 design and the exact adding-up reduction in design B (below); after the round-2 review the theorem is claimed with \(q=\dim\beta\) or a documented exact reduction, and a numerical cutoff is only a safeguard (remark after P3.5(b2)). Since the Batch B correction P3.5(b2)(ii) requires either a positive definite \(\Sigma\) (ii-a) or a reduced parameterisation with exactly zero error in the null directions (ii-b). If \(\hat\Sigma\) were singular while the estimation error had a component in its null space, those directions would not be controlled. The numerical rank (eigenvalues above \(10^{-10}\) times the largest) equals \(\dim\beta\) in every replication of every 2×2 design, from \(\dim\beta=3\) to 401.

\begin{itemize}
\tightlist
\item
  Conditioning: the smallest-to-largest eigenvalue ratio is 0.005--0.04 for the pooled fits, and falls with the number of districts, to a minimum of \(2.6\times10^{-8}\) at 200 × 10. That is still far above the rank threshold.
\item
  Design B fits all three columns (\(\dim\beta=18\)) and has \(q=12\) in every replication. The six null directions are the adding-up constraint: summed across columns, the coefficients are the regression of 1 on the design, which has no error. So the estimation error lies in the range of \(\hat\Sigma\), and nothing is left uncontrolled.
\item
  On the \(C-1\) free columns (Remark P3.4.1, as the package fits), \(\hat\Sigma\) has full rank 12 in every replication.
\item
  No replication had a numerical rank drop, so there is no rank-drop subsample whose coverage could be compared.
\end{itemize}

\hypertarget{pE-model-holds-m1-or-m0-without-shifts}{%
\subsubsection{8.1 Model holds (M1, or M0 without shifts)}\label{pE-model-holds-m1-or-m0-without-shifts}}

\begin{longtable}[]{@{}
  >{\raggedright\arraybackslash\hspace{0pt}}p{(\columnwidth - 18\tabcolsep) * \real{0.1000}}
  >{\raggedright\arraybackslash\hspace{0pt}}p{(\columnwidth - 18\tabcolsep) * \real{0.1000}}
  >{\raggedright\arraybackslash\hspace{0pt}}p{(\columnwidth - 18\tabcolsep) * \real{0.1000}}
  >{\raggedright\arraybackslash\hspace{0pt}}p{(\columnwidth - 18\tabcolsep) * \real{0.1000}}
  >{\raggedright\arraybackslash\hspace{0pt}}p{(\columnwidth - 18\tabcolsep) * \real{0.1000}}
  >{\raggedright\arraybackslash\hspace{0pt}}p{(\columnwidth - 18\tabcolsep) * \real{0.1000}}
  >{\raggedright\arraybackslash\hspace{0pt}}p{(\columnwidth - 18\tabcolsep) * \real{0.1000}}
  >{\raggedright\arraybackslash\hspace{0pt}}p{(\columnwidth - 18\tabcolsep) * \real{0.1000}}
  >{\raggedright\arraybackslash\hspace{0pt}}p{(\columnwidth - 18\tabcolsep) * \real{0.1000}}
  >{\raggedright\arraybackslash\hspace{0pt}}p{(\columnwidth - 18\tabcolsep) * \real{0.1000}}@{}}
\toprule\noalign{}
\begin{minipage}[b]{\linewidth}\raggedright\hspace{0pt}Design
\end{minipage} & \begin{minipage}[b]{\linewidth}\raggedright\hspace{0pt}\(q\)
\end{minipage} & \begin{minipage}[b]{\linewidth}\raggedright\hspace{0pt}DD
\end{minipage} & \begin{minipage}[b]{\linewidth}\raggedright\hspace{0pt}Plug-in: cover / empty / width
\end{minipage} & \begin{minipage}[b]{\linewidth}\raggedright\hspace{0pt}Conservative, \(B\)
\end{minipage} & \begin{minipage}[b]{\linewidth}\raggedright\hspace{0pt}CLR, \(B\)
\end{minipage} & \begin{minipage}[b]{\linewidth}\raggedright\hspace{0pt}Projection, \(B\)
\end{minipage} & \begin{minipage}[b]{\linewidth}\raggedright\hspace{0pt}Conservative, \(\tilde B\)
\end{minipage} & \begin{minipage}[b]{\linewidth}\raggedright\hspace{0pt}CLR, \(\tilde B\)
\end{minipage} & \begin{minipage}[b]{\linewidth}\raggedright\hspace{0pt}Projection, \(\tilde B\)
\end{minipage} \\
\midrule\noalign{}
\endhead
\bottomrule\noalign{}
\endlastfoot
A: 2×2, M1, 10 districts × 93 & 21 & 0.574 & 0.982 (0.004) / 0.000 / 0.287 & 1 / 0.696 & 1 / 0.538 & 1 / 0.662 & 1 / 0.596 & 1 / 0.432 & 1 / 0.554 \\
A: 2×2, M1, 93 × 30 & 187 & 0.567 & 0.656 (0.015) / 0.002 / 0.160 & 1 / 0.694 & 1 / 0.477 & 1 / 0.964 & 1 / 0.633 & 1 / 0.414 & 1 / 0.932 \\
C: 2×2 pooled, truth near the boundary & 3 & 0.325 & 0.661 (0.015) / 0.076 / 0.031 & 1 / 0.087 & 0.999 (0.001) / 0.057 & 1 / 0.059 & 0.996 (0.002) / 0.041 & 0.994 (0.002) / 0.039 & 0.999 (0.001) / 0.041 \\
\textbf{E: 2×2 pooled, Slater fails} & 3 & 0.271 & \textbf{0.267 (0.014) / 0.503} / 0.038 & 1 / 0.249 & 0.998 (0.001) / 0.207 & 1 / 0.251 & 0.990 (0.003) / 0.179 & \textbf{0.944 (0.007)} / 0.143 & 0.995 (0.002) / 0.194 \\
B: 3×3 pooled, loyalty P→P & 12 & 0.835 & 1 / 0.000 / 0.533 & 1 / 0.675 & 1 / 0.659 & 1 / 0.679 & 1 / 0.592 & 1 / 0.557 & 1 / 0.581 \\
B: 3×3 pooled, defection P→O & 12 & 0.635 & 1 / 0.000 / 0.348 & 1 / 0.447 & 1 / 0.431 & 1 / 0.445 & 1 / 0.399 & 1 / 0.369 & 1 / 0.395 \\
B: 3×3 pooled, demobilisation & 12 & 0.533 & 1 / 0.000 / 0.283 & 1 / 0.381 & 1 / 0.363 & 1 / 0.378 & 1 / 0.333 & 1 / 0.304 & 1 / 0.329 \\
B: 3×3 pooled, mobilisation N→P & 12 & 0.884 & 1 / 0.000 / 0.528 & 1 / 0.683 & 1 / 0.646 & 1 / 0.661 & 1 / 0.619 & 1 / 0.559 & 1 / 0.595 \\
\end{longtable}

Design E puts the truth on the boundary so that Slater fails. Compositions span \([0.02,0.98]\) with both ends observed. The mean rate of group \(b\) is exactly 0 at the right vertex and that of group \(w\) is exactly 0 at the left vertex. The lower and upper constraints on \(k\) then meet at the truth, so the identified set of the model-implied rate is a single point, and no \(k\) satisfies every constraint strictly. Design C is the earlier boundary design: the rate is close to 0 at one vertex, so Slater holds, but only just. The CLR interval for \(B\) with the Popoviciu allowance covered in all 1000 replications of every row. Its widths are 0.530, 0.472, 0.086, 0.221 and 0.651/0.428/0.363/0.641 (rows in order); full output in \path{phase3_montecarlo_output.txt}.

\hypertarget{pE-many-small-districts-r2-m1-22-10-polls-per-district}{%
\subsubsection{8.2 Many small districts (R2): M1, 2×2, 10 polls per district}\label{pE-many-small-districts-r2-m1-22-10-polls-per-district}}

\begin{longtable}[]{@{}
  >{\raggedright\arraybackslash\hspace{0pt}}p{(\columnwidth - 20\tabcolsep) * \real{0.0909}}
  >{\raggedright\arraybackslash\hspace{0pt}}p{(\columnwidth - 20\tabcolsep) * \real{0.0909}}
  >{\raggedright\arraybackslash\hspace{0pt}}p{(\columnwidth - 20\tabcolsep) * \real{0.0909}}
  >{\raggedright\arraybackslash\hspace{0pt}}p{(\columnwidth - 20\tabcolsep) * \real{0.0909}}
  >{\raggedright\arraybackslash\hspace{0pt}}p{(\columnwidth - 20\tabcolsep) * \real{0.0909}}
  >{\raggedright\arraybackslash\hspace{0pt}}p{(\columnwidth - 20\tabcolsep) * \real{0.0909}}
  >{\raggedright\arraybackslash\hspace{0pt}}p{(\columnwidth - 20\tabcolsep) * \real{0.0909}}
  >{\raggedright\arraybackslash\hspace{0pt}}p{(\columnwidth - 20\tabcolsep) * \real{0.0909}}
  >{\raggedright\arraybackslash\hspace{0pt}}p{(\columnwidth - 20\tabcolsep) * \real{0.0909}}
  >{\raggedright\arraybackslash\hspace{0pt}}p{(\columnwidth - 20\tabcolsep) * \real{0.0909}}
  >{\raggedright\arraybackslash\hspace{0pt}}p{(\columnwidth - 20\tabcolsep) * \real{0.0909}}@{}}
\toprule\noalign{}
\begin{minipage}[b]{\linewidth}\raggedright\hspace{0pt}Districts
\end{minipage} & \begin{minipage}[b]{\linewidth}\raggedright\hspace{0pt}\(q\)
\end{minipage} & \begin{minipage}[b]{\linewidth}\raggedright\hspace{0pt}DD
\end{minipage} & \begin{minipage}[b]{\linewidth}\raggedright\hspace{0pt}Plug-in: cover / empty
\end{minipage} & \begin{minipage}[b]{\linewidth}\raggedright\hspace{0pt}CLR, \(\tilde B\)
\end{minipage} & \begin{minipage}[b]{\linewidth}\raggedright\hspace{0pt}Projection, \(\tilde B\)
\end{minipage} & \begin{minipage}[b]{\linewidth}\raggedright\hspace{0pt}Thm P3.8, \(\tilde B\)
\end{minipage} & \begin{minipage}[b]{\linewidth}\raggedright\hspace{0pt}CLR, \(B\)
\end{minipage} & \begin{minipage}[b]{\linewidth}\raggedright\hspace{0pt}Projection, \(B\)
\end{minipage} & \begin{minipage}[b]{\linewidth}\raggedright\hspace{0pt}Thm P3.8, \(B\)
\end{minipage} & \begin{minipage}[b]{\linewidth}\raggedright\hspace{0pt}Conservative, \(B\)
\end{minipage} \\
\midrule\noalign{}
\endhead
\bottomrule\noalign{}
\endlastfoot
30 & 61 & 0.571 & 0.305 (0.015) / 0.241 & 0.999 (0.001) / 0.557 & 1 / 0.967 & 1 / 1.000 & 1 / 0.746 & 1 / 1.000 & 1 / 1.000 & 1 / 1.000 \\
93 & 187 & 0.566 & 0.042 (0.006) / 0.463 & 0.995 (0.002) / 0.467 & 1 / 1.000 & 1 / 1.000 & 0.998 (0.001) / 0.575 & 1 / 1.000 & 1 / 1.000 & 1 / 0.930 \\
200 & 401 & 0.566 & 0.005 (0.002) / 0.717 & 0.996 (0.002) / 0.416 & 1 / 1.000 & 1 / 1.000 & 1 / 0.491 & 1 / 1.000 & 1 / 1.000 & 1 / 0.890 \\
\end{longtable}

Theorem P3.8 also covered in all 1000 replications of designs A 10 × 93 and 93 × 30, with median clipped width 1.000 in both, for \(\tilde B\) and \(B\).

\hypertarget{pE-riding-shocks-m0-against-m1-22-20-districts-50-polls}{%
\subsubsection{8.3 Riding shocks: M0 against M1, 2×2, 20 districts × 50 polls}\label{pE-riding-shocks-m0-against-m1-22-20-districts-50-polls}}

The district shifts are added to both groups' rates, and depend on the district's mean composition \(m_d\):

\begin{itemize}
\tightlist
\item
  \textbf{S0:} none;
\item
  \textbf{S1:} affine in \(m_d\), plus noise (as in design A);
\item
  \textbf{S2:} a common shock \(0.25\cos(2\pi m_d)\), non-affine in \(m_d\): the reviewer's common-shock example.
\end{itemize}

\begin{longtable}[]{@{}
  >{\raggedright\arraybackslash\hspace{0pt}}p{(\columnwidth - 16\tabcolsep) * \real{0.1111}}
  >{\raggedright\arraybackslash\hspace{0pt}}p{(\columnwidth - 16\tabcolsep) * \real{0.1111}}
  >{\raggedright\arraybackslash\hspace{0pt}}p{(\columnwidth - 16\tabcolsep) * \real{0.1111}}
  >{\raggedright\arraybackslash\hspace{0pt}}p{(\columnwidth - 16\tabcolsep) * \real{0.1111}}
  >{\raggedright\arraybackslash\hspace{0pt}}p{(\columnwidth - 16\tabcolsep) * \real{0.1111}}
  >{\raggedright\arraybackslash\hspace{0pt}}p{(\columnwidth - 16\tabcolsep) * \real{0.1111}}
  >{\raggedright\arraybackslash\hspace{0pt}}p{(\columnwidth - 16\tabcolsep) * \real{0.1111}}
  >{\raggedright\arraybackslash\hspace{0pt}}p{(\columnwidth - 16\tabcolsep) * \real{0.1111}}
  >{\raggedright\arraybackslash\hspace{0pt}}p{(\columnwidth - 16\tabcolsep) * \real{0.1111}}@{}}
\toprule\noalign{}
\begin{minipage}[b]{\linewidth}\raggedright\hspace{0pt}Design
\end{minipage} & \begin{minipage}[b]{\linewidth}\raggedright\hspace{0pt}Fit
\end{minipage} & \begin{minipage}[b]{\linewidth}\raggedright\hspace{0pt}Plug-in: cover / empty
\end{minipage} & \begin{minipage}[b]{\linewidth}\raggedright\hspace{0pt}CLR, \(\tilde B\): cover / empty / width
\end{minipage} & \begin{minipage}[b]{\linewidth}\raggedright\hspace{0pt}Projection, \(\tilde B\)
\end{minipage} & \begin{minipage}[b]{\linewidth}\raggedright\hspace{0pt}Conservative, \(\tilde B\)
\end{minipage} & \begin{minipage}[b]{\linewidth}\raggedright\hspace{0pt}CLR, \(B\)
\end{minipage} & \begin{minipage}[b]{\linewidth}\raggedright\hspace{0pt}Projection, \(B\)
\end{minipage} & \begin{minipage}[b]{\linewidth}\raggedright\hspace{0pt}Pooling test: rejection at 5\%
\end{minipage} \\
\midrule\noalign{}
\endhead
\bottomrule\noalign{}
\endlastfoot
S0 no shifts & M0 & 1 / 0.000 & 1 / 0.000 / 0.329 & 1 / 0.336 & 1 / 0.342 & 1 / 0.429 & 1 / 0.436 & \\
& M1 & 0.999 (0.001) / 0.000 & 1 / 0.000 / 0.468 & 1 / 0.704 & 1 / 0.667 & 1 / 0.571 & 1 / 0.808 & \textbf{0.090 (0.009)} (size; df 38) \\
S1 affine shifts & M0 & 1 / 0.000 & 1 / 0.000 / 0.335 & 1 / 0.344 & 1 / 0.351 & 1 / 0.435 & 1 / 0.444 & \\
& M1 & 0.938 (0.008) / 0.000 & 1 / 0.000 / 0.437 & 1 / 0.672 & 1 / 0.638 & 1 / 0.539 & 1 / 0.773 & 1.000 (0.000) \\
S2 common non-affine shock & M0 & \textbf{0.002 (0.001) / 0.987} & \textbf{0.009 (0.003) / 0.963} / 0.023 & \textbf{0.013 (0.004)}, empty 0.950 & \textbf{0.015 (0.004)}, empty 0.964 & \textbf{0.155 (0.011)}, empty 0.585 & \textbf{0.206 (0.013)}, empty 0.478 & \\
& M1 & 0.997 (0.002) / 0.002 & 1 / 0.000 / 0.328 & 1 / 0.509 & 1 / 0.541 & 1 / 0.429 & 1 / 0.611 & 1.000 (0.000) \\
\end{longtable}

\hypertarget{pE-reading-the-results}{%
\subsubsection{8.4 Reading the results}\label{pE-reading-the-results}}

\begin{enumerate}
\def\labelenumi{\arabic{enumi}.}
\tightlist
\item
  \textbf{No proved interval under-covers.}

  \begin{itemize}
  \tightlist
  \item
    Where its theorem applies, the projection interval covers at least 0.995 (design E, MC SE 0.002) and the conservative interval at least 0.990 (design E, SE 0.003).
  \item
    Theorem P3.8 covers in every replication of every design where it was computed, including R2.
  \item
    The only figures below 0.95 for a correctly specified model are the plug-in (not a CI) and CLR for \(\tilde B\) in design E.
  \end{itemize}
\item
  \textbf{CLR (a simulation finding) covers at least 0.994 in every design where Slater holds}, including R2 with 200 districts (0.995--0.999 for \(\tilde B\), 0.998--1 for \(B\)).

  \begin{itemize}
  \tightlist
  \item
    \textbf{At the Slater failure (design E), CLR for \(\tilde B\) covers 0.944 (SE 0.007).} This is within one standard error of the nominal 0.95, so it is not evidence of under-coverage.
  \item
    It is the informative case: the identified set is a point, so the interval has no slack, and coverage falls from near 1 to the nominal level. The 100\% elsewhere reflects truths in the interior of the identified set (§4).
  \item
    With the residual allowance (target \(B\)) CLR covers 0.998 in design E.
  \end{itemize}
\item
  \textbf{The plug-in fails where Theorem P3.4 predicts, and badly at the Slater failure.}

  \begin{itemize}
  \tightlist
  \item
    In R2 its coverage falls with the number of districts: 0.305, 0.042 and 0.005 at 30, 93 and 200 districts. It is empty 24\%, 46\% and 72\% of the time.
  \item
    In design E it covers 0.267 and is empty half the time: with a point identified set, any sampling noise empties the plug-in.
  \end{itemize}
\item
  \textbf{The projection interval is useful only when \(q\) is small.} Its critical value grows like \(\sqrt q\): 2.8 at \(q=3\), 4.6 at \(q=12\), 14.8 at \(q=187\) and 21.2 at \(q=401\), while CLR's is about 2--3.

  \begin{itemize}
  \tightlist
  \item
    Pooled designs (\(q=3\), 12): it is only 0.01--0.03 wider than CLR.
  \item
    M1 with 10 districts (\(q=21\)): 0.554 against CLR's 0.432 for \(\tilde B\).
  \item
    M1 with 93 districts: 0.93 at 30 polls per district, and \([0,1]\) at 10.
  \item
    So in grouped designs with many districts CLR is the only informative interval, and its support there is simulation only (§0.5).
  \end{itemize}
\item
  \textbf{Theorem P3.8 remains uninformative.} Its clipped width is 1.000 in every M1 design, including 10 districts × 93 polls, where CLR is 0.43 and DD is 0.57. This confirms §5.1 at 1000 replications.
\item
  \textbf{In R2, CLR narrows as districts are added.} For \(\tilde B\) its width is 0.557, 0.467 and 0.416 at 30, 93 and 200 districts; for \(B\), 0.746, 0.575 and 0.491. With 200 districts of 10 polls it is narrower than DD (0.566). The earlier statement that ``with 10 units per district the valid CI is no narrower than DD'' holds at 93 districts but not at 200.
\item
  \textbf{Riding shocks (R-37(i)).}

  \begin{itemize}
  \tightlist
  \item
    \textbf{Common non-affine shock (S2):} every M0 interval fails, for \(B\) and for \(\tilde B\). CLR for \(\tilde B\) covers 0.9\%, and is empty (ends crossing, so M0 is rejected) 96\% of the time. For \(B\), CLR covers 16\% and projection 21\% (the residual allowance widens them), and the conservative interval 4\%; again most failures are empty intervals. M1 covers in every replication, and the pooling test rejects in every replication.
  \item
    \textbf{Affine shifts (S1):} the pooling test rejects in every replication, because M0 is false for the observed design (P3.6(c)), but M0's intervals still cover. Shifts affine in composition change the pooled conditional mean only through terms the pooled model absorbs (2b Prop. D4(b)), approximately, so M0's bias is small here. A rejection of the pooling test therefore says that M0 is wrong, not that its intervals are.
  \item
    \textbf{Size of the pooling test (S0): 0.090 (SE 0.009) at nominal 0.05} (this Monte Carlo, 1000 replications, χ² calibration; the separate 2000-replication run in \path{benchmarks/pooling_test_size.R} gives 0.078 for χ², and 0.043 with the wild bootstrap now used by the package, P3.6 calibration table), with 38 restrictions and 50 polls per district. The HC3 Wald statistic has mean 39.7 and variance 100, against 38 and 76 for \(\chi^2_{38}\). The F form (\(W/\text{df}\) against \(F_{38,959}\)) gives 0.079. This is consistent with P3.6, whose \(\chi^2\) calibration is established for fixed dimension only, but the test over-rejects in a moderate-dimension R1 design. It errs towards M1, the safer rung. \emph{(Follow-up, 0.3.0: the restricted wild bootstrap rejects 0.043 (SE 0.005) here; P3.6, calibration paragraph.)}
  \end{itemize}
\item
  \textbf{The residual term is a large part of the width.} Targeting \(\tilde B\) instead of \(B\) narrows CLR by 0.06--0.10 in the 3×3 design and 0.06--0.11 in the 2×2 M1 designs with 30 or more polls per district; with 10 polls per district, by 0.08--0.19.
\item
  \textbf{Bhatia--Davis versus Popoviciu.} In the boundary design C, Bhatia--Davis narrows CLR for \(B\) from 0.086 to 0.057. In the Slater design E it narrows it from 0.221 to 0.207. Elsewhere the plausible range of rates includes 0.5, the allowance equals Popoviciu's, and the interval is up to 0.015 wider (the α share spent on the variance bound). It remains the package default.
\item
  \textbf{Real data} (NZ benchmark, \path{benchmarks/results/ei_datasets_nz_summary.md}; run with eiPartialID 0.2.0 during WS6, before the 0.3.x changes to the tests and the \texttt{proj\_actual} variance bound, which can move those columns slightly).

  \begin{itemize}
  \tightlist
  \item
    Per electorate, with about 64 polling places each: M0's CLR and projection intervals covered all 4302 cells, at about DD width.
  \item
    Pooling each election's 62--69 electorates and fitting M1 with electorates as groups: the M1 test (against M2) rejects M1 in all seven elections (\(p\approx0\)), and the pooling test rejects M0 against M1 (\(p\approx0\)). The plug-in M1 estimate is empty in all seven elections, but that is a sample diagnostic: the emptiness test does not reject M1 in any election (it rejects pooled M0 in 2002 and 2020 only).
  \item
    There, CLR intervals are rejected (ends crossing) in 21\% of cells and miss the truth in 8\%. The projection interval never misses, but it is as wide as DD.
  \item
    So where the model is wrong, the tests flag it, and a non-empty CLR interval is not reassurance.
  \end{itemize}
\end{enumerate}

\begin{center}\rule{0.5\linewidth}{0.5pt}\end{center}

\hypertarget{pE-practical-recipe-input-to-phase-4}{%
\subsection{9. Practical recipe (input to Phase 4)}\label{pE-practical-recipe-input-to-phase-4}}

\begin{enumerate}
\def\labelenumi{\arabic{enumi}.}
\tightlist
\item
  Fit the rung's regression (M0 / M1 / M2) on \(C-1\) columns with HC3 covariance (no clustering, §0.4), and the sample hulls.
\item
  Build the LP (2a Prop. 7 / 2b Thm 2) and report the plug-in bound, labelled as an estimate of the identified set, not a confidence interval.
\item
  Report the CLR confidence interval for \(\tilde B\) (default) and for \(B\), labelled algorithmic (Conjecture P3.5(b3); for \(B\) its Bhatia--Davis allowance is P3.5(c2)). Report the projection interval (P3.5(b2)(ii), with \(q=\dim\beta\) or an exact reduction) alongside it as the interval with an asymptotic guarantee; for \(B\) its Bhatia--Davis allowance must take \(\bar m\) from the projection interval on the rung alone (P3.5(c1)), or use Popoviciu. Intersect both with Duncan--Davis, which is valid because \(B\in DD\) always.
\item
  Theorem P3.8 gives a finite-sample guarantee in every regime. In the designs tried, its width is that of the logical range, so it is not yet worth reporting by default. With many small districts (R2), say that the CLR interval's coverage there rests on simulation (§0.5).
\item
  Report the emptiness test and, for M1, the test of its testable part (P3.6(b)) and the pooling test (P3.6(c)). The \(\chi^2\) version of the pooling test over-rejects with many restrictions (9\% at nominal 5\% with 38 restrictions, §8.4); since 0.3.0 the package calibrates both tests by a restricted wild bootstrap (4.3\% and 5.4\% in design S0; P3.6). A rejection says M0 is wrong for this design, not that its intervals fail (§8.3, S1).
\item
  Report the sensitivity curves in \(\eta\) (and \(\delta\)), with the breakdown value of the headline conclusion.
\item
  With few units per district, apply P3.7 selection, or prefer M1(\(\delta\)) or coarser groups, and say so.
\end{enumerate}

\hypertarget{pE-open-items}{%
\subsection{10. Open items}\label{pE-open-items}}

\begin{itemize}
\tightlist
\item
  \emph{Closed after review (Batch B, Nov 2026):} (1) P3.5(b2)(ii) now requires a positive definite \(\Sigma\) or a reduced full-rank parameterisation; a singular \(\hat V\) in the package's projection interval needs a guard (rank check, or \(q=\dim\beta\) with a warning). (2) The residual step of P3.5(a) is completed by the Chebyshev/Lindeberg subsequence argument. Both repairs leave every default and every allowance unchanged.
\item
  \emph{Added after review (round 2, Nov 2026):} uniform (in the design sequence) limit theory for the plug-in endpoint when candidates are within \(O(n^{-1/2})\) of binding without a convergent localisation (Proposition P3.4′ covers convergent localised slacks).
\item
  Conjecture P3.5(b3): validity of collecting dual vertices by simulation, with \(k_p\) computed over the collected set. A proved reduction to (b1) is given; recovery and the \(k_p\) effect are open.
\item
  Boundary cases where Slater fails (implied rates exactly 0 at a vertex): rate and limit of the endpoint estimator.
\item
  The rate of the hull error \(\varphi(\beta_0;\hat A)-\varphi(\beta_0;A)\), and inference for the population-hull endpoint.
\item
  Inference for the superpopulation target \(B^*\), which needs design-sampling variation.
\item
  Uniform validity near point identification for the CLR version (Stoye-type adaptation).
\item
  Asymptotic theory for M1/M2 in R3 (high-dimensional Gaussian approximation; §0.5) and in R2.
\item
  \textbf{Informative finite-sample inference in R2.} Theorem P3.8 is valid but uninformative in the designs tried. Candidates:

  \begin{itemize}
  \tightlist
  \item
    bound the endpoint directly rather than every constraint, e.g.~a finite-sample version of (b2) with a data-dependent dual;
  \item
    variance-adaptive bounds (empirical Bernstein);
  \item
    Bhatia--Davis;
  \item
    uneven \(\alpha\) splits.
  \end{itemize}
\item
  CLR when Slater fails: coverage fell to the nominal level (0.944, SE 0.007) in design E of §8, as expected without slack. Whether it holds uniformly near point identification is the Stoye-type question above.
\item
  Calibration of the pooling and M1 tests in moderate and high dimension. The HC3 Wald \(\chi^2\) test over-rejects (9\% at 5\% with 38 restrictions, §8.4). \emph{Partly addressed (0.3.0):} both tests now use a restricted wild bootstrap, which is close to nominal in design S0 (P3.6, calibration paragraph). Open: a proof for a growing number of restrictions.
\item
  Inference under R2 with a hierarchical model for the district shifts.
\item
  Dependence within districts beyond the district effect (spillovers; §0.4).
\item
  Joint relaxation for population turnover.
\item
  Coverage of the plug-in under misspecification, with the menu's dials.
\end{itemize}


\clearpage
\section{PHASE5a: the sign module (M-SIGN)}\label{app:phase5a}
\textit{Source: \texttt{docs/extension/PHASE5a\_sign\_module.md}.}\par
% GENERATED by tools/make_appendices.py from docs/extension/PHASE5a_sign_module.md
% Do not edit by hand; edit the source and regenerate.
\emph{Proof-register items R3, R2 and R10-union from \texttt{PACKAGE\_DESIGN.md} §8. Notation and justification tags follow 2a §0.3, 2c §0.3 and P3 §0.3. Each item ends with its numerical check. Status: R3 (§1), R2 (§2) and R10-union (§3) written and checked. An independent review pass (fresh session) found one error, in R10(d), and several gaps; all are fixed below and listed in §4. Two external referee reports (Batch B, Nov 2026: \texttt{reviews/B\_partial\_identification\_report.md}, \texttt{reviews/B\_ecological\_inference\_report.md}) found further errors and gaps, ledgered in \texttt{reviews/B\_PI\_RESPONSE\_PLAN.md} (rows P-B15 to P-B27, P-B36 to P-B38, P-B43; P-E6, P-E12, P-E17, P-E19, P-E21, P-E22, P-E28). Each correction below is marked ``Corrected after review, Batch B''. \textbf{Scope:} every statement here is for the \textbf{pooled} LCE model (M0). The sign and CR rows are valid linear rows under the district rungs too (2c M5(a), as corrected), but R3 and R10 for M1, M1(δ), M2 and L4 are proof item \textbf{R16}, not proved.}

\emph{Round 2 (Nov 2026): the second reports of the partial-identification and EI referees (\texttt{reviews/B\_partial\_identification\_report\_round2.md}, \texttt{reviews/B\_ecological\_inference\_report\_round2.md}) asked for the R10 scope to match the pooled scope, a narrower ``finite LP'' sentence, the pairwise-CR pooling, and the coverage figure of each plug-in procedure; these are marked ``round 2''. Checks are tagged as in 2c: \textbf{{[}F{]}} fails on an earlier statement, \textbf{{[}V{]}} verifies a corrected or new statement, \textbf{{[}I{]}} illustration.}

\emph{Extra tags used here:}

\begin{longtable}[]{@{}
  >{\raggedright\arraybackslash}p{(\columnwidth - 2\tabcolsep) * \real{0.5000}}
  >{\raggedright\arraybackslash}p{(\columnwidth - 2\tabcolsep) * \real{0.5000}}@{}}
\toprule\noalign{}
\begin{minipage}[b]{\linewidth}\raggedright
Tag
\end{minipage} & \begin{minipage}[b]{\linewidth}\raggedright
Meaning
\end{minipage} \\
\midrule\noalign{}
\endhead
\bottomrule\noalign{}
\endlastfoot
{[}Weier{]} & a continuous function on a nonempty compact set attains its maximum \\
{[}pert{]} & for an LP, the perturbation function \(v(\beta)=\max\{\varphi(k):k\in F,\ a^\top k\le\beta\}\) is concave and nondecreasing on its domain. Its superdifferential at \(\beta\) is the set of optimal dual multipliers of the row \(a^\top k\le\beta\), and for a concave function on \(\mathbb R\) it is \([v'_+(\beta),v'_-(\beta)]\) (Rockafellar 1970, Thms 23.2, 24.1, 29.1†; Bertsimas \& Tsitsiklis 1997, §5.2†) \\
\end{longtable}

\begin{center}\rule{0.5\linewidth}{0.5pt}\end{center}

\hypertarget{pF-r3-sharpness-of-lce-plus-restrictions-on-k}{%
\subsection{\texorpdfstring{1. R3: sharpness of LCE plus restrictions on \(k\)}{1. R3: sharpness of LCE plus restrictions on k}}\label{pF-r3-sharpness-of-lce-plus-restrictions-on-k}}

\hypertarget{pF-setting}{%
\subsubsection{1.1 Setting}\label{pF-setting}}

\begin{itemize}
\item
  The design \((N_i,X_i)_{i\le p}\) is fixed, and \(A\supseteq\operatorname{supp}P_X\) is a polytope (2a §0).
\item
  LCE rung, \textbf{pooled} (M0): \(H=Q(P)+K\) with \(K\in\mathcal K\), and \(k\in\mathbb R^d\) its free coordinates (2a Prop. 7). The district rungs are not covered (proof item R16; 2c M5).
\item
  \(\Theta_I=\Theta_I(P;A)\) is the sharp set (2a Thm 1). \(\Theta_O=\Theta_O(P;A)\) is the regression-only set that the LP computes (2a Prop. 7). Both are treated as subsets of \(\mathbb R^d\) through \(k\).
\item
  \textbf{Target.} A model-implied linear target \(\varphi(k)=\varphi_0+h^\top k\), for example \(\tilde B_{rc}\) or \(\tilde D=\tilde B_1-\tilde B_2\) (2a Prop. 7; 2c §0.2). (The letter \(\gamma\) is reserved for the variance ratio \(\gamma_r\) of (d).)

  \begin{itemize}
  \tightlist
  \item
    \(\tilde B\) uses the design weights \(N_i\). Fixing the design permits the auxiliary \textbf{convention} of giving its units equal probability, so that the finite-design target is represented as 2a's \(B^*\) for that empirical law. Conditioning alone does not equate this target with an arbitrary superpopulation rate. \emph{(Corrected after review, Batch B, Nov 2026: earlier worded as an identity.)} In the population, \(ER\) computed from \(T\) and from \(\tau=E[T\mid X]\) coincide by iterated expectations, so restrictions written with the \(\tau\)-based \(ER^*\) satisfy (C2).
  \item
    The actual-population rate \(B\) is \emph{not} covered: \(B-\tilde B\) is a weighted sum of unit residuals that the data never pin down, and it needs the P3 allowance.
  \end{itemize}
\item
  \textbf{Admissible restriction.} A set \(C=C(P)\subseteq\mathbb R^d\) such that

  \begin{itemize}
  \tightlist
  \item
    (C1) \(C\) is closed and convex;
  \item
    (C2) \(C\) depends on \(P\) only through the design and \(Q(P)\).
  \end{itemize}

  (C2) is stronger than (a) needs; (b) uses it, because there the restriction must be the same for every law with the same regression and design. (C1)--(C2) admit non-polyhedral sets (a ball, for instance), so sharpness alone does not make the computation a finite LP. \emph{(Corrected after review, Batch B, Nov 2026.)}

  The examples below are each a finite set of rows \(\{Gk\le g\}\), so for them the regression-only computation, optimisation over \(\Theta_O\cap C\), is a finite LP. The full-law set \(\Theta_I\cap C\) of R3(a) need not have a finite representation even for polyhedral \(C\). \emph{(Corrected after review, round 2, Nov 2026: earlier ``for them every computation is a finite LP''.)}

  \begin{itemize}
  \tightlist
  \item
    any sign on \(\delta^*_B(r,s)\), \(\delta^{r*}_B\), \(\Delta^*_r\) or \(\delta^{r*}_W\) (2c M5(a), pooled form);
  \item
    mean-level ordering on \(A\) (2c M5(b));
  \item
    the M-LOC rows \((x^\top H)_r-(x^\top H)_s\ge0\) on \(A\), which reduce to the vertices because they are linear in \(x\) (2a L3(a)).
  \end{itemize}
\item
  \textbf{Union restriction.} \(C=C_+\cup C_-\) with \(C_\pm\) admissible. Contextual reinforcement for a pair \((g_B,g_W)\) of affine functions of \(k\) is \(C_\pm=\{\pm g_B\ge0,\ \pm g_W\ge0\}\).
\end{itemize}

\hypertarget{pF-statement}{%
\subsubsection{1.2 Statement}\label{pF-statement}}

\textbf{Theorem R3.} Assume 2a (A1)--(A5), standing integrability, \(P(X\in A)=1\), and that \(\operatorname{affdim}(A)=R-1\). \emph{(Corrected in the review response: ``nonempty relative interior'' is too weak, see 2a Remark 5.2. Under A5 the condition holds automatically.)} The dimension condition is sufficient, not necessary: \(R\) affinely independent points of \(A\) bound every entry of \(K\) (2a L3(b)). \(P(X\in A)=1\) is used twice: the deterministic witness of (b) and the full target both use the region on which the mean is specified. With trimming (\(P(X\in A)<1\)) the outside-\(A\) terms need 2a Prop. 7′.

\begin{enumerate}
\def\labelenumi{(\alph{enumi})}
\item
  \textbf{Sharp interval.} For admissible \(C\), the sharp identified set for \(\varphi\) under LCE \(\cap\,C\) is
  \[\varphi(\Theta_I\cap C)=\Big[\min_{\Theta_I\cap C}\varphi,\ \max_{\Theta_I\cap C}\varphi\Big],\]
  or \(\emptyset\).
\item
  \textbf{What the LP delivers.} \(\Theta_O\cap C=\bigcup_{P'\in\mathcal P(P)}\big(\Theta_I(P')\cap C(P')\big)\), where \(\mathcal P(P)\) is as in 2a Prop. 2.
\end{enumerate}

\begin{itemize}
\tightlist
\item
  So the interval \(\varphi(\Theta_O\cap C)\) is sharp for the information in the regression \(E[T\mid X]\) and the design. It is a finite LP when \(C\) is polyhedral (all package restrictions are). \emph{(Corrected after review, Batch B, Nov 2026: earlier called ``the LP interval'' for every admissible \(C\).)}
\item
  It equals the sharp interval of (a) whenever \(\Theta_I=\Theta_O\) (2a Cor. 2, Prop. 3).
\item
  Otherwise it is a weakly outer interval, and M-SHARP can tighten it.
\item
  \textbf{Remark R3.1 (full law versus regression; added after review, Batch B).} Let \(X=1/4,1/2,3/4\) with equal probability and equal sizes, and \(T\) a fair binary total at every composition. The full law forces both row rates to equal \(T\) in every unit, so both row means are \(.5\) everywhere and \(k=0\). The regression alone (\(\tau=.5\)) admits \(m_1=.5+k(1-x)\), \(m_0=.5-kx\). For \(C=\{k\in[.2,.4]\}\) the full-law identified set (a) is \textbf{empty}, while the regression-only gap set (b) is \([1/6,1/3]\). This is exactly the distinction (a)--(b) draw. Example due to the referee; check S6.
\end{itemize}

\begin{enumerate}
\def\labelenumi{(\alph{enumi})}
\setcounter{enumi}{2}
\tightlist
\item
  \textbf{Unions.} For \(C=C_+\cup C_-\), the identified set is \(\varphi(\Theta\cap C_+)\cup\varphi(\Theta\cap C_-)\) for \(\Theta\in\{\Theta_I,\Theta_O\}\). That is at most two closed intervals.
\end{enumerate}

\begin{itemize}
\tightlist
\item
  If one piece is empty, the set is the other piece, a single interval, and the hull is sharp.
\item
  If both are non-empty, the hull \([\min,\max]\) is sharp exactly when the two intervals intersect.
\item
  \textbf{Certificate:} if \(\Theta\cap C_+\cap C_-\neq\emptyset\), they intersect. For \(\Theta_O\) with polyhedral \(C_\pm\) this is one feasibility LP. The certificate is sufficient only (§1.4).
\end{itemize}

\begin{enumerate}
\def\labelenumi{(\alph{enumi})}
\setcounter{enumi}{3}
\tightlist
\item
  \textbf{One-vs-rest CR is one LP.} Let \((g_B,g_W)=(\delta^{r*}_B,\delta^{r*}_W)\) (W1), and assume \(\operatorname{Var}_p(X_r)>0\). Put:
\end{enumerate}

\begin{itemize}
\tightlist
\item
  \(\gamma_r=\operatorname{Var}_p(X_r)/(\mu_r(1-\mu_r))\in(0,1]\);
\item
  \(ER^{r*}=\operatorname{Cov}_p(X_r,\tau(X))/\operatorname{Var}_p(X_r)\), with \(\tau(x)=(x^\top Q_cx)_c\);
\item
  \(D^{r*}=\tilde B_r-\tilde B_{\neg r}\).
\end{itemize}

Then
\[C_+\cup C_-=\big\{k:\ D^{r*}(k)\in\operatorname{hull}\{\gamma_rER^{r*},\ ER^{r*}\}\big\},\]
which is two linear rows in \(k\). For \(R=2\) this is the only CR, so in \(2\times2\) CR is a single LP. Under it, \(\operatorname{sign}\tilde D=\operatorname{sign}ER^*\). Ties (\(ER^{r*}=0\), or \(\gamma_r=1\)) can leave both branches feasible; the slab is unaffected. The slab can have an empty intersection with the LCE set (Remark R3.2).

\begin{enumerate}
\def\labelenumi{(\alph{enumi})}
\setcounter{enumi}{4}
\tightlist
\item
  \textbf{Refutation is joint.} If \(P\) is generated by a model in LCE \(\cap\,C\) with skew part \(K_0\), then \(K_0\in\Theta_I\cap C\subseteq\Theta_O\cap C\). An empty population LP therefore rejects the \textbf{conjunction} LCE \(\cap\,C\); emptiness alone does not attribute the failure to either part.
\end{enumerate}

\begin{itemize}
\item
  A belief that also has a model-free meaning can sometimes be tested on its own. Example: on \(x=1/4,1/2,3/4\) let both row means be \(1.5-2x\) and the aggregate equal this function. LCE holds on \(A=[1/4,3/4]\), \(ER=-2\), and every feasible gap is at least \(-1\), so a positive within association is impossible even in the model-free continuous table problem. For a restriction defined only on the LCE parameter \(k\), the meaning of \(C\) already presumes LCE.
\item
  \begin{enumerate}
  \def\labelenumi{(\alph{enumi})}
  \setcounter{enumi}{4}
  \tightlist
  \item
    is a population statement. Infeasibility of the LP at \(\hat\beta\) is a sample diagnostic; a level-controlled rejection uses R10(d)'s rule.
  \end{enumerate}
\end{itemize}

\emph{(Corrected after review, Batch B, Nov 2026: earlier ``never \(C\) alone''. The referee's examples show both directions: a belief can be refutable without LCE, and LCE can exclude a belief that is never refutable model-free (Remark R3.2).)}

\hypertarget{pF-proof}{%
\subsubsection{1.3 Proof}\label{pF-proof}}

\emph{(a)}
1. A value \(t\) is in the identified set iff some \(\mathbb Q\in\mathcal M(Q(P)+K;A)\) with \(\mathbb Q\circ W^{-1}=P\) has \(k\in C\) and \(\varphi(k)=t\). {[}Def identified set; Def \(\Theta_I\){]}
2. \(C\) is evaluated at the true \(P\) for every such \(\mathbb Q\), because it depends on \(\mathbb Q\) only through its \(W\)-marginal \(P\), via the design and \(Q(P)\). So the condition in 1 is \(k\in\Theta_I\cap C\). {[}C2; 1{]}
3. The identified set is therefore the image \(\varphi(\Theta_I\cap C)\). {[}2{]}
4. \(\Theta_I\) is convex and compact. \(C\) is closed and convex. So \(\Theta_I\cap C\) is convex and compact, or empty. {[}2a T1; C1{]}
5. \(\varphi\) is affine and continuous, so the image of a nonempty compact convex set is a compact interval with the stated endpoints. {[}lin; IVT{]} ∎(a)

\emph{(b)}
6. ⊇: for \(P'\in\mathcal P(P)\), \(Q(P')=Q(P)\) and \(P'_{(N,X)}=P_{(N,X)}\), so \(C(P')=C(P)\). Also \(\Theta_I(P')\subseteq\Theta_O(P')=\Theta_O(P)\). {[}Def \(\mathcal P\); C2; 2a P1, P2 step 2{]}
7. ⊆: take \(k\in\Theta_O\cap C\). The construction in 2a Prop. 2 (steps 4--10) gives \(P'\in\mathcal P(P)\) with \(k\in\Theta_I(P')\). Then \(k\in C(P)=C(P')\). {[}2a P2; 6{]}
8. The statements about the LP interval follow from (a) applied to each \(P'\), and from \(\Theta_I=\Theta_O\) in the cases cited. {[}7; 2a Cor. 2, Prop. 3{]} ∎(b)

\emph{(c)}
9. \(\varphi(\Theta\cap(C_+\cup C_-))=\varphi(\Theta\cap C_+)\cup\varphi(\Theta\cap C_-)\). Each piece is a closed interval or empty, by (a), or by (b) for \(\Theta_O\). {[}image of a union; (a), (b){]}
10. If one piece is empty the union is the other piece. A union of two non-empty closed intervals is an interval iff they intersect. A point \(k_0\in\Theta\cap C_+\cap C_-\) gives \(\varphi(k_0)\) in both. {[}order on \(\mathbb R\){]} ∎(c)

\emph{(d)}
11. Apply 2c Lemma M1(c) with \(m(X_i)=M_H(X_i)\) in place of \(\beta_i\), and \(\tau(X_i)\) in place of \(T_i\). The identities there use only \(T=\sum_sX_s\beta_s\) and the simplex, and \(x^\top M_H(x)=\tau(x)\) holds. This gives
\[\delta^{r*}_W=\operatorname{Var}_p(X_r)\,(ER^{r*}-D^{r*}),\qquad \delta^{r*}_B=\mu_r(1-\mu_r)\,(D^{r*}-NM^{r*}),\]
where \(NM^{r*}\) is the neighbourhood gap computed from \(\tau\). {[}2c L M1(c); 2a L1(a); alg{]}
12. \(NM^{r*}=\operatorname{Cov}_p(X_r,\tau)/(\mu_r(1-\mu_r))=\gamma_rER^{r*}\). {[}Def; alg; 2c M3(d) for the collapsed population{]}
13. Both \(\ge0\) ⇔ \(\gamma_rER^{r*}\le D^{r*}\le ER^{r*}\). Both \(\le0\) ⇔ \(ER^{r*}\le D^{r*}\le\gamma_rER^{r*}\). The union of these is \(D^{r*}\in\operatorname{hull}\{\gamma_rER^{r*},ER^{r*}\}\). {[}11--12; positive factors{]}
14. \(D^{r*}\) is affine in \(k\), and \(ER^{r*}\) and \(\gamma_r\) depend only on the design and \(Q\). So the set is two rows of an admissible \(C\). The sign statement follows because \(\gamma_r>0\). {[}2c M5(a); C2; 13{]} ∎(d)

\emph{(e)}
15. \(K_0\in\Theta_I\) by definition, and \(k_0\in C\) by hypothesis. Then apply 2a P1. {[}Def; 2a P1{]} ∎

\hypertarget{pF-what-fails-and-where}{%
\subsubsection{1.4 What fails, and where}\label{pF-what-fails-and-where}}

\begin{itemize}
\tightlist
\item
  \textbf{Pairings other than W1.} For CR on \((\delta^{r*}_B,\Delta^*_r)\) or \((\delta^*_B(r,s),\Delta^*_r)\), (d) fails. The two sign conditions are on \emph{different} affine functions of \(k\), so both branches can be non-empty and their images can be disjoint (S5). For the pairwise version only \(\delta^*_B(r,s)\) is pairwise; \(\Delta^*_r\) pools every original group, including those outside \(\{r,s\}\). So an argument about \(r\) against \(s\) alone supports the between term, not the contextual term. \emph{(Added after review, round 2, EI referee.)}
\item
  \textbf{Not covered:}

  \begin{itemize}
  \tightlist
  \item
    N2 shape restrictions, which do not restrict \(k\);
  \item
    N1 signs on the realised \(\delta\), which are unit-level, not mean-level;
  \item
    any target that mixes \(\tilde B\) with \(B\) (marked ⚠ in the matrix).
  \end{itemize}
\item
  \textbf{The common-witness certificate is only sufficient} (example added after review, Batch B). In the square \([-1,1]^2\) with \(g_B=u+.5\) and \(g_W=u-.5\), the two branches are the disjoint strips \(u\ge.5\) and \(u\le-.5\). For the target \(u\) the images are \([.5,1]\) and \([-1,-.5]\), a gap of length one. For the target \(v\) both images are \([-1,1]\), identical, although the certificate LP \(\{g_B=g_W=0\}\) is infeasible.
\item
  \textbf{Remark R3.2 (LCE can exclude CR although LCE fits; added after review, Batch B).} Let \(X=0,1/4,1\) with equal probability, \(A=[0,1]\), \(m_0=1-x\), \(m_1=(49/55)x\), and deterministic \(T\). The regression is \((1-x)^2+(49/55)x^2\). Its Bernstein cross coefficient is 0, so the vertex rows force \(k=0\): LCE fits, with a single point. There \(ER=0\) and \(\tilde D=-261/1925\). CR needs \(\tilde D\in\operatorname{hull}\{\gamma ER,ER\}=\{0\}\), so \textbf{both LCE branches are empty}, although LCE alone is not rejected. In the model-free continuous table problem (2c M3) CR is still feasible (\(\beta_{i1}=\beta_{i2}=T_i\)). Any ``CR cannot be refuted'' statement must name its base system. Example due to the referee; check S7.
\item
  \textbf{Degenerate hull.} If \(\operatorname{affdim}(A)<R-1\), \(\Theta_O\) can be unbounded (2a L3(b), Remark 5.2). A singleton \(A\) has nonempty relative interior and is still degenerate; under A5 this case cannot arise. For polyhedral \(C\), (a) and (b) then hold with a possibly infinite endpoint, and the LP reports unboundedness. For non-polyhedral closed convex \(C\) the image of an unbounded closed convex set need not be closed, so an endpoint may be a supremum that is not attained. Added rows can restore a finite endpoint (V5, near-1d support).
\end{itemize}

\hypertarget{pF-numerical-check-verify_r3.py-output-verify_r3_output.txt}{%
\subsubsection{\texorpdfstring{1.5 Numerical check (\texttt{verify\_R3.py}, output \texttt{verify\_R3\_output.txt})}{1.5 Numerical check (verify\_R3.py, output verify\_R3\_output.txt)}}\label{pF-numerical-check-verify_r3.py-output-verify_r3_output.txt}}

\begin{longtable}[]{@{}
  >{\raggedright\arraybackslash}p{(\columnwidth - 2\tabcolsep) * \real{0.5000}}
  >{\raggedright\arraybackslash}p{(\columnwidth - 2\tabcolsep) * \real{0.5000}}@{}}
\toprule\noalign{}
\begin{minipage}[b]{\linewidth}\raggedright
Check
\end{minipage} & \begin{minipage}[b]{\linewidth}\raggedright
Result
\end{minipage} \\
\midrule\noalign{}
\endhead
\bottomrule\noalign{}
\endlastfoot
S1: constructive sharpness \emph{for the LP set}, i.e.~part (b) (it uses the 2a Prop. 2 construction with \(x^\top\beta=\tau\)). 3×2, LCE with the true signs of \(\delta^{0*}_B\) and \(\delta^{0*}_W\). For 1000 random \(t\in[L,U]\), solve for \(k\), build \(\beta_i=M_H(X_i)\), and check simplex, accounting, rows and \(\varphi=t\) & worst violation \(1.8\times10^{-15}\) \\
S2: 2×2 CR branches (2000 designs). \emph{Relabelled after review (Batch B): each branch is classified as empty, a point, or a positive-width interval; the old code called ``no positive-width branch'' ``both degenerate'', which also admitted two empty branches} & see output: no case with two positive-width branches and none with two empty branches. Frequencies on these draws are descriptive (the design generator is reused) \\
S3: 3×2 W1 CR, row 0 (1000 designs) & both branches non-empty 0 times, as (d) predicts \\
S4: 3×2 CR with the W2 or pairwise pairings & both non-empty in 357/1000 and 614/1000; no gaps \\
S6 (added after review, Batch B): full law versus regression, Remark R3.1 & regression gap set \([1/6,1/3]\); full-law set empty \\
S7 (added after review, Batch B): LCE excludes CR although LCE fits, Remark R3.2. \emph{Fails on the old §3.5 ``unrefutable'' as applied there} & LCE: \(k=0\); \(ER=0\); \(\tilde D=-261/1925\); both branches empty \\
S5: adversarial search (3 bases × 2 pairings × 3 targets; 3600 problems) & both non-empty 951; certificate held 945; \textbf{gaps 3}, all where the certificate failed; gap ≈ 2\% of hull width on this seed. The reviewer's rerun on a fresh seed found 14 gaps, up to ≈ 31\% of hull width, each confirmed by checking that both branch LPs are infeasible at the gap midpoint \\
\end{longtable}

\hypertarget{pF-consequences-for-the-package}{%
\subsubsection{1.6 Consequences for the package}\label{pF-consequences-for-the-package}}

\begin{enumerate}
\def\labelenumi{\arabic{enumi}.}
\tightlist
\item
  For the \textbf{pooled} linear rung (M0) plus any sign, ordering or local rows, \texttt{eiBounds()} may label the LP interval \textbf{sharp for the regression}. It is the sharp identified set when \(\Theta_I=\Theta_O\), and outer otherwise. Under district rungs (M1, M1(δ), M2, L4) the rows are correct but sharpness is not proved (R16): the label should read ``outer for the regression'', and the note should say the sign is one pooled restriction across groups. \emph{(Corrected after review, Batch B, Nov 2026.)}
\item
  One-vs-rest CR (the default, W1) and \(2\times2\) CR are implemented as two rows in \(k\), so identification needs no union. Inference is not quite ordinary: the rows' constants are a min/max of estimated quantities. The conservative CI inflates each piece before the max (R10(b)); the CLR CI takes the union of the two branch CIs (R10(c)). With the package's simulated candidates, the CLR version is algorithmic (Conjecture P3.5(b3)).
\item
  For other CR pairings the package:

  \begin{itemize}
  \tightlist
  \item
    runs the certificate LP. The certificate is sufficient, not necessary, for the two branch images to meet. In S5, 3 cases without a certificate had no gap, and every gap had a failing certificate. R3(c) allows a gap; the certificate is not a hypothesis;
  \item
    if the certificate fails, solves both branches and reports both intervals when they are disjoint.
  \end{itemize}
\item
  Emptiness messages always name the \emph{combination} (LCE + restriction) that has been rejected, say whether the rung alone is rejected, and distinguish three events: an empty plug-in set at \(\hat\beta\) (a sample diagnostic), an empty population set (R3(e)), and a rejection by the calibrated test (R10(d)). A statement that a belief ``cannot be refuted'' must name its base system (Remark R3.2). \emph{(Corrected after review, Batch B, Nov 2026.)}
\end{enumerate}

\begin{center}\rule{0.5\linewidth}{0.5pt}\end{center}

\hypertarget{pF-r2-when-does-an-added-restriction-tighten-the-bound}{%
\subsection{2. R2: when does an added restriction tighten the bound?}\label{pF-r2-when-does-an-added-restriction-tighten-the-bound}}

\hypertarget{pF-setting-1}{%
\subsubsection{2.1 Setting}\label{pF-setting-1}}

\begin{itemize}
\tightlist
\item
  \(F=\{k:Gk\le g\}\) is the current feasible set: the rung, plus any rows already imposed. It is compact by 2a L3(b) when \(\operatorname{affdim}(A)=R-1\), which A5 implies (2a Remark 5.2). R2's proof uses this compactness (step 3). Nonemptiness is \emph{assumed}: R3 does not deliver it, and an empty \(F\) is a rejection (R3(e)).
\item
  \(\varphi(k)=\varphi_0+h^\top k\) is a linear target, with \(U=\max_F\varphi\), \(L=\min_F\varphi\), and optimal faces \(F^*_U=\arg\max_F\varphi\) and \(F^*_L=\arg\min_F\varphi\).
\item
  \(N\) is the added restriction: an admissible set as in §1.1, closed and convex, with \(F\cap N\ne\emptyset\). \(U_N\) and \(L_N\) are the endpoints after adding it.
\item
  Three possible verdicts for \(N\) relative to \(F\):

  \begin{itemize}
  \tightlist
  \item
    \textbf{implied:} \(F\subseteq N\);
  \item
    \textbf{cuts the set only:} \(F\not\subseteq N\), but \([L_N,U_N]=[L,U]\) for the target in question;
  \item
    \textbf{tightens the target:} \([L_N,U_N]\subsetneq[L,U]\).
  \end{itemize}
\end{itemize}

\hypertarget{pF-statement-1}{%
\subsubsection{2.2 Statement}\label{pF-statement-1}}

\textbf{Proposition R2.}

\begin{enumerate}
\def\labelenumi{(\alph{enumi})}
\item
  \textbf{Face criterion.} \(U_N<U\) iff \(F^*_U\cap N=\emptyset\), and \(L_N>L\) iff \(F^*_L\cap N=\emptyset\) (given attainment on both nonempty sets, which compactness supplies). Each is one feasibility LP when \(N\) is polyhedral; the LP must carry the target intercept (step 4). In words: the restriction moves an endpoint iff it removes \emph{every} point that attains that endpoint.
\item
  \textbf{One row \(a^\top k\le b\).} \(U_N<U\) iff \(\min_{F^*_U}a^\top k>b\). This is also equivalent to: every optimal dual multiplier of the new row is \(>0\) (that is, \(v'_+(b)>0\)).
\end{enumerate}

\begin{itemize}
\tightlist
\item
  A positive multiplier returned by the solver at \emph{one} optimum does \textbf{not} certify tightening. Example (added after review, Batch B): maximise \(x\) on \([0,1]\) with the row \(x\le1\) duplicated; every split \((1-\lambda,\lambda)\) of the unit dual cost is optimal, so the new row can carry multiplier 1 without moving the endpoint. A positive multiplier certifies tightening only if its uniqueness, or a positive lower bound on all optimal multipliers, has been established separately.
\item
  A zero multiplier, or a row that is slack at an optimum, certifies that the row does not tighten.
\item
  ``Active'' (the row holds with equality at the optimum) does \textbf{not} mean ``tightening''.
\end{itemize}

\begin{enumerate}
\def\labelenumi{(\alph{enumi})}
\setcounter{enumi}{2}
\item
  \textbf{Implied = useless for every target.} \(N\) is implied iff it tightens no linear target. For general closed convex \(N\) the witness comes from a hyperplane separating a point of \(F\setminus N\) from \(N\); for polyhedral \(N\), any non-implied row of \(N\) serves. For a row \(a^\top k\le b\), it is implied iff \(\max_Fa^\top k\le b\), and when it is not implied the target \(\varphi=a^\top k\) is a witness that tightens. So the ``already implied'' table answers ``can this assumption buy anything at all?''. Whether it buys something \emph{for your target} is a separate question, answered by (a). The witness objective quantifies over all linear functions of \(k\); it need not be a reported cell-rate target.
\item
  \textbf{\(2\times2\) under LCE.} Assume \(\operatorname{Var}_p(X)>0\) and \(\gamma<1\). Then \(k=w_1\) is scalar and \(\tilde D\) is strictly monotone in \(w_1\), with slope \(-(1-\gamma)\) times a positive factor.
\end{enumerate}

\begin{itemize}
\tightlist
\item
  Each of the four sign rows is either implied or tightens \(\tilde D\), and any target that is non-constant in \(w_1\); for such targets the middle verdict cannot occur. A target constant in \(w_1\) is never tightened. \emph{(Corrected after review, Batch B, Nov 2026: the qualifications were implicit.)}
\item
  The closed forms are those of G2:

  \begin{itemize}
  \tightlist
  \item
    \(\delta^*_W\ge0\): \([L_J,\min(U_J,ER^*)]\); implied iff \(ER^*\ge U_J\); empty iff \(ER^*<L_J\).
  \item
    \(\delta^*_W\le0\): \([\max(L_J,ER^*),U_J]\).
  \item
    \(\delta^*_B\ge0\): \([\max(L_J,NM^*),U_J]\).
  \item
    \(\delta^*_B\le0\): \([L_J,\min(U_J,NM^*)]\).
  \end{itemize}

  Here \(ER^*\) and \(NM^*\) are computed from \(\tau\) (§1.2(d)), and \([L_J,U_J]\) is the JKST interval for \(\tilde D\). The formulas are intersections of \([L_J,U_J]\) with a half-line; reversed endpoints mean the set is empty.
\end{itemize}

\begin{enumerate}
\def\labelenumi{(\alph{enumi})}
\setcounter{enumi}{4}
\tightlist
\item
  \textbf{Several rows.} For \(N=\bigcap_jN_j\), (a) applies to \(N\) as a whole. A row's contribution depends on what is already imposed.
\end{enumerate}

\begin{itemize}
\tightlist
\item
  The package reports, for each row:

  \begin{itemize}
  \tightlist
  \item
    its gain when added first (to the rung alone);
  \item
    its gain when added last (dropped from the full stack).
  \end{itemize}
\item
  A row whose last-added gain is zero has \textbf{zero incremental gain for this target}, given the others, even if it tightens on its own. It need not be implied by the others, and it may still tighten other targets. Example: on \([0,1]^2\) with target \(x\) and rows \(x+y\le.8\), \(x\le.5\), adding \(x+y\le.8\) last gives zero gain for \(x\), yet it still removes \((.5,1)\) and lowers the upper endpoint of \(x+y\) from 1.5 to .8 (check T4b). A row is redundant \textbf{for the feasible set} only if the implication test (c) says so. \emph{(Corrected after review, Batch B, Nov 2026: earlier ``redundant given the others'', which reads as ``implied''.)}
\end{itemize}

\hypertarget{pF-proof-1}{%
\subsubsection{2.3 Proof}\label{pF-proof-1}}

\emph{(a)}
1. \(F\cap N\subseteq F\), so \(U_N\le U\). {[}nest{]}
2. If \(k\in F^*_U\cap N\), then \(U_N\ge\varphi(k)=U\). So \(U_N=U\). {[}Def; 1{]}
3. If \(F^*_U\cap N=\emptyset\): \(F\cap N\) is nonempty and compact, so \(\varphi\) attains \(U_N\) at some \(k'\in F\cap N\). Since \(k'\notin F^*_U\), \(\varphi(k')<U\). {[}C1; compact; Weier; Def \(F^*_U\){]}
4. The lower endpoint is the same argument applied to \(-\varphi\). If \(N=\{Mk\le m\}\), then \(F^*_U\cap N=\{Gk\le g,\ h^\top k\ge U-\varphi_0,\ Mk\le m\}\) is a feasibility LP. {[}sym; Def{]} ∎(a) \emph{(Corrected after review, Batch B, Nov 2026: the row was printed as \(h^\top k\ge U\), omitting the target intercept \(\varphi_0\). On \([0,1]^2\) with \(\varphi=2+x\) that row is \(x\ge3\), which rejects the whole optimal face and so reports ``moves'' for every added row. The package's face test (\texttt{.rxc\_face\_moves}) subtracts the intercept and was not affected; \texttt{verify\_R2.py} used zero intercepts and could not see the error. Check T1c.)}

\emph{(b)}
5. \(F^*_U\) is a nonempty compact face of \(F\), so \(\beta^*:=\min_{F^*_U}a^\top k\) is attained, and \(F^*_U\cap\{a^\top k\le b\}=\emptyset\) iff \(\beta^*>b\). The first claim follows from (a). {[}Weier; (a){]}
6. Let \(v(\beta)\) be the perturbation function. Then \(v(\beta)=U\) for \(\beta\ge\beta^*\), by (a) applied at \(\beta\), and \(v(\beta)<U\) for \(\beta<\beta^*\). {[}(a); Def \(v\){]}
7. If \(b\ge\beta^*\), \(v\) is constant on \([b,\infty)\), so \(v'_+(b)=0\) and \(0\) is an optimal multiplier. {[}6; pert{]}
8. If \(b<\beta^*\): by concavity, \(v'_+(b)\ge(v(\beta^*)-v(b))/(\beta^*-b)>0\). So every optimal multiplier, which lies in \([v'_+(b),v'_-(b)]\), is positive. At the left end of the domain, \(b=\min_Fa^\top k\), the multiplier set is \([v'_+(b),\infty)\) (\(v'_-=+\infty\)); the equivalence is unaffected. {[}pert; 6{]}
9. A solver returns \emph{one} multiplier from \(\partial v(b)\), and when \(\partial v(b)=[0,\lambda]\) with \(\lambda>0\) it may return a positive one even though the row does not tighten. This is the degenerate case of a row through an optimal vertex: it is active, but it removes nothing at the optimum (T1). Suppose the row is slack, \(a^\top k<b\), at an optimum \(k\) of the \emph{augmented} LP. Near \(k\) the augmented set coincides with \(F\), so \(k\) is a local maximiser of the linear \(\varphi\) on the convex \(F\), hence a global one. Then \(k\in F^*_U\cap N\), and by 2 the row does not tighten. {[}pert; local = global for concave maximisation; 2{]} ∎(b)

\emph{(c)}
10. If \(F\subseteq N\), then \(F\cap N=F\), and no target changes. {[}Def{]}
11. If \(a^\top k\le b\) is not implied, then \(\max_Fa^\top k>b\). For \(\varphi=a^\top k\), \(U_N\le b<U\). {[}Def; 1{]}
11′. General closed convex \(N\) with \(F\not\subseteq N\): take \(k_1\in F\setminus N\). Strict separation of the point \(k_1\) from the closed convex set \(N\) gives \((a,b)\) with \(a^\top k\le b\) on \(N\) and \(a^\top k_1>b\). For \(\varphi=a^\top k\): \(U_N\le b<a^\top k_1\le U\). {[}sep. hyperplane; 11{]} ∎(c)

\emph{(d)}
12. By 2c M3 / §1.2(d) with \(R=2\): \(\delta^*_W=\operatorname{Var}_p(X)(ER^*-\tilde D)\) and \(\delta^*_B=\mu(1-\mu)(\tilde D-NM^*)\), with \(\tilde D\) affine in \(w_1\) with slope proportional to \(-(1-\gamma)\), so non-constant iff \(\gamma<1\). {[}§1.3 step 11; G2 step 3{]}
13. Each sign row is therefore a half-line in \(\tilde D\), and \(F\) maps to \([L_J,U_J]\). Intersecting gives the closed forms. {[}12; R3(b) and 2a §6.1, since \(F\) is the LP set{]}
14. Since \(\tilde D\) is a bijection between \(F\) and \([L_J,U_J]\), cutting \(F\) means cutting \([L_J,U_J]\). So ``cuts the set only'' is impossible for \(\tilde D\), and also for any target that is non-constant in \(w_1\). {[}12; bij{]} ∎(d)

\emph{(e)} This is (a) applied to \(N\) and to each sub-stack. Order dependence follows from (a), since \(F^*_U\) changes as rows are added. ∎

\hypertarget{pF-numerical-check-verify_r2.py-output-verify_r2_output.txt}{%
\subsubsection{\texorpdfstring{2.4 Numerical check (\texttt{verify\_R2.py}, output \texttt{verify\_R2\_output.txt})}{2.4 Numerical check (verify\_R2.py, output verify\_R2\_output.txt)}}\label{pF-numerical-check-verify_r2.py-output-verify_r2_output.txt}}

\begin{longtable}[]{@{}
  >{\raggedright\arraybackslash}p{(\columnwidth - 2\tabcolsep) * \real{0.5000}}
  >{\raggedright\arraybackslash}p{(\columnwidth - 2\tabcolsep) * \real{0.5000}}@{}}
\toprule\noalign{}
\begin{minipage}[b]{\linewidth}\raggedright
Check
\end{minipage} & \begin{minipage}[b]{\linewidth}\raggedright
Result
\end{minipage} \\
\midrule\noalign{}
\endhead
\bottomrule\noalign{}
\endlastfoot
T1: face test vs brute force. 3×2 LCE base, random targets; 400 random rows and 400 rows through the solver's optimal vertex & agree 800/800. Of the 400 active-but-redundant rows, the solver reported a \textbf{positive dual in 133} \\
T1b (added after review; random targets give a vertex face almost surely, so T1 never tests a larger face): target normal = a facet normal, so the optimal face is that facet; added rows cut through the facet's interior, cut all of it, or miss it & face test = brute force in every case (see output) \\
T2: implied ⇔ the witness target \(a^\top k\) does not tighten (a sanity check, close to a tautology since \(U_N=\min(U,b)\) for this target) & 333/333 \\
T3: 2×2, four sign rows, 2000 designs & closed forms = LP (error \(1.7\times10^{-14}\)); not implied ⇔ tightens \(\tilde D\) in 7409/7409; ``cuts set only'' 0 \\
T4: 3×2, true-sign rows on \(\delta^{0*}_B\) and \(\delta^{0*}_W\), targets \(\tilde B_0,\tilde B_1,\tilde B_2\) (1800 = 300 designs × 2 rows × 3 targets) & 147 implied (8\%); \textbf{494 cuts set but not this target (28\%)}; 1153 tightens (64\%); 6 not classified, solver status recorded (base solve status 4, numerical difficulty). \emph{Corrected after review, Batch B: unsuccessful base solves were skipped silently; the status is now printed. The referee's rerun on other library versions classified 153/494/1153 of 1800} \\
T1c (added after review, Batch B): target with intercept, \(\varphi=2+x\) on \([0,1]^2\); rows \(x+y\le1.2\) (keeps part of the optimal facet) and \(x+y\le.8\) (removes it). \emph{Fails on the old step 4} & old face LP (\(h^\top k\ge U\)) reports ``moves'' for both rows, wrong for the first; corrected LP (\(h^\top k\ge U-\varphi_0\)) agrees with brute force for both \\
T4b (added after review, Batch B): square, target \(x\), rows \(x+y\le.8\), \(x\le.5\). \emph{Fails on the old (e) sentence read as ``implied''} & added-last gain 0 for \(x\), 0.7 for \(x+y\); not implied by the other row \\
\end{longtable}

\hypertarget{pF-consequences-for-the-package-1}{%
\subsubsection{2.5 Consequences for the package}\label{pF-consequences-for-the-package-1}}

\begin{enumerate}
\def\labelenumi{\arabic{enumi}.}
\tightlist
\item
  \textbf{The redundancy diagnostic uses the face test, not solver duals.} For each added module and each reported target, the package runs one feasibility LP per endpoint. Dual values are shown only as supplementary sensitivity information.
\item
  \textbf{Three verdicts per module, reported separately:}

  \begin{itemize}
  \tightlist
  \item
    ``already implied by the data and the rung'' (target-free; the ``already implied'' table);
  \item
    ``narrows the set but not this quantity'';
  \item
    ``narrows this quantity by X''.
  \end{itemize}

  In \(2\times2\) the middle verdict cannot occur for targets that are non-constant in \(w_1\) (R2(d)); with three or more groups it is common (T4: 28\%). Each verdict is relative to a nonempty convex base, a feasible added restriction, one target and one order (first or last), and the output should name all four.
\item
  \textbf{Added-first and added-last gains} are reported for stacks, so users can see which beliefs are doing the work. A zero added-last gain means ``no gain for this target'', not ``implied'' (R2(e)).
\item
  \textbf{Statistical caveat (feeds R10).} In samples, ``implied'' and ``tightens'' are estimated sharp-edged verdicts. Near a tie (e.g.~\(ER^*\approx U_J\)), report the width change with uncertainty rather than a binary verdict. Inference on a width gain needs \textbf{joint} treatment of the base and augmented endpoints (and of any estimated threshold, such as \(ER^*\)); an interval for either endpoint alone is not enough. A one-sided test of \(\max_Fa^\top k\le b\) is an LP-value test of the kind P3.5 covers. \emph{(Corrected after review, Batch B, Nov 2026: earlier ``the width change with its confidence interval'', without the joint requirement.)}
\end{enumerate}

\begin{center}\rule{0.5\linewidth}{0.5pt}\end{center}

\hypertarget{pF-r10-union-inference-with-added-rows-eitheror-restrictions-and-near-empty-sets}{%
\subsection{3. R10-union: inference with added rows, either/or restrictions, and near-empty sets}\label{pF-r10-union-inference-with-added-rows-eitheror-restrictions-and-near-empty-sets}}

\hypertarget{pF-setting-2}{%
\subsubsection{3.1 Setting}\label{pF-setting-2}}

P3 §0, with the regime table of P3 §0.5: the asymptotic results are theorems for M0 in every regime and for M1/M2 with a fixed number of districts (R1); M1/M2 with growing districts (R3) is a conjecture \textbf{before} any question of candidate collection; R2 has only the finite-sample interval P3.8 and simulation findings. Everything below is for the pooled M0 model unless stated (R16 for the district rungs). \emph{(Corrected after review, Batch B, Nov 2026: earlier ``regimes R1/R3'', which does not settle increasing dimension.)} \(\hat\beta\) is the regression coefficient vector with \(\sqrt n(\hat\beta-\beta_0)\Rightarrow N(0,\Sigma)\) and \(\hat\Sigma\) consistent. The design is conditioned on (P3-S2).
* Base rows: \(s_j(\beta)+a_j^\top k\ge0\) (P3 §1 step 2).
* Target: \(\varphi=c(\beta)+h^\top k\).
* \textbf{Linear-in-\(\beta\) added row:} \(s_j(\beta)+a_j^\top k\ge0\), with \(s_j\) linear in \(\beta\) and \(a_j\) fixed by the design. Every sign, ordering and local row in §1.1 has this form. The reason: \(\delta^*\), \(\tilde B\), \(ER^*\) and the orderings are linear in \(m=M_{Q+K}(X_i)\) with design weights, and \(Q\) is linear in \(\beta\) {[}2c M5 step 2; P3 §1{]}.
* \textbf{Parallel disjunction:} a row of the form \(\max_{l\le L}s_l(\beta)+a^\top k\ge0\), equivalently \(\bigcup_l\{s_l(\beta)+a^\top k\ge0\}\). The pieces share the same \(k\)-normal \(a\).
* \textbf{General union:} \(C_+\cup C_-\), each branch a finite set of linear-in-\(\beta\) rows.
* \(\mathrm{CI}(\mathcal R)\) denotes the P3.5 interval (version (a) or (b)) computed with the base rows plus the row set \(\mathcal R\), at total level \(\alpha=\alpha_U+\alpha_L(+\alpha_E)\). \textbf{Emptiness convention}, one rule per version: for P3.5(a), \(\mathrm{CI}=\emptyset\) iff the inflated set \(\hat\Theta^+\) is empty; for P3.5(b), \(\mathrm{CI}=\emptyset\) iff the computed ends cross, \(\hat L>\hat U\). (Applying both rules to (b) and combining them by a union bound would double the \textbf{bound} in (d); the actual rejection probability need not double. \emph{Reworded after review, Batch B.})

\hypertarget{pF-statement-2}{%
\subsubsection{3.2 Statement}\label{pF-statement-2}}

\textbf{Theorem R10.} Assume P3.5's conditions \emph{for the augmented row system} (base rows plus added rows), and that the restrictions hold at the truth, i.e.~\(k_0\) satisfies them. \textbf{Pooled M0.} The fixed-district baseline intervals of P3 (M1/M2 in regime R1, P3 §0.5) are background: they cover the rung alone. The added-row, slab and union statements below are \textbf{not} claimed for district rungs; that is proof item R16. \emph{(Corrected after review, round 2, Nov 2026: the statement earlier also named ``M1/M2 in regime R1'', contradicting the pooled scope of §1.1 and the header.)}

\textbf{Status by interval version} \emph{(replaces the earlier status note; corrected after review, Batch B, Nov 2026, following the referee's version table)}:

\begin{longtable}[]{@{}
  >{\raggedright\arraybackslash}p{(\columnwidth - 4\tabcolsep) * \real{0.3333}}
  >{\raggedright\arraybackslash}p{(\columnwidth - 4\tabcolsep) * \real{0.3333}}
  >{\raggedright\arraybackslash}p{(\columnwidth - 4\tabcolsep) * \real{0.3333}}@{}}
\toprule\noalign{}
\begin{minipage}[b]{\linewidth}\raggedright
Interval version (P3.5)
\end{minipage} & \begin{minipage}[b]{\linewidth}\raggedright
What R10 claims
\end{minipage} & \begin{minipage}[b]{\linewidth}\raggedright
Status
\end{minipage} \\
\midrule\noalign{}
\endhead
\bottomrule\noalign{}
\endlastfoot
(a) simultaneous row inflation (conservative) & (a)--(d), including the direct parallel slab (b), under a valid joint row-and-constant approximation for the design sequence & Theorem \\
(b1) CLR over all dual vertices & (a), (c), (d), under the \textbf{full augmented} CLR conditions, including one-sided validity. The added rows change the dual feasible set, its vertices, the gradients, the selection step and the quantile; appending row standard errors to the old CLR maximum is not enough & Theorem under those conditions (cited) \\
(b2)(i) fixed simultaneous superset & the same branch-union and crossing arguments, with the fixed-superset bounds and the allocated endpoint levels & Theorem \\
(b2)(ii) projection & the joint coefficient event controls both ends and any data-dependent collection of dual-feasible candidates; spend the coefficient-error budget once, and account separately for the residual of actual targets. Needs \(\Sigma\) positive definite or a reduced full-rank parameterisation (P3.5(b2)(ii), corrected) & Theorem \\
(b3) CLR calibrated on simulated candidates (the package default for CR: the CLR branch union) & the true-branch inclusion and the crossing implications are deterministic; their nominal probabilities are conjectural & \textbf{Algorithm} (Conjecture P3.5(b3)); simulation coverage cannot promote it \\
\end{longtable}

Further limits:
* \textbf{Regime.} M1/M2 with growing districts (R3) is a conjecture in P3 §0.5 independently of candidate collection.
* \textbf{CLR Condition P is not verified by a finite vertex set.} Their Condition P needs a fixed compact index representation along the sequence, gradients uniformly bounded away from zero, a well-conditioned covariance and a covariance-estimation rate. Nonzero gradients for one design are not enough.
* \textbf{Zero-variance candidates.} A candidate endpoint with \(\hat\sigma=0\) is a different matter from an exactly known row intercept. Treat it as exact only if its gradient is exactly zero; otherwise check it separately.
* In the package, the default CR interval (CLR branch union) and, with \texttt{ci\ =\ "clr"}, the ``already implied'' intervals and the near-tie margin are therefore algorithmic. The gated Monte Carlo coverage (0.990) is a simulation finding.
* Version (a) needs no Slater-type condition. Statements of P3.2/P3.4 type do, and rows that bind at the truth (e.g.~the CR edge) can break it.

\begin{enumerate}
\def\labelenumi{(\alph{enumi})}
\tightlist
\item
  \textbf{Added rows.} For a finite set \(\mathcal R\) of linear-in-\(\beta\) rows, P3.5(a), P3.5(b) and P3.6(a) hold for the augmented system, version by version as in the table above, for \(\tilde B\); for \(B\) they hold with P3-S1a, the weight condition and the residual allowance carried over. The new rows join the Gaussian maximum that sets the critical value. A row whose \(s_j\) does not depend on \(\beta\) has \(\sigma_j=0\): it is imposed exactly and left out of the maximum.
\end{enumerate}

\begin{itemize}
\tightlist
\item
  \textbf{Actual targets that are contrasts.} For an actual gap \(B_1-B_0\) the residual is \(\sum_i(w_{i1}e_{i1}-w_{i0}e_{i0})\), and each unit term has range at most \(w_{i1}+w_{i0}\), so \(\tfrac12\{\sum_i(w_{i1}+w_{i0})^2\}^{1/2}\) is a valid standard-deviation envelope under conditional unit independence. The single-row \(S_1\) need not bound it. The sum of the per-cell envelopes (which the package uses) is also valid, by Minkowski's inequality, and more conservative. \emph{(Corrected after review, Batch B, Nov 2026.)}
\end{itemize}

\begin{enumerate}
\def\labelenumi{(\alph{enumi})}
\setcounter{enumi}{1}
\tightlist
\item
  \textbf{Parallel disjunction.} Inflate each piece separately, then take the maximum:
  \[\max_l\big(s_l(\hat\beta)+\hat k\hat\sigma_l\big)+a^\top k\ge0,\]
  with all pieces included in the critical-value maximum. \textbf{For version (a) (conservative) only}, this covers at the nominal level. The row is not linear in \(\beta\), so Lemma P3.1 fails (an endpoint becomes a max of minima) and CLR does not apply. For CLR, use (c): in W1/2×2 each branch is a set of linear-in-\(\beta\) rows.
\end{enumerate}

\begin{itemize}
\tightlist
\item
  W1 CR and \(2\times2\) CR are of this form, by §1.2(d):

  \begin{itemize}
  \tightlist
  \item
    lower row: \(h_D^\top k+\max(c_D-ER^*,\ c_D-\gamma_rER^*)\ge0\);
  \item
    upper row: \(-h_D^\top k+\max(ER^*-c_D,\ \gamma_rER^*-c_D)\ge0\).
  \end{itemize}
\item
  No branch is selected, so the sign of \(\widehat{ER}\) never enters.
\item
  The proof uses \textbf{one} simultaneous event: every population row intercept is at most its inflated estimate, and the fitted target constant is within its two-sided allowance. It needs no strict feasibility. At a tie, inflating all pieces and then maximising still dominates the true piece; taking the gradient of a winning estimated piece would be a different procedure. \emph{(Added after review, Batch B.)}
\end{itemize}

\begin{enumerate}
\def\labelenumi{(\alph{enumi})}
\setcounter{enumi}{2}
\tightlist
\item
  \textbf{General union.} \(\mathrm{CI}_\cup:=\mathrm{CI}(C_+)\cup\mathrm{CI}(C_-)\), with \emph{each} branch at the full level \(\alpha\), satisfies
  \[\liminf_nP(\theta_0\in\mathrm{CI}_\cup)\ \ge\ 1-\alpha.\]
\end{enumerate}

\begin{itemize}
\tightlist
\item
  No Bonferroni split across branches is needed.
\item
  The hull of \(\mathrm{CI}_\cup\) is also valid; it is wider or equal.
\item
  The emptiness test ``reject if both branch sets are empty'' has asymptotic level \(\le\alpha_U+\alpha_L\).
\item
  This is \textbf{point coverage} of the true target. It gives no simultaneous coverage of both branch images, and it does not license choosing a branch from \(\widehat{ER}\). If designs form a triangular array and the true branch switches with \(n\), branch validity is needed along each subsequence on which that branch is true (there are only two labels). \emph{(Added after review, Batch B.)}
\end{itemize}

\begin{enumerate}
\def\labelenumi{(\alph{enumi})}
\setcounter{enumi}{3}
\tightlist
\item
  \textbf{Near-empty sets.}
\end{enumerate}

\begin{itemize}
\tightlist
\item
  Under the model, \(P(\mathrm{CI}=\emptyset)\to\le\alpha_U+\alpha_L\), including when the identified set is a point or nearly one.
\item
  Coverage holds unconditionally, so reporting ``empty'' is part of the allowed error, not an extra one.
\item
  Conditioning on passing the emptiness test can only \emph{raise} coverage: covering implies non-empty, so \(P(\text{cover}\mid\ne\emptyset)=P(\text{cover})/P(\ne\emptyset)\ge P(\text{cover})\).
\item
  Conditioning on other data events (a near-rejection, a CI that came out narrow) is not covered by the theorem.
\item
  If the restrictions are false, no coverage statement holds. In particular, a non-empty, narrow CI is no evidence that they are true (G3). The EI referee's experiment makes this concrete: with correct linear means and a false CR, the conservative CR interval at 600 precincts covered 0 of 400 times under a large departure and was never empty (§3.4, V12).
\item
  All statements are \textbf{pointwise} in the DGP; no uniform guarantee over sequences of nearly infeasible laws follows.
\item
  Conditioning on non-emptiness offers no protection when the researcher selects a belief because it gives an attractive result. \emph{(Added after review, Batch B.)}
\end{itemize}

\begin{enumerate}
\def\labelenumi{(\alph{enumi})}
\setcounter{enumi}{4}
\tightlist
\item
  \textbf{Plug-in restrictions are not valid.} A rule that estimates the CR bracket \(\operatorname{hull}\{\gamma_r\widehat{ER},\widehat{ER}\}\) and imposes it without inflating it can under-cover when the truth is on the bracket's edge, and the under-coverage does not vanish as \(n\) grows. For a fixed repeating three-point design, \(\limsup_nP(\tilde D_0\in\mathrm{CI}_{\rm naive})\le\Phi(\tfrac7{12}z_{.995})=0.933525<0.95\) (proof below). \emph{(Corrected after review, Batch B: earlier supported by a Monte Carlo and an asymptotic sketch only.)}
\end{enumerate}

\hypertarget{pF-proof-2}{%
\subsubsection{3.3 Proof}\label{pF-proof-2}}

\emph{(a)}
1. Under the model, \(k_0\) satisfies the base rows and the rows in \(\mathcal R\) at \(\beta_0\). {[}R3(e); hypothesis{]}
2. Each added \(s_j(\hat\beta)\) is a linear map of \(\hat\beta\), so it is jointly asymptotically normal with the base terms. Let \(\mathcal E\) be P3.5(a)'s event, with the maximum taken over all rows. Then \(P(\mathcal E)\to\ge1-\alpha_U-\alpha_L\). {[}OLS; lin; P3.5(a) step 1{]}
3. On \(\mathcal E\), \(k_0\in\hat\Theta^+\). P3.5(a) steps 3--7 and P3.6(a) then go through verbatim, since they use only \(k_0\in\hat\Theta^+\) and the residual term. {[}1--2; P3.5, P3.6{]}
4. For (b) of P3.5, the endpoint is still a minimum of finitely many linear functions of \(\beta\), because Lemma P3.1 needs only rows that are linear in \(\beta\) with fixed \(a_j\). CLR applies to the \textbf{augmented} system if its conditions hold there: the added rows change the dual feasible set, its vertices, the gradients, the selection and the quantile, so the conditions must be checked anew (status table). It inherits P3's open item on collecting candidates. {[}P3.1; CLR{]} ∎(a)

\emph{(b)}
5. Let \(l^*\) be a piece that holds at the truth: \(s_{l^*}(\beta_0)+a^\top k_0\ge0\). {[}hypothesis; Def disjunction{]}
6. On \(\mathcal E\), which now includes every piece, \(s_{l^*}(\beta_0)\le s_{l^*}(\hat\beta)+\hat k\hat\sigma_{l^*}\). So \(\max_l(s_l(\hat\beta)+\hat k\hat\sigma_l)+a^\top k_0\ge0\). {[}Def \(\mathcal E\); 5{]}
7. For given data this is one linear row in \(k\), and \(k_0\) satisfies it on \(\mathcal E\). Proceed as in step 3. {[}6; (a){]}
8. For W1/2×2 CR: §1.2(d) gives \(D^{r*}\ge\min(ER^*,\gamma_rER^*)\iff\max(D^{r*}-ER^*,\ D^{r*}-\gamma_rER^*)\ge0\), and similarly for the upper side. Each piece is linear in \(\beta\) apart from the common \(h_D^\top k\). {[}§1.3 steps 11--13; alg{]} ∎(b)

\emph{(c)}
9. Fix a branch \(j^*\in\{+,-\}\) with \(k_0\in C_{j^*}\). It is non-random, because \(k_0\) is a population quantity. {[}hypothesis{]}
10. The model ``rung \(\cap\,C_{j^*}\)'' holds, so by (a), \(\liminf P(\theta_0\in\mathrm{CI}(C_{j^*}))\ge1-\alpha\). {[}9; (a){]}
11. \(\mathrm{CI}(C_{j^*})\subseteq\mathrm{CI}_\cup\subseteq\operatorname{hull}\mathrm{CI}_\cup\). {[}Def{]}
12. The other branch's CI is never used for coverage, so its level is irrelevant. That is why no split is needed. {[}10--11{]}
13. Both sets empty implies the \(j^*\) set is empty, which has probability \(\to\le\alpha_U+\alpha_L\). {[}P3.6(a) for version (a); the crossing argument of (d) step 15 for version (b); both applied to branch \(j^*\){]} ∎(c)

\emph{(d)}
14. For P3.5(a): \(\mathrm{CI}=\emptyset\) requires \(\hat\Theta^+=\emptyset\), which requires \(\mathcal E^c\). {[}P3.5(a) step 2{]}
15. For P3.5(b), take \(\theta_0=\tilde B\) (or \(\tilde D\)). A crossing \(\hat L>\hat U\) requires \(\hat U<\theta_0\) or \(\hat L>\theta_0\). These have asymptotic probabilities \(\le\alpha_U\) and \(\le\alpha_L\) respectively, by CLR's one-sided validity, whatever the width of the identified set. For the target \(B\), the E-allowance only widens the interval, so a crossing of the \(B\)-interval implies a crossing of the \(\tilde B\)-interval, and the same bound holds. {[}CLR; union; P3 allowance{]}
16. In both cases the empty event is contained in the non-coverage event whose probability is bounded. Neither step uses the width of the identified set. {[}14--15{]}
17. \(\{\theta_0\in\mathrm{CI}\}\subseteq\{\mathrm{CI}\ne\emptyset\}\), so \(P(\text{cover}\mid\ne\emptyset)=P(\text{cover})/P(\ne\emptyset)\ge P(\text{cover})\). If the restrictions are false, step 1 fails. {[}Def cond. prob.; 1{]} ∎(d)

\emph{(e)} \textbf{Fixed-design proof} (the partial-identification referee's argument, adopted with credit; it replaces the earlier Monte Carlo sketch).
18. Repeat each of \(X=1/4,1/2,3/4\) equally often, with equal sizes. Take \(m_0=.4\) and \(m_1=.55\) constant, and draw the group rates independently across units and groups from bounded laws with these means and positive variance (e.g.~beta draws of fixed concentration). Every base probability row has strict slack, OLS has a full-rank fixed design and a positive definite asymptotic covariance, and HC3 is consistent. The positive CR branch has \(\tilde D_0=ER=.15\) on its upper edge. {[}Def; OLS{]}
19. In the basis \(((1-x)^2,2x(1-x),x^2)\), for \textbf{any} coefficient vector \(q\): \(ER(q)=q_{11}-q_{00}\) and \(c_D(q)=\tfrac7{12}(q_{11}-q_{00})\), because \(E[X^2]/E[X]=7/12\) and the middle coefficient cancels by symmetry. So \(\operatorname{se}(c_D)/\operatorname{se}(\widehat{ER})=7/12\) exactly. {[}alg{]}
20. With probability tending to one, the plug-in sign is positive and the base rows allow the \(k\) that attains the plug-in upper edge (strict slack and continuity). The naive upper endpoint is then \(\widehat{ER}+k_{\rm crit}\operatorname{se}(c_D)\). {[}19; continuity{]}
21. Its critical value is calibrated on eight one-sided base rows and two target tails. By the union bound, the population Gaussian critical value satisfies \(k_{\rm crit}\le z_{.995}\) at level .95, whatever the correlations. {[}union{]}
22. Non-coverage is \(\{\widehat{ER}-ER<-k_{\rm crit}\operatorname{se}(c_D)\}\), so \(\limsup_nP(\tilde D_0\in\mathrm{CI}_{\rm naive})\le\Phi(\tfrac7{12}z_{.995})=0.933525<0.95\). {[}20--21; OLS{]}
23. With 4,000 critical-value draws (as in \texttt{verify\_R10.py}), each draw exceeds \(z_{.996}\) with conditional probability at most .04 by the same union bound; an empirical .95 quantile above that threshold needs at least 200 exceedances. So the limiting coverage is at most \(\Phi(\tfrac7{12}z_{.996})+P\{\mathrm{Bin}(4000,.04)\ge200\}=0.940092\). ∎ {[}union; binomial tail{]}

\emph{Remarks.} (i) The far-edge design of §3.4 also satisfies the condition in its design limit: with \(X=.1+.8U\), \(U\sim\)Beta\((2,2)\), symmetry gives \(c_D(q)=(.5+2\operatorname{Var}X)(q_{11}-q_{00})\) with \(\operatorname{Var}X=.032\), a ratio of .564 and a limiting bound of about .927 with population critical values (referee's derivation; the ratio is recomputed in V9). (ii) The Monte Carlo of §3.4 (coverage 0.917, Monte Carlo SE about 0.011, and the reviewer's 0.914, 0.909, 0.903 at \(n=150\), 600, 2000) is now an illustration, not the argument. Check: V9.

\textbf{Uniformity.} The statements are pointwise in the DGP, as in P3. Uniform validity near point identification (a Stoye-type adaptation) remains P3's open item. (d) only says that emptiness is controlled pointwise, including at point-identified DGPs.

\hypertarget{pF-monte-carlo-verify_r10.py-output-verify_r10_output.txt}{%
\subsubsection{\texorpdfstring{3.4 Monte Carlo (\texttt{verify\_R10.py}, output \texttt{verify\_R10\_output.txt})}{3.4 Monte Carlo (verify\_R10.py, output verify\_R10\_output.txt)}}\label{pF-monte-carlo-verify_r10.py-output-verify_r10_output.txt}}

Setup:
* 2×2, \(n=150\) units, 600 replications.
* Unit rates are Beta draws with linear conditional means, so LCE and CR hold.
* HC3 covariance, conservative CI (P3.5(a)) at 95\%, target \(\tilde D\).

\begin{longtable}[]{@{}
  >{\raggedright\arraybackslash}p{(\columnwidth - 8\tabcolsep) * \real{0.2000}}
  >{\raggedright\arraybackslash}p{(\columnwidth - 8\tabcolsep) * \real{0.2000}}
  >{\raggedright\arraybackslash}p{(\columnwidth - 8\tabcolsep) * \real{0.2000}}
  >{\raggedright\arraybackslash}p{(\columnwidth - 8\tabcolsep) * \real{0.2000}}
  >{\raggedright\arraybackslash}p{(\columnwidth - 8\tabcolsep) * \real{0.2000}}@{}}
\toprule\noalign{}
\begin{minipage}[b]{\linewidth}\raggedright
Design
\end{minipage} & \begin{minipage}[b]{\linewidth}\raggedright
JKST
\end{minipage} & \begin{minipage}[b]{\linewidth}\raggedright
CR slab (b)
\end{minipage} & \begin{minipage}[b]{\linewidth}\raggedright
CR union (c): coverage / occupied length (hull span; unmerged sum; runs with overlap)
\end{minipage} & \begin{minipage}[b]{\linewidth}\raggedright
CR naive (plug-in bracket)
\end{minipage} \\
\midrule\noalign{}
\endhead
\bottomrule\noalign{}
\endlastfoot
Truth on the CR edge (\(\delta^*_W=0\)), \(ER^*=0.15\) & 1.000 / 1.106 & \textbf{0.992} / 0.359 & 0.992 / 0.357 (0.357; 0.386; 134/600) & \textbf{0.917} / 0.256 \\
Truth on the edge, \(ER^*=0.005\) (near-point CR set, branch ambiguous) & 1.000 / 1.223 & 1.000 / 0.242 & 1.000 / 0.242 (0.242; 0.430; 591/600) & 1.000 / 0.147 \\
Interior truth & 1.000 / 1.232 & 1.000 / 0.480 & 1.000 / 0.480 (0.480; 0.480; 0/600) & 1.000 / 0.388 \\
\end{longtable}

Entries are coverage / median length. No CI was empty in any design. \emph{(Corrected after review, Batch B, Nov 2026: the union column earlier reported the sum of the two branch lengths, 0.386 / 0.430 / 0.480, which double-counts overlaps. It now reports the occupied (Lebesgue) length, with the hull span and the old unmerged sum beside it; the numbers reproduce the EI referee's correction exactly. The apparent length penalty of the branch union in the near-zero design disappears; the headline slab comparison, 0.359 against 1.106, is unaffected. The package merges overlapping pieces (\texttt{.rxc\_union}) and was not affected. Check V11.)}

\textbf{What these runs do not test:} the CLR procedure (all four methods use the conservative version), coverage of the actual-rate target (residual allowance), the crossing emptiness rule, or a growing-dimension regime. The design is narrow: fixed composition and sizes, Beta rates, and cancelling slopes chosen to put the truth on the edge.

\textbf{A false CR (EI referee's experiment, reproduced as V12).} Same generator, with \(m_0=.40-sx\) and \(m_1=.50+s-sx\): the aggregate regression is \(.40+.10x\) for every \(s\), the between sign is positive throughout, and for \(s>0\) both contextual slopes are negative, so CR is false. The design is fixed across \(s\).

\begin{longtable}[]{@{}
  >{\raggedleft\arraybackslash}p{(\columnwidth - 12\tabcolsep) * \real{0.1429}}
  >{\raggedleft\arraybackslash}p{(\columnwidth - 12\tabcolsep) * \real{0.1429}}
  >{\raggedleft\arraybackslash}p{(\columnwidth - 12\tabcolsep) * \real{0.1429}}
  >{\raggedleft\arraybackslash}p{(\columnwidth - 12\tabcolsep) * \real{0.1429}}
  >{\raggedleft\arraybackslash}p{(\columnwidth - 12\tabcolsep) * \real{0.1429}}
  >{\raggedleft\arraybackslash}p{(\columnwidth - 12\tabcolsep) * \real{0.1429}}
  >{\raggedleft\arraybackslash}p{(\columnwidth - 12\tabcolsep) * \real{0.1429}}@{}}
\toprule\noalign{}
\begin{minipage}[b]{\linewidth}\raggedleft
Precincts
\end{minipage} & \begin{minipage}[b]{\linewidth}\raggedleft
\(s\)
\end{minipage} & \begin{minipage}[b]{\linewidth}\raggedleft
true \(\tilde D\)
\end{minipage} & \begin{minipage}[b]{\linewidth}\raggedleft
JKST covered / 400
\end{minipage} & \begin{minipage}[b]{\linewidth}\raggedleft
CR slab covered / 400
\end{minipage} & \begin{minipage}[b]{\linewidth}\raggedleft
median JKST width
\end{minipage} & \begin{minipage}[b]{\linewidth}\raggedleft
median CR width
\end{minipage} \\
\midrule\noalign{}
\endhead
\bottomrule\noalign{}
\endlastfoot
150 & 0 & .100 & 400 & 399 & 1.262 & .306 \\
150 & .08 & .170 & 400 & 336 & 1.256 & .304 \\
150 & .25 & .318 & 400 & 2 & 1.254 & .303 \\
600 & 0 & .100 & 400 & 398 & 1.074 & .194 \\
600 & .08 & .169 & 400 & 110 & 1.069 & .196 \\
600 & .25 & .317 & 400 & 0 & 1.066 & .195 \\
\end{longtable}

No interval was empty. A false restriction gives narrow, nonempty and reproducibly wrong intervals; precision makes the wrong answer more reproducible, not more credible. This is the price R10(d) states, made visible. \emph{(Added after review, Batch B.)}

\hypertarget{pF-consequences-for-the-package-2}{%
\subsubsection{3.5 Consequences for the package}\label{pF-consequences-for-the-package-2}}

\begin{enumerate}
\def\labelenumi{\arabic{enumi}.}
\tightlist
\item
  \textbf{Every added row goes through the same CI machinery.} Sign, ordering and local rows join the critical-value maximum. Nothing new is needed beyond appending rows (R10(a)).
\item
  \textbf{CR (W1 default and 2×2) uses the slab form (b) for the conservative CI} (one row per side, each piece inflated before the maximum) \textbf{and the branch union (c) for CLR}, the default CI. That default is algorithmic (Conjecture P3.5(b3)); the conservative version and the projection interval carry the guarantee. In the edge design of §3.4 the \textbf{conservative slab} has median width 0.359 (coverage 0.992) against 1.106 for JKST alone; the default \textbf{CLR branch union}, in the package's gated run, had median width 0.230 and coverage 0.990 (§5). It is a large gain, and it rests on a belief whose refutability depends on the base system: in the model-free continuous table problem (2c M3, no linear rung) CR is always feasible; under LCE it can be rejected even when LCE alone fits (Remark R3.2); with integer counts it can be rejected outright (2c Remark M2.1). A false CR gives narrow wrong intervals (§3.4, V12). \emph{(Corrected after review, Batch B, Nov 2026: earlier ``it rests on an unrefutable belief (2c M3)'', which imported M3 into the LCE base; and the quoted numbers did not name their method.)}
\item
  \textbf{Other CR pairings use the union of the branch CIs, each at the full level (c).} The package reports the union, which can be two intervals, and offers the hull as an option.
\item
  \textbf{Never impose an estimated bracket as if it were known.} The EGGHM bracket \(\operatorname{hull}\{\gamma\,\widehat{ER},\widehat{ER}\}\) intersected with a JKST CI is not a confidence interval (e). Direct argument: in the design of (e), the upper endpoint of that intersection cannot exceed \(\widehat{ER}\) with probability tending to one, so its limiting coverage of \(\tilde D_0=ER\) is at most one half. \emph{(Corrected after review, Batch B, Nov 2026: earlier ``the MC naive interval contains bracket ∩ JKST CI, so the under-coverage transfers a fortiori''. That containment is unproved and false in general: dropping an empty pre-expansion intersection and expanding endpoints do not commute. With raw base interval \([0,.1]\), bracket \([.15,.2]\) and target allowance .1, the raw intersection is empty although the expanded base meets the bracket. Check V10.)} \textbf{Which procedure the numbers describe} \emph{(added after review, round 2, EI referee; check V13)}. In the edge design of §3.4 (\(n=150\), 600 replications), the \textbf{literal} intersection of the fitted bracket with the completed JKST interval covered the true \(\tilde D\) in \textbf{304/600} (50.7\%) replications, in line with the at-most-one-half limit above. The \texttt{CR-naive} procedure of \texttt{verify\_R10.py}, which adds the target's own inflation after imposing the bracket, covered \textbf{550/600} (91.7\%). The miss rate of about 8--10\% quoted in the guide and vignette refers to \texttt{CR-naive} only; the literal intersection misses about half the time. When the regression and the target use different weights, the two sample ERs must be distinguished. (Our \(\widehat{ER}^*\) is \(\tau\)-based. It equals the \(T\)-based EGGHM \(\widehat{ER}\) under \(N\)-weighted OLS, because \(X_r\) and \(1\) lie in the span of the quadratic monomials. P3 uses unweighted OLS, so the two differ slightly in samples.)
\item
  \textbf{Emptiness is reported as a test result, not as an error.} Output includes:

  \begin{itemize}
  \tightlist
  \item
    the rejection flag;
  \item
    the combination that was rejected (R3(e));
  \item
    a warning that a narrow CI after a near-rejection is not reassurance (d).
  \end{itemize}
\end{enumerate}

\begin{center}\rule{0.5\linewidth}{0.5pt}\end{center}

\hypertarget{pF-independent-review-log}{%
\subsection{4. Independent review log}\label{pF-independent-review-log}}

A fresh-session adversarial review of §§1--3 and the three scripts. Verdicts: R3, R2 and R10 all \textbf{accept with fixes}. Every fix below has been applied.

\begin{longtable}[]{@{}
  >{\raggedright\arraybackslash}p{(\columnwidth - 8\tabcolsep) * \real{0.2000}}
  >{\raggedright\arraybackslash}p{(\columnwidth - 8\tabcolsep) * \real{0.2000}}
  >{\raggedright\arraybackslash}p{(\columnwidth - 8\tabcolsep) * \real{0.2000}}
  >{\raggedright\arraybackslash}p{(\columnwidth - 8\tabcolsep) * \real{0.2000}}
  >{\raggedright\arraybackslash}p{(\columnwidth - 8\tabcolsep) * \real{0.2000}}@{}}
\toprule\noalign{}
\begin{minipage}[b]{\linewidth}\raggedright
\#
\end{minipage} & \begin{minipage}[b]{\linewidth}\raggedright
Item
\end{minipage} & \begin{minipage}[b]{\linewidth}\raggedright
Class
\end{minipage} & \begin{minipage}[b]{\linewidth}\raggedright
Finding
\end{minipage} & \begin{minipage}[b]{\linewidth}\raggedright
Fix
\end{minipage} \\
\midrule\noalign{}
\endhead
\bottomrule\noalign{}
\endlastfoot
1 & R10(d) & ERROR & It claimed coverage conditional on non-emptiness can fall below \(1-\alpha\). In fact cover ⊆ non-empty, so conditioning raises coverage & (d) and step 17 rewritten \\
2 & R10(b) & GAP & The slab row is not linear in \(\beta\), so CLR is not covered & Slab restricted to version (a); CLR goes through (c) \\
3 & R10 §3.1 & GAP & Two emptiness rules for version (b) would double the level & One rule per version \\
4 & R10 step 15 & GAP & Wrong for \(\theta_0=B\) (the E-allowance) & Argument run with \(\tilde B\) \\
5 & R10(a) & GAP & Rows with \(\sigma_j=0\); conditions on the augmented system; Slater & Stated \\
6 & R10(e) & GAP & Shown by MC only & Asymptotic sketch added \\
7 & R10 minors & MINOR & Step-13 citation; ``wider or equal''; a-fortiori remark; \(\tau\)- vs \(T\)-based \(ER\) & Fixed \\
8 & R3(c) & GAP & ``Sharp iff intersect'' is false when one branch is empty & Case split \\
9 & R3 setting & GAP & Meaning of \(\tilde B\) and of \(ER^*\) under (C2) & Note added \\
10 & R3 §1.6.2 & GAP & Contradicted R10(b): CR rows need inflation & Rewritten \\
11 & R3 §1.4 & MINOR & Non-polyhedral images may be non-closed & Stated \\
12 & R3 minors & MINOR & S1 tests (b), not (a); certificate is an LP only for \(\Theta_O\); ``weakly outer''; \(\gamma\) used for two things (target normal now \(h\)); S5 gap sizes reach ≈ 31\% on other seeds & Fixed \\
14 & R2(c) & GAP & Proved for one row only & Separation step 11′ \\
15 & R2 step 9 & GAP & Which optimum is meant & Local ⇒ global argument \\
16--17 & R2 minors & MINOR & Domain boundary in {[}pert{]}; citations; nonemptiness of \(F\) is an assumption & Fixed \\
18 & R2 checks & GAP & T1 never tests a non-singleton face; T2 nearly tautological & T1b added; T2 relabelled \\
\end{longtable}

Numerical re-checks by the reviewer: the identities in R3 step 11 hold for every \(r\) to \(1.5\times10^{-16}\); the S5 holes are genuine; the T4 28\% ``cuts set only'' rate is not a tolerance artefact (shallowest cut \(3.8\times10^{-6}\)); \texttt{verify\_R10.py} reproduces exactly.

\textbf{External review (Batch B, Nov 2026).} This internal review covered §§1--3 and three scripts. It missed the face-LP intercept (R2 step 4), the set-versus-target meaning of redundancy (R2(e)), the LCE example that excludes CR (Remark R3.2), the double-counted union lengths (§3.4), and the scope of the ``unrefutable'' sentence; and it left R10(e) to simulation. All are corrected above; the full ledger is \texttt{reviews/B\_PI\_RESPONSE\_PLAN.md} §2. Process lesson: a pre-screen must also cover the scope document a proof relies on (here 2c), not only the new proofs.

\begin{center}\rule{0.5\linewidth}{0.5pt}\end{center}

\hypertarget{pF-implementation-status-build-step-1-branch-claudemsign-build}{%
\subsection{\texorpdfstring{5. Implementation status (build step 1, branch \texttt{claude/msign-build})}{5. Implementation status (build step 1, branch claude/msign-build)}}\label{pF-implementation-status-build-step-1-branch-claudemsign-build}}

\textbf{Implemented} (\texttt{package/eiPartialID}):

\begin{longtable}[]{@{}
  >{\raggedright\arraybackslash}p{(\columnwidth - 4\tabcolsep) * \real{0.3333}}
  >{\raggedright\arraybackslash}p{(\columnwidth - 4\tabcolsep) * \real{0.3333}}
  >{\raggedright\arraybackslash}p{(\columnwidth - 4\tabcolsep) * \real{0.3333}}@{}}
\toprule\noalign{}
\begin{minipage}[b]{\linewidth}\raggedright
Spec item
\end{minipage} & \begin{minipage}[b]{\linewidth}\raggedright
Where
\end{minipage} & \begin{minipage}[b]{\linewidth}\raggedright
Notes
\end{minipage} \\
\midrule\noalign{}
\endhead
\bottomrule\noalign{}
\endlastfoot
Sign rows on \(\delta_B\) (W1, pairwise) and \(\delta_W\) (W1, W2 \(=\Delta_r\)), §1.6.1, R10(a) & \texttt{R/rxc\_sign.R} (\texttt{.rxc\_mlin}, \texttt{.rxc\_assoc\_W}, \texttt{.rxc\_assoc\_form}); rows added in \texttt{.rxc\_extra\_rows} (\texttt{R/rxc\_core.R}) & \texttt{restrictions\$assoc\ =\ list(list(type\ =\ "dB"/"dW",\ row,\ col,\ sign,\ version,\ other))}. One design-weighted (\(N_i\)) row per restriction, pooled over groups; unit \(i\) uses its own group's \(H\) \\
CR, W1 / 2×2 slab, §1.6.2, R10(b) & \texttt{.rxc\_cr\_parts} (slab = two disjunctive rows), \texttt{.rxc\_fix\_disj}; \texttt{.rxc\_ci\_conservative} (\texttt{R/rxc\_infer.R}) & Point bounds: slab at \(\hat\beta\). Conservative CI: every piece in the critical-value maximum, each inflated by its own \(\hat k\hat\sigma_l\), then the max. \texttt{.rxc\_ci\_clr} refuses disjunctive rows \\
CR, W2 / pair, and CLR for every CR, §1.6.3, R10(c) & \texttt{.rxc\_systems(cons,\ "branch")}, \texttt{.rxc\_estimate} & Union of branch CIs, each at the full level; branches with empty / crossed / out-of-range intervals drop out; all empty = rejected. Several CRs give \(2^m\) systems \\
Two intervals per target, hull option & \texttt{.rxc\_union}, \texttt{.rxc\_union\_summary}; table columns \texttt{n\_pieces}, \texttt{*\_gap\_lower/upper}; \texttt{fit\$pieces}; \texttt{eiBounds(union\ =\ "pieces"/"hull")} & Printed as \texttt{{[}a,\ b{]}\ U\ {[}c,\ d{]}} \\
R2 redundancy diagnostic, §2.5 & \texttt{eiRedundancy()} (exported; stored in \texttt{fit\$sign} by default when \texttt{assoc}/\texttt{cr} are given) & Face test via one feasibility LP per endpoint (\texttt{.rxc\_meet}, \texttt{.rxc\_face\_moves}); duals never used. Verdicts implied / cuts set only / tightens, added first and added last, with width gains. If the base is itself a union (another non-W1 CR), the verdict falls back to a brute-force width comparison and says so (\texttt{method}) \\
``Already implied'' table & \texttt{.rxc\_implied\_table} & Every \(\delta\) (W1, W2, pairwise) under the rung alone: point interval and CI, signs flagged \\
CR certificate, §1.6.3 & \texttt{.rxc\_cr\_table} & \(\Theta\cap\{g_B=0,g_W=0\}\neq\emptyset\) (feasibility LP), branch intervals, number of pieces \\
Emptiness, §1.6.4, R10(d) & \texttt{.rxc\_empty\_test} (conservative; reject iff every system's inflated set is empty; also reports whether the rung alone is rejected); notes in \texttt{eiBounds} & Messages name the combination ``rung (M) + restriction labels''; 2-row note that \(\delta_B\) / CR are never refuted on their own; ``narrow interval after a near-rejection is not reassurance'' \\
CLR no longer falls back to the conservative CI & \texttt{.rxc\_collect(relax\ =\ TRUE)}, \texttt{.rxc\_ci\_clr} & When no draw is feasible, candidates come from rows relaxed by the smallest uniform slack. The dual feasible set does not depend on the row constants, so these vertices are dual-feasible for the original rows and bound the endpoint by weak duality (P3.5(b2) step 1). Coverage of the CLR critical value over the collected set is Conjecture P3.5(b3). The CLR rule (ends cross ⇒ empty) is kept, so the two emptiness rules are never mixed \\
Descriptive \(\widehat{ER}\), \(\widehat{NM}=\gamma\widehat{ER}\), \(\gamma\), bracket, §3.5.4 & \texttt{.rxc\_descriptive}, printed under ``DESCRIPTIVE ONLY'' & \(\tau\)-based; equals the EGGHM \(T\)-based value under \texttt{weights\ =\ "n"} (tested to 1e-10). Never imposed \\
Integration & \texttt{ladder(sign\_rungs\ =\ ...)}, \texttt{print}/\texttt{summary}, \texttt{sensitivity()}, \texttt{breakdown()} (pass \texttt{union}) & Point-identified rungs (L5, L6) ignore sign/CR rows, with a note \\
Docs & roxygen (\texttt{eiBounds}, \texttt{eiRedundancy}, \texttt{ladder}), vignette \texttt{eiPartialID\_rxc.Rmd} section ``Signs on associations and contextual reinforcement'', \texttt{PLAIN\_LANGUAGE\_GUIDE.md} §9 & Applied framing as in §3.5 \\
\end{longtable}

\textbf{Tests} (\texttt{tests/testthat/test-msign.R}; reference numbers from \texttt{docs/extension/make\_msign\_reference.py}, which calls \texttt{verify\_R3.py}'s \texttt{assoc\_3x2} / \texttt{assoc\_alt} and writes \texttt{msign\_example\_3x2.csv}, \texttt{msign\_gap\_3x2.csv}, \texttt{msign\_reference\_3x2.csv}):

\begin{itemize}
\tightlist
\item
  all association forms, ER, NM, γ, D and \(\tilde B\) match Python to 1e-9 on a noise-free 3×2 example and on the gap example;
\item
  the W1 identities hold for random \(k\) to 1e-12;
\item
  3×2 CR point bounds (W1, W2, pair; three targets) equal the Python branch unions; W1 slab = union of branches; the gap example (near-1-d support plus one ordering row, found with the S5 search) gives two pieces for W2 and pair, certificate false, and the hull option;
\item
  noise-free gap example: both CI versions return the two pieces;
\item
  2×2 closed forms of R2(d) for all four sign rows and CR (voter\_vax and two simulated designs, both regression weightings); 2×2 CR slab = the single non-empty branch;
\item
  MonotoneEI \texttt{voter\_vax} (bundled as \texttt{voter\_vax\_MonotoneEI.csv}, MIT, attribution in \texttt{voter\_vax\_MonotoneEI.README}): CR bounds for the gap \([-0.4699, -0.0502]\), ER \(=-0.470\), NM \(=-0.050\), N-weighted;
\item
  redundancy verdicts = brute force for 24 sign rows × 3 targets (both ``tightens'' and ``cuts the set, not this target'' occur); the middle verdict never occurs in 2×2; added-last gains = width differences;
\item
  joint refutation and its messages; CR CIs contain the estimate and are narrower than the rung's; an empty CLR branch is dropped via crossing ends, not by fallback;
\item
  \texttt{sensitivity}, \texttt{breakdown}, \texttt{ladder(sign\_rungs)} run with sign/CR rows; input checks;
\item
  gated coverage MC (\texttt{EIPARTIALID\_SLOW=1}): truth on the CR edge, \(n=150\), 200 reps: conservative slab \textbf{1.000} (median width 0.346), CLR branch union \textbf{0.990} (median width 0.230).
\end{itemize}

\textbf{Run:} R 4.3 (Ubuntu apt) in the build session: full \texttt{testthat} suite 0 failures (skips: \texttt{ei.Datasets} not installed, slow MC unless enabled); the slow MC was run once with the result above; \texttt{R\ CMD\ check\ -\/-no-manual\ -\/-no-vignettes}: no errors (pre-existing notes: unused imports \texttt{MASS}/\texttt{eco}, unavailable Suggests; a locale warning from the container). The vignette was rendered with \texttt{rmarkdown::render}.

\textbf{Release 0.3.0 (response to the Batch B referees, Nov 2026).} Changes since the table above (\texttt{package/eiPartialID/NEWS.md}; tests in \texttt{tests/testthat/test-batchB.R}):

\begin{longtable}[]{@{}
  >{\raggedright\arraybackslash}p{(\columnwidth - 2\tabcolsep) * \real{0.5000}}
  >{\raggedright\arraybackslash}p{(\columnwidth - 2\tabcolsep) * \real{0.5000}}@{}}
\toprule\noalign{}
\begin{minipage}[b]{\linewidth}\raggedright
Item
\end{minipage} & \begin{minipage}[b]{\linewidth}\raggedright
Status in 0.3.0
\end{minipage} \\
\midrule\noalign{}
\endhead
\bottomrule\noalign{}
\endlastfoot
Proof status of \texttt{assoc} and \texttt{cr} & \textbf{experimental}; \texttt{print()} shows the EXPERIMENTAL flag \\
Grouped rungs (M1, M1(δ), M2, L4 with groups) & enabled, with a fit-time note: each sign is one pooled restriction across groups; sharpness proved for the pooled model only (R16 open). \texttt{print()} says ``outer set; sharpness under grouped rungs not yet proved'' \\
Refutation wording & ``always feasible'' only in the model-free problem with continuous tables; under the rung the combination can be rejected (R3(e), Remark R3.2). Next to a joint rejection, the output says whether the rung alone is rejected. An empty plug-in is ``a sample diagnostic, not a test'' \\
CR confidence interval & CLR branch union by default, labelled \textbf{algorithmic} on its line; the projection branch union (proved, P3.5(b2)) is printed beside it \\
Reporting & every union piece printed (three or more included); \texttt{base\_*} columns with the rung-only estimate and CIs; weight concentration per signed association in \texttt{eiRedundancy()\$concentration} (largest unit weight, effective number of units) \\
Redundancy verdicts & ``empty at the point estimate (see the emptiness test)''; ``cuts the set, not this target'' means zero gain for this target, not that the row is implied (R2(e)) \\
W2 & labelled ``pooled over groups''; a per-group version is planned (decision U8(a)) \\
Projection and CLR guards & candidates whose gradients use a numerical null direction of \(\hat V\) are dropped from the projection interval. \emph{Corrected after review, round 2:} this is a \textbf{safeguard}, not case (ii-b) of P3.5(b2): a numerical rank cutoff does not establish the exact identities (ii-b) needs, and the theorem is claimed with \(q=\dim\beta\) (decision: the package sets \(q=\dim\beta\); P3 §5 remark). A zero-variance CLR candidate is exact only if its gradient is zero \\
CI status by fit (package labels, 0.3.x; round 2) & \texttt{print()} states the CI status per fit. \textbf{M0:} P3.5(b2)(ii) for the projection interval and R10 for added rows (pooled). \textbf{Grouped rungs without M-SIGN rows:} P3's fixed-district results, regime R1 only. \textbf{Grouped rungs with M-SIGN rows:} no theorem (R16). The default CLR intervals are algorithmic throughout (Conjecture P3.5(b3)) \\
Slack \(\eta\) on \(\delta\) & implemented (\texttt{sensitivity(dial\ =\ "slack")}) \\
Not yet & the drop-one display; the L5/L6 sign check at the fitted point; per-group W2 rows \\
\end{longtable}

\textbf{Open / not done:}
* Uniformity near point identification remains P3's open item; the CLR candidate collection is unchanged apart from the relaxation above.
* With several non-W1 CRs, systems grow as \(2^m\); no pruning.
* The redundancy verdicts are point-level; the R2 §2.5.4 ``width change with its CI'' near ties is not yet reported (the ``already implied'' table does give CIs for the associations).
* Out of scope as instructed: N1, N2, M-LOC (R1), R×C rank (R5), DiD.
* \emph{Corrected after review (Batch B, Nov 2026):} this list earlier also named ``slack \(\eta\) on \(\delta\)''. That was stale: a per-restriction slack in probability units is implemented (\texttt{sensitivity(dial\ =\ "slack")}, \texttt{slack} in each restriction, \texttt{tests/testthat/test-msign-dial.R}). The stale line misled the EI referee.


\clearpage
\section{PACKAGE\_DESIGN: audit, new results and proof register (excerpts)}\label{app:pkgdesign}
\textit{Source: \texttt{docs/extension/PACKAGE\_DESIGN.md}, \S\S2, 5 and 8 (the other sections are design notes).}\par
% GENERATED by tools/make_appendices.py from docs/extension/PACKAGE_DESIGN.md
% Do not edit by hand; edit the source and regenerate.
\hypertarget{pG-audit-of-the-gpt-conversation}{%
\subsection{2. Audit of the GPT conversation}\label{pG-audit-of-the-gpt-conversation}}

\hypertarget{pG-claims-and-verdicts}{%
\subsubsection{2.1 Claims and verdicts}\label{pG-claims-and-verdicts}}

\begin{longtable}[]{@{}
  >{\raggedright\arraybackslash}p{(\columnwidth - 6\tabcolsep) * \real{0.2500}}
  >{\raggedright\arraybackslash}p{(\columnwidth - 6\tabcolsep) * \real{0.2500}}
  >{\raggedright\arraybackslash}p{(\columnwidth - 6\tabcolsep) * \real{0.2500}}
  >{\raggedright\arraybackslash}p{(\columnwidth - 6\tabcolsep) * \real{0.2500}}@{}}
\toprule\noalign{}
\begin{minipage}[b]{\linewidth}\raggedright
\#
\end{minipage} & \begin{minipage}[b]{\linewidth}\raggedright
Claim
\end{minipage} & \begin{minipage}[b]{\linewidth}\raggedright
Verdict
\end{minipage} & \begin{minipage}[b]{\linewidth}\raggedright
Basis
\end{minipage} \\
\midrule\noalign{}
\endhead
\bottomrule\noalign{}
\endlastfoot
G1 & JKST bounds on \(w_1\) map to bounds on \((\delta_W,\delta_B)\), and re-imposing those adds nothing about \(D\) & \textbf{Correct.} In 2×2 the pair is a line segment, \(\delta_W/\operatorname{Var}(X_N)+\delta_B/\operatorname{Var}(X)=D_{ER}-D_{NM}\). Re-imposing the marginals changes the \(D\) interval by \(<10^{-15}\). & M5(a); {[}V1{]} \\
G2 & Sign restrictions on \(\delta_W,\delta_B\) and CR give \([L_J,\min(U_J,D_{ER})]\), \([\max(L_J,D_{NM}),U_J]\), and \([L_J,U_J]\cap[\min,\max](D_{ER},D_{NM})\) & \textbf{Correct} (proved below) & EGGHM Props 4 and 6 hold for mean-level quantities {[}V1{]} \\
G3 & ``Incompatible ⇒ evidence against the assumptions'' & \textbf{Correct only jointly with LCE.} Model-free, the \(\delta_B\) sign and CR are never refutable (M3). Under LCE, a wrong \(\delta_B\) sign gave an empty set in 223 of 2000 designs; but in 1777 it gave a non-empty interval that excludes the truth {[}V2{]}. Emptiness is a one-sided test with low power. & M3; {[}V2{]} \\
G4 & Pointwise ordering \(c_1-w_1+d_1x\ge0\) on the hull & \textbf{Already implemented} as \texttt{order} (mean level, hull) & PHASE4; M5(b) \\
G5 & Pointwise contextual monotonicity \(w_1+d_1x\ge0\) & \textbf{Correct algebra; new as a module.} It is EGGHM's local within term (N2.3), \(x m_1'+(1-x)m_0'\), written under LCE. It is weaker than the existing \texttt{sign} rows and (given C1) stronger than global \(\delta_W\ge0\). & P-N; {[}V3{]} \\
G6 & Unidentified dimension \((C-1)\binom R2\) & \textbf{Correct} & 2a Lemma 1 \\
G7 & R×C gives ``more opportunity'' & \textbf{Correct in a precise sense:} P-R and P-S. The caveats: W1 vs W2 for \(R\ge3\); N2 no longer implies \(\delta_W\) for \(R\ge3\); a known table pins only \((R-1)(C-1)\) directions. & §5; {[}V5{]}, {[}V7{]} \\
G8 & 2×3: adding-up gives \(\sum_ct_c=0\) & \textbf{Correct but not new:} it is built into \(K\). The gain comes from joint restrictions across columns, which are strictly tighter than the best single-column restriction in 18\% of designs {[}V6{]}. M4 reported about half, but under a different, model-free set-up. & M4; {[}V6{]} \\
G9 & Joint sampling uncertainty & \textbf{Correct concern; handled} if all rows are written in \(k\) and enter one inference stack. Proved for the conservative interval and the projection interval in fixed-dimension regimes (P3 §0.5); the CLR interval as implemented is algorithmic (Conjecture P3.5(b3)) & P3.5(a), (b2); (b3) \\
\end{longtable}

\hypertarget{pG-proof-of-g2-22-any-mean-level-model}{%
\subsubsection{2.2 Proof of G2 (2×2, any mean-level model)}\label{pG-proof-of-g2-22-any-mean-level-model}}

\begin{enumerate}
\def\labelenumi{\arabic{enumi}.}
\tightlist
\item
  EGGHM Prop. 4: \(D_{ER}-D=\delta_W/\operatorname{Var}_p(X)\), and EGGHM Prop. 6: \(D_{NM}-D=-\delta_B/(\mu(1-\mu))\). Both are accounting identities for any \(m\) (PHASE2c Lemma M1), so they hold for the model-implied quantities. {[}M1; checked to \(3\times10^{-16}\) in V1{]}
\item
  With \(\operatorname{Var}_p(X)>0\) and \(0<\mu<1\): \(\delta_W\ge0\iff D\le D_{ER}\), and \(\delta_B\ge0\iff D\ge D_{NM}\). {[}1{]}
\item
  Under LCE, \(D\) is affine and non-constant in \(w_1\) when \(\gamma<1\), so the JKST set maps to \([L_J,U_J]\). Intersecting with the half-lines from step 2 gives G2. CR is the union of the two same-sign cases, i.e.~\(D\) lies between \(D_{NM}\) and \(D_{ER}\). ∎
\end{enumerate}

\textbf{How often each adds information} {[}V2{]}, over 2000 random LCE designs with true signs imposed:

\begin{longtable}[]{@{}
  >{\raggedright\arraybackslash}p{(\columnwidth - 6\tabcolsep) * \real{0.2500}}
  >{\raggedright\arraybackslash}p{(\columnwidth - 6\tabcolsep) * \real{0.2500}}
  >{\raggedright\arraybackslash}p{(\columnwidth - 6\tabcolsep) * \real{0.2500}}
  >{\raggedright\arraybackslash}p{(\columnwidth - 6\tabcolsep) * \real{0.2500}}@{}}
\toprule\noalign{}
\begin{minipage}[b]{\linewidth}\raggedright
Restriction
\end{minipage} & \begin{minipage}[b]{\linewidth}\raggedright
Strictly tightens JKST
\end{minipage} & \begin{minipage}[b]{\linewidth}\raggedright
Redundant
\end{minipage} & \begin{minipage}[b]{\linewidth}\raggedright
Median width ratio
\end{minipage} \\
\midrule\noalign{}
\endhead
\bottomrule\noalign{}
\endlastfoot
\(\delta_W\) sign & 83\% & 17\% & 0.70 \\
\(\delta_B\) sign & 89\% & 11\% & 0.67 \\
CR, sign unknown (union) & 80\% of the 1060 designs where CR holds & --- & 0.69 \\
\end{longtable}

\textbf{Redundancy diagnostic.} ``Redundant'' happens exactly when the JKST-implied interval for that \(\delta\) excludes 0. The package should print this table (``signs already implied by the data'') before any sign is imposed.

\begin{center}\rule{0.5\linewidth}{0.5pt}\end{center}

\hypertarget{pG-new-results-for-the-extension-modules}{%
\subsection{5. New results for the extension modules}\label{pG-new-results-for-the-extension-modules}}

\hypertarget{pG-sandwich-lemma-p-s-any-r-one-column}{%
\subsubsection{5.1 Sandwich lemma (P-S), any R, one column}\label{pG-sandwich-lemma-p-s-any-r-one-column}}

Write \(q_{rs}(x):=(Qx)_r-(Qx)_s\), which is identified, and \(g_{rs}(x;k):=(Kx)_r-(Kx)_s\), which is linear in \(k\). Define
\[O_{rs}(x):=m_r(x)-m_s(x)=(Hx)_r-(Hx)_s,\qquad L_{rs}(x):=\sum_tx_t\,\partial_{e_r-e_s}m_t(x)=(x^\top H)_r-(x^\top H)_s.\]
\(O\) is the between-group gap. \(L\) is the composition-weighted contextual derivative when a resident of group \(s\) is replaced by one of group \(r\).

\textbf{Proposition P-S (proved).}
* (a) \(O_{rs}(x)+L_{rs}(x)=2q_{rs}(x)=\partial_{e_r-e_s}\tau(x)\). This is the local Monotone-EI decomposition.
* (b) \(O_{rs}(x)=q_{rs}(x)+g_{rs}(x;k)\) and \(L_{rs}(x)=q_{rs}(x)-g_{rs}(x;k)\).
* (c) Hence \(O_{rs}\ge0\) and \(L_{rs}\ge0\) at \(x\) \(\iff\) \(\lvert g_{rs}(x;k)\rvert\le q_{rs}(x)\). This is non-empty only if \(q_{rs}(x)\ge0\), which is a data test.

\emph{Proof.}
1. \(\tau(x)=x^\top Hx\), so \(\nabla\tau=(H+H^\top)x=2Qx\). {[}def{]}
2. \(x^\top H=(H^\top x)^\top\) and \(H^\top=Q-K\), since \(K\) is antisymmetric. {[}def{]}
3. Substitute \(H=Q+K\) into \(O\) and \(H^\top=Q-K\) into \(L\). This gives (b); adding them gives (a). {[}alg{]}
4. (c) follows from (b). ∎

Checked to \(4\times10^{-16}\) {[}V4{]}. With \(R=2\), \(x=\) the group-1 share and \(k=w_1\), this is exactly GPT's pair \(c_1-w_1+d_1x\ge0\) and \(w_1+d_1x\ge0\).

\hypertarget{pG-rank-lemma-p-r-binary-outcome-or-per-column}{%
\subsubsection{5.2 Rank lemma (P-R), binary outcome or per column}\label{pG-rank-lemma-p-r-binary-outcome-or-per-column}}

For a pair \((r,s)\), \(g_{rs}(x;k)=\sum_t(K_{rt}-K_{st})x_t\). As \(x\) ranges over a full-dimensional subset of the simplex, the functionals reached are \(K_{rs}\) and \(K_{rt}-K_{st}\) for \(t\ne r,s\): that is, \(R-1\) independent directions.

\textbf{Proposition P-R.}
* (a) One pair restricts an \((R-1)\)-dimensional subspace of \(\mathcal K_c\). \textbf{Proved} (above).
* (b) If the pairs form a star tree \(\{(0,s)\}\), all \(\binom R2\) directions are reached. \textbf{Proved}: \(K_{0s}\) comes directly, and then \(K_{st}=K_{0t}-(K_{0t}-K_{st})\).
* (c) The same holds for any spanning tree. \textbf{Checked} for chain trees with \(R=3,4,5\) {[}V7{]}; general proof by induction on the tree (R5).
* (d) Combined with P-S(c), M-ORD ∩ M-LOC on a spanning tree bounds \(K_c\) with no Duncan--Davis rows. \textbf{Follows from (b) and P-S} for star trees.

Rank is necessary, not sufficient, for tightening: it says which directions can be cut, and the signs and the data decide whether they are. {[}V5{]} shows the practical effect in 3×2. With true restrictions only, ordering and local rows tighten about 2/3 of group-rate targets (median width 0.77 of DD under both). With near-1-D support they tighten more often (≈94\%, median 0.47), because DD alone is weaker there.

\hypertarget{pG-nesting-in-22-under-lce-p-n}{%
\subsubsection{5.3 Nesting in 2×2 under LCE (P-N)}\label{pG-nesting-in-22-under-lce-p-n}}

Let \(a=\mu\operatorname{Var}_{w^1}(X)\), \(b=(1-\mu)\operatorname{Var}_{w^0}(X)\), and \(x^*=a/(a+b)\).
* (a) Signed slopes (\(w_1\ge0\), \(w_1+d_1\ge0\)) ⇒ weighted-local on the hull. \textbf{Proved}: \(w_1+d_1x=(1-x)w_1+x(w_1+d_1)\).
* (b) \(\delta_W=(a+b)(w_1+d_1x^*)\). \textbf{Proved} (M5(c); asserted in V3). Hence weighted-local on the hull ⇒ \(\delta_W\ge0\) whenever \(x^*\in[x_L,x_U]\), and the implication is strict when \(d_1\ne0\) and \(x^*\) is interior.
* (c) \textbf{Conjecture C1:} \(x^*\) always lies in the composition hull. Evidence: 2000 random designs plus an adversarial search over 50,000 discrete designs with 2--4 support points and extreme weights found no violation beyond \(10^{-6}\) {[}V3{]}. A pairwise representation \(\operatorname{Var}_w(X)=\tfrac12\sum w_iw_j(X_i-X_j)^2\) writes \(x^*\) as a weighted mean of pair ratios, but individual ratios can leave the hull, so a proof needs a cancellation argument (R1).

\begin{center}\rule{0.5\linewidth}{0.5pt}\end{center}

\hypertarget{pG-proof-register}{%
\subsection{8. Proof register}\label{pG-proof-register}}

Status: \textbf{P} proved · \textbf{C} checked numerically · \textbf{J} conjecture · \textbf{O} open.

\begin{longtable}[]{@{}
  >{\raggedright\arraybackslash}p{(\columnwidth - 8\tabcolsep) * \real{0.2000}}
  >{\raggedright\arraybackslash}p{(\columnwidth - 8\tabcolsep) * \real{0.2000}}
  >{\raggedright\arraybackslash}p{(\columnwidth - 8\tabcolsep) * \real{0.2000}}
  >{\raggedright\arraybackslash}p{(\columnwidth - 8\tabcolsep) * \real{0.2000}}
  >{\raggedright\arraybackslash}p{(\columnwidth - 8\tabcolsep) * \real{0.2000}}@{}}
\toprule\noalign{}
\begin{minipage}[b]{\linewidth}\raggedright
ID
\end{minipage} & \begin{minipage}[b]{\linewidth}\raggedright
Claim
\end{minipage} & \begin{minipage}[b]{\linewidth}\raggedright
Status
\end{minipage} & \begin{minipage}[b]{\linewidth}\raggedright
Strategy / check
\end{minipage} & \begin{minipage}[b]{\linewidth}\raggedright
Blocks
\end{minipage} \\
\midrule\noalign{}
\endhead
\bottomrule\noalign{}
\endlastfoot
R1 & C1: \(x^*\in\) hull (2×2), hence weighted-local ⇒ \(\delta_W\ge0\) under LCE & J (C {[}V3{]}) & pairwise variance representation plus a cancellation argument; try two-point designs first & the M-LOC ⇒ SIGN-W cell \\
R2 & An added restriction moves an endpoint iff it removes the whole optimal face (one feasibility LP, with the target intercept); for one row, iff every optimal dual is positive (a solver's dual alone does not certify); implied ⇔ tightens no target; 2×2 closed form = G2; ``redundant'' = zero gain for this target & \textbf{P} (5a §2; \texttt{verify\_R2.py}; reviewed internally and externally (Batch B), corrections applied) & --- & redundancy diagnostic \\
R3 & \textbf{Pooled} LCE ∩ any closed convex restriction on \(k\) is sharp for the model-implied target (full law, and for the regression); CR is a union, a single slab for W1/2×2, and may have a gap otherwise; the conjunction can be rejected under LCE (Remark R3.2) & \textbf{P} for the pooled model (5a §1; \texttt{verify\_R3.py}; reviewed internally and externally (Batch B), corrections applied). Grouped rungs: R16 & --- & M-SIGN reporting \\
R4 & Refutability classification for \(R\ge3\) (which signs, orderings and local rows can produce emptiness, model-free and under LCE) & O & construct feasible completions; counterexamples & §4.2 codes, tests \\
R5 & P-R(c): any spanning tree of pairs reaches all directions; behaviour with lower-dimensional support & P (star); C (chain) {[}V7{]} & induction on tree leaves; support rank & step 4 \\
R6 & When joint column restrictions in 2×C are strictly tighter than the best single one; ordinal cumulative version & C {[}V6{]} & LP geometry in the \((C-1)\)-dimensional \(t\)-space & step 3 \\
R7 & Global N2 LP: outer validity with estimated τ (prespecified degree); binned version via 2c Lemma N2.4 & O (Lemma N2.4 proved; sharp version Conjecture N2.2′) & uniform confidence band for τ, then a monotone envelope & step 8 \\
R8 & B1(M), κ, CAR-δ: each is a valid outer bound, monotone in its dial; B1 rank condition & partly (NA) & per module & steps 6, 7, 10 \\
R9 & Consistency of the M-SHARP series bound; never wider than the current outer set. It conditions on \(X\) only, so it is ``tighter than \(\Theta_O\)'', not sharp, when size is informative (2c Remark N2.2.1) & O & sieve argument & step 1b \\
R10 & CLR validity with added rows, unions (min/max over LPs), near-emptiness, and the actual-rate allowance with new rows & By version (5a §3.2 table): conservative (a) \textbf{P}; full-candidate CLR (b1) cited, under the augmented Condition P; fixed-superset and projection (b2) \textbf{P}; CLR with simulated candidates (b3) \textbf{algorithm} (conjecture). (e) naive bracket under-covers: \textbf{P} (analytic, referee B). Regime: fixed dimension only (P3 §0.5). Package 0.3.0: projection rank guard and zero-variance-candidate check implemented & --- & all code \\
R11 & DiD: DiD-S sharp within FD; FD4 ∩ Δ4; LCE-levels ∩ sign-Δ; union inference & DiD-S validity P; rest O & as in 2d & §7 \\
R12 & P-S (sandwich) & P {[}V4{]} & --- & --- \\
R13 & G2 table (JKST + signs) & P (§2.2) {[}V1, V2{]} & --- & --- \\
R14 & Single-index dimension \(C-1\) & J (NA check) & generic-rank argument & T2.4 \\
R15 & Integer unit tables: N1 and Δ4 with integer cells (discrete target sets; integer branch feasibility). 2×2 with margins only: lattice of step \(1/(N\mu(1-\mu))\) inside the continuous interval (2c Remark M2.1) & O (2×2 margins-only case P) & integer programming; small tables first & integer-exact reporting \\
R16 & Grouped rungs: R3′ (regression sharpness under M1/M1(δ)/M2 via 2b Thm 1′, with \(A_d\) and all free coordinates) and R10′ (inference, with the regime caveat) & O & 2b Thm 1′; tests on noise-free grouped designs & removing ``experimental'' for grouped M-SIGN \\
R17 & P3.4′ (round 2, Nov 2026): limit of the plug-in endpoint along a triangular array of designs (changing hulls), under a localisation condition on the gradients and the slacks \(\sqrt n(\theta_\lambda-\varphi_U)\); P3.4 restated for a fixed constraint system & \textbf{P} (P3 §4; \texttt{verify\_phase3.py} ``P3.4′'') & continuous mapping after Lemma P3.1 & uniform limit theory (P3 §10, open) \\
R18 & P3.5(c) (round 2, Nov 2026): variance-adaptive (Bhatia--Davis) residual allowance. (c1) valid when \(\bar m\) is bounded by a P3.5(a) or (b2)(ii) interval on the rung's own constraints and \(\max_iw_i^2/\bar s^2\to0\); (c2) with a CLR interval, algorithmic & (c1) \textbf{P}; (c2) inherits (b3) (P3 §5; \texttt{verify\_phase3.py} ``P3.5(c)'') & union bound; Chebyshev/Lindeberg relative to \(\bar s\) & the projection column for \(B\) (package: \(\bar m\) projection interval on the rung alone; note near 0 or 1) \\
\end{longtable}

\textbf{Standard for ``proved'':}
1. The proof is written in the numbered-step format used in PHASE2c, with a justification tag on each step.
2. It is checked against brute force or the LP, including adversarial designs.
3. An independent fresh-session review re-derives it (S12).
4. The §4 tables are then updated.

Until then, the package labels any module that depends on the claim as \texttt{experimental} in \texttt{meta\$proof\_status}, and \texttt{ladder()} prints that label.

\endgroup


\end{document}
